% =====================================================================
%  H. Brøchner
%  On the Religious in its Unity with the Human.
%  A Positive Supplement to "The Problem of Faith and Knowledge".
%  Copenhagen: P. G. Philipsen, 1869.   [IV] + 211 pp.
%
%  ENGLISH TRANSLATION
%  Status: IN PROGRESS. Translated from transcription.tex (the Danish
%  diplomatic transcription, COMPLETE), never from the scan.
%  Structure mirrors transcription.tex 1:1: the same divisions and heads,
%  the same paragraphing, page markers % --- p. N --- at the same joints,
%  footnotes at the same anchors with the same per-printed-page reset, and
%  the transcribed centred rules at the same eighteen places.
%
%  This is the positive sequel to `problemet-tro-viden' (1868), which is
%  complete in this collection in both languages. The conventions below are
%  INHERITED from that volume's translation so the pair reads as a pair;
%  where this book forces a new decision it is marked NEW below.
%
%  RIGHTS: clear. Brøchner died 1875; long in the public domain.
%
%  ---------------------------------------------------------------
%  TRANSLATION CONVENTIONS
%  ---------------------------------------------------------------
%  Standing method: ../../../TRANSLATION-PLAYBOOK.md.
%
%  INHERITED from problemet-tro-viden (do not diverge without recording it):
%   * Erkjendelse -> "cognition"; Viden -> "knowledge". The distinction is
%     load-bearing in both books -- the sibling's title is "The Problem of
%     Faith and Knowledge" -- and the two words are never swapped.
%   * ophæve / Ophævelse -> "abolish" / "abolition", or "supersede" where
%     the sense is a superseding rather than an annulling. NOT "sublate":
%     the sibling volume avoids it, and so does this file.
%   * Quotation marks: Danish „…“ -> English curly doubles ``…''.
%   * Emphasis: every \emph{} of the Danish is carried over as \emph{}.
%   * Latin -> \textit{}; Greek glyphs copied verbatim from the Danish and
%     the surrounding prose translated (textalpha is loaded for this).
%   * The old id-est mark ɔ: -> "i.e."
%   * Em-dash ---. Proper names unchanged.
%   * Register: scholarly, moderately literal but readable. Brøchner's
%     syntax is long and nominalised; compound coinages are kept rather
%     than dissolved, to match the sibling.
%
%  NEW, forced by this book:
%   * ANTIQUA LATIN. The Danish distinguishes SPACED antiqua
%     (\emph{\textit{…}}) from TIGHT antiqua (\textit{…}) -- "causa sui" is
%     spaced on p. 68, "modus" tight on p. 83, "conditio sine qua non" on
%     p. 59 half and half. In English, italic no longer marks the Latin off
%     from the surrounding Fraktur, because there is no Fraktur; italic now
%     does double duty for Latin AND for emphasis. The distinction is kept
%     mechanically all the same -- \emph{\textit{…}} where the Danish has
%     it, \textit{…} where it does not -- so the two files stay
%     comparable, and the loss of the fount contrast is recorded here
%     rather than silently absorbed.
%   * ANTIQUA GUILLEMETS. Where the printing puts «…» round antiqua Latin
%     rather than the usual „…“ -- («Deum pati»), p. 41 -- the guillemets
%     are KEPT as printed, per the playbook's rule for foreign-language
%     guillemet terms. They are the one exception to ``…''.
%   * LETTERSPACING THAT BREAKS MID-WORD. The compositor's spacing is
%     inconsistent across line breaks and sometimes starts or stops inside
%     a word (Aands-functioner p. 6, \emph{Tids}bestemmelsen p. 18,
%     Væsen\emph{stotaliteten} p. 19, \emph{Livsforhold i Be}tydning
%     p. 32). Where the English word divides at a comparable point the
%     partial run is reproduced; where it cannot, the emphasis is placed on
%     the whole English word and a comment records what the printing does.
%     Either way the RUN COUNT is preserved, so the parity check holds.
%   * HEADS. The book's Indhold and its printed heads disagree in eight
%     places (listed in transcription.tex §3). Several of the divergences
%     are Danish-internal and vanish in English -- "Det"/"De" at p. 90,
%     the enumerator styles 1)/1. -- so they are recorded in comments at
%     the Indhold and at the head, not manufactured in the English.
%   * INDHOLD. The book's own contents leaf is translated as a list rather
%     than replaced by \tableofcontents (which is what the sibling volume
%     does), because here it is a bibliographic object in its own right:
%     it carries six rubric lines that appear nowhere in the body, two body
%     heads that it omits, and a page range that contradicts itself. All of
%     that is reproduced.
%
%  BUILD: pdflatex (libertinus + libertinust1math), like the Danish.
%  Sandbox compile-test per TRANSCRIPTION-PLAYBOOK.md §5.
% =====================================================================

\documentclass[12pt,a4paper]{book}
\usepackage[utf8]{inputenc}
\usepackage[T1]{fontenc}
\usepackage{libertinus}
\usepackage{libertinust1math}
\usepackage{textalpha}   % Greek (τέλος etc.), copied verbatim from the Danish
\usepackage{graphicx}
% The Danish needs \reflectbox{c} for the id-est mark ɔ:. The English
% renders that mark "i.e.", so the declaration is not needed for the body;
% it is kept so the file compiles unchanged if a quotation ever restores it.
\DeclareUnicodeCharacter{0254}{\reflectbox{c}}
\usepackage[english]{babel}
\usepackage[protrusion=true,expansion=true]{microtype}
\usepackage{setspace}
\usepackage[margin=1.2in]{geometry}
\usepackage{enumitem}
\usepackage{fancyhdr}
\setlength{\headheight}{28pt}
\addtolength{\topmargin}{-13.5pt}
\usepackage{hyperref}

\hypersetup{
  pdftitle={On the Religious in its Unity with the Human},
  pdfauthor={H. Brøchner},
  hidelinks
}

% The book prints *) and **), restarting on every printed page. The counter
% is reset at the page marker of every note-bearing page, exactly as in the
% Danish. (footmisc[perpage] is wrong here: it resets on the TYPESET page.)
\renewcommand{\thefootnote}{%
  \ifcase\value{footnote}\or *)\or **)\or ***)\else \arabic{footnote})\fi}

\pagestyle{fancy}
\fancyhf{}
\fancyhead[LE,RO]{\thepage}
\fancyhead[RE]{\textit{On the Religious in its Unity with the Human}}
\fancyhead[LO]{\textit{\leftmark}}

\setstretch{1.3}

\title{\textbf{On the Religious\\[2pt] in\\[2pt] its Unity with the Human}\\[10pt]
  \large A Positive Supplement to\\ ``The Problem of Faith and Knowledge''\\[10pt]
  \normalsize\textit{Om det Religiøse i dets Enhed med det Humane}}
\author{Dr.\ H.\ Brøchner\\[4pt]
  \small Professor of Philosophy\\[4pt]
  \small Translated from the Danish}
\date{Copenhagen: P.\ G.\ Philipsen, 1869}

\begin{document}
\maketitle
\thispagestyle{empty}
\newpage

% =====================================================================
%  CONTENTS  (printed pp. [III]--IV)
%  The book's own contents leaf, translated. Its printed page ranges are
%  reproduced EXACTLY as printed, including the one that contradicts
%  itself (entry 2 given 167--201 while its own sub-entry a) is given
%  165--167), and the six rubric lines that appear here and nowhere in the
%  body. Top-level entries are bold Fraktur in the print; the rubric lines
%  are letterspaced; the Greek enumerators α) β) γ) are antiqua italic.
% =====================================================================
\chapter*{Contents}
\addcontentsline{toc}{chapter}{Contents}
\markboth{Contents}{Contents}

\begingroup
\setstretch{1.15}
\begin{list}{}{%
  \setlength{\leftmargin}{2.2em}\setlength{\itemindent}{-2.2em}%
  \setlength{\itemsep}{2pt}\setlength{\parsep}{0pt}}

\item[] \hfill Page

\item \textbf{Introduction} \dotfill 1--4.

\item \textbf{I. On the Relation of Unity between True Cognition and True
  Action} \dotfill 5--29.

\item \textbf{II. On the Presence of the Essential Marks of the Religious
  in its True Sense in True Cognition and True Action, and on their
  Relation of Unity to the Former} \dotfill 30--63.

\item[] \hspace{1.2em} 1. The Essence of the Religious \dotfill 30--43.
\item[] \hspace{1.2em} 2. Cognition's Relation to the Religious according
  to its Ideal Determinateness \dotfill 43--51.
\item[] \hspace{1.2em} 3. The Relation of the Will and of Action to the
  Religious according to its Ideal Determinateness \dotfill 51--56.
\item[] \hspace{1.2em} 4. The Reconciliation of the Human Spirit in the
  Ethical-Human Life Resting upon True Cognition \dotfill 56--63.

\item \textbf{III. The Central Ethical-Religious Problems, Apprehended
  from the Standpoint of Knowledge} \dotfill 64--201.

\item[] \hspace{1.2em} \emph{Preparatory Developments of Concepts:}
\item[] \hspace{1.2em} 1. The Concept of the Absolute \dotfill 64--89.
% The Danish Indhold reads „De menneskelige Væsens Begreb“ where the
% printed head on p. 90 reads „Det menneskelige Væsens Begreb“. The
% divergence is a Danish article form and has no English counterpart; it is
% recorded here and at the head rather than manufactured in the English.
\item[] \hspace{1.2em} 2. The Concept of the Human Being, Determined in
  Relation to the Absolute and to the Totality of Existence \dotfill 90--97.
\item[] \hspace{1.2em} \emph{Comparison of the Conception Developed with}
% Indhold: „1) Ludvig Feuerbachs“ / „2) den Hegelske Philosophies“ -- both
% genitive and both carrying an enumerator, where the run-in heads on
% pp. 97 and 126 have neither. The genitive hangs on the rubric line above.
\item[] \hspace{1.2em} 1) Ludvig Feuerbach's \dotfill 97--126.
\item[] \hspace{1.2em} 2) that of the Hegelian philosophy \dotfill 126--128.
\item[] \hspace{1.2em} \emph{Development of the Problems:}
\item[] \hspace{1.2em} 1. On the Freedom of the Will \dotfill 128--143.
\item[] \hspace{1.2em} 2. On the Essence and Possibility of Evil \dotfill 144--152.
\item[] \hspace{1.2em} 3. The Self-Dissolution of Evil as the Possibility
  of Reconciliation \dotfill 152--159.
\item[] \hspace{1.2em} \emph{Further Developments towards the Problem of Evil:}
% Indhold: without „i det Onde“, which the printed head on p. 159 carries.
\item[] \hspace{1.2em} 1. The Relation between the Moment of Necessity and
  the Moment of Freedom \dotfill 159--164.
\item[] \hspace{2.4em} a) in the Realisation of Evil \dotfill 159--162.
\item[] \hspace{2.4em} b) in the Abolition of Evil \dotfill 162--164.
% --- p. IV ---
% As printed: entry 2 is given 167--201 while its own sub-entry a) is given
% 165--167, and the printed head stands on p. 164. Not in the Trykfeil.
% Reproduced as printed.
\item[] \hspace{1.2em} 2. Reconciliation, its Source and Goal \dotfill 167--201.
\item[] \hspace{2.4em} a) that which is operative in Reconciliation \dotfill 165--167.
\item[] \hspace{2.4em} b) that which is to be attained in Reconciliation \dotfill 167--168.
\item[] \hspace{2.4em} \emph{Conditions for the Completeness of Reconciliation:}
% Indhold: „Skyldens Ophævelse“; the printed head on p. 168 reverses the
% word order, „Ophævelsen af Skylden“. Both render as "the abolition of
% guilt"; the variance is recorded, not manufactured.
\item[] \hspace{2.4em} α) The Abolition of Guilt \dotfill 168--176.
\item[] \hspace{2.4em} β) The Ethical as a Determination of Totality for
  Human Life \dotfill 177--182.
\item[] \hspace{2.4em} γ) The Eternity-Determinateness of the Individual
  Human Life Answering to the Concept of the Ideal \dotfill 182--190.
\item[] \hspace{1.2em} \emph{Comparison of the Conception Developed of
  Reconciliation with}
\item[] \hspace{1.2em} 1) the abstract metaphysical conception of
  Reconciliation \dotfill 190--192.
\item[] \hspace{1.2em} 2) the Christian doctrine of Reconciliation \dotfill 192--201.

\item \textbf{Conclusion} \dotfill 202--210.

\end{list}
\endgroup
% NOTE, as in the Danish: the Indhold has NO entry for the two body heads
% „a) Opfattelsen af Synden og af dens Virkning.“ (p. 193) and
% „b) Opfattelsen af Forsoningen og af Ophævelsen af Syndens Virkninger.“
% (p. 195). It runs straight from 2) 192--201 to Slutning. Both heads are
% real and are set in the body, here as there.

% =====================================================================
%  BODY
% =====================================================================

% The body of printed pp. 1--4 carries NO printed head whatever: the text
% begins directly, with a two-line ornamental initial. „Indledning“ is the
% Indhold's word for it, supplied editorially -- as in the Danish, where the
% head is bracketed for the same reason.
\chapter*{[Introduction]}
\addcontentsline{toc}{chapter}{[Introduction]}
\markboth{Introduction}{Introduction}

% --- p. 1 ---
% p. 1 carries NO printed head of any kind: the text begins directly with a
% two-line ornamental initial N (set here as an ordinary letter, per the
% book's conventions). The signature numeral „1“ at the foot is not
% transcribed. „Indledning“ is the Indhold's word only, supplied editorially
% by the skeleton above, and bracketed there as in the Danish.
The present work, whose publication has been delayed by circumstances
independent of the author's will, attaches itself most closely to the
historical-critical treatise: Problemet om Tro og Viden, in which
% The Danish sets the title PLAIN -- no letterspacing, no other device -- and
% so does the English. Italicising it here would add an emphasis run the
% printing does not have and break parity with the Danish.
reference is made to it in several places, and reproduces, like that treatise,
in all essentials unchanged, the content of a university lecture course. It
forms the positive supplement to that treatise, whose result, from the nature
of the case, had essentially to be negative; and taking its starting point
from this negative, it seeks, presupposing the developments by which that
result was established, the definitive determination of faith's relation to
knowledge.
% The Danish prints „ved hvilket“ where one would expect „ved hvilke“. It is
% a real word in the language and so stands as printed and is not logged as a
% defect; the disagreement is Danish-internal and has no English counterpart.

The result, of a negative kind, to which all the developments given in the
earlier work (in the ``preliminary determinations of concepts'', in the
critique of theology's standpoint and of the attempt to bring about a
reconciliation by a distinction of principle) have led, may be summed up in
the following points:

% Displayed enumeration, printed with the numeral at the left margin and the
% text hanging: „1.“ and „2.“, both wholly on p. 1.
\begin{enumerate}[label=\arabic*., leftmargin=*, labelsep=1em,
                  topsep=0.35\baselineskip, itemsep=0.35\baselineskip]
\item Between faith in the sense of positive religion, i.e. faith as ``faith
  in revelation'', and knowledge there arises --- whether the standpoint be
  taken in the positively religious or in science --- in virtue of their
  different content and their different source, an unabolishable incongruence,
  an irreconcilable opposition.
% The Danish sets „Incongruents“; the initial is the Fraktur I/J sort,
% normalised to I per the house convention (ABBYY reads it „Jncongruents“).
\item It therefore becomes impossible that faith (in the sense indicated) and
  knowledge should be united without discord in the same consciousness, should
  preserve their validity for the same personality.
\end{enumerate}

% --- p. 2 ---
This negative result now converts itself immediately into a positive one
differing from it only formally. If the opposites are irreconcilable in
consciousness, then it becomes a necessity that a \emph{choice} be made
% letterspaced: „v æ l g e s“ only; the following „disse“ is set tight. The
% English run falls on "choice", the word the Danish spaces.
between these opposites, of which one must be the expression of the human
spirit's highest energy, must enclose within itself the goal of its highest
striving.

This choice, whose necessity under that presupposition is not to be evaded, is
then, with respect to its character, to be formally determined thus: that in
accordance with the demand which lies in the unity of the human being and of
human self-consciousness, it must be such a choice as one in which the
spirit's fundamental activities are in harmony within themselves and with one
another, in which therefore the being, undivided and unmutilated, is in unity
with itself in its cognising and its willing.

When this determination is taken as the norm, it becomes clear what direction
the choice can and must take, what can and must become the object of the
choice.

For from the developments given in the earlier work it emerges that faith, as
faith in the sense of the positive religions, \emph{cannot} be chosen if the
demand set up above is to be fulfilled. That it cannot be so chosen lies
partly in its relation to the practical, partly in its relation to the
theoretical. It cannot be chosen because, in the practical respect, it cannot
% Both occurrences of „Theoretiske“ on this page are set TIGHT in the Danish,
% not letterspaced; the BATCH-AGENT calibration set wrongly lists one as
% spaced. No \emph{} on "theoretical" here, accordingly.
become the forming principle for the wholeness of an ethical human life,
cannot form human nature. For on the standpoint of positive religion the
source of the ethical demand becomes the supranatural, transcendent subject of
revelation, distinguished from the human being, whose will is absolute
arbitrariness; and the demands of this will cannot therefore be congruent with
the essential demand of the human spirit. And when positive religion has been
clarified into the highest form of a religion of spirit, then its ``ethical''
demand requires the denial of the natural, that is, of something fundamental
in the human being, of something that \emph{must} assert itself in its
essentiality. --- If we look to the theoretical, then faith cannot be chosen
because it cannot round itself off into a totality, because, while it cannot
make itself inde\-%
% --- p. 3 ---
\setcounter{footnote}{0}% (no space: the word is broken across pp. 2/3, as in the Danish)
pendent of knowledge, it cannot take up into itself what is essential in
knowledge, cannot see through and explain knowledge\footnote{Cf.\ Problemet om
Tro og Viden, pp.\ 195; 210 f.}.
% printed mark: *) --- the only note on p. 3, set below a short rule.
% The work-title reference is left in Danish, per the playbook: it names the
% sibling volume, which this collection catalogues under its Danish title.
% The signature „1*“ at the foot of p. 3 is not transcribed.
Knowledge must then, as an essential determination in human nature,
inevitably assert itself as justified over against faith, and bring discord
into the consciousness that was to be filled by faith. A relation of
subordination on knowledge's part is excluded by the kind of the opposition
and by the essence of knowledge. When faith is chosen, therefore, the being's
harmony and its inner reconciliation cannot be attained.

But when the demanded goal cannot be attained through the choice of faith,
then there shows itself, through what was developed in the previous work, a
possibility that it can be attained through a choice of knowledge. For in that
knowledge is chosen, that is chosen which, by being an expression of the
universality of the human being, has the capacity to unfold itself as a
totality resting in itself with an inner infinity. In that knowledge is
chosen, that is chosen in which the being in its unity becomes manifest to
itself, that which in the indivisible human spirit holds the ideal principate,
and which therefore can become the norm for a true ethical \emph{wholeness of
life}. In that knowledge is chosen, there is finally chosen that which, as
enclosing the cognition of the spirit's \emph{whole} life and development of
life, can see through the essence of the finite forms of religion and
comprehend them in their springing from human nature in its unfolding, and
therefore does not retain a disturbing barrier of consciousness within them;
and which therefore also, in separating out within the historical forms of
religion the transient and finite from the abiding and essential, can stand in
harmony with and take up into itself this essential element, which is
precisely the religious according to its true and eternal significance.

How now, from this standpoint of the necessity of a choice, and of a choice so
determined, a \emph{reconciling} solution of the problem of the relation of
faith and knowledge becomes possible; how the right and validity of knowledge
can be asserted without faith being violated in its truth, without the
% --- p. 4 ---
religious that is demanded by a need of the human spirit being denied --- this
it will be the task of the developments that follow to make clear. Their
course is determined by the nature of the task. \emph{First} it will be
necessary to establish that there is present, conditioned by the unity of
essence, a relation of unity between true cognition and true action, ethical
action. \emph{Thereafter} it shall be shown that the essential marks of the
religious in its true and eternal sense --- faith and self-surrender --- are
inseparable from true cognition and true action, and that these are therefore
in harmony with the religious and bear it within themselves, so that human
life, according to its harmonious unity with itself in the totality of its
fundamental functions, can show itself as penetrated by the religious. And
\emph{finally} the possibility shall be shown that out of philosophy's
% „Og“ is set tight in the Danish; only „e n d e l i g“ is letterspaced, and
% it is broken „en=/delig“ across the printed line.
concepts of the concretely determined Absolute --- the highest determination
of knowledge --- and of the human being in its relation to the principle and
to real existence, there can be developed a view of life that can satisfy
ethically and religiously, and that there can be attained those essential
determinations of the ethical and religious relation of which it has been
asserted that they can be derived only from a principle of will and of faith
absolutely heterogeneous with the principle of knowledge. Here, then, the
central ethical-religious problems --- the problems of the freedom of the
will, of the possibility and actuality of ethical evil, and of reconciliation
and restitution --- will come up for discussion.

In connection with these developments there will at the same time be made
clear the relation of the view of life indicated to those conceptions which
offer real or apparent points of contact with it, and the critique of them
given, partly directly, partly indirectly.

\begin{center}\rule{2cm}{0.4pt}\end{center}
% RULE, p. 4: the Introduction closes with a short centred DOUBLE rule --- two
% parallel hairlines --- standing below the last line of type, the rest of the
% leaf blank. The same weight as the division-closing rules of pp. 29, 63, 201
% and 210, and unlike every other rule in the book. Set with the single \rule
% this edition uses throughout, as in the Danish.

\chapter*{I.\\ On the Relation of Unity between True Cognition and True Action.}
\addcontentsline{toc}{chapter}{I. On the Relation of Unity}
\markboth{I. On the Relation of Unity}{I. On the Relation of Unity}

% --- p. 5 ---
% p. 5 is the opening page of division I. The \chapter* head is in the skeleton
% and is not repeated here; the short centred rule below it is head ornament
% and is not transcribed at all. p. 5 carries no folio (by design) and opens
% with a two-line ornamental initial G, set as an ordinary letter.
%
% NOTE on letterspacing across a line break, which happens repeatedly in this
% division: where a letterspaced word runs over a line ending, the compositor
% sometimes spaced only the part standing on the first line and set the
% continuation tight, so \emph{} sometimes ends inside a word. The opposite
% also occurs. The practice is genuinely inconsistent and is not normalised in
% the Danish; in the English the emphasis falls on the corresponding element
% and the RUN COUNT is preserved, with a comment wherever the English word
% cannot divide where the Danish one does.
The fundamental relation between cognition and will, as primitive forms of
expression of the human being, has been indicated in the earlier work
(pp.\ 134 ff.), and to the determinations there given I may at this place
refer back. The conception there asserted was this: that the human being,
whose indivisible unity is posited by the concept of subjectivity and impresses
itself in the unity of self-consciousness, is according to its fundamental
concept a selfhood-energy, and that this determination of the selfhood-energy
is the pervading universal which gives itself its peculiar forms of activity
in cognition and will. In these peculiar forms of expression, then, the
universal determinateness of the being reveals itself; they become primitive
functions of the selfhood-energy; and while their originality encloses their
peculiarity and their self-validity, their relation to the universal
determination of the selfhood-energy --- the determination inherent in the
particular --- encloses this: that they are governed by the unity of the
being, that they reflect one another's formal determinatenesses, and that the
essential energy therefore preserves its unity in the twofold form of
appearance. By the \emph{equal} relation in which cognition and will have been
shown to stand to the determination of the energy, that confusion of concepts
is removed which, through an arbitrary and unclear extension of the concept
\emph{drive}, identifies drive with the energy and thereby procures for the
will an illusory precedence in the human being, without being able, then, to
explain the origin of cognition, or of consciousness generally;
% --- p. 6 ---
\setcounter{footnote}{0}% p. 6 carries TWO notes, printed *) and **)
without being able to remove the confusion that ``the drive'' is itself to
posit what the drive, in the \emph{proper} sense of the word, strives to hold
fast \emph{or} to remove; and without being able to account for the fact that
the first expression of the selfhood-energy, as unconscious organic forming,
continues \emph{after} the will has been actualised, and continues
\emph{independently} of the will. Through the critique of this one-sided and
unclear conception the primitivity both of cognition and of will has been
asserted; and when the question of the ideal principate was raised, this was
assigned to cognition in virtue of its relation to that fundamental
determination in the human being, \emph{essential universality}, by which the
human being categorically distinguishes itself from the other stages of the
development of life\footnote{Cf.\ Problemet om Tro og Viden, p.\ 137; Et Svar
til Prof.\ Nielsen, pp.\ 12 f.}.
% printed mark: *) --- first of the two notes on p. 6, below a short rule.
This ideal principate of cognition reveals itself also in
\emph{consciousness's} relation to the particular functions of spirit. As
functions of the human selfhood-energy they are \emph{accompanied by
consciousness}, and are precisely thereby functions \emph{of spirit}; but
% letterspacing: „A a n d s =“ at the line end is spaced, „functioner,“ on the
% next line is set tight. The English compound will not divide at the same
% point, so the run falls on the corresponding element, "of spirit".
consciousness stands according to its concept in another relation to cognition
and thought than, for example, to the will; it becomes the \emph{abstract,
formal,} expression of cognition, that which can be unfolded into an
\emph{appropriating cognition}, into \emph{comprehending} thinking, into the
transparency of spirit to itself in its content; and the fact that it enters
into relation to the different functions of spirit, in becoming the form of
their being-for-the-self, points precisely to the ideal principate of
cognition, and to the possibility of a transparency of the self to itself
through cognition in the self's various peculiar forms of
expression\footnote{With respect to the charge which a critic has brought
against me of confounding what is essentially different, I remark only --- what
every thoughtful reader of the earlier work will surely have seen --- that I do
not identify the \emph{immediate forms of consciousness} with unfolded
\emph{scientific cognition}, but maintain only that what comes forward in those
forms of immediacy is partly such as can be unfolded in clear cognition, and
partly a \emph{resulting} immediate, and more precisely such a one as has
resulted from a preceding \emph{cognition}. Just as little do I confound the
\emph{subjective} with the \emph{objective}; rather I maintain that the former,
as an expression of the subject's \emph{being}, has the same transparency for
cognition, i.e. for a \emph{theoretical} cognition, as has that immediate.}.
% printed mark: **) --- the second note on p. 6. Its text RUNS OVER onto p. 7
% in the print, where it fills the last six lines below the rule; it is set
% here in full at its call, and the page marker below stands at the beginning
% of p. 7's body text, as in the Danish.
% --- p. 7 ---
When this fundamental point --- the being's undivided unity in cognition and
will, and the ideal principate of cognition --- has already been won as a
result in the earlier work, then at this place the \emph{relative} opposition
is to be more closely elucidated which must come forward in virtue of the
peculiarity of the fundamental functions and in virtue of the conditions of
the development of spirit, but through which, in virtue of the fundamental
relation, harmony --- the unity grounded in essence --- must actualise itself
as governing the opposition.

\emph{The relative opposition} in which cognition and will come forward in
their independence is to be apprehended first as an opposition
\emph{of concept}, and thereafter as an opposition \emph{of facticity}.
% „B e g r e b s modsætning“ and „F a c t i c i t e t s modsætning“: in both
% compounds only the first element is letterspaced in the Danish and the
% second set tight. Two lines below, the whole word is spaced throughout.

\emph{The opposition of concept} between the two primitive functions has
already been indicated in ``Problemet om Tro og Viden'' (pp.\ 135 f.). It
showed itself as coming forward immediately in the difference of formal
determinateness; but the difference of form was seen not to constitute an
abstract opposition, since the determinations which expressed it pointed to
one another and developed out of one another, so that the fundamental
functions reflected one another's formal determinatenesses.
% PRINTER'S DEFECT here in the Danish: „hinandes“ as printed, for
% „hinandens“ --- a dropped n, fully legible at 300 dpi and read so by both
% OCR witnesses, with „hinandens“ set correctly on p. 5. Orthographic, and so
% not reproducible in English; recorded here.
Thereby it became possible that beneath the difference of form the
\emph{essential} unity of the content can be preserved: a unity which is
required by the fact that the same undivided being gives both forms their
content. When thus, in the ethical resolution, the cognition of duty is
converted into a willing of duty --- when therefore that which as an object of
cognition has the form of possibility is posited in the form of actuality in
the I --- then this \emph{determinate} content is in both forms
\emph{essentially} unchanged, and must be so already for this reason, that
otherwise the fulfilment of duty in the action would \emph{not} be a
fulfilment of this determinate thing cognised as duty; and in the preceding
% --- p. 8 ---
cognition the universal determination of possibility is so punctualised in
relation to the existing subject, and the object of cognition so brought into
relation to the cognising \emph{being}, so seen as an essential determination
which is required to be posited as a determination of reality, that the
conversion into a content of will cannot abolish the essential unity of the
content.

That the difference between cognition and will does not transform itself into
an \emph{abstract} opposition shows itself also indirectly when we elucidate
the attempt to assert such an opposition through the opposition between the
``theoretical'' and the ``practical consciousness'', between ``theoretical''
and ``practical cognition'' (thinking).

With respect to this theory we touch only upon what we have shown elsewhere:
the lack of clarity in the application of the categories of which it makes
itself guilty, in that the concepts consciousness, cognition and thinking are
used indiscriminately without clear delimitation, and cognition, which relates
itself to action, is confounded with the act of will itself, the resolution of
the will. Disregarding this confused application of the concepts, and
distinguishing what has been called practical cognition from the immediate
``practical consciousness'' --- which is partly a consciousness of the
subject's state in the practical respect, of the subject's resolution of will
and of the motive, the factually moving, and partly a consciousness of the
ethically a priori, of the practical demand --- we investigate how the sought
opposition between a theoretical and a practical cognition is to be
determined.

It will then be immediately evident here that there cannot be two essentially
distinct kinds of cognition, two absolutely heterogeneous forms of cognition,
since in that case the determination of cognition would transform itself into
a purely contentless name and in its duality could not be the property of the
same cognising being. An opposition which springs from the ground of the same
being must preserve within itself the essential unity. And therefore it is to
be held fast that cognition, \emph{as} cognition, in every form and in every
application, must preserve its essential determinateness.
% --- p. 9 ---
The opposition between a ``theoretical'' and a ``practical'' cognition then
becomes a merely relative opposition between a cognition which is determined
solely by the laws of thought and which has theoretical insight as its result,
and a cognition which is determined also by motives of will (laws of will) and
which has action, resolution, as its result.

The designation of practical cognition as determined by motives of will points
partly to this cognition's presuppositions and starting points (premisses),
partly to its being affected by the will in the course of its advance. Its
starting points are practical incitements: partly promptings of drive, partly
effects of an ethical a priori (grounds of conscience). From these starting
points, which stand as something given, ``practical'' cognition then advances
in the same way as the theoretical, which draws consequences from its
premisses, partly from unresolved ones and partly from transparent ones. ---
Cognition's being affected in the course of its advance may proceed from a
more clearly or a more obscurely conscious will, which either supplements
incomplete premisses, or corrects one-sided theories, or hinders true
cognition. In the first case the will supplies a starting point which, in
order to be used, must not stand in conflict with what is theoretically
cognised. The relation here forms an analogy to that in which, in theoretical
reasoning, an unproved current opinion (a prejudice) is used. In the second
case the will inserts a corrective which cognition acknowledges as its own,
and which forms the starting point for a new inference. Here too the
theoretical domain offers perfectly corresponding analogies (cf.\ Problemet om
Tro og Viden, pp.\ 144 f.). --- In the last case the relation is more
complicated. During cognition's advance the will works partly by bringing
forward instances that can prevent a conclusion, and partly by spacing out the
unfolding of the premisses from the conclusion, the resolution, whereby the
clarity of cognition can be obscured, and contradicting instances, specious
grounds and pretexts can take on the appearance of vali\-%
% --- p. 10 ---
dity. To this hindering there is formed, in the theoretical domain, the
analogy of the working of prejudices upon cognition.

The \emph{result} of ``practical cognition'' is partly the acceptance of a
doctrine which relates itself to the subject's action, to its practical life,
and partly an inference in the direction of action, which is converted into a
resolution. If the result is the acceptance of a doctrine, then this
acceptance comes forward as an inference which proceeds with necessity from
its premisses; but among these premisses there may be found, according to what
has been developed above, such as have been inserted by the will. If the
result is resolution, as resolution to the determinate action, then the
preceding act of cognition --- to which the resolution in the proper sense
always points back --- is, like the cognition upon which the acceptance of the
practical doctrine depended, and like the cognition which drew consequences
from practical premisses, in its form not essentially different from the
theoretical; but this would be impossible if an absolute heterogeneity
separated ``practical'' and ``theoretical'' cognition from one another. For
cognition is only what it is by its form; but by form is then of course
understood something other and more than the empty syllogistic form applied
\emph{externally} to an arbitrary content.

The result with respect to the cognition which relates itself to the practical
is therefore this: that in virtue of its form and in virtue of its relation to
the universal concept of cognition, and likewise in virtue of the relation of
the practical --- of the starting points and of the result --- to the human
\emph{being}, it must be essentially one with the cognition designated as
% letterspacing: „Menneske“ is set tight in the Danish, only „v æ s e n e t ,“
% is spaced. The English compound divides at the same point.
theoretical. All cognition is \emph{as} cognition theoretical. That cognition
can be applied in a practical direction, and in this application be affected
by the will, proves nothing for an essentially distinct practical cognition;
for a unity of the will with cognition comes forward \emph{everywhere}, its
reciprocal action with cognition becomes manifest in the domain of
indisputably theoretical cognition as well, and theoretical cognition can
\emph{as such} become an object of the will.
% --- p. 11 ---
That the relation is of this kind will become still more evident when the
\emph{marks of character} are examined by which the alleged abstract
opposition between theoretical and practical cognition is \emph{described}:
% The Danish sets „abstrakte“ with k here, where the book's usual form is
% „abstract“ (pp. 8, 13, 15). Orthographic, and not reproducible in English.
they will all show themselves as leading back from the relative difference to
the essential unity.

When the opposition is first characterised thus, that the \emph{theoretical}
consciousness (cognition) is \emph{impersonal} and the \emph{practical
personal}, then this opposition, as an abstract and exclusive opposition,
shows itself untenable already for this reason: that theoretical cognition, as
the expression of the unity of that being whose fundamental determination is
personality, must not merely be an expression of personality in general, but
must have for its final goal a cognition of personality in its connection with
the whole of existence and in its subordination under the principle of
existence. In this cognition of the personal self, the cognition of objective
real existence --- of that world of corporeality to which personality stands
in relation through its side of reality and its conditions of reality --- will
then also show itself as having ``personal'' significance. That determination
of theoretical cognition as impersonal encloses the consequence that
\emph{theoretical cognition, although a primitive expression of the self,
cannot become self-cognition}, inasmuch as it finds its limit in the very self
which is the source of this cognition. If, in order to maintain the assertion
of theoretical cognition's impersonal character, one were to lay the emphasis
on this, that the cognition of the self too, as theoretical cognition, must
remain in the form of the universal, and that therefore the personal self can
through it be cognised only in an impersonal manner --- then it would have to
be urged against this that the universal in such cognition is not held fast in
the original form of abstraction, but is constantly more and more punctualised
towards concretion, so that what remains ``impersonal'' in theoretical
cognition would designate only what remains relative in the punctualisation of
the universal; and this relative element is moreover such
% --- p. 12 ---
as is not merely abiding in theoretical cognition, but is and must be so in
the resolutions of the will as well, precisely because these are conscious and
are expressions of spirit, and as such are raised above that form of
singularity which belongs to natural existence and to natural determination by
drive.

The object of impersonal, theoretical cognition is determined as the
\emph{objectively valid}, that of personal, practical cognition as the
\emph{subjectively valid}, and an abstract opposition is made between the
objectively universally valid and the subjectively universally valid. When,
however, that self-contradicting limit upon theoretical cognition is removed,
it shows itself to comprehend also that which, as a ``subjectively valid'',
was set in opposition to the objectively valid, and to comprehend not merely
what, in virtue of its connection with cognised human nature, has validity in
its universality for all subjects, but also what, as determined in its
concretion by co-operating universal grounds, has validity for this determinate
subject.

Only in virtue of this relation can there be any talk of a ``subjectively
universally valid'', whereas the abstract opposition between the subjective
and the objective must lead to letting \emph{individual} subjectivity as such
determine the valid as it \emph{wills} to determine it, and so to letting the
subjective grounds become one with individually accidental grounds, promptings
of drive and the like --- to letting ``grounds of will'' become such as an
arbitrary will makes into grounds.

If the opposition between the two kinds of cognition is then determined thus,
that theoretical cognition is moved by \emph{necessity} while practical
cognition ``\emph{is grounded in freedom}'' --- the former, as lost in the
universal, being compellingly determined by incontrovertible grounds of
knowledge, the latter, as conditioned by regard for the self, resting on
grounds of conscience --- then this opposition too shows itself to be
dialectical. Theoretical cognition is governed by necessity where the
presuppositions of the inference are fully transparent, where the proof
acquires stringency through clear and complete pre\-%
% --- p. 13 ---
misses. But just as fully is ``practical cognition'' determined with necessity
where these conditions are present, where the presuppositions are given clear
and complete. What produces the appearance of an abstract opposition is
\emph{the difference in concretion} between the object of practical cognition
and certain spheres of the theoretical. In the domain of the concretely
ethical a seemingly unabolishable difference may persist, precisely on account
of the incompleteness of the premisses, in the same way as within the
theoretical, in the more concrete domains. The difference can maintain itself
because the manifold presuppositions are not clarified and are therefore
partially supplemented in an individual manner. Thus, under the same
presuppositions, the necessity is the same. And the necessity which determines
theoretical cognition is, when the cognition is \emph{comprehending} and as
such \emph{appropriating}, an expression of \emph{the lawfulness of the
cognising being itself}, and thus converts itself into \emph{freedom}. In the
form of unfree necessity theoretical cognition comes forward only where what
is given in sensation is determinative in its foreignness. But to this stage
there corresponds that stage in the development of practical ``cognition''
where the regard for the self which determines it is regard for something
obscurely drive-like, by which the will is brought under natural necessity.
Only as the expression of something essential has the subject's need, which
determines that ``cognition'', any warrant; but as an expression of the being
this need bears at once the stamp of true necessity and of freedom, while the
so-called ``freedom'' of determination by drive is unfree necessity, just as
the necessity of sensation is. In theoretical and in practical cognition
alike, therefore, \emph{both the} determinations come forward which were to
% letterspacing: „b e g g e  d e“ ends the line and is spaced in the Danish;
% „Bestemmelser,“ opening the next line is set tight. The English breaks at
% the same point.
constitute the sharp opposition; the real relation becomes obscure only
through this, that the individual who acts ``arbitrarily'' does not become
conscious of being, precisely in this action, unfree.

When finally the opposition is determined in this way, that theoretical
cognition is to be \emph{disin}\-%
% --- p. 14 ---
\setcounter{footnote}{0}% p. 14 carries one note, printed *)
\emph{terested, contemplative, without responsibility}, while the practical is
% the letterspacing of „u i n t e r = / e s s e r e t“ IS carried across the
% page break in the Danish; both halves are spaced on the image. The English
% divides the word at the same joint.
\emph{interested}, moved by \emph{fear and hope}, an object of
\emph{responsibility} --- then this opposition stands as the previous ones do.
Theoretical consciousness is \emph{disinterested} when ``interested'' is taken
in the sense of egoism; it is on the contrary interested in the highest degree
when interest is taken in the sense of a striving towards an essential end,
and this end for theoretical cognition is precisely the self's clarity about
itself, its being for itself as self-conscious according to its total
concretion. That spirit theoretically finds rest in any result whatever to
which the consequences of the presuppositions lead is correct only in so far
as spirit can halt momentarily at that result; but in the untrue or only
half-true result there is always a stimulus which again awakens cognition's
interest for a new investigation and again drives it forward, precisely
because the untrue result cannot come into accord with, or comprehend, that
totality of cognition which cognition according to its essence must strive to
work out\footnote{Theoretical cognition's interest in its result comes
strikingly forward in its supplementing of the incomplete premisses of the
inductive inference.}.
% printed mark: *) --- the only note on p. 14, below a short rule.

Theoretical cognition is to be \emph{contemplative}, absorbed in the grounding
of the objective; but it is so only in the sense limited by the earlier
determinations. From the contemplative, self-forgetting absorption in the
universal and objective it seeks back to the punctual and to that
self-cognition towards which an irresistible inner need drives it; and in
striving towards its goal there is, in cognition's \emph{struggle} for the
appropriation of actuality, in the \emph{true} cognition's struggle against
the untrue, something corresponding to the will's conflict with existence, to
its struggle for its realisation, to the good will's struggle against evil.

Theoretical cognition is finally to be \emph{without}
% --- p. 15 ---
\emph{personal responsibility}, the practical accompanied by the consciousness
of such. This opposition too is untenable in the form in which it is set up.
That with respect to the consciousness of responsibility there is a difference
between the \emph{act of will} and the activity of thought, when regard is had
only to this, that the former intervenes determiningly in reality while the
latter does not, at any rate not in a conscious manner, goes without saying.
But the activity of cognition does not therefore fall outside responsibility.
The act of cognition, which is subject to \emph{the norm of the idea of
truth}, has in that norm an ``ought'', in relation to which there is a
consciousness of responsibility. When responsibility in general points to
this, that what the individual is responsible for is the individual's own
deed, that it proceeds from the individual's being, and that it stands in
relation to a higher, more comprehensive law, then responsibility becomes
common to what proceeds from the activity of cognition and from that of the
will; and it must become so all the more, since the normal activity of
cognition is, as will appear in what follows, the condition of the normal
activity of the will. Cognition's connection with responsibility is therefore
not to be apprehended merely thus, that the will's relation to cognition can
bring cognition under responsibility, but thus, that cognition \emph{in and
for itself}, by its relation to the individual's being and to the higher law
of truth, comes under it.

Through the consideration of the marks of character by which an abstract
opposition between ``theoretical'' and ``practical cognition''
(consciousness), and thereby between cognition and will, was to be maintained,
we are thus led to the \emph{same re}sult as above: that the \emph{abstract
% letterspacing: „s a m m e  R e =“ closes the line spaced in the Danish,
% „sultat“ opens the next line tight. The English divides at the same joint.
opposition} yields its place to an \emph{essential unity} which
\emph{governs} the \emph{relative} difference, and \emph{must} do so, on
account of the relation of the primitive expressions of the being to the unity
of the being. This result will now receive a new confirmation when we consider
the \emph{opposition of facticity} between cognition and will, the opposition
into which they can enter towards one another in the course of the individual
development of the life of spirit.

A \emph{factual} opposition comes forward at the imperfect stages of
% --- p. 16 ---
\setcounter{footnote}{0}% p. 16 carries TWO notes, printed *) and **)
development of cognition and will: partly under the form of a one-sided
development of cognition or of will; partly under the form of an indifference
of the will towards, or an independence of it from, cognition; partly under
the form of the act of will's opposition to cognition.

The \emph{first} of these forms of divergence finds its
explanation\footnote{Cf.\ Problemet om Tro og Viden, pp.\ 136 f.} in the
sporadic element which comes forward, and must come forward, at the
subordinate stages of development, but which at the highest stage of
development --- the stage at which cognition and will are determined by reason
in their concrete content --- must be joined together in unity and harmony.
And unity is not merely the goal of the development: in the course of the
diverging development there is a pointer towards it in this, that the
one-sided development of cognition is defective not merely with respect to
will but also with respect to cognition itself, just as the one-sided
development of will is defective not merely with respect to cognition but with
respect to will as well --- so that the belonging-together and reciprocal
conditioning of what has been separated becomes plain. Cognition developed
one-sidedly, and not united with a development and firmness of will in the
personality, will always show itself to be an abstract and uncarried-through
cognition; a will developed one-sidedly, and not united with the richness and
clarity of cognition, will show itself to be a will not yet purified through,
and therefore not liberated from the instinctive, the fanatical and the
violent, and thereby from the egoistic.

In the \emph{second} form of the difference an effect of the sporadic
likewise\footnote{The same work, pp.\ 142 ff.} makes itself felt, as where one
acts from a mere prompting of drive or in blind passion; or else the
resolution has only an incomplete cognition for its presupposition, and one
acts from the immediacy of the concentrated personality. With respect to this
last case there has been shown, in the earlier work (p.\ 144), the form
% --- p. 17 ---
\setcounter{footnote}{0}% p. 17 carries TWO notes and RESTARTS at *)
under which, in such a resolution, a preceding cognition that has ``gone back
into immediate being in spirit'' is at work\footnote{A more thorough
psychological observation will show that the concept of a determinateness of
the will proceeding from a cognition that has gone back into immediate being
--- cognition taken in the widest sense of the concept --- finds application
also in the will's first expressions in the child, since the ``pronouncedness''
of these expressions follows step by step the pronouncedness of the life of
representation.}. The relation between the resolution and the incomplete
presupposition of cognition has its exact parallel within the theoretical, in
the inference which, with an incomplete empirical foundation --- an experience
not exhausted --- for its presupposition, states the universal not
hypothetically but categorically.

In the \emph{third} form of divergence there makes itself felt partly\footnote{Problemet
om Tro og Viden, p.\ 144.} a working of the ethically a priori in the will's
determining of itself in opposition to one-sided theory, and partly a
counter-working of the will against the true, through an affecting of
cognition and through a hindering of cognition's realisation. With respect to
the activity which the ethically a priori exercises in consciousness,
``Problemet om Tro og Viden'', p.\ 145, has pointed to the parallel with the
working which the a priori exercises within the theoretical, and it has been
established that the \emph{ethically} a priori, as an \emph{essence}-determination
and as an expression of the \emph{undivided} being for which \emph{knowledge}
too is an expression, must also be a \emph{knowledge}-determination and be
transparent to cognition, so that its
% letterspacing: „V æ s e n s =“ and „V i d e n s =“ close their lines spaced
% in the Danish; „bestemmelse“ opens the following line tight in both cases.
% The English compounds divide at the same joint.
working is a working of that which finds its place within knowledge and enters
into the system of knowledge's fundamental concepts. --- As regards the will's
conflict with what is cognised as true, this will show itself to take place
only in relation to an \emph{imperfect} cognition, and therefore to be
confined to those stages of development at which the power of the sporadic is
unconquered. A sharper psychological and ethical observation will confirm that
clear, \emph{concrete} cognition is never powerless against the will, that it
is impossible to act against
% The signature numeral „2“ at the foot of p. 17 is not transcribed.
% --- p. 18 ---
such an \emph{actual} cognising. Could the will withstand what was in this way
cognised as the \emph{being's} determination, as the being's demand, then it
would act merely in virtue of the essenceless, or the being would have to be
split into two heterogeneous parts; but the former is self-contradictory, and
such a split in the being, according to which the will would be something
irrational and cognition something rational, is an empty representation, which
becomes no less empty if, instead of a split, a ``defectiveness'' in the being
is spoken of by way of postulate. In such cases as seem to speak for the will's
unconditional power to withstand cognition, cognition will partly show itself
to be incomplete in such a way that in its abstraction it as it were gives
room for the conflicting action; and partly it will appear that cognition, as
not yet fully unfolded, is by an act of will --- more clearly or more obscurely
conscious --- which affects the direction of thought, brought outside the focus
of consciousness, moved off into a distance of time for consciousness, and
thereby obscured and hindered, so that it can be paralysed by a conflicting
illusory cognition, an excuse, a pretext; and only through such a spacing-out
does the conflicting action become possible. The clear cognition of what is
unethical in an action, of its opposition to a preceding better cognition,
sets in only \emph{after} the action. Here the \emph{time}-determination
acquires a decisive significance.
% letterspacing: „T i d s“ is spaced and „bestemmelsen“ runs on tight, within
% one and the same line, in the Danish. The English divides at the same joint.

While thus the factual divergence which at the imperfect stages of development
can take place between cognition and will does not show itself to be of such a
kind as to abolish the inner unity, this inner unity --- in which cognition
and will stand as expressions of the same being --- finds its impress in that
parallelism between the universal stages of development of cognition and of
will which comes forward in spite of the unlike development the two functions
of spirit may have in the individual. Just as cognition develops from its
first form, sensation, through the involuntary formation of universal
representations, to a conscious and progressive separating out of the
universal in the concept, so the will develops from the first, immediate,
single determinateness by drive,
% --- p. 19 ---
\setcounter{footnote}{0}%
through a determinateness by involuntarily formed \emph{practical universal
representations} --- produced by a spontaneously arising momentary and partial
reflection, such as make themselves felt in an abiding inclination --- to a
determinateness by consciously formed principles, maxims, out of which the
unity of a view of life may then proceed.

The harmony and unity to which there was a pointer even in the divergence
itself then comes forward unveiled at that stage of development of cognition
and will at which the human spirit has attained to clarity about itself
according to its concrete essential determination. In that the human spirit
here becomes conscious in cognition of the content of its being, it expresses
this content of cognition in the determinateness of the will: the will makes
that content of cognition, in which the being becomes manifest to
\emph{itself}, into \emph{its own} content, realises the being in its own
form, in the resolution; and it does so in virtue of the necessity lying in
the unity of the individually-universal, ideally-real being, in virtue of that
law of the being which lets the being manifest itself in its \emph{total}
formal determinateness, and therein manifest itself in harmony with
itself\footnote{Cf.\ Problemet om Tro og Viden, p.\ 142.}.
% printed mark: *) --- the only note on p. 19, below a short rule.
In the resolution the will is determined by the necessity of the being; but
inasmuch as this necessity is the necessity of its \emph{own} being, the will
is \emph{self}-determining through freedom.
% letterspacing: the line closes with the spaced „s e l v =“ and
% „bestemmende“ opens the next line tight. The English divides at the same
% joint.
In order that the being may realise itself in an act of will, in the proper
sense of that concept, it must first have become manifest to itself in
cognition; but the will's self-determining is maintained through that relation
to the essence-\emph{totality}.
% letterspacing: the spacing begins MID-WORD in the Danish, at the linking
% „s“ --- „Væsen“ is tight and „s t o t a l i t e t e n“ is spaced, all on one
% line. The English hyphenates the compound so the run can begin at the same
% point in the word.

\begin{center}\rule{2cm}{0.4pt}\end{center}
% A short centred rule stands here on p. 19, between the two paragraphs; it is
% a thought-break in the running text, not a head. A SINGLE hairline, like
% every rule in the book but the five doubles of pp. 4, 29, 63, 201 and 210.

If, after having elucidated the universal relation between the theoretical and
the practical, we seek the closer determination of \emph{cognition's true
form} --- of the form which is the goal of its development --- then it is to
be found from the fundamental relation of consciousness. In the relation given
in the human spirit between self-consciousness and consciousness there lies
% --- p. 20 ---
the necessity of a striving to bring all content of consciousness under the
unity of self-consciousness, to abolish the foreignness of the object of
consciousness and to make it transparent for spirit in a comprehending
cognising, in which the self has its ideal unfolding and its self-revelation,
in that it ideally appropriates actuality. Only in such a comprehending
cognising, in which the content of consciousness forms for spirit a unity
answering to the unity of self-consciousness, is that contradiction abolished
which lies in the relation of consciousness: that the object is in my
consciousness as a determination of \emph{my} consciousness, and yet
according to its content is something foreign to me; and only in such
cognising is that contradiction likewise abolished, that spirit --- which in
self-consciousness relates itself to itself in an infinite manner and is in
unity with itself --- is at the same time foreign to itself through the
foreign content which the self comprehends.
% letterspacing: only „m i n“ is spaced in the Danish (the first „min“ four
% words earlier is tight, on the same line); „Bevidstheds Bestemmelse“ opens
% the next line tight.
But when such a leading back of the content of consciousness to transparent
unity with self-consciousness is cognition's demand, determined by the
fundamental psychological relation, then such a leading back and such a unity
become possible only on condition of a cognition of the principle which is the
principle of unity of the self and of existence; and cognition's true form
then becomes a principle-determined unity of cognition, in which the self, in
cognising itself as an organic member ordered into the totality governed by
the principle, can cognise itself according to its concretion. Such a form of
cognition, in its essence like that which philosophy according to its concept
is to give, must --- in virtue of its relation to the original demand of
consciousness, in virtue of its universal principle, in virtue of the kind of
thought it presupposes, in virtue of the consciousness it encloses of the
essence and fundamental relation of thought, in virtue of its combining of the
a priori and the empirical --- claim for itself the place of cognition's
τέλος.
% Greek, antiqua italic with accent in the printing, read off the image;
% neither OCR witness reads it. Copied verbatim, per the cluster convention.

If we seek \emph{action's true} form, then this must be determined from the
will's fundamental relation to the essence of the self-conscious spirit. Since
man is a self-conscious being and as self-conscious is a willing being, that
working which in the animal, as unconscious and unfree fulfilment of drive, is
only a
% --- p. 21 ---
deed acquires the character of \emph{action}, i.e. of a \emph{conscious}
working proceeding from \emph{will}. More closely, the nature of action is to be
determined thus: that it is not merely accompanied by the consciousness
\emph{that} one is acting, but is determined by the consciousness of
\emph{why} one is acting. Action as action, in the proper sense of the word,
is directed towards a goal, an \emph{end}, is guided by this, is a conscious
working for the attainment of this end. Action's goal, the consciously willed
end, may, seen from the standpoint of the objective, show itself as an
immediate or as a mediate goal: as that which stands directly in relation to
the will and is willed for its own sake, or as that which stands in relation
to the will only through something else and is willed for the sake of
something else, as a means to the attainment of the end. But when, with
respect to the objective, we speak of a willing of it for its own sake, this
is nevertheless to be understood only in comparison with the objective means,
not so that the subject's relation to itself is excluded, or the agent's
subjective interest in the objective forgotten. The subject of action relates
itself, as subject, at every point to itself; in determining the objective as
an end, it determines it as \emph{its} end, and relates itself \emph{to
itself} in the realisation of the end, wills \emph{itself} in willing the
objective. Through the relation to the objective we are therefore led back to
the relation to the subject; and that we are led back to it lies in the very
essence of subjectivity, in its fundamental determination: to be for itself in
its other. The end of action becomes then everywhere the subject itself, and
this holds also of the action through which the subject gives itself up to the
higher; for through self-surrender the higher self-affirmation is won.

But when the determination given holds of \emph{all} action, then within it
falls the opposition between \emph{true} and \emph{untrue} action.
% letterspacing: „a l“ closes its line spaced in the Danish, „Handlen“ opens
% the next tight.
Action's truth is determined by the truth of the will (of the will's
determinateness); but the will's truth is determined by whether the willing
subject, in willing \emph{itself}, wills itself \emph{according to its truth},
% --- p. 22 ---
i.e. according to its \emph{essential determination}, and more closely,
according to its \emph{full} essential determination. \emph{True action} is
therefore that in which the will makes the realisation of the being --- of
which the will is a primitive expression --- into its end, in which the will
therefore works as \emph{essence-willing}.

In opposition to this true form of action, which we anticipatively designate
as \emph{ethical} action, the \emph{untrue} form of action becomes --- setting
aside for the present the unethical action, whose essence will be determined
later --- the \emph{not yet} ethical, imperfect action, which proceeds from
the subject's willing of itself according to its particular determinateness of
immediacy, according to its finitude and contingency, according to its
phenomenal existence. Such a willing becomes a willing by the subject of
itself according to its ungoverned side of sensibility --- sensibility being
taken here in the widest compass of the concept --- as an unconsciously
egoistic willing of enjoyment. In this form of action, and in the will from
which it proceeds, there lies an inner contradiction which reveals its untruth.
When the self, as sensibly determined, wills enjoyment, it wills \emph{itself}
through enjoyment, in that in enjoyment it wills its self-affirmation as
sensible. But in willing enjoyment the self wills that which is
\emph{independent} of the self's willing and which according to its nature is
\emph{changeable} --- that, therefore, which \emph{cannot} be governed by the
will and which in its incommensurability with the self's universal essence
\emph{cannot} be held fast by the will. And in willing itself in the ununified
manifoldness and essenceless changeableness of enjoyment, the self wills
itself in a determinateness --- the determinateness of sensible externality
--- in which its unity is abolished; it wills itself, therefore, in that in
which, while it seeks a self-affirmation, it \emph{loses itself}; it wills
itself as not-self. The inner contradiction is abolished only when the will
--- which as an expression of the being must according to its concept be
essence-willing --- becomes a willing of that which in truth \emph{can} be
willed, because it is one with the essence of the willing subject: a willing
of that in which the self affirms itself according to its unity and abiding
essentiality, according to the infinity of its essence. To such a willing of
itself
% p. 22 carries NO footnote --- the page was on the brief's list of
% note-bearing pages, but the image shows an unbroken type block.
% --- p. 23 ---
there is, at the imperfect stages of development of the will, a pointer, as to
the goal at which the will finds rest, in that the will realises what was
obscurely striven for at those lower stages of development. The goal of will
and action, which is determined by the self and in which the self relates
itself to itself, must from the beginning be covertly determined from
\emph{the self's unity} --- a determining which betrays itself through the
working-forth of the universal --- and as the self becomes clear to itself
according to its essence, it must \emph{consciously} and \emph{freely} be
determined from \emph{the self's content-rich unity}, so that the self
according to its \emph{total} essential determinateness, according to its true
universality and infinity, becomes the contradiction-free τέλος of will and
action. But such a willing by the self of itself according to its truth,
according to the infinity of its essence, is \emph{ethical} willing.
% Greek: a second τέλος, at the head of a line on p. 23, likewise read off the
% image; ABBYY drops the word entirely.

For this ethical, essential willing a twofold determination has been indicated
in what immediately precedes, and the two stand in the closest inner
connection. The ethical is constituted in that the self wills itself
\emph{according to its essence}, and in that --- this being the condition of
the former --- it draws itself out of its immediate being with that being's
unfreedom, finitude and changeableness, and by \emph{freedom} determines
itself \emph{as spirit} according to its \emph{inner infinity}. At the stage
which lies before the ethical, the self willed that which as finite is
something vanishing and is independent of the will; it willed \emph{itself},
therefore, in what is something vanishing and something foreign to the self,
that is to say, it lost itself. As it becomes conscious of the transitoriness
of what is willed and of the contradiction lying therein, the opposition comes
forward between the transitory finite and the infinite and unconditioned;
there shows itself a \emph{choice}; and as the will, in despair of the finite,
frees itself from the finite and chooses the infinite, the ethical breaks
through. But the infinite is chosen precisely in that the self chooses itself
according to its essence; for that means: according to the
\emph{determination of eternity} within it, according to its \emph{eternal
validity}. In this \emph{resolution of choice}, which is a \emph{self-acquiring
of freedom}, the oppositions in the choice receive a new elucidation and a new
% --- p. 24 ---
\setcounter{footnote}{0}%
determination: the determination of the \emph{good} and the \emph{evil} ---
concepts which first belong in the ethical domain; and in the ethical choice
the good is chosen.

In choosing itself according to its essence and through the choice willing its
essence, the self has therefore made a movement of \emph{infinity}: its
\emph{immediate} being according to its determinateness by \emph{nature} is
converted into a being of \emph{spirit} through \emph{itself}, through its
\emph{freedom}.
% letterspacing: „N a t u r“ is spaced and „bestemthed“ runs on tight in the
% same line; „ved s i g / s e l v“ carries the spacing across the line break.
In that the self, through that infinitising choice, chooses and wills itself,
it wills itself \emph{absolutely}, i.e. \emph{as an absolute}; and \emph{only
itself}, i.e. \emph{its own essence}, can the self will absolutely; but this
willing is possible only in that it wills the \emph{Absolute} in the self.
\emph{The self} which chooses itself by the absolute willing of the ethical
resolution has a \emph{twofold} significance: it is the \emph{ideal, universal
self}, the self's absolute essence, that is chosen; but the choosing self,
which chooses that essence as \emph{its own}, is the \emph{individual self},
which in choosing that other self chooses itself in its individuality. This is
to say: it \emph{takes over} its own development, lying prior to the ethical
choice; but in taking it over in the choice of the ideal self, it takes it
over as that development shows itself in the ideal self's light. In other
words: it passes the ethical judgement of \emph{repentance} upon that
development; it \emph{transforms} itself in principle in the ethical
resolution, in virtue of its fundamental essence. In this repentant choosing
the self gathers itself together in its whole finite concretion; and while in
the choice it as it were takes itself out of finitude, it is in the most
absolute continuity with finitude and has itself as its task in its ethically
chosen concrete actuality --- an actuality which is to be worked through by,
and penetrated with, the universal, so that the actual self is determined by
the ideal self which it has within it, and expresses at one and the same time
the individual and the universally human\footnote{Cf.\ the developments in
\emph{Kierkegaard}'s Enten --- Eller, II.}.
% printed mark: *) --- the only note on p. 24, below a short rule. The name
% „K i e r k e g a a r d s“ is letterspaced Fraktur, not antiqua; the run
% falls on the name here too, the genitive standing outside it.

For \emph{ethical willing}, whose essence and task are thus determined, there
now shows itself --- both in its first constituting of itself and in its
unfolding in action's con\-%
% The Danish prints „deres“ (for „dens“) here; transcribed as printed there,
% and having no English counterpart, recorded only.
% --- p. 25 ---
crete determinateness --- \emph{true cognition} as an essential and necessary
\emph{condition}.

For the immediate working of the ethically a priori, by which such a
presupposition of cognition might seem to be replaced, shows itself
insufficient. That the a priori has an activity which reveals itself in
conscience is incontrovertible; but in order that this activity may be raised
above a merely enveloped (implicit) presence in the form of the merely
instinctive and drive-like, it presupposes a development of cognition, just as
the logically a priori presupposes such a development in order to be freed
from what hinders and obscures it. In that first working the ethically a
priori works immediately together with promptings of drive, works itself under
the form of a prompting of drive --- for example under the form of sympathy,
of natural love --- but is not clarified into ethical freedom. Therefore the
ethically a priori in its immediate working is insufficient as a regulative
for action. In the immediate prompting, which is not yet clarified through by
cognition, it remains constantly uncertain how far what prompts is something
in truth a priori, and as such a true essential determination, or whether it
is not something untrue, an expression of something unwarranted, something
phenomenal and contingent. And since the a priori \emph{as} a priori has the
stamp of the abstract, it becomes, when it is not unfolded into concretion by
cognition, insufficient for the determination of the concrete action. When
action is called for, the subject will come to seek the determinateness of the
action outside the abstract norm, in its own unclarified, nature-determined
immediacy, in which it is not conscious of itself according to its essence;
and there will then arise a dividedness and doubleness in the will, since the
universal and the particular, the norm and the concrete content of the action,
cannot be willed as one. This contradictory doubleness can vanish only when
the essential demand unfolds itself in concrete determinateness for cognition.

The diversity of what at different times, at the different stages of human
development, has been held
% --- p. 26 ---
to be morally right confirms the insufficiency of the immediate prompting by
the a priori. One need take as an example only the manner in which the feeling
of reverence for the divine, or the immediate feeling of piety towards
parents, has found its expression.

In virtue of the ethically a priori there may be an immediate consciousness of
an ethical obligation, as an obligation towards something higher; but the
clear consciousness of \emph{that which obliges} is conditioned by true
cognition, and likewise the \emph{assured} consciousness of the concrete
content of the obligation is conditioned by such cognition --- by the concrete
cognition of that which obliges and of the \emph{essential determinateness}
and \emph{essential content} of the one obliged.
% Printer's defect here in the Danish: the page prints „I n h o l d“
% (letterspaced) for „Indhold“ --- the „d“ simply absent, verified on the image
% and read so by both OCR witnesses. Orthographic, and so not reproducible in
% English; recorded here.
Instinctively one may live in accordance with the moral, in that the
objectively moral --- custom and tradition --- immediately determines the
acting individual, who in a relation of piety subordinates itself to this
objective as to a higher power in which it feels something akin to itself, and
whose law it therefore follows, not as a foreign law. Instinctively the
immediate consciousness may grasp the good, may act in the spirit of morality,
in that it, unconsciously affected by the surrounding moral element of life,
acknowledges something universal and ideal as the norm of action, gives up the
egoistic, and in the particular action acts with the will of resignation,
guided by its determinateness in principle. But custom and tradition, the
surrounding moral atmosphere of life --- whether they be the resultant of the
immediate working of the a priori or of a preceding cognition which has as it
were been deposited in immediate being --- can be the secure norm only when
they can again be freed, in living cognition, from the immediate form in which
something contingent and foreign may cling to them. The clear, free, assuring
consciousness of the moral and of the inner connection within it can be
attained only through such a cognition as separates out the contingency of
immediacy, and more determinately through that kind of cognition by which what
is cognised can acquire inner unity and necessity --- there\-%
% --- p. 27 ---
fore through the principle-determined cognition whose essence has been
developed above. Such a cognition has, in its incompleteness --- inseparable
from the conditions of human development --- and in its relative abstraction,
the essential truth and the possibility of as it were growing out from its
centre and filling out its forms without their being displaced. With such
cognition for its presupposition the ethical will acquires its true form,
precisely because the being is in that cognition clear to itself; and in its
liberation from finite entanglement, from the obscurely drive-like which in
itself is the changeable and which, even in its direction towards the good,
knows no limit, the ethical will acquires that true resoluteness, firmness and
purity which is conditioned by the being-for-cognition of what is willed. The
subject's concrete cognition of its essence and of the true demand of that
essence acquires here its decisive practical significance.

From what has been developed it becomes clear that cognition's significance
for the ethical already comes forward where the ideal self is to be made
\emph{with consciousness} into the object of the will, where something
\emph{absolute} and \emph{unconditioned} announces itself and is to be
separated out from what veils and obscures it. But it comes forward still more
clearly when the idea of the good is to be realised in the \emph{concretion}
of the moral will, to be worked out in the actuality of a wholeness of life
--- when, what becomes identical with this, the self's willing of itself
according to its ideal essence is to attain the concrete form of actuality of
the ethical action, the realisation of the human being through free human
action, action in truth agreeing with itself and with its own essence. For
only through the true and concrete cognition of the self's essence and of its
fundamental relation does free action become in its determinateness action
\emph{determined by reason}, does the ethical unfold its rational content
\emph{clearly and unmixed}. Since the \emph{ethically a priori} in its
abstraction meets with something concretely given, with something
\emph{empirical}, then in this meeting --- which is to lead to a determining
of the latter by the former, to an unfolding of the abstract ethical
consciousness of essence into an essential consciousness of actuality --- only
the true form of cognition, which acknowledges the
% --- p. 28 ---
\setcounter{footnote}{0}%
a priori and realises its essential unity with the empirical --- a cognition,
therefore, of the kind determined above --- can be guiding in the last
instance. Only through it can the self's willing of its essence become a
willing of the essence according to its complete concrete determinateness;
only through it can the self in its acting truly ``will one
thing''\footnote{Cf.\ Section III (on the realisation of the ideal).}. ---
% printed mark: *) --- the only note on p. 28, below a short rule. „III“ is
% set in antiqua, as all figures in this book are; not italicised.

An expression of this significance of cognition for action is the fact that
\emph{ethics}, the scientific cognition of the ethical, must acquire its
\emph{essential} significance for that ethical. If \emph{scientific} cognition
is cognition's clearest and most complete unfolding, then what has been said
of true cognition must find its immediate application to it.

Since ethical action is determined as action's true form on the ground that
the will from which it proceeds is an \emph{essence}-willing, a willing of the
self according to its essential determinateness, ethical action acquires its
\emph{norm} in the \emph{human being}, its \emph{law} in the being's own law
of activity, and its task becomes to express the total actual personality.
% letterspacing: „V æ s e n s“ closes spaced and „villen“ runs on tight, in
% the same line. The English compound divides at the same joint.
Herewith there is once more a pointing back to \emph{the opposition to the
positively religious}. This opposition shows itself here partly as an
opposition between the being's indwelling law and the supranatural,
arbitrarily determined law of action, partly as an opposition between an
ethical life which comprehends the total, real existence and a life which, as
obeying the supranatural demand, flees away from the natural, relates itself
negatively to the natural through asceticism, and expresses the spirit-side of
the being in an abstractly spiritualistic way.

\begin{center}\rule{2cm}{0.4pt}\end{center}
% A short centred rule stands here on p. 28, between „spiritualistisk.“ and
% the closing summary: a thought-break, not a head rule. A SINGLE hairline.

The result of the developments given in the foregoing we sum up in the
following propositions:

\begin{enumerate}[label=\arabic*., leftmargin=*, labelsep=1em,
                  topsep=0.35\baselineskip, itemsep=0.35\baselineskip]
\item True cognition, which is chosen in the choice between \emph{knowledge}
  and \emph{faith}, becomes the condition of the clearly con\-%
% --- p. 29 ---
  scious willing of the ethical, and of the ethical's concrete unfolding into
  a wholeness which comprehends the life of the actual personality.
\item The personal life can be clarified through by true principle-determined
  cognition, and this cognition can, as it advances in concretion, win such a
  punctuality that the concrete ethical action can find in it its
  \textit{ratio sufficiens}.
% antiqua Latin on p. 29, upright in the print; \textit{} per the house
% convention taken from the 1868 volume, and TIGHT antiqua, so no \emph{}.
\item The true unity of theory and practice, of the ``theoretical'' and the
  ``practical'' consciousness, is present in this relation, in which the
  subject's consciousness of its action and of its ``grounds of will'' is in
  harmony with its theoretical cognition of its essence and of the essence's
  demand; --- and it is present \emph{only} in this.
\item In virtue of this relation between cognition and ethical action, since
  no splitting into incommensurable forms takes place within cognition, all
  that is truly ethical must be a rational ethical, just as all true ethics is
  a rational ethics. An antirational ethical is a self-contradiction, an
  antirational ethics an absurdity.
\end{enumerate}

\begin{center}\rule{2cm}{0.4pt}\end{center}
% p. 29 is the last page of division I and is a short page: the type block
% ends after item 4 and a short centred rule closes the division. The rule is
% printed as a DOUBLE rule (two thin lines); set here with the same single
% \rule used for the other short rules in this book.

\chapter*{II.\\ On the Presence of the Essential Marks of the Religious in its True Sense in True Cognition and True Action, and on their Relation of Unity to the Former.}
\addcontentsline{toc}{chapter}{II. On the Presence of the Essential Marks}
\markboth{II. On the Presence}{II. On the Presence}

% --- p. 30 ---
\setcounter{footnote}{0}%
% p. 30 carries the division head II (already in the skeleton), then a short
% centred rule (not transcribed, per the book's convention for rules under
% display heads), then the bold centred sub-head below. The text opens with an
% enlarged initial capital, set normally.
\section*{1. The Essence of the Religious.}

The expression for the religious relation which was used in the ``preliminary
determinations of concepts'' in ``Problemet om Tro og Viden'' to designate that
relation according to its ideal determinateness becomes here our starting
point, and through the analysis of it we shall be able to reach the more
concrete unfolding of the concept of the religious.

We there determined the religious relation as ``\emph{the absolute
relation}\footnote{If it should seem to ``inverted'' dialecticians
contradictory to speak of an \emph{absolute relation}, because relation
expresses \emph{relativity}, we need presumably only recall the Absolute's
relation to itself, and spirit's infinity-relation to itself in the concretion
of self-consciousness.} \emph{to the Absolute}'' --- by an expression,
therefore, which not merely in form approaches S.\ Kierkegaard's usual
designation of the stage of a religiousness not yet Christian, which he
characterises as ``the absolute direction towards the absolute τέλος'', but
which has already been used by him where he sets the religious absolute
relation beside the relative relations, which with him comprehend the ethical
relations as well, and in which that absolute relation is not to exhaust
itself or become concrete. While we thus appropriate his
% printed mark: *) --- the only note on p. 30, below a short rule.
% Greek: τέλος, antiqua italic with acute in the print, read off the image; the
% ABBYY layer carries no Greek at all. The second quotation, „den absolute
% Retning mod det absolute τέλος“, is NOT letterspaced, though the first is;
% the English follows.
% --- p. 31 ---
designation of the concept, it will nevertheless appear through the
developments that follow how in our conception this designation acquires
another significance than it had with Kierkegaard --- already where he speaks
of the pre-Christian religious stage, and still more determinately where he
speaks of the paradoxically religious.

If we analyse the determination of the religious relation as the absolute
relation to the Absolute, then the analysis leads us to the following
consequences, through which the concept acquires its more concrete
determinateness.

a) \emph{With respect to the Absolute}, which is the object of the relation,
the fundamental determination encloses the following:
% the enumerators a)--h) on pp. 31--41, and the Greek α) β) on p. 31, are
% antiqua italic in the print; set plain here, as in the Contents.

α) The ``Absolute'' which is the object of the religious relation can be an
absolute \emph{for us} only when it is \emph{our} Absolute, i.e. when it stands
in a relation of \emph{essential unity} to us. In an opposition of essence
there would lie immediately a determination irreconcilable with absoluteness.
Against this unity of essence the ethical ``absolute distinction'' between God
and guilty man, so strongly accentuated by Kierkegaard, does not tell; for the
determination of guilt lies, as will be more closely developed under III,
within the relation of immanence, and therefore makes no definitive breach in
the unity of essence.
% letterspacing: „V æ s e n s e n h e d s =“ closes the line spaced and
% „Forhold til os.“ opens the next tight; the capital F shows this is a real
% hyphenated compound, not a broken word. The English run falls on the same
% element.

In the positive religion's conception of the relation between God and man
there is a doubleness present. On the one side likeness is accentuated, out of
the immediate religious need: man is posited as created in God's image; but on
the other side the reflection determined by the form of representation leads to
accentuating the opposition between the transcendent God and man, and this
opposition fixes itself as an unabolishable one (cf.\ Section III, on
reconciliation).

β) The object of the religious relation can be the Absolute for us only when it
is \emph{the in truth infinite}. Only the in truth infinite can be absolutely
willed, absolutely sought through the need of the \emph{whole} human being; and
only what we will absolutely, absolutely strive towards, is for
% --- p. 32 ---
us an absolute. A willing of something finitely determinate is a finite
willing, and the relation to it a finite, relative relation. This
determination --- that the Absolute must be the in truth infinite --- is
inseparable from the other, that it must stand in a relation of essential
unity to us. For it can be the in truth infinite only when it is not separated
from us by the barrier of a represented, and therefore finitely determined,
transcendence.

Only these more formal determinations of the Absolute are necessary in the
present connection; the more concrete determinations of the concept will be
given under III, Chapter 1.

b) \emph{With respect to the subject of the relation: man}, it follows from
the determination: the absolute relation, that the subject must be present
% QUOTES, p. 32: the printing sets „af Bestemmelsen: det absolute Forhold, at
% Subjectet maa være tilstede“ with NO quotation marks --- the colon alone marks
% the quoted determination off. An earlier draft supplied a pair; removed, since
% quotation marks are under the parity rule and none stands here.
in the relation according to its essence's \emph{inner infinity}, that is to
say, according to its \emph{whole} essential determinateness, according to the
totality of its \emph{universal} essence. The religious relation cannot,
therefore, be a relation which relates itself to a single side of the human
spirit or to a single fundamental function, even if this were thought to
enclose the others within itself as moments --- whereby the others would enter
the relation precisely only \emph{in the determinateness} they could have as
moments; but it must be a relation into which the fundamental functions of the
human spirit, as totalised within themselves, enter equally. Just as it
therefore becomes a one-sidedness in the \emph{Hegelian} philosophy of
religion that it apprehends the religious relation exclusively as a relation
of \emph{cognition}, so it becomes just as fully a one-sidedness to determine
it as a mere relation of \emph{will}, or to designate it as what has been
called, with a somewhat unclear expression, a relation of \emph{life}, in
\emph{opposition} to a relation of cognition. The religious relation is a
\emph{life-relation in the sen}se of a relation which comprehends \emph{the
whole of human life}; but human life, the life of spirit, is the life of
cognition just as much as of will and of feeling; it is therefore apprehended
in a one-sided and imperfect manner when it is apprehended only as opposed to
cognition. But since the relation comprehends the totality of human life, the
relation of will acquires on the other side a just as decisive and primitive
significance, and it
% letterspacing on p. 32 breaks in the middle of compounds and runs past the
% word: „E r k j e n = d e l s e s“ spaced with „forhold“ tight; the same for
% „V i l l i e s forhold“ and „L i v s forhold“; and „L i v s f o r h o l d  i
% B e =“ is spaced while „tydning“ on the next line is not. The English
% reproduces the last of these by breaking "sense" at the same joint.
% --- p. 33 ---
\setcounter{footnote}{0}%
can step into the foreground where reconciliation is \emph{sought}. The will
in the relation to the Absolute is, when the relation is in truth, the
religiously determined, \emph{ethical} will, which actively transforms
existence through the relation. Through this will the relation becomes a
relation of \emph{conscience}. That the relation is a \emph{relation of
personality} points back equally to cognition and to will\footnote{Cf.\ the
preceding Section.}: what is a fundamental expression of the being is at the
same time a fundamental expression of the personality; otherwise the
determination of personality would stand in a merely contingent relation to
the being, which is immediately self-contradictory.
% printed mark: *) --- the only note on p. 33, below a short rule.

That towards which spirit infinitely strives through the absolute relation so
apprehended is \emph{reconciliation}, and the religious life becomes a life
for spirit's reconciliation \emph{within itself and with its ground of
essence}. This is as yet more a \emph{metaphysical} determination of the
religious relation, and it is reached already through the concept of spirit as
that which according to its essence is infinite, and which, in order to find
rest in itself, strives to realise this infinity within itself, but first
realises it through the absolute relation to the in truth infinite, to its own
and to existence's ground of essence. The specifically ethical-religious
determination of the relation --- the determination of reconciliation as
reconciliation of \emph{guilt} --- will be developed only later.

Since the equal entry of the \emph{whole} human spirit into the religious
relation is thus accentuated, the divergent conceptions are to be touched upon
which emphasise the one-sidedly limited determination of a \emph{life-relation}.
We may here take account partly of \emph{Grundtvig}, partly of
\emph{Martensen}.

With the \emph{former} an emphasis is laid on the relation of \emph{life} in
the sense of a relation of \emph{immediacy}, which is so apprehended that the
immediacy acquires the aesthetic character of sensible presence and comes to
enclose within itself the moment of the magical. In the conception of the
working of the ``living Word'', of the power of the sacrament,
% signature numeral „3“ at the foot of p. 33 (first page of the gathering);
% not transcribed.
% --- p. 34 ---
spirit's relation of freedom to spirit here gives room for a magical working;
and while the religious consciousness, which by a postulate cuts reflection
off at a determinate point, at the symbol, seeks rest in an unshakable,
immovable objectivity, it exchanges the in truth religious relation for an
aesthetic, poetic relation, in which the personality has not itself made the
movement of infinity through true inwardisation, in which it does not know
infinite resignation, but relates itself to the divine without knowing the
seriousness of the relation, without remembering that the possibility of doubt
becomes at every moment that which must be overcome if the relation is to be
in truth, and that the anxiety of uncertainty becomes the ground upon which
the victorious certainty of the relation of faith is to rest. Since, with a
half-conscious shyness of thought, there is unclear talk of a relation of life
into which the life of thought is not taken up, it is overlooked at this
standpoint that a relation through cognition, through thought --- precisely
because cognition is the being's expression, because it is the form of
realisation for the relation of essence --- becomes in the true sense a
relation of life, and that it leads to \emph{that} immediate unity which can
be thought in the domain of the life of spirit.

With \emph{Martensen} there comes forward an endeavour to bring the
\emph{whole} of human life into the religious relation: \emph{faith} is
determined as \emph{a unity of religious feeling, cognition and will}. But in
virtue of the subordinate and governed position which he lets cognition occupy
in relation to faith, and in virtue of the conception he asserts of the
\emph{will} as ``the concluding moment in religious consciousness'', this
endeavour is not realised, and the relation of cognition is pressed back.
Since the religious God-relation is separated from the speculative and the
aesthetic, which are determined as relations ``at second hand'', while the
religious becomes an \emph{existential relation}, a ``consciousness of the
God-relation which is one with the personal living and being in this
relation'', then in this form of separation what is essential in the relation
of cognition is pushed aside; the one-sided distinction between the impersonal
and the personal is
% --- p. 35 ---
asserted also for that true, concrete cognition in which the self seeks a
satisfaction for its deepest need, grounded in the essence of its personality.
And this one-sided pressing back of the moment of cognition betrays itself
also in the closer developments, in which the religious relation is
characterised as a \emph{holy} relation, a relation of \emph{conscience}; for
through these determinations it is to be expressed that the relation of
religion is a relation to the ``\emph{personal}'' God, in \emph{that}
conception of the concept which raises it above the ``logical'' and makes the
\emph{will} the organ of the God-relation.

c) With respect to the \emph{relative} relations it follows from the
determination of the religious relation as the \emph{absolute} one that this
relation cannot fall outside them, but must comprehend them, penetrate them
and make them absolute through their unity with it. In this unity --- a unity
conditioned by the ethical becoming something all-comprehending --- the
relative relations must acquire a higher significance, so that they can become
adequate expressions of the religious, which has its concrete unfolding in
them.
% unlike a) on p. 31 and b) on p. 32, only „r e l a t i v e“ is letterspaced
% here; „Med Hensyn til de“ is set tight. The English follows.

This determination receives its sharper elucidation by a comparison with S.\
\emph{Kierkegaard}'s conception of the relation between the absolute religious
τέλος and the relative ones. As he seeks the expression for that form of
religiousness which is not yet that of the paradoxically religious, the
\emph{first} expression of religious existence becomes for him \emph{the
absolute direction (respect) towards the absolute} τέλος, \emph{expressed in
action in the transformation of existence}. In the relation to the absolute
τέλος \emph{all} the moments of finitude are to be once and for all reduced,
in action, to what must be given up in relation to eternal blessedness; only
on condition of this exceptionless surrender is the existing individual in
actuality to relate himself to the absolute τέλος, to will it absolutely. As
the absolute τέλος for the willing individual who wills to strive absolutely,
it is to be so abstract that, aesthetically regarded, it is the poorest
representation; the willing individual is not even to will to know anything of
this τέλος beyond that it exists; for
% --- p. 36 ---
\setcounter{footnote}{0}%
as soon as he comes to know something of it, he is already to begin to be
delayed in his course. The absolute τέλος is then not to be possessed at rest
in existence; the whole time, the whole of existence, becomes the time of
expectation. Mediation is to have no place here, because the existing
individual is prevented from gaining such a footing outside existence that
from it he could mediate what moreover, by being in becoming, withdraws itself
from completedness. In relation to the absolute τέλος, that mediation takes
place is to mean precisely that the absolute τέλος is reduced to a relative
τέλος. It is not to be true that the absolute τέλος becomes concrete in the
relative; for resignation's absolute distinction is at every moment to secure
it against all fraternising. The absolute relation is to be able to demand the
renunciation of \emph{all} the relative ones. But for that very reason the one
who relates himself to the absolute τέλος is nevertheless to be very well able
to be in the relative ones; it will be precisely in renouncing that the
absolute relation is practised: the relation which rests on the absolute
distinction, the relation to that which is won only by absolutely venturing
all\footnote{\emph{Kierkegaard}: Afsluttende Efterskrift, pp.\ 295 ff.}.

In the first determination of the religious relation thus more closely
developed, there are already indicated in Kierkegaard the consequences which
come forward decisively where at last, within Christianity, the demand of the
religious is raised so high that it consumingly burns up all finitude and with
it all those ethical relations which, seen from this height, come under the
determination of finitude. The religious relation here acquires the utmost
abstraction: it can become the power which with the force of infinity draws
the existing individual out of finitude and constantly gives him the direction
towards the eternal; but it cannot win shape in the existing individual,
cannot form his concrete actuality; for, as I have shown in another
place\footnote{Problemet om Tro og Viden, p.\ 221.}, it becomes only a
\emph{demand} that the individual should be able at one and the same time to
relate himself absolutely to the absolute τέλος and relatively to the relative
ones,
% printed marks on p. 36: *) and **), below a short rule. The name in the
% first note is letterspaced Fraktur: „K i e r k e g a a r d“.
% --- p. 37 ---
in which the former is not to become concrete. But by this abstraction, which
unavoidably attaches itself to the religious relation, that relation acquires
precisely the character of \emph{relativity} and loses the possibility of
expressing the highest energy of the undivided human being; and in the finite
which it is to be able to demand be given up there is also that which, as
something essential, is something that cannot be given up: the ethical
relations as a totality. A corrective against that conception lies already in
this, that with Kierkegaard the fundamental ethical relation comes into being
only as a relation of will to \emph{the eternal}, and that the determination
of the eternal in the domain of the ethical is not to be distinguished from
the determination of it at that stage of the religious which we have
characterised. If this is held fast, then the ethical cannot be apprehended as
a stage of transition in the sense that there could be an ethical life which
was not yet religious; rather the ethical becomes ethical only by being from
the first religious. And precisely for that reason the religious stands in
such a relation to the ethical that it must always find its expression in the
ethical --- so that there cannot be thought a surrender, warranted by the
religious, of the totality of the ethical, of all the ethical relations, but
the renunciation of an ethical relation must always be in virtue of a higher
and more comprehensive ethical, and religious renunciation must therefore
always be borne by ethical action.

d) The first, essential determination of the religious relation was therefore
that it was an \emph{absolute, infinite} relation. And in the concept of the
human spirit, in its essential determination, lies the \emph{necessity} of
such a relation for the human spirit. Infinity is its essence, its true
element of life; only in a relation of infinity can it therefore come to rest,
and so that it stands in an absolute relation --- i.e. in a relation of
infinity --- not merely according to a single one of its fundamental functions
but according to their totality. Even when the content of spirit is converted
into the form of \emph{feeling}, the relation's absoluteness can be preserved,
since feeling, in totalised unity with cognition and will, becomes that
intensive form of immediacy of spirit whose content has the infinite richness
of unfolded possibility, and which at every
% --- p. 38 ---
moment can become mobile and be clarified through in cognition and will.

e) That first determination of the religious relation, as the absolute
relation to the Absolute, now acquires a \emph{more concrete} form in that the
limitation of \emph{existence}, the determinateness by \emph{finitude}, of the
human being that enters the religious relation is accentuated. At this point,
however, only the general determinations to which this standpoint leads are to
be indicated; the closer consequences --- namely with respect to the
significance of the temporal relation for the ethical-religious life, for
guilt and its abolition --- will be shown in the last Section.

The new determinations that come in, and that express that the absolute
relation to the Absolute is the absolute relation to the Absolute of a human
spirit which is \emph{according to its essence infinite} but which
\emph{exists in determinateness by finitude}, are the determinations of
\emph{faith} and of \emph{self-surrender}. Of these, when they are
distinguished, the first points rather to cognition, the second to will --- but
so that faith's cognition must according to its truth be determined as being
in inner, harmonious unity with the will, and self-surrender's will as
standing in the same relation of unity to cognition.

f) The determination of \emph{faith} comes forward in that the finite
consciousness is, out of its determinateness by finitude, to grasp the
infinite and eternal in the form of certainty and to hold it fast as that
which penetrates and determines finitude. Faith is a postulate of cognition,
an inference incommensurable with its premisses: in the relation to the
eternally present it leaps over the intermediate links of grounding cognition
and grasps the eternally present as this present; it transforms the
uncertainty of what is to come into the certainty of something present. In
this relation faith shows itself to stand at one and the same time in relation
to the being's determinateness by finitude and to its determinateness by
infinity. Faith is only as that certainty which is born of uncertainty, which
as it were unwinds itself out of uncertainty; and through the uncertainty
which at every moment is to be conquered, finitude betrays itself. But that
certainty can be won through
% --- p. 39 ---
\setcounter{footnote}{0}%
uncertainty, that uncertainty \emph{can} be conquered, shows the infinity of
spirit. That which warrants faith's postulate, that which fills out the
distance between its premisses and its conclusion, is \emph{the essential
determination of the human spirit}, which with the necessity of the
\emph{a priori} points to \emph{the Absolute} as that without which
consciousness would lose its hold upon itself. It is the being's infinity
which, through the enclosure of finitude, compels consciousness to acknowledge
the infinite and eternal\footnote{Immediately there lies in this as yet only
such an acknowledgement as can be united with an averted disposition, as when
it is said of the demons in the New Testament that they believe in God and
tremble. But when the infinite is to be acknowledged as the in truth infinite,
then this acknowledgement is conditioned by a faith which acknowledges the
infinite and eternal in its validity as the ideal of human life, by a faith
which is in unity with self-surrender. In that relation in which faith stands
beside an averted disposition, the divine becomes, through the
\emph{opposition} to the self's abstract infinity, only an abstract infinite,
in reality therefore only a finite, with the form of negative infinity. The
rest of infinity which spirit seeks, because infinity is its essential
determination, is therefore found only where faith and self-surrender are in
unity, i.e. where faith is according to the completeness of its essence.}. And
through the infinity of the \emph{being} there works the infinity of the
\emph{principle}, by which the being has its being. That postulate of faith
has therefore at once a human and a divine origin; it is an expression of the
unity of the working of the human and of the divine spirit.
% printed mark: *) --- the only note on p. 39, below a short rule.

As such a postulate of cognition faith is at work in that, in the ethical
choice, man's living is determined by the infinite and eternal; and as such it
is at work in that the infinity of life is held fast in the particular action
which fills out human life and which, as \emph{ethical}, confirms that
primitive choice's constituting of the character of the life of spirit.

In the ethical choice and in the ethical action which confirms the choice
there reveals itself the \emph{true form of the unity of cognition and will}
in the relation of faith. The choice, and the action which rests upon the
choice, is at one and the same time
% --- p. 40 ---
\setcounter{footnote}{0}%
the expression of that energy of cognition with which cognition breaks through
the finitude that veils infinity and with the power of its own essence grasps
infinity, and the expression of the will's resolution. The will, which through
cognition is clear about itself, demands as the goal in which it, as
essence-willing, can come to rest, the infinite, which is the condition of the
infinity of essence; it demands the infinite in order to be able to be
itself\footnote{The will's possible resistance to the divine will be
elucidated under Section III.}. And just as the will demands the Absolute in
order to rest in it, so also does feeling first find rest in it; and feeling
preserves the Absolute's absoluteness in that at every moment of its intensive
concentration it as it were conceals within itself the memory of cognition and
the tension towards will.
% printed mark: *) --- the only note on p. 40, below a short rule.

g) The determination of \emph{self-surrender} comes forward in that the
particular finite will is to break through the determinateness by finitude in
which it immediately has its life, in order to will infinitely, and in this
willing --- which is a willing of the true infinity of its own essence --- to
enter into the absolute relation, to set itself in unity with the in truth
infinite, to let itself be determined by the in truth infinite as a member of
that totality, that community of spirit, which the infinite governs.

Self-surrender is first the self-surrender of \emph{the choice}, the choosing
will's decisive surrender in principle of the determinateness by finitude in
order to be able to grasp infinity. But the moment of the constituting choice
is to be unfolded so as to comprehend a \emph{whole} life, in whose particular
actions the choice, formed into an abiding resolution, repeats itself, through
the confirmation of the will's first surrender of the finite content of will
and of its own selfishness. In the choice and in the repetition of the choice,
self-surrender acquires on the one side the form of resignation, of dying away
from immediacy (S.\ Kierkegaard), of inwardisation, of ``self-annihilation''
in the sense of the abolition of the selfish; and on the other side --- since
the goal is a regaining by the self of itself according to its determinateness
by eternity and in concrete actuality --- the form of a self-assertion
according to the self's
% --- p. 41 ---
\setcounter{footnote}{0}%
truth. But since such a regaining, such a self-asserting, is possible only
through the unity with the divine, which is the principle of the being,
self-surrender acquires the form of a surrender to the working of the divine
(«\textit{Deum pati}») --- a surrender which is not merely the purification of
the finite will but, as its infinitising, the expression of its highest
energy: it becomes a suffering which is one with the highest
acting\footnote{With S.\ \emph{Kierkegaard} religious suffering and
self-annihilation are set in opposition to the ethical self-assertion in which
the existing individual, by ``repenting through'', is to find himself again.
This opposition hangs together with the conception of the ethical stage as a
stage of transition. If the ethical is apprehended in that original unity with
the religious which has been shown above, then already the ethical
self-assertion becomes a self-assertion through unity with the self's absolute
principle, and the opposition between the religious and the ethical relation
is abolished.}.
% printed mark: *) --- the only note on p. 41, below a short rule.
% ANTIQUA LATIN, p. 41: «Deum pati» is set in upright antiqua inside FRENCH
% GUILLEMETS -- the guillemets are antiqua too, and are the one place in the
% book where the usual „…“ is not used. Latin -> \textit{} per the house
% convention; the guillemets are KEPT as printed rather than converted to
% ``…'', per the playbook's rule for foreign-language guillemet terms.
% Tight antiqua, so no \emph{}.

Its \emph{definitive} determination self-surrender acquires in that the
subject which surrenders itself to the divine determines its own life in the
determinateness by finitude as \emph{guilt} --- in that self-surrender
therefore becomes one with consciousness's repentant acknowledgement of guilt,
and with its striving after that guilt's reconciliation in God through a
divine working. This determination is here only touched upon; its closer
development will be given under Section III, as will the demonstration of its
relation to the determinations of the Christian.

In self-surrender, in which the will's working comes forward primitively, this
will is, as emerges from the earlier developments, conditioned by and in unity
with cognition. And just as fully as self-surrender beneath the infinite must
be determined as the will's goal, so must it also be determined as cognition's
and as feeling's; for when the cognising spirit, through the surrender of its
finitude, finds rest in self-surrender to the infinite, spirit as feeling
finds rest in that self-surrender which is the self-surrender of the feeling
of absolute dependence.

h) In the foregoing it has already been indicated how
% --- p. 42 ---
the moment of faith and the moment of self-surrender in the religious relation
stand in the closest connection with one another. Their concepts presuppose
one another, and in their full and true significance they cannot be thought
without one another. Faith presupposes, as was shown (p.\ 39 note),
self-surrender in order to be able to be faith in the in truth infinite; it
presupposes self-surrender in order to be able to grow its own growth; and it
bears self-surrender within itself, since as a grasping and a holding fast of
the eternal --- which in faith's certainty is lifted out of the uncertainty of
doubt --- it is a self-giving of spirit to the eternal truth in which human
life as a whole finds rest. And just as faith presupposes self-surrender, so
self-surrender presupposes faith, precisely in order to be able to be
self-surrender to the \emph{true} divine, i.e. in order in truth to be able to
be self-surrender; and that action of self-surrender by which the self wins
its life through the surrender of its life's selfishness is such that life can
be regained in its truth and eternity only through faith in the eternal, to
which life can thus be surrendered.
% letterspacing on p. 42: „s a n d e“ only --- „Guddommelige“ that follows it
% is tight. The English run falls on the same element.

i) If we sum up the given determinations of the essence of the religious
relation, then religion, according to its ideal significance, becomes the
absolute relation to the Absolute as the true infinite standing in a relation
of essential unity with us. It becomes a relation which, as grounded in the
human being itself, comprehends the whole of human life, all the fundamental
activities of spirit. It becomes --- since man is seen as a unity of infinity
and finitude --- a relation of faith and of self-surrender; a relation through
which the human spirit, entangled in finitude and bound in guilt, finds
reconciliation within itself and with its own origin.
% the enumerators i) on p. 42 and a)--e) on pp. 44--51 are antiqua italic in
% the print, exactly like a)--h) on pp. 31--41; set plain here, following the
% convention already established for pp. 30--41 and in the Contents.

\begin{center}\rule{2cm}{0.4pt}\end{center}
% A short centred rule stands here on p. 42, between „med sit Ophav.“ and the
% paragraph on the positive religions; a thought-break in the running text,
% not a head. Same treatment as the rules on pp. 19 and 28. A SINGLE hairline.

The relation between religion in the sense indicated and the \emph{positive}
religions is indicated in its main features in the ``preliminary
determinations of concepts'' in the earlier work; it will emerge more
determinately from the developments that follow.

The relation of the conception developed above to S.\ \emph{Kierkegaard}'s
determinations of the essence of the religious has
% --- p. 43 ---
\setcounter{footnote}{0}%
already been partly touched upon. It shows itself at first glance that this
conception has essential points of connection in his determination of that
stage of religiousness which is not yet that of the paradoxically religious;
but an incisive difference came forward in the difference of conception of the
religious's relation to the ethical, of the absolute relation's relation to
the relative ones. From \emph{Kierkegaard}'s determination of the religious as
the \emph{paradoxically} religious that conception must of course diverge
still more decisively. What sharpens the difference here is the opposition
between immanence and abstract transcendence; between the view of the human
\emph{being}'s imperturbability and man's \emph{essence}-transformation
through sin; between the maintaining of cognition in the religious relation
and its annihilation through the paradox, the absolutely incomprehensible;
between the acknowledgement of the real's essential validity and essential
unity with the ideal in existence's principle, in existence's wholeness and in
the human being\footnote{Cf.\ the following Section (III).}, and
Christianity's absolutely spiritualistic conception, which with Kierkegaard
becomes manifest in the negative consequences with which he concludes. With
these oppositions an essential difference in the conception of reconciliation
will also show itself to stand in connection.
% PRINTER'S DEFECT, p. 43: the Danish prints „Tilintetgjørrelse“, with a
% doubled r. Orthographic, and so not reproducible in English; recorded here.
% letterspacing across the line break, p. 43: the compositor spaces only the
% SECOND half of „Menneske-væsenets“ and only the FIRST half of
% „Væsens-forvandling“. The English compounds are hyphenated so that each run
% can fall on the same half of the word.
% printed mark: *) --- the only note on p. 43, below a short rule. The Roman
% numeral „(III)“ in it is antiqua, as all figures and Roman numerals in this
% book are; not italicised.

\section*{2. Cognition's Relation to the Religious according to its Ideal Determinateness.}
% Bold centred head in two lines, standing mid-page on p. 43 above a short
% centred rule (rule not transcribed). Agrees with the Contents exactly.

In order that cognition's relation to the religious, in the ideal sense of
this concept, may show itself as a relation of unity, it must be established
that cognition, by its relation to the human being and by the form in which it
expresses the being, can become the medium for a relation of infinity, an
absolute relation; that it can have absolute significance for the human
spirit; and that it encloses within itself that moment of faith and
% --- p. 44 ---
of self-surrender which showed itself to be an essential mark of the
religious.

a) The condition for cognition's being able to become, for the human spirit,
the form of an absolute relation is that it can preserve the object's true
infinity within the relation. That it must be able to preserve the
determinateness by infinity lies already, so far, in cognition's form, in that
this form is determined by the universal; and this universal --- whose being
as universal for human consciousness is conditioned by the being's inner
infinity --- stands in essential connection with the infinite. But in order
that through this form of cognition the infinite may be preserved as the in
truth infinite, cognition's form of universality must be in concrete unity
with the other fundamental forms of the concept. In order that it may be this,
it is presupposed that cognition in the relation to the Absolute does not
stand in abstract opposition to what is just as fully a fundamental function
of spirit, and that it does not become a one-sided form of spirit's activity,
isolated from what supplements and totalises it. As such it would, for
example, unavoidably have to be apprehended if it were separated from the will
by an absolute heterogeneity of principle; for in this separation the will
would appropriate the determination of singularity as its fundamental form and
let cognition's form of universality become that of abstract universality,
while the will's form of singularity, in its separation from the universal,
would become precisely only the form of immediate natural singularity. In this
abstract opposition neither of the functions of spirit would be apprehended
according to its truth, i.e. as an expression of the undivided, total human
spirit, and neither of them would come in truth to relate itself to the
infinite, because the infinite in its true infinity could not be for either of
them.

With respect to the relation between the formal determinatenesses of the
primitive functions of spirit, what is needful has been developed in the
previous Section and in the earlier work, pp.\ 134 ff. It has been shown in
what relation cognition's determinateness by universality stands to the
particular and the concrete, and how the will's determinateness by singularity
--- precisely because the \emph{will's} singularity is ``singularity of
spir\-%
% --- p. 45 ---
it'' --- unfolds the determinateness of the universal out of itself. Of what
nature the will's determinateness by singularity is shows itself already
through the fact that the will has its limit in the organism against the
singularity and immediacy of nature-determined reality; and the will's
psychological relation to that determinateness by drive out of which it
unfolds itself in the course of the development of the life of spirit, and
into which it is unconsciously converted as acting upon the self's natural
side, its organism --- as for example when an act of will calls forth a
momentary or a mechanically continued movement --- points to a limit of the
will as will, by which its form of singularity is determined. If, however, in
order to maintain the absolute heterogeneity, one were to let the will's
universal set itself up as something absolutely heterogeneous from cognition's
universal, then such a doubling of the universal into two absolutely
heterogeneous kinds of universal, yet designated by the same common
determination, would immediately enclose a self-contradiction, an absurdity.
From what has been developed in the foregoing it therefore emerges that for
that cognition which, as an essential expression of the human spirit undivided
within itself, cannot be held fast in abstract separation from what is just as
fully a primitive expression of spirit, and which is therefore not
one-sidedly determined by a single fundamental form of the concept --- for
such cognition the Absolute can be an object according to its truth, as the in
truth infinite; and its infinity can be preserved under the relative
abstraction. It can therefore be present for cognition in the form in which it
is the object of religion according to religion's ideal significance, in which
it is the object of faith as that which is demanded by the human spirit, as
that goal of spirit in which the being's claim finds satisfaction, in which
spirit first finds rest.

b) That cognition can have \emph{absolute} significance for the human spirit
--- that significance, therefore, which it must be able to have in order to be
able to express the religious relation --- shows itself first through this,
that cognition, in virtue of its relation to the human being, must be \emph{an
end in and for itself}. It becomes an end in that it is a realisation of a
fundamental determination in the human spirit, and more closely, of the
determination which is according to
% letterspacing p. 45: „a b s o l u t“ (not „Betydning“), and the whole phrase
% „Ø i e m e d  i  o g  f o r  s i g“; the second „Øiemed“, two lines below,
% is tight. The English follows in both places.
% --- p. 46 ---
the concept principal in spirit; but this means: it is an end as the being's
realisation in one of the being's fundamental forms. But since it is a form
for the realisation of the universal being, it is the form for the individual
spirit's expansion so as ideally to take up the whole of infinite existence;
it is the energy which ideally binds the individual spirit together with this
whole of existence. Thereby cognition's absolute significance shows itself
from a new standpoint; for in realising spirit's determinateness by
universality, and in ideally binding spirit together with the whole of
existence, it binds spirit together with existence's principle, the Absolute,
which as such is cognition's principle and moving power, that upon whose
cognition all true cognition rests. Cognition therefore acquires, from this
standpoint, absolute significance because it realises that relation of spirit
to the Absolute which is laid down in the human being; and since this
realising of the human being's deepest relation is seen as being only through
the working of the indwelling Absolute, cognition's absolute significance
shows itself once more in a new light, which makes it still clearer how spirit
can relate itself religiously in cognition and to cognition.

From the standpoints emphasised, cognition has as yet been directly regarded
only according to its significance in the \emph{theoretical} respect, through
which it already shows itself as a self-valid end for that being which in
cognition realises itself according to a fundamental determination within
itself, and which precisely only in relating itself to cognition as an end can
be present in cognition with spirit's whole power, with that love of cognition
by which alone cognition becomes productive and creative, animated by the
idea. But since now the human being --- which in cognition realises itself and
realises its ideal unity with existence and with existence's principle --- as
a being which comprehends a \emph{totality} of functions, enters as an active
\emph{member} into a whole of existence and becomes conscious of itself as
such, cognition's absolute significance is also to be seen from the standpoint
of \emph{action}. How it is to be
% letterspacing p. 46: „t h e o r e t i s k“ (not „Henseende“),
% „T o t a l i t e t“, „L e d“ (which stands at the head of its line, „activt“
% before it tight), and „H a n d l i n g e n s“ (not „Synspunkt“).
% --- p. 47 ---
apprehended from this standpoint emerges from what was developed in the
previous Section. Only through cognition does that action become possible in
which the individual works according to its absolute determination, which is
the being's realisation through the will's freedom. Cognition showed itself to
become the presupposition of such an action, partly by its clarifying,
liberating, purifying influence upon the will, partly as determining the
content of the true will. From this standpoint --- where cognition shows
itself as a function within an organic wholeness of life, while it works as a
condition, as a means, for the ethical --- it is at the same time to be seen
as an ethical \emph{end}, and as working towards its own realisation in
working towards the realisation of that unity of life in which it has its
truth. And since cognition exercises upon the other activities of spirit, upon
the life of feeling and the life of imagination, the same purifying and
clarifying influence which it exercised upon the will, it shows itself in its
relation to them too in its significance as realising the inner unity and
harmony of the total human being; and this standpoint thus leads back to the
first, from which cognition was seen as a realisation of the human being.

c) That there is in cognition a moment of \emph{faith} --- not in the Jacobian
sense of an immediate certainty of being, but in a sense in which the
fundamental determination of the religious relation of faith is preserved ---
shows itself at \emph{every} stage of cognition's development. This moment of
faith is covertly present in the lower forms of cognition everywhere where a
universal is asserted on the ground of the particular that is incongruent with
it. Where a universal law asserts in virtue of the incompletely given
particular, where the incomplete premiss of induction is therefore given
validity as a complete one, there the universal is brought forward in that the
empirical is supplemented in virtue of something a priori, that is, in virtue
of a claim of the \emph{being} to unity and lawfulness. And what thus already
shows itself where in a lower sphere a \emph{relative universal} comes
forward, shows itself still more plainly where the \emph{universal in the
highest sense}, the \emph{total existence's}
% letterspacing p. 47: „Ø i e m e d“ (not „ethisk“); „T r o e n s“;
% „e t h v e r t“, with „Trin“ on the next line tight --- as printed;
% „V æ s e n e t s“, with the following „For=“ tight;
% „r e l a t i v t  A l m e n t“, spaced through the line break; and
% „h ø i e s t e  F o r s t a n d  A l m e n e“ and „t o t a l e
% T i l v æ r e l s e s“ as two runs, the „den“ between them tight. The run
% stops at the page break: „universelle Princip“ on p. 48 is tight.
% --- p. 48 ---
universal principle, is grasped by cognition out of incomplete presuppositions.
For since one goes beyond experience's unavoidable incompleteness, the
unity-principle of an existence not exhausted by experience is grasped, not as
that whose certainty experience secures, but as that which is reached by a
necessary postulate out of the individual consciousness, which without it
would lose its hold upon itself, become uncertain of its own unity, be
confused within itself. This postulate of the principle --- which as
existence's principle is the human spirit's principle of life --- comes
forward also in that cognition which is the presupposition of the primitive
ethical choice by which the self constitutes itself as an ethical self; and
just as, therefore, the cognition upon which this choice rests encloses the
moment of faith, so this moment is present also in that cognition for which
the incomplete premisses of the particular ethical action acquire validity as
complete ones, as an action's \textit{ratio sufficiens}.

In this --- that it is a demand proceeding with necessity from the being's own
source, and a demand which encloses the immediate acknowledgement of what is
demanded --- the faith which cognition encloses agrees with the faith which
has been determined as the distinguishing mark of the religious. But the faith
which is present in cognition itself must have the character of the
\emph{rational}. In demanding the Absolute, its demand rests upon the thinking
human being's consciousness of itself; it does not come into conflict with
that consciousness, but confirms it. Faith as rational faith becomes a faith
in that which is higher than the individual, in that it is the principle which
grounds the individual spirits and the whole of existence; but which, in being
determined as higher, is at the same time determined as of the same essence,
precisely because it is the \emph{principle} of the human being and of
existence. And since the being from which the religious demand of faith
proceeds is an indivisible unity, there lies in this incontrovertibly that
this demand, in order in truth to express the being's claim, cannot stand in
conflict with any true expression of the being, and therefore not with
thought, with cognition. Faith in the sense of the in truth reli\-%
% antiqua on p. 48: „ratio sufficiens“, upright antiqua in the print; set
% \textit{} per the house convention, tight, so no \emph{}.
% letterspacing p. 48: „R a t i o n e l l e s“ (not „Charakteer“) and
% „P r i n c i p“ in „Menneskevæsenets og Tilværelsens P r i n c i p.“
% --- p. 49 ---
gious can therefore not be a faith which stands in principle opposed to
knowledge; by this disharmony in principle it would show itself not to have
its source in the essential. This faith must, in virtue of the being's unity,
stand in harmony in principle with knowledge; but in virtue of the human
spirit's determinateness by finitude, faith must, as \emph{anticipating the
result}, acquire an \emph{abiding} significance as something necessary for the
human spirit. It is the form of the indispensable anticipation of certainty;
it is the demand grounded in essence, which has the being's own necessity; it
is therefore not like an uncertain knowledge, but like an inference without
premisses which nevertheless has incontrovertible certainty, because the
premisses are hidden in the being itself. But precisely therefore, precisely
because faith draws its certainty from this source in the being, it cannot
come into conflict with what springs clearly and transparently from the same
source. Where, therefore, that which is posited as a determination of faith
incontrovertibly conflicts with knowledge --- and it does so when it conflicts
not merely with this or that determination of knowledge, but with what makes
knowledge knowledge, with \emph{knowledge's principle, the principle of
grounding and of necessity} --- there it cannot have truth. But in such
conflict with knowledge's very principle stands, for example, faith in the
miracle, when the miracle is taken in the sense in which positive religion
takes it, not in the effaced and falsified sense in which it is taken by a
half-theological inverted speculation.

When ``faith'' at the stage of the positive religions appears with a claim to
unconditional validity, then it must support this claim either upon its divine
origin or upon its springing from what is innermost and most essential in the
human spirit. If it supports the claim upon its springing from the
transcendent divine --- and only when the divine is so apprehended do the two
standpoints fall outside one another --- then faith with its supernatural
origin and its supernatural content must be posited as raised above all
knowledge and at the same time as standing in
% letterspacing p. 49: „f o r e g r i b e n d e  R e s u l t a t e t“ (spaced
% through the line break), „b l i v e n d e“ (with „Betydning“ tight), and the
% run „V i d e n s  P r i n c i p ,  B e g r u n d e l s e n s  o g
% N ø d v e n d i g h e d e n s“ --- the second „Princip,“ that closes the
% phrase is tight. The English follows in each case.
% --- p. 50 ---
\setcounter{footnote}{0}%
opposition to all knowledge. But since knowledge is seen as an expression of
that human nature which has its springing from God, there arises not merely a
discord between the divine's opposed forms of revelation, in faith and in
thought, but there arises an unabolished and unabolishable discord in the
human spirit, which is to deny one expression of essence in order to
acknowledge an effect proceeding from the same source from which the human
being itself has sprung. If then, in order to maintain the necessity of
acknowledging ``revelation'', reference is made to the demand of personality,
to a norm within personality, then the last of the standpoints indicated above
has hereby been put in the place of the first. But when faith then rests upon
the being's inescapable claim, then what was said above holds again: that this
being's unity must preserve the harmony between the expressions of the being;
and a harmony is then required between faith's content and what is essential
in knowledge. As surely as there have been, and constantly can be, false forms
of the positively religious --- forms of doctrine and cultus which most deeply
offend the moral consciousness\footnote{It is here sufficient to recall the
Near Eastern religions.} --- and as surely as conscience can be entangled and
unclear in the individual, so surely is a norm required which cannot be sought
in something absolutely heterogeneous with knowledge, but must be sought in
knowledge's negative corrective, in that veto of knowledge which must retain
its validity; because, when knowledge is denied, that regulative is lost by
which a distinction can be drawn between the true faith grounded in the being
and that which is only an expression of finite craving and desire, of
something unessential and particular in human nature. Knowledge's relation to
the proposition of faith may be such that knowledge must halt at a
\textit{non liquet}; but knowledge's relation can never become such that it
must approve what indisputably contradicts its \emph{principle}, annihilates
it \emph{as} knowledge.

d) In cognition, as cognition is according to its truth, there is finally
present that moment of \emph{self-surrender} which characterised the in truth
religious relation. In all true cognition there is a purification of thought,
a surrender of individual
% antiqua on p. 50: „non liquet“, upright antiqua; neither OCR witness marks
% it, and it is the second Latin tag on this stretch besides „ratio
% sufficiens“ (p. 48). Tight antiqua, so no \emph{}.
% letterspacing p. 50: „P r i n c i p“ and „s o m“ in „modsiger dens
% P r i n c i p, tilintetgjør den s o m Viden.“, and „H e n g i v e l s e n s“.
% printed mark: *) --- the only note on p. 50, below a short rule.
% --- p. 51 ---
entanglement, opinion and prejudice --- in one word, of thought's selfishness;
there is in it a self-surrender to truth's absolute power. And this
self-surrender, which is present in all true cognition, shows itself most
clearly in that relation in which cognition, in virtue of the being's demand
of faith, surrenders itself to the infinite that satisfies the being's need.
In this surrender, which is conditioned by an inner necessity of the being,
cognition is in unity with the will; and in being itself infinitised in the
surrender, cognition clarifies that will whose choice is determined by the
same power of the being.

e) The result of what has been developed is therefore that cognition in its
truth stands in an essential relation of unity to the religious in its ideal
determinateness, to the religious faith which is the expression of the human
being's infinite demand. Under the conditions of finitude of the human spirit
there is indeed as it were an interval between faith and knowledge, but an
interval which must be able to be traversed by knowledge. A discord in
principle, unabolishable, becomes an impossibility, because that faith which
abolished knowledge's essential determination would precisely thereby show
itself as not according with faith's own essence. Under the changing of its
historical forms faith has a development, and its essential connection with
knowledge shows itself in this, that this development of faith is not torn
loose from that of knowledge, is not without the liberating knowledge.
% p. 51: a short centred rule stands between this paragraph and the sub-head
% below --- a rule above a display head, and so not transcribed, as at p. 30.

\section*{3. The Relation of the Will and of Action to the Religious according to its Ideal Determinateness.}
% Bold centred head in two lines, p. 51. Agrees with the Contents exactly.

In determining the relation of the will and of action to the religious, the
same standpoints will have to be brought to bear as were emphasised in the
foregoing with respect to cognition.

a) There is then first to be brought out that the possibility of an absolute
relation to the Absolute through the will is present, in that the will, in its
harmony with cognition and in its being clarified through by the universal,
has the capacity to pre\-%
% No letterspacing anywhere on p. 51 --- checked line by line at 260 dpi.
% --- p. 52 ---
serve its object's true infinity --- a capacity which would have to be denied
if one let an absolute heterogeneity separate the will from cognition, or if,
on the ground of an unclear representation of the being's ``defectiveness'',
an \emph{abiding} irrational element in the will were postulated. The Absolute
is present as absolute in the will's form as \emph{the good}, and man's
willing of the good is in unity with the cognition of it. As this willing of
the Absolute as the good impresses itself in action, this action acquires ---
as will be more closely developed in what follows --- the shape of a
realisation of the good in the form of the human being, governed by the will's
unconditional submission to the Absolute's demand; so that this realising,
through the unconditional submission, becomes the expression of the
realisation of essential unity with the divine (ὁμοιωθῆναι θεῷ), and therefore
the form of an absolute relation of will to the Absolute.
% GREEK on p. 52: „ὁμοιωθῆναι θεῷ“, antiqua italic, broken „ὁμοιω-/θῆναι“
% across the line in the print. Neither OCR witness reads any of it, and p. 52
% is not in the recon pass's Greek list; it exists only because it was caught
% by eye. Copied verbatim from transcription.tex, not retyped.
% The parenthesis „hvad der i det Efterfølgende nærmere vil blive udviklet“ is
% printed with plain commas, not parentheses; set with dashes here for English
% readability, the punctuation carrying no emphasis either way.

b) The \emph{absolute significance} of true willing and of true acting for the
human spirit shows itself immediately, in that the will as
\emph{essence}-willing conditions the unfolding of the human being's endowment
in real actuality, and constitutes the true relation to the Absolute. As shown
above, this twofold standpoint comes together into an inseparable unity.

When, in the development of the ethical's coming into being in the individual,
the expression was used that the self in the ethical choice chooses and wills
\emph{itself absolutely}, there was a pointer to this: that such choosing is
possible only through the self's standing in a \emph{relation of essential
unity to the Absolute}; for only in virtue of such a relation can it be chosen
absolutely. That expression must be taken as the starting point, because, if
we determined the object of the choice as the absolute idea of the good, the
choice of this as that which is to be realised in \emph{human} action would be
identical with the choice of the \emph{humanly good}, of the \emph{ideally
human}. In this the absolute idea is for man in \emph{the determinateness of
his essence}, and in being realised through the free will, actuality is given
to the idea. Beyond the determination of the ideally human
% letterspacing p. 52: „b l i v e n d e“ (not „Irrationelt“); „d e t  G o d e“
% at its FIRST occurrence only --- the two that follow are tight;
% „a b s o - / l u t e  B e t y d n i n g“, spaced through the line break;
% „V æ s e n s“ with „villen“ tight, i.e. only the first half of the compound;
% „s i g  a b s o l u t“; „V æ s e n s = / e n h e d s f o r h o l d  t i l
% d e t  A b s o l u t e“; „m e n n e s k e l i g“ with „Handling“ tight;
% „m e n n e s k e l i g t  G o d e“; „i d e a l t  M e n n e s k e l i g e“;
% and „d e t s  V æ s e n s  B e s t e m t h e d“. The last line of the page
% repeats „det idealt Menneskelige“ TIGHT --- as printed, not regularised, and
% the English repeats it without \emph{} accordingly.
% --- p. 53 ---
as the task of human action we cannot come. We can realise only that Absolute
which is \emph{our} absolute essence, actualise only that ideal which is
\emph{our} ideal. From the determination of the ethical task as the
realisation of the idea of the good as such we are at once led back to the
limit indicated, by the fact that what realises must be congruent with what is
to be realised; the Absolute can then be \emph{positively} realised by us only
in so far as it expresses itself in the determinateness of our essence. But
this limitation of the positive expression for the idea's realisation through
the will does not exclude, but precisely presupposes, another negative form of
realisation, in which the idea makes itself felt \emph{as such}. If we see the
human individual ordered as a member into a principle-determined organised
totality, then it follows indeed from its determinateness as a single organic
member that it can positively express the principle only through the active
unfolding of the individuality's own essence; but in the form of negation
(self-surrender) it expresses the whole as whole, the principle as principle,
in virtue of the whole's and the principle's presence within it in the
abstract form of possibility. Through that determination of the ethical choice
and the ethical task which was our starting point there is therefore a pointer
to the relation to the Absolute as that by which the self in its ideality can
become an absolute task to itself, and which acquires presence in the
ethically unfolded self; and the absolute significance of ethical willing and
acting therefore shows itself from a twofold standpoint.

To the same result we are led when we start from the other expression for the
ethical choice: that the self in that choice chooses \emph{itself according to
its essential determinateness}. For the choice of the self according to its
essence is a choice of the being \emph{according to its total concretion}, and
therefore according to its determinateness \emph{as an organic member in a
totality of existence}. In cognising itself in this determinateness, the self
is from the beginning in relation to that principle which is the governing
principle of life of the total existence. The I's self-assertion in the
ethical choice is only through the principle's working through the I,
% letterspacing p. 53: „v o r t“ twice; „p o s i t i v t“; „s o m
% s a a d a n.“ (spaced through the line break); „s i g  e f t e r  s i n
% V æ s e n s b e s t e m t h e d“; „e f t e r  d e t s  t o t a l e
% C o n c r e t i o n“; and „s o m“ at the end of a line followed by
% „o r g a n i s k  L e d  i  e n  T i l v æ r e l s e n s  T o t a l i t e t.“
% --- ONE emphasis carried across the break, not two. The English keeps it as
% one run accordingly.
% --- p. 54 ---
and in that the I asserts itself, but asserts itself in its being ordered into
the totality, its self-assertion is only as self-surrender to the principle.
And what holds of the I's relation in the ethical choice holds of its relation
in the particular ethical action, for which this choice forms the foundation.
As the ethical wins concretion, and ethical willing forms itself as that
willing of the being according to its completeness which has been clarified by
true cognition, the relation is at every point posited to that totality into
which the self as ethical freely orders itself, and to the totality's
principle. The relation which was posited in the ethical choice, when the
ethically a priori first acquired validity for the will, is preserved as this
a priori wins fullness in the concrete action. The twofold standpoint for the
ethical's absolute significance thus gathers itself once more into unity.

When the relation to the Absolute which is constituted in the ethical choice
is posited as preserved everywhere in ethical action, this leads to an
essential consequence --- already drawn above from another standpoint --- for
the conception of the \emph{particular} ethical relations. From the fact that
in these relations in their singularity there is at every point a relation to
the Absolute working through them, it follows that the particular ethical
relations, which in their isolation would have to be determined as merely
relative relations, become relations with absolute significance, expressions
of the absolute relation. The absolute relation is seen as comprehending all
the ethical relations, as indwelling in them, and as unfolding through them a
concrete content and winning the shape of actuality. Thereby the particular
ethical relations acquire a higher significance; and in apprehending them in
this significance the subject relates itself religiously to the ethical, which
now in its articulation acquires absolute validity, becomes an absolute end
for the subject.

It is only another expression for the same thing when it was said that the
ethical can never as a whole be given up for a religious τέλος which as an
absolute should raise itself above everything ethical. It cannot, because the
ethical is an expression of essence, because through its relation to the
Absolute it has absolute
% GREEK on p. 54: „τέλος“, antiqua italic. tesseract renders it „7é4o0ç“ and
% the ABBYY layer drops it altogether; p. 54 is not in the recon pass's Greek
% list. Copied verbatim from transcription.tex.
% letterspacing p. 54: only „s æ r l i g e“, at its FIRST occurrence; the
% second occurrence eight lines below is tight, and „ethiske Forhold“ is tight
% in both. As printed, and so in the English.
% --- p. 55 ---
significance, and because it stands in indissoluble connection with the
religious. In the collision in actuality between the different forms of the
ethical, the ethical from one sphere may give way to the ethical from another
which comes forward with the right of the higher universal; but in resigning a
single one of the ethical relations --- which, as members in the ethical
totality, are according to their essence absolute relations --- the
resignation of the particular relation is in virtue of the ethical totality,
which in the resignation \emph{integrates itself}. The resignation of the
particular ethical relation is therefore in virtue of the ethical itself;
there arises no absolute relation outside the ethical, but the Absolute
becomes omnipresent within the ethical.

c) In the ethical thus apprehended --- the expression of true, essential
willing and acting --- the moment of faith has already been shown to be
present. It has been shown that it comes forward not merely in the ethical
choice's first grasping of the Absolute, but everywhere in the choice's
confirmation in repetition, in every ethical resolution in which the
individual determines itself to action in spite of cognition's incompleteness
and the uncertainty of realisation, of the outcome; and it comes forward
likewise in that the individual, under the incomplete realisation of the
ethical demand, acknowledges in the indwelling, through-working principle ---
with which it is in unity through love's self-surrender --- that which
ethically completes.

d) The moment of self-surrender in the ethical has already been touched upon,
in that the negative form for the realisation of the absolute good was
indicated, and in that it was shown that the self's assertion of itself, as an
assertion of the self according to its essence's \emph{complete}
determinateness --- according to which it is a moment in a whole governed by
the principle --- became one with a self-surrender of love and of resignation
to that principle by which the self's essence is, and in which it first finds
itself again according to the truth of its essence. This holds with respect to
the first ethical choice and to each particular ethical action. Since now
ethical self-assertion is self-assertion through self-surrender, \emph{love}
is the principle of the ethical; and since the self is seen in relation to the
Absolute and to the spir\-%
% letterspacing p. 55: „i n t e g r e r e r  s i g  s e l v.“;
% „f u l d s t æ n d i g e“ (with „Bestemthed“ tight); „K j æ r l i g h e d e n“.
% The section letters „c)“ and „d)“ are antiqua, as „a)“/„b)“ on pp. 42--52.
% --- p. 56 ---
itual wholeness of life into which it is ordered, love's object becomes the
Absolute, the principle, and man \emph{as} man (the neighbour). In the
relation of love to man the relation of love to the Absolute is revealed: in
the ethical-religious relation to humanity there is a relation to the
Absolute, which is humanity's ground of life. In this love of the Absolute and
of man \emph{as} man, self-love is transfigured and man's innermost essence
revealed. ``Only that man is something who loves something, and what a man is
shall be known by what he loves''. (Feuerbach).
% The Feuerbach quotation, checked at ~490 dpi in the Danish: low „ opening,
% high “ closing, and the full stop stands AFTER the closing mark, not before
% it. „(Feuerbach).“ runs on as the last, short line of the same paragraph; it
% is not set off. The English keeps the stop outside the quotation mark and
% the attribution running on, as printed.

e) Thus the result with respect to the will's relation to the religious
becomes the same as with respect to cognition's. The relation between the will
in its truth --- the essence-willing at work in ethical action --- and the
religious according to its ideal determinateness is a relation of unity, and
must be so, because a fundamental function of the being cannot contradict what
expresses the human being's deepest claim, its infinite demand.
% A short centred rule stands between this paragraph and the sub-head below ---
% a rule above a display head, and so not transcribed, as at pp. 30 and 51.

\section*{4. The Reconciliation of the Human Spirit in the Ethical-Human Life Resting upon True Cognition.}
% Bold centred head in two lines, p. 56; agrees with the Contents. The line
% break falls after „sande“; „ethisk=humane“ is the ordinary Fraktur hyphen.

When human life is unfolded in harmonious unity of true cognition and true
action; when cognition in its truth bears within itself faith according to the
ideal significance of that concept; when the ethical action, in which the
human being according to its concretion finds its realisation, has as ethical
action the true stamp of the religious; and when finally the human being
through cognition and action stands in relation to an existence and to a
principle of existence from which all dualistic dividing has been removed ---
then the conditions are present without which spirit cannot find rest and
reconciliation. What proceeded immediately from the essence of the human
spirit was a striving after the realisation of its inner infinity as a
% letterspacing p. 56: only „s o m“, twice, in „Mennesket s o m Menneske“ ---
% in both places the neighbouring „Mennesket“/„Menneske“ is TIGHT. The
% compositor spaced the little word alone; the English does the same.
% --- p. 57 ---
true infinity; and this realisation can take place only on the conditions that
spirit's activities are not divided, that they do not follow heterogeneous
laws, and that spirit in its relation to existence and to the principle is not
hindered by the finitude-barrier of a dualism. But in the true cognition of
the Absolute, of the totality of existence and of the self; in the true
ethical willing, which equally accentuates the self's self-assertion and its
ordering into the totality, its spirit-side and its nature-side; in the true
religious relation, which as the absolute relation penetrates the relative
ones and in them wins concretion --- in these this concrete infinity is
present; and since cognition, morality and religiousness are harmoniously
bound together, the condition is brought about under which spirit's original
striving can come to rest, under which spirit's reconciliation with itself in
its actuality can take place --- a reconciliation which must be lacking so
long as spirit is divided by conflicting powers, so long as consciousness is
split within itself, at every moment torn asunder by doubt about cognition's
truth and about action's validity. Only through that harmonious realisation of
all the powers of the life of spirit does human life also become in truth a
\emph{human} life according to the full significance of that concept; and the
character of humanity reveals itself in the \emph{universal} love of humanity,
of the human \emph{as} human.

When at this point of the development --- where \emph{guilt}, the breach in the
ethical, and the possibility of restitution in spite of that breach, have not
yet been treated --- reconciliation is spoken of, then of course this concept
is partly to be taken in a more abstract, more metaphysical sense: in the
sense it must acquire when one starts from infinity as an essential
determination, and when the form is sought under which the striving after
infinity that has proceeded from the ground of the human being can find its
satisfaction. And partly the actuality of reconciliation remains as yet only
\emph{hypothetical}, in that only the general conditions are indicated
\emph{without which it does not become} thinkable. But when in what follows
reconciliation is developed in more concrete determinateness as \emph{the
reconciliation of guilt},
% letterspacing p. 57: „h u m a n t“ (Liv tight); „u n i v e r s e l l e“;
% „s o m“ alone in „det Menneskelige s o m menneskeligt“, the same idiom as on
% p. 56; „S k y l =/d e n“, spaced through the line break, with „Bruddet“
% tight; „h y p o t h e t i s k“; „u d e n  h v i l k e  d e n  i k k e
% b l i v e r“ with „tænkelig.“ tight; and „S k y l d e n s  F o r =“ running
% over onto p. 58, where „s o n i n g,“ is spaced too --- ONE emphasis carried
% across the page turn, so the English \emph{} likewise spans the page marker.
% --- p. 58 ---
then this more concrete determination must have that more abstract one for its
presupposition and condition; it will not abolish it, but fill it out.

In the sense in which reconciliation is spoken of at this point of the
development, this concept designates first \emph{spirit's reconciliation
within itself}, as grounded in spirit's inner harmony, in the inner unity of
all the functions of spirit, in the being's undivided working. This conscious
harmony, resting upon true cognition, and the reconciliation which it
encloses, differ essentially from the immediate harmony of beauty of \emph{the
Greek consciousness} with its illusory reconciliation. In the Greek spirit's
harmony of beauty the oppositions were brought to unity only in an illusory
manner, precisely because the form of union was beauty's immediacy, the
\emph{immediate} connection of the \emph{originally different}. As soon as
this harmony, which was there for intuition, was to be determined by thought,
the oppositions parted; and Greek thought's determinateness by the form of
intuition let the parted oppositions become fixed in the separation without
being able to bring about a higher reconciliation. Therefore Greek
consciousness hovers between the immediate unity of beauty and the
unreconciled opposition which unavoidably succeeds it; and because that unity
of beauty in its immediacy does not have the true form of the life of spirit,
spirit lacks in it the deeper inwardisation, and precisely therefore that
inner infinity which is the condition of the in truth ethical; and likewise
there is lacking the possibility that life's concrete content can unfold
itself for the ethical consciousness as a principle-determined unity. Since
that unity which we have determined as the unity of the in truth human life
--- as a freely acquired unity, transparent for cognition, resting upon
another determination of the fundamental relation in the human and in
existence --- is secured against dividedness, it encloses the possibility of
that reconciliation which the Greek spirit, in virtue of its form of
consciousness, had to lack.

To the determination of the human spirit's reconciliation within itself there
attaches itself then the determination of its \emph{reconciliation with its
% --- p. 59 ---
divine principle}; and this determination of reconciliation shows itself to
stand in necessary connection with the first, in that only in spirit's
reconciliation with its indwelling principle can its inner reconciliation
become complete --- just as, conversely, the reconciliation with the
principle, which is the undivided spirit's reconciliation, can be accomplished
only as spirit wins the inner reconciliation. Reconciliation with the
principle takes place in that the unity of essence becomes conscious and finds
its --- negative and positive --- expression; it takes place in that man's relation
to the Absolute becomes an absolute relation, through the unconditionedness of
faith and of self-surrender. For the reconciled consciousness the Absolute is
the through-working power of life which, in virtue of the unity of essence,
can unconditionally govern human thought and determine human will, and in
which the self in the highest sense can find itself again. How this
finding-again becomes possible in spite of guilt it will be the task of the
following Section to show.

The human spirit's reconciliation is finally to be apprehended as its
\emph{reconciliation with existence and within existence}. It acquires this
determination in that man becomes conscious of himself as ordered into an
organic existence whose principle is the human being's principle, whose
governing laws are reason's inviolable laws, and in that he joins himself in
love with that wholeness of life in which he constitutes a member.

When it has now become clear from the foregoing that what is the
\textit{conditio} \emph{\textit{sine qua non}} for spirit's true
reconciliation, for its reconciliation in actuality, is present in that form
of religiousness which we may call \emph{the religiousness of humanity}, then
it will also become clear that it cannot be found within the positive
religions.

What prevents this is, in one word, the limitation by \emph{finitude} in the
positively religious. The striving after the in truth infinite, which is the
distinguishing mark of the religious in the ideal sense, is also the moving
power in the historically appearing positive forms of religion; but
% ANTIQUA LATIN on p. 59: the whole of „conditio sine qua non“ is antiqua, and
% within it „s i n e  q u a  n o n“ is ALSO letterspaced while „conditio“ is
% tight. Latin -> \textit{} per the house convention; the letterspacing is
% carried by the extra \emph{} round the spaced half, exactly as in the Danish.
% This is the clearest case in the book of why the spaced/tight antiqua
% distinction is worth keeping mechanically: in English the two halves would
% otherwise be indistinguishable.
% Other letterspacing p. 59: „g u d d o m m e l i g e  P r i n c i p;“ (the
% run-over from p. 58); „F o r s o n i n g  m e d  T i l v æ r e l s e n  o g
% i  T i l v æ =/r e l s e n.“; „H u m a n i t e t e n s  R e =/l i g i ø s i t e t“;
% and „E n d e l i g =/h e d s“ with „begrændsningen“ TIGHT --- the same
% half-spaced compound as „Aands“functioner on p. 6.
% --- p. 60 ---
\setcounter{footnote}{0}%
this striving acquires, as has been shown in another place\footnote{Cf.\
Problemet om Tro og Viden, p.\ 33.}, its limitation by finitude through the
peculiar formal determinateness of the theoretical and practical spirit at the
stage of these religions.
% printed mark: *) --- the only note on p. 60, below a short rule. The
% asterisk stands after „Sted“, before „er viist“.

That the human spirit, at the stage of the positive religions, must lack the
conditions for reconciliation shows itself from the various standpoints from
which the concept of reconciliation was apprehended above.

The human spirit at the stage of the positive religions \emph{cannot} attain
the \emph{inner reconciliation}, spirit's harmonious unity with itself. It
cannot, because the functions of spirit are in conflict with one another and
within themselves, and because this conflict is such as cannot be abolished.
Cognition and will, both of which are held fast in a determinateness by
finitude, must abidingly diverge; between knowledge and faith there arises an
insoluble contradiction, and likewise between the action determined by the
true cognition of the being's law and the action which is fulfilment of the
supranatural principle's arbitrary command. This finds its expression also in
the fact that the positive religion's determinateness by finitude makes the
religious relation --- the relation to the transcendent divine --- into a
\emph{particular} relation, which separates itself from the other essential
and, as such, warranted human relations, and negates them in laying claim to
spirit's whole power. The reconciliation which spirit was to attain in the
exclusive God-relation would --- if a reconciliation could be thought where
something essential is sought to be negated --- at all events not be able to
become a reconciliation \emph{within actuality}, not the reconciliation within
itself of the spirit existing in the concretion of actuality. It would become
a reconciliation in a faith which, as relating itself to the supranatural,
loses knowledge from sight, and in an action in which the ethical is absorbed
into an abstract asceticism or into an arbitrary statutory element.

Nor can the human spirit at the stage of the positive religions attain the
full reconciliation \emph{with its principle}. This shows itself whether we
look to the human spirit's
% letterspacing p. 60: „i k k e“ IS spaced in „Trin i k k e naae den
% i n d r e  F o r s o n i n g“, but „naae den“ between them is TIGHT --- an
% emphasis the compositor set in two pieces. The parallel sentence twelve
% lines below reads „ikke heller naae den fulde Forsoning“ ALL tight, with
% only „m e d  s i t  P r i n c i p.“ spaced, which is why the first „ikke“ is
% a positive reading and not a justification artefact. The English follows.
% --- p. 61 ---
\setcounter{footnote}{0}%
conditions or to the determinateness of the divine. If, with respect to the
human spirit's conditions, we leave aside what has already been touched upon
--- the opposition of the functions of spirit, which must prevent spirit from
entering into the in truth absolute relation --- then there is a particular
hindrance to such a relation in the unclarified drive-like element which, in
virtue of the determinateness by finitude of cognition and will, and in virtue
of the divergence of will and faith from cognition, remains over at this
standpoint. For the drive-like as such, which is preserved in the affect of
fear and hope in the form in which that affect comes forward in the religions,
conceals unconsciously, in seeking its satisfaction, a residue of egoism,
which can be removed only by the clarification of the drive through true
cognition, by the bringing about of spirit's inner unity, transparent to
itself. --- If we look to the conception of the divine at the stage of the
positive religions, then it has been shown\footnote{Cf.\ Problemet om Tro og
Viden, pp.\ 35 f.} how the form of representation brings an abiding
determinateness by finitude into that conception. And this determinateness by
finitude in the transcendent divine then unavoidably entails that man's
relation to the divine cannot become in truth absolute. The unconditional
self-surrender, freed from all egoism, must be lacking; for it can be found
only where man relates himself to the in truth infinite. At the stage of the
positive religions the true relation of subjectivity can find no place either.
The subject is here, in the religious relation, determined by an other for the
subject, which therefore --- although itself determined as subject, as
personality --- becomes, on account of the transcendence, of the separation of
essence, something objective for the subject. Only when the relation is
determined by immanence does the relation of subjectivity become in the
fullest sense a relation of subjectivity.

Finally, the human spirit at the stage of the positive religions cannot attain
reconciliation \emph{with and within existence}. What hinders this in this
respect is first and foremost the representation of \emph{the miracle}. By
this represen\-%
% printed mark: *) --- the only note on p. 61, below a short rule.
% letterspacing p. 61: only „m e d  o g  i  T i l v æ r e l s e n.“ (with
% „ikke naae Forsoningen“ tight) and „U n d e r e t“.
% --- p. 62 ---
\setcounter{footnote}{0}%
tation, which, as has been shown in another place\footnote{Cf.\ Problemet om
Tro og Viden, p.\ 35.}, finitises the divine and brings a discord into God,
there is brought also a discord and a contradiction into existence; and this
inner contradiction of existence --- its essencelessness, posited by the
miracle, which yet must at every moment as it were be revoked\footnote{The
same work, pp.\ 79 ff.}, if only to be posited anew --- prevents spirit from
finding rest in this existence, which shows itself as in a confusing double
light for thought and will. And were the religious consciousness, out of this
doubleness, to grasp the representation of the natural's essencelessness, then
spirit would be referred to seeking its life's goal through a flight from
natural existence, through a denial of it; but of a reconciliation within
natural existence --- back to which the human being's nature-side
\emph{must} constantly draw spirit --- there could be no question.
% printed marks on p. 62: *) and **), below a short rule.

From the standpoint here brought to bear it therefore becomes the
\emph{provisional} result that what is the \emph{indispensable} condition for
the possibility of reconciliation is present only in that form of
religiousness which is the highest expression of the human. So far, then, as
this provisional result is not overturned by the \emph{definitive}
determination of the conditions of reconciliation, this form of religiousness
will show itself as the realisation of religion's concept, as the adequate
expression from that which is striven after in the positive religions.
% ERRATUM, p. 62 line 24: the page prints „fra det i de positive Religioner
% Tilstræbte“; the Trykfeil leaf (printed p. 211) asks for „for det“. The
% Danish is transcribed AS PRINTED, per the editorial stance, and the English
% follows it: "expression FROM that which is striven after", not "for". This
% is the author's own and only erratum, and it is not applied in either file.
% letterspacing p. 62: „f o r e l ø b i g e“ (Resultat tight);
% „u u n d =/v æ r l i g“ (Betingelse tight); „d e f i n i t i v e“; and „m a a“
% alone at the end of a line, with „drage det tilbage“ on the next line TIGHT.
% That last is the one doubtful call on the page --- it sits at a line end,
% where a compositor may space to justify --- and the transcription sets it as
% printed while recording the alternative reading. The English follows the
% transcription.

\begin{center}\rule{2cm}{0.4pt}\end{center}
% A short centred rule stands here on p. 62, between the two paragraphs (not
% under a head), as at pp. 19 and 28. A SINGLE hairline.

By this result there is a pointer, then, to our last and most difficult task:
through a treatment of the central ethical-religious problems --- the problems
of the freedom of the will, of the possibility and actuality of evil, and of
evil's abolition in reconciliation --- to maintain the life-significance, the
actuality-significance, of the standpoint indicated. From a starting point in
knowledge there is to be sought a solution of these problems which satisfies
ethically and religiously; and the given conception of
% --- p. 63 ---
the fundamental relation of the human is to be made more concrete and deepened
in such a way that, in fulfilling the demand of the knowing spirit, it
fulfils the deepest demand of the moral and religious spirit. While the breach
in the life of the human spirit through guilt is acknowledged, the in truth
absolute relation to the Absolute is to show itself as capable of being
preserved, and the possibility of the true and full reconciliation is to be
maintained --- and maintained as something that can become an object of
cognition. Reconciliation, whose possibility was shown from the human spirit's
universal essential determinateness, is to be maintained also from its most
concrete determinateness in actuality.

\begin{center}\rule{2cm}{0.4pt}\end{center}
% p. 63 is the last page of division II and is a short page: the type block
% ends after eleven lines and a short centred rule closes the division,
% exactly as at p. 29. The rule is printed as a DOUBLE rule and is genuine
% type, not show-through. Set here with the same single \rule used throughout.
% No letterspacing, no footnote and no Greek anywhere on p. 63.

\chapter*{III.\\ The Central Ethical-Religious Problems, Apprehended from the Standpoint of Knowledge.}
\addcontentsline{toc}{chapter}{III. The Central Problems}
\markboth{III. The Central Problems}{III. The Central Problems}

% --- p. 64 ---
% p. 64 opens division III. Order on the page: the bold centred division head
% (already in the skeleton), a short centred rule (not transcribed), an
% ornamental two-line initial (set normally), the introductory paragraph
% below, a second short centred rule (not transcribed), then the bold centred
% sub-head.
% Letterspacing on this page is asymmetric across the parallel phrases:
% „Begrebet om det Absolute“ is spaced throughout, but in the second member
% only „det menneskelige Væsen“ is (the words „og Begrebet om“ are set close).
% Likewise „Det Absolutes Begreb“ has the spacing on „Absolutes“ alone. As
% printed; the English runs fall on the same members.
The presupposition and foundation for the solution of these problems from the
standpoint indicated must be the development of the mutually conditioning
concepts: \emph{the concept of the Absolute}, and the concept of \emph{the
human being}. The determination of the concept of the Absolute must have a
footing in the determination of the human as that which is the highest form of
finite existence, and particularly in the determination of that form of
subjectivity which appears in liberated form in human self-consciousness. The
human then becomes, from this standpoint, \textit{principium cognoscendi} for
the divine. But when the human being is definitively to be determined
according to its peculiarity and its fundamental relations, then the concept
of the Absolute as the human being's \emph{principle} must become the
condition --- so that what as \textit{principium essendi} has the priority
definitively acquires it also as \textit{principium cognoscendi}, since for
thought, which reaches the Absolute from finitude, the Absolute becomes that
which, as resting in itself and having certainty in itself, explains what has
being through it. The concept of the \emph{Absolute} becomes therefore the
fundamental one.
% antiqua Latin on p. 64: principium cognoscendi (twice), principium essendi.
% All three are set in UPRIGHT antiqua at normal letterspacing --- the
% calibration case for the letterspaced antiqua found on p. 68 below. Tight,
% so \textit{} without \emph{}.

\section*{1. The Concept of the Absolute.}

For the concrete determination of the concept of the Absolute, finite
existence as a whole becomes the foundation, since
% --- p. 65 ---
the Absolute is to be apprehended as this existence's principle, and is to be
so determined that it explains the fact of the existence grounded by it, while
it is cognised as that which, as principle, explains itself, in that it is its
own ground and rests in itself.

That the Absolute is apprehended as existence's principle acquires partly the
significance that it is apprehended as the \emph{principle of being} for
existence, i.e. as existence's springing-forth, its ground; and partly the
significance that it is apprehended as existence's \emph{goal}, its law and
its norm. And as existence's τέλος the Absolute is again determined partly as
the teleologically working, partly as the ethically determining (as ethical
ideal). But these different ways of viewing the apprehension of the principle
are to be seen in their inner unity.
% letterspacing broken across the line break: „V æ r e n s -“ ends the line
% spaced, „princip“ opens the next close-set. The English compound divides at
% the same joint.
% Greek on p. 65: τέλος, antiqua italic with accent. Neither OCR witness reads
% it; caught by eye at 300 dpi. Copied verbatim.

Since the Absolute is to be determined as \emph{principle of being}, there
comes forward first the relation between the \emph{ideal and the real} in
existence. The existence which is immediately given for human consciousness,
and which provides the filling for the self's first consciousness of itself in
time, is real, natural existence. This existence is to be explained by the
Absolute, and the springing-forth of its being is to be sought in the
Absolute's determinations. But the determination of reality points back,
beyond itself, to the ideal. If the determination of reality is held fast in
its isolation, and in the form which is necessary for the explanation of
actuality's concrete manifoldness, then it leads us to the concept of
\emph{the atoms}, as the original independent realia which in their
unchangeableness are what lies at the foundation of all existence's
manifoldness. But precisely this concept --- which immediately encloses a
contradiction in that the atom as independent, eternal and unchangeable is
posited in a being absolute and closed within itself, while in its
manifoldness and its finite determinateness by quality it is posited as
something relative --- leads, when the relations into which the atom enters
are thought through, beyond the merely real to something ideal. In order that
the phenomena of existence may be explained, it must be laid down that the
atoms enter into constantly changing connections, and that through mutual
% „Realia“ is Fraktur in the Danish, not antiqua --- the book sets only Latin
% phrases in antiqua, and this naturalised plural is not one. So no \textit{}.
% --- p. 66 ---
communication of motion they become members in that infinite chain of causes
and effects which comprehends all real existence. But in the atom's nature as
closed within itself and unchangeable no impulse to such relations ---
relations enclosing a change --- can be thought to lie; for even if the atoms
are equipped with original forces, the atom's unchangeableness will condition
the unchangeableness of the working of force. The constant change in real
existence, which as a fact is indisputable, must then unavoidably be thought
to have its source outside the atoms as single ones; but this would --- since
at this standpoint the thought of a supranatural factor must fall away --- be
in the totality of the really existent, in nature as a whole, so that the
change would come about through an activity of the whole. But as soon as the
determination of totality is brought in, as the expression of a real unity and
not merely of the aggregate's apparent unity, which could not explain the
fact, then one goes beyond the \emph{merely} real and posits something ideal;
for the totality \emph{as tota}lity is \emph{ideally} determined in that it is
and works in its parts, its moments, \emph{as totality}, and thus relates
itself through them to itself, and in the self-doubling within this relation
--- which is the fundamental form of ideality --- asserts itself. The fact of
change therefore points to the totality with its ideal unity, for which the
atoms --- which as the original principles of real existence had to be posited
as independent and unchangeable --- become serving moments. --- To the same
result we are led also through the consideration of the relation of causality,
which in spite of the atoms' independence and unchangeableness acquires
validity for them, in that the atoms, as communicating motion to one another,
are apprehended as functioning members in a chain of causes and effects, even
though the relation thereby posited becomes as it were a superficial one, not
touching the atom's interior. The concept of cause is then here reduced to the
concept of the exclusively real cause, the thing-cause, the finite cause. The
cause so determined must have the impulse to activity from outside; its
working is cause of another in that it is effect of an
% „kun Reale“: only „kun“ is spaced, „Reale“ is not. „som Tota-“ is spaced at
% the line end and „litet“ is not. Both as printed, and the English divides
% "totality" at the same joint.
% --- p. 67 ---
other: it is cause and effect with respect to something different.

Since now what is acted upon is, in virtue of the conceptually necessary
connection between action and reaction, seen as reacting, the concept of cause
leads to the concept of a \emph{reciprocal action}; but at the standpoint
indicated the relation of reciprocal action, like the relation of cause, must
become a superficial relation, in which the members of the reciprocal action
--- even if the substrates of the activity are represented as in inner unity
with the activity --- are not to be absorbed: they are to preserve a being as
given outside the reciprocal action, are to enter with only one side of their
existence into the determinate relation of reciprocal action. As soon,
however, as a reciprocal action is laid down --- which must necessarily happen
as soon as the concepts of cause and effect are laid down, however far the
origin of the activity remains inexplicable on the presupposition --- it
encloses consequences which unavoidably lead beyond the finite relations and
the finite thing-causes, to an apprehension of the concept of cause answering
to that concept of totality which stepped into the place of the concept of the
atom as the concept of what is primitive in existence. What can enter into a
relation of reciprocal action with another points, by that very relation, to
an essential connection so far as the reciprocal action reaches. But when now
every thing whatever is only what it is in its relations --- when a holding
fast of it as a thing in and for itself, as a relationless thing, makes it a
contentless abstraction --- then the thing must according to its \emph{whole}
essence enter into a comprehensive active relation, into a \emph{total
reciprocal action}. And since reciprocal action thus comprehends the whole of
real existence, there is through the comprehensive reciprocal action a pointer
to a comprehensive connection of essence; and the single thing, the single
real cause, which at every point is at once conditioning and conditioned,
presupposes its own condition, standing in essential connection with it, just
as it is presupposed by that condition as a condition essentially bound up
with it. Since now none of the real things standing in essential connection,
as mutually presupposing one another, becomes \emph{primitive} cause, the
relation of cau\-%
% --- p. 68 ---
sality points back, through the reciprocal action of the conditioned finite
thing-causes, to something which comprehends the essence of them all,
conditions them all and determines their reciprocal action, as primitive
cause. In the relation of subordination to this primitive cause, which works
through them as \emph{indwelling}, the finite real causes become
\emph{secondary} causes, \emph{intermediate causes}, causes in a
\emph{derived} sense. As the principle for the active totality of the
essentially connected reals --- for that totality which, as being in its
moments, is itself --- the primitive cause becomes \emph{ideal} cause; it
becomes cause in relation to itself (\emph{\textit{causa sui}}) and end for
itself (\emph{\textit{causa finalis}}). By its relation to the conditioned
causes the ideal cause asserts itself as the \emph{unconditioned}, i.e.
\emph{that which conditions itself}, that which posits its own conditions: the
system of the real causes. Comprehended by the totality, as its members, these
become, in another sense than in the finite reciprocal action, at once cause
and effect: the series of causes and effects is bent back into itself in such
a cycle that the single members, in relating themselves as moments of the
totality to one another --- not in relation to something different, but in
relation to the same, namely to the totality --- are at once causes and
effects, conditions for the unconditioned, but conditions posited by the
unconditioned, which through them relates itself to itself and is
\emph{principle for its own teleological realisation} in that it is
\emph{principle for all actuality}.
% antiqua Latin on p. 68 --- causa sui, causa finalis --- is itself
% LETTERSPACED, unlike principium cognoscendi/essendi on p. 64, which is set
% close in the same upright antiqua. Encoded \emph{\textit{...}}: \textit{}
% for the antiqua sort, \emph{} for the letterspacing. This is exactly the
% distinction the file header warns is invisible in English, and is why it is
% kept mechanically.

When thus real existence, even when apprehended one-sidedly from the
standpoint of reality, pointed back with necessity to something ideal as that
which unites and grounds it, then its explanation is to be sought in an
\emph{ideally} determined Absolute; but in order that this may be \emph{the
ground of real existence}, the point of attachment for the real must be shown
in the Absolute's own determination of ideality, so that the ideal principle
is seen as an \emph{ideally-real} principle. This connection between the
determination of ideality and that of reality in the Absolute shows itself
first \emph{\textit{in abstracto}}
% PRINTER'S DEFECT, p. 68: the Danish prints „Asolutes“ for „Absolutes“ --- a
% dropped b, verified at 600 dpi and not in the author's Trykfeil list.
% Orthographic, and so not reproducible in English; recorded here.
% „in abstracto“ is antiqua and letterspaced, like causa sui above.
% --- p. 69 ---
through this, that the determination of ideality itself encloses the relation
\emph{to the other}, through which the ideal is in relation to itself. If,
then, the ideal's other, its opposite, is \emph{the real}, then we must posit
this determination --- from which we \emph{cannot abstract}, and which we saw
as pointing back to the ideal --- as \emph{essential to the Absolute itself},
as demanded by it, and as an original determination in it. But when in this
way, through the concept of \emph{otherness}, which encloses the
\emph{necessity of oppo}sition, we come to the determination of the real, then
this opposition between ideal and real is as yet \emph{only formal}; the
real's \emph{concrete determinateness} is not yet explained, and only when
that too is explained by the ideal principle's fundamental determinations is
the principle in actuality the ground of explanation for real existence. The
explanation of the real's essential determinations we find in accentuating
that the ideal principle, as such, in asserting itself as \textit{principium
sui} through its other, is \emph{energy}. The determination of energy, which
lies in the determination of ideality, presupposes, like that determination,
\emph{the other}, through the relation to which energy asserts itself as
energy. This inner opposition of the ideal energy, which is grounded in its
\emph{essence}, is the \emph{passive}; but this determination of the passive,
grounded in energy, does not denote absolute passivity, but a passivity which
has the possibility of activity within it, which is like a slumbering
activity, like an activity which, having gone to rest in passivity, bears
energy within itself as real possibility. A passivity of this kind is demanded
by energy, is the condition of energy's actuality, and is only through energy,
is thought only through it. But such a passivity, such a hidden energy veiled
beneath inertia, is the fundamental stamp of the \emph{material}, of that
which mechanism determines as \emph{dead being, as passive matter}. Dead,
passive in an absolute sense, an existent cannot be; for all being requires,
in order to \emph{assert} itself, \emph{being, working and activity}. But the
concept \emph{passive matter} can acquire its validity in the sense that
activity is present only as possi\-%
% antiqua Latin on p. 69: principium sui, set close, NOT letterspaced --- so
% \textit{} without \emph{}, unlike causa sui on p. 68.
% „Modsætningens Nød-“ is spaced at the line end and „vendighed“ is not; the
% English divides "opposition" at the same joint.
% --- p. 70 ---
bility while its actuality is conditioned by what falls outside the material
so determined. And the concept of the material coincides with that of the
\emph{totality-less}, since totality --- also as \emph{relative totality} ---
by its ideal relation to itself, by its being for itself, points immediately
to the determination of energy, so that the passive can only be that which
lacks totality. But the lack of the ideal relation to itself which lies in the
determination of the totality-less encloses in turn the determination of
\emph{externality}, a being bound by externality's forms, governed by
\emph{the separation of space and time}. As determined by these forms the
material exists \emph{as} material in the separation of the objectively single
elements of matter; and the being of these as isolated is not merely for the
abstracting thought which disregards the conditions of concrete existence and
separates the element from the totality; rather it has its relative validity
in the objective, the real, in virtue of the fundamental relation in the
principle itself.

Thus there shows itself, in the principle's own determination of ideality, a
point of attachment for the determination of reality in existence. How this
determination of reality is now governed by the ideal, is \emph{idealised},
but \emph{without losing its essentiality}, and how \emph{the real} as a
determination \emph{in the Absolute} is related to \emph{the real in
existence}, will emerge from developments in what follows.

At this point the conception put forward is to be more closely elucidated
through comparisons. From those conceptions which from the outset postulate
ideality and reality as primitive determinations in the Absolute, but which
\emph{merely demand} this union, or restrict themselves to substituting the
concept of reality for the otherness demanded by ideality, the conception
indicated distinguishes itself in this: that it deduces the \emph{peculiar}
determinations of the real from the determination of ideality, and that it can
therefore lead to the cognition of a unity in which the real is governed by
the grounding ideal, without being compelled to postulate a contradictory
uniting medium between oppositions originally separated in their peculiarity
--- separated by the lack of a conceptual connection --- and irreconcilable in
their immediate heterogeneity.
% „idealiseres,“ is spaced, the following „men“ is not, and „uden at tabe sin
% Væsentlighed,“ is spaced again. As printed.
% --- p. 71 ---

To the representation of a \emph{creation} the conception indicated stands in
this relation, that it appropriates from that representation the truth
\emph{that the real is grounded in the ideal, the spiritual, and determined by
it;} while it sees what is untrue in the representation of creation in this,
that it does not let the real be of essential significance for the ideal's
being, does not let the ideal Absolute in its unconditionedness condition
itself through the real, and does not acknowledge this real's essential unity
with the ideal. It holds fast that the unconditioned is unconditioned only as
that which conditions itself through a system of real conditions; that the
Absolute is \textit{principium sui}, \textit{causa sui}, in that it is
principle for an other, and that this other, as the principle's other, can be
only that ideal principle's \emph{own} determination.

Already when we started from the side of reality --- and more determinately,
from the abstractly apprehended side of reality in existence --- we were
therefore referred to the determination of ideality in the Absolute, and the
immanent development of this concept led to more concrete determinations. But
these come forward \emph{directly} when existence is seen from the standpoint
of the \emph{highest forms of existence}, from the standpoint of \emph{the
living}, which is inexplicable by the merely real, and more closely: from the
\emph{standpoint of the psychical and the spiritual}. Since life acquires its
highest expression in human life, and human life becomes \emph{self-conscious}
life, self-consciousness --- when its \emph{ground} is to be sought --- points
back to the Absolute's determination as \emph{absolute subjectivity}. To this
determination there has already been a pointer above, through the
determination of the principle as \textit{causa sui}, as the unconditioned
which, in relating itself to itself through the totality of the real
conditions, stands in a relation of freedom to itself. Of this determinateness
of the principle, the determination of \emph{self-consciousness} is an
impress, but in the \emph{form of fini}tude. This determinateness by finitude
of self-consciousness shows itself from various standpoints. It shows itself
in that the determination of self-consciousness encloses and pre\-%
% antiqua Latin on p. 71: principium sui, causa sui (twice) --- all set close,
% NOT letterspaced, unlike p. 68.
% „Ende-lighedens Form“: here the spacing starts only in the SECOND half of
% the divided word --- the mirror image of the usual case. The English divides
% "finitude" so the run can begin at the same point.
% --- p. 72 ---
supposes that of consciousness, and thereby the I's separation from the object
and from the other I's as the limiting other. It shows itself in that
self-consciousness is in its actuality conditioned by the \emph{individual
organism}, and thus becomes a determination of \emph{individuality}. And it
shows itself in that self-consciousness's being-for-itself, since it rests
upon a concrete but individually limited foundation, is an \emph{abstract}
doubling of the self, a being-for-itself in the form of abstract universality.
Self-consciousness, as we know it in \emph{man} --- and that is, as we know it
\emph{at all} --- thus shows itself as the form of limited spiritual
individuality. \emph{Immediately}, therefore, we cannot transfer the
determination of self-consciousness to the Absolute without finitising the
Absolute; the concept must undergo a necessary \emph{metamorphosis}, which
happens in that the finitising determinations are removed; and these are
removed in that the concept of \emph{absolute subjectivity} designates the
Absolute as that which, as freed from self-consciousness's finite limitation,
can be \emph{principle for self-consciousnesses}. We can give the concept of
absolute subjectivity its content by starting from the relation between
reason-determined objective existence and the finite subjective. What in the
latter comes forward in the form of reflection, in the relations of the
relatively abstract determinations of concepts, comes forward in the objective
as the objective connection of the really concrete; what comes forward in the
ideal --- which is according to its essence universal and exists in the form
of individuality --- in the form of the I's abstract universality, has its
concrete counterpart in the objective in the real totality. But in this real
totality's \emph{principle of unity} this objective real, which encloses
within itself the whole of the objective relations, must be reflected in a
form of ideal reflection which differs in a twofold respect from the form in
the human. For since this ideal reflection does not rest upon an
\emph{individual} concrete which, in excluding the rest of the concrete, could
give it only the form of negativity and of relative abstraction,
% --- p. 73 ---
but since it is the form for the principle which bears the whole of the
concrete within itself as its own determination, posited by it, the principle
has in that reflection a being-for-itself in the form of the absolutely
\emph{concrete universality}, and is not merely reproducing in the form of
abstraction but positing, producing, in the form of actuality; it does not
have the opposition within itself as an abstract, formal opposition, and
outside itself as a concrete, real one, but comprehends within itself, in
infinitely concrete being-for-itself, the opposition of otherness unfolded in
the manifoldness of real existence. In the Absolute as absolute subjectivity
we therefore see that power which comprehends, bears and works through all
existence and with it all the finite existences of spirit; in whose infinitely
concrete, energetic being-for-itself the source of self-consciousnesses lies,
and in whose eternal being the infinity of space, of time and of matter is
gathered together into ideality's living, content-rich unity. In this
gathering into ideality's content-rich unity the Absolute is \emph{absolutely
knowing subjectivity}, as a knowing of \emph{itself in all things}.

In the Absolute so determined we seek the explanation of the finite
self-consciousness's essence, taking as our starting point the relation
developed above between the determination of ideality and the determination of
reality. Since existence was apprehended as the \emph{universal totality}
determined by the ideal principle, there lies in this concept the possibility
of \emph{relative totalities} within the universal one, in that the unity of
the universal totality, as articulating itself, posits \emph{particular}
determinations which, precisely as \emph{determinations of totality}, can
\emph{themselves} have the determination of \emph{totality} within them. This
concept of the relative totality, which expresses the determination of the
\emph{organic} and forms the middle term between the concept of the
totality-less, the inorganic, and that of the universal totality, finds its
\emph{most adequate} expression in \emph{human nature}, which by its inner
\emph{universality} has the possibility of becoming the foundation for that
being-for-itself of the ideal which is in \emph{self-consciousness, in the
form of liberated universality}. In virtue of the being's \emph{universality}
% --- p. 74 ---
the individual human soul has the possibility of self-consciousness, and
through it the possibility of a consciousness of existence and of its
principle; but since the human individual is a \emph{relative} totality, the
universality which is for-itself in consciousness acquires the form of
\emph{abstraction}, though of an abstraction \emph{constantly integrable},
precisely because in the relative totality, as an organic member in the
universal one, the universal totality and the remaining members can be present
only in the form of \emph{possibility}, not in the form of unfolded reality,
in which they are only \emph{as} universal totality and as those determinate,
really formed members. Therefore self-consciousness becomes an
\emph{abstract} image of the absolute subjectivity's infinitely concrete
being-for-itself: it has a determinateness by finitude within it which is
explained by the principle but cannot be immediately transferred to the
principle.

Through the determination of \emph{absolute subjectivity} the Absolute shows
itself, as already touched upon, as that which, in reaching out beyond the
objective, relates itself to itself \emph{and rests in itself as the
unconditioned which is its own ground}. And it is in this determination of the
Absolute that the connection is to show itself between what has been wished to
be separated as a \emph{merely logical-metaphysical and an ethical}
determination of it. The ethical determination must have its foundation in the
metaphysical, and comes forward as an ethical-metaphysical determination. With
respect to this, the connection between the \emph{twofold} determination of
\emph{the Absolute as existence's} τέλος acquires its significance, since the
determination of the Absolute as τέλος points partly to the standpoint of the
\emph{organic}, partly to \emph{the standpoint of the ethical}; and precisely
in the \emph{organic} apprehension of existence lies the starting point for
that determination in which \emph{human freedom} is grounded.

In the determination of the Absolute as \emph{absolute subjectivity} and as
\emph{ideal-real principle} there lies enclosed the determination of that
which posits reality as determined, in its real being, by the ideal. Since the
emphasis here falls upon reality's being \emph{as} reality, as being in
% --- p. 75 ---
the fundamental forms of real existence, its determinateness by the ideal
becomes a \emph{teleological} determinateness, and the Absolute is seen from
this standpoint as teleologically working. It is seen as τέλος for that
development in real existence through which existence lets ideality acquire
actuality within it. But in this determination the point of attachment is
given for a higher one. In the teleological conversion of the determination of
reality into that of ideality the concept of the \emph{relative totality}
acquires its significance: it becomes the middle term through which the
conversion, the leading back, is realised. But the concept of the relative
totality first finds its \emph{completion} in the concept of the relative
totality \emph{determined by essential universality}, i.e. in the
\emph{self-conscious} human individual existence of spirit. In this the
teleological working of the ideal becomes \emph{conscious and free working}: a
conscious idealising, in progressive development of thought, of the real that
has been taken up, i.e. a \emph{theoretical} working; and a conscious free
forming, under progressive development of will, of reality by ideality, i.e. a
\emph{practical} working. The determination \emph{with freedom to give
ideality actuality in the real} is the \emph{ethical's abstract} expression,
% ONE run in the Danish („E t h i s k e s  a b s t r a c t e“, with „Udtryk“
% tight), not two: an earlier draft split it and the parity check caught it.
but an expression which does not yet distinguish the ethical from
``practical'' and ``poietic'' activity in general. In order to express the
ethical's \textit{differentia specifica} the more abstract determination must
acquire a more concrete form. Ethical action is essentially a realisation of
the ideal in \emph{the personality's side of reality} --- in the drive
determined by natural energy, which acquires the form of the ideal ---
through a \emph{free resolution}; and the realisation of the ideal outside the
personality is ethical only by being the expression of the personality's inner
acting. Without this connection the ethical becomes something historical, or
something practical in the ordinary sense. Since reality's being determined by
ideality is posited through ethical resolution and action, and since the will,
in order to be able to realise this determining, must go back to its own
infinity and thereby to \emph{its principle}, the Absolute, the
\emph{logical-metaphysical} concept of the Absolute determines itself into an
\emph{ethical-metaphysical} one:
% Greek on p. 75: τέλος, antiqua italic. Latin: differentia specifica,
% antiqua, set close (not letterspaced), so \textit{} without \emph{}.
% --- p. 76 ---
the Absolute is apprehended under the determination \emph{the idea of the
good}.

The \emph{abstract} foundation for this determination of the idea is
\emph{ideality's relation to itself as the goal of its own energy}. But this
abstract foundation acquires a more concrete shape in that the ideal's
relation to itself is seen as a relation \emph{through the real}, and in that
the \emph{real} is apprehended according to its \emph{concrete} essential
determinateness. The Absolute is seen as \emph{the good in the metaphysical}
sense, in that it is seen as that towards which the whole of real existence
strives as towards \emph{the ideal goal} by which this striving itself is; in
that it is seen as that which in essential communication gives to the existent
a being which is being only through the relation to the ideal. It is seen as
\emph{the good in the ethical-metaphysical} sense, in that it is seen as
\emph{the goal and norm of ethical action}, as that which in the form of
individuality is to be realised through free action, as that which, as
something unconditionally valid for the will, is to govern and fill out the
will. That the relation to the \emph{human will} is reflected in the
Absolute's relation to itself is the condition for the Absolute's being seen
as \emph{the idea of the good in ethical} determinateness. This relation comes
forward whether we determine the Absolute as that which is the norm, the law,
for the human will, as that which is to be expressed through it, or whether we
determine it as that which in its eternal self-realisation grounds the human
wills in which it gives itself shape. Given this relation to the will, it lies
near to transfer the will's determination to the divine. It has however been
shown in the earlier work (pp.\ 165 f.) what metamorphosis must take place
with this concept, as with that of knowledge, in order that it may find
application to the eternal and infinite. Just as in the eternal we sought a
ground of self-consciousness, so we seek in it also a ground of will; but we
seek the ground of will in something freed from finitude.

With the determination of will there comes forward the determination of the
Absolute as \emph{absolute personality}. What
% --- p. 77 ---
\setcounter{footnote}{0}%
is sought to be maintained through this determination on the side of
\emph{theology} has already been shown\footnote{Problemet om Tro og Viden,
pp.\ 68 f.}. It is on the one side the separation, in representation, from the
human; on the other side the unboundness by the universal law of necessity.
What is aimed at on the side of \emph{philosophy}, when this concept is
applied, is partly the same; but partly it is also this: to obtain in the
divine an ideal of will answering to the demands of the personal life; and
partly, finally, this: to keep the eternally complete divine personality out
of the world's becoming and development.
% printed mark: *) --- the only note on p. 77, below a short rule. Its text is
% plain Fraktur, neither letterspaced nor in quotation marks.

When there is demanded in the divine an ideal of will answering to the claims
of the personal life, this demand is more closely understood thus: that since
the good is to be realised in the individual's personal life, and the personal
life has the form of singularity, there is to be in the ideal of the good a
corresponding form of singularity; and that the good is to be apprehended as a
divine person with determinations of will which express personality's essence,
heterogeneous with the determinations of knowledge. But in understanding the
demand thus, one has partly brought to bear that abstract opposition between
cognition's determinateness by universality and the will's determinateness by
singularity whose falsity has already been shown; and partly it has been
overlooked that the carrying through of the determination of \emph{absolute}
personality, by making God the \emph{absolutely singular} being, would
precisely annihilate the determination of God as the will's ideal; for the
determinateness of the absolutely singular being can have validity for no
other being than precisely for that being itself. To have recourse to a
combination of the determination of omnipotence with that of the ideal of
will, and to apprehend the divine will as the omnipotently commanding will
whose commands, as omnipotence's commands, were to become laws for the finite
will, could be of no help; for omnipotence's \emph{arbitrary} commands cannot
become the absolute ideal for the will's \emph{essence-determined} relation of
freedom.

In the application of the determination of personality to the divine, both in
theology and in the various philosophical
% „Saavel i Theologiens som i de forskjellige philosophiske“ is set close
% throughout --- unlike the letterspaced „Theologiens“/„Philosophiens“ earlier
% on the page. So no \emph{} here.
% --- p. 78 ---
systems, the necessary \emph{metamorphosis} which that determination must
undergo in being transferred to the Absolute is either wholly forgotten or
carried through only inconsistently. It has already been partly elucidated
through the earlier developments, and in what follows it will be indicated
what significance the determinations of \emph{holiness}, \emph{righteousness}
and \emph{love}, which are posited as predicates of the absolute personality,
can acquire in their application to the Absolute, on the presupposition of
that metamorphosis.
% The three nouns are letterspaced but the „og“ between them is set close;
% checked at 200 dpi. Compare p. 82, where the „og“ inside a run IS spaced.

When, in the determination of the Absolute as personality, a particular
emphasis is laid on this, that by this determination God as the eternally
existent is to be kept out of the becoming and development of natural
existence and of the finite spirit --- which, however, from the nature of the
case, none of the philosophical systems that have accentuated that
determination has been able to carry through --- then it will be shown in what
immediately follows how, starting from the given determination of
\emph{absolute subjectivity}, it will be possible, by laying down the relation
of \emph{immanence}, to free the Absolute from becoming's imperfection and to
maintain it in its absoluteness, while avoiding those contradictions which are
inseparable from the determination of absolute personality. This will emerge
from an investigation of \emph{the essence of pantheism}, of its
\emph{truth} and of its \emph{one-sidedness} in the ordinary conception.
% PRINTER'S DEFECT, p. 78: in „Immanentsforholdet“ a thin full-height vertical
% stroke stands between the two m's --- a risen space or quad from the
% letterspacing, taking ink. Not a sort; not transcribed in either file, and
% logged here. Letterspacing: „Immanents“ is spaced and „forholdet“ is not ---
% the compositor's usual habit of dropping the spacing part-way through a
% compound. The English compound divides at the same joint.

To the name \emph{pantheism} there clings as it were a stain, and the
immediate consciousness, and particularly the religious consciousness, has an
instinctive shyness of it. What frightens one back is, on the one side, the
individual's resistance to being absorbed into the Absolute --- an absorption
which, even when the Absolute is determined as absolute subjectivity, shows
itself to consciousness as a losing of oneself in the

self-less; and on the other side it is this, that the divine, in having to
comprehend all existence, seems to be degraded, to be drawn into the finite's
lowliness, to go together with the transitory, even with evil. But as regards
the \emph{latter},
% --- p. 79 ---
the pantheistic conception has never --- not even where it appeared at the
lower stages of development of spirit, as in the Orient --- been to be
understood in such a way that God should be all, in the sense of the sum
total, the collective sum, of all the finite, single things. When it
apprehended God as all, what was thereby understood was that God was things'
abiding essence, in relation to which the finite was only as vanishing
accidents which eternally alternate while the essence abides. They were to
proceed from the substance and to go back into it, and through their vanishing
they were to show its exaltedness and the untruth of finite being as finite.
And when this essence was to find an expression in the finite, it was the most
glorious thing in each kind of the finite that became the image, the symbol,
for it: the excellent stands nearest to the essence and was therefore used for
its designation. When it was said, then, that God is all, what was meant was
not: God is the finite things; for these have, as finite, as accidents, no
being: in their transitoriness their essence is revealed, and the substance
shows itself as what abides in them. But by this --- God is all --- it was
expressed that God is existence's universal essence, the absolutely universal
which is indwelling in the single existences, but in such a way that to these
as single no real reality belongs. Even at that standpoint, then, where spirit
did not yet know how to distinguish its own essence from nature's, pantheism
was not a doctrine of God's identity with the sum of the finite, but the
doctrine of God's identity with the universal power of nature, over against
which no finite thing has subsistence, no imperfect thing a right to be called
existent. God is therefore not finitised, not drawn down into what is lowly;
rather the finite is abolished over against the divine: the world, not God,
vanishes. And what holds when God is apprehended as the universal power of
nature, the substance of nature, holds also of the higher forms of pantheism
that have appeared within the domain of philosophy, which determine God as
\emph{subject}, as \emph{spirit}, but which let the finite spirits be taken up
and abolished as moments in his essence as absolute spirit. In these
% --- p. 80 ---
forms too --- whose originators, by an \emph{arbitrary} limitation of
pantheism's concept, seek to keep them outside it --- God is not drawn down
into the finite, but the finite vanishes in God. And as regards God's relation
to \emph{evil}, it will be attempted in what follows to show that, even though
the historically given pantheistic systems have remained standing at the
dilemma of either abolishing evil's actuality or making it a determination in
the divine, there is a possibility of holding fast evil's actuality
\emph{within a pantheistic view}, without letting it abolish the determination
of the Absolute as the ideal of the good.

As regards the \emph{first} of the motives emphasised above for a shyness of
pantheism --- the fear of letting the single being be absorbed into the
divine's unity --- it shows itself, on closer consideration, that this shyness
is conditioned not by anything universal and pervasive in human nature, but by
something peculiar, developed under determinate historical conditions. When
the Indian longs for liberation from the barriers of single existence, and in
order to attain it lets himself be crushed under Siva's car, or swallowed by
the waves of the Ganges, or sinks his soul into the formless nothing, then
that shyness is not present. When Persian poets develop a pantheistic
mysticism, then self-surrender to the infinite One is a source of enthusiasm.
And if we go from the Orient to Greece, we find, among that people whose
peculiarity was a stamping out of the beautiful individuality, that virtually
all the greatest representatives of its own peculiar thinking --- including
such as whose systems do not have pantheism's character --- let the individual
vanish, without loss and without craving, over against the eternal divine
thought which animates the universe. When the Eleatics posit the One ---
whether it be determined as God or as the existent --- as existence's
wholeness, they find precisely rest and thought's repose in this thought, and
Parmenides becomes a poet in uttering it. When Aristotle seeks an immortality
for the human spirit, it is only for that
% --- p. 81 ---
highest form of spirit which is in individuality-less unity with the eternal
divine thought resting in a νόησις τῆς νοήσεως. When the Stoics in their
striving after unity let the Aristotelian Νοῦς become the world's indwelling
force, and let this rational force go together with matter, then within this
pantheistic unity they let individuals stand governed by a world-law which
they willingly follow and under which they bow with resignation's repose. And
when the Neoplatonists strive towards ecstatic unity with the divine One, then
so far is it from being the case that the abolition of individual existence by
this unity has anything deterrent for them, that it is on the contrary the
goal of the whole life's burning longing. And within Christianity the same
holds of the mystics' suffering of God, of their striving after a union in
which the soul ``makes itself simple'' so as to have no determination
separating it from the divine.
% GREEK on p. 81, two instances, both antiqua italic with accents, both read
% off the image at 600 dpi; neither OCR witness reads either. Neither is
% letterspaced. p. 81 was on no list of Greek pages. Copied verbatim.

It is essentially in the modern age that we find that longing after the
preservation of individuality --- often appearing in sickly forms --- which
has the instinctive abhorrence of pantheism for its consequence; and this
longing has, in the immediate consciousness, its ground in that emphasis on
subjectivity peculiar to the modern age, which from the first assumes that
form, in that the spiritualistically determined individual is apprehended as
infinite in its being torn loose from its natural conditions. And once such an
emphasis has been laid on the preservation of individuality, the craving after
it comes forward most strongly precisely in those times and among those
peoples where the individual --- its finitude being infinitised by egoism's
interest --- becomes the centre of existence, and where the infinite is pushed
aside as incomprehensible or indifferent. The period of enlightenment was thus
hostile to pantheism, and among the Romance nations the antipathy to it has
been strongest.

When thus that instinctive shyness of pantheism has shown itself to be
grounded not in a necessity in human nature but in a determinate development
conditioned by peculiar historical presuppositions, then we can, unaf\-%
% --- p. 82 ---
fected by it, investigate pantheism's \emph{essence and truth}. And when we do
not limit pantheism's concept to the \emph{imperfect and one-sided} forms of
it, then we must acknowledge the pantheistic conception as \emph{essential and
necessary for philosophy}. A pantheistic element there must be in all
philosophy, because no philosophy can think the infinite as being separated
from and outside the finite, but --- if the infinite is to preserve its
infinity --- must think the finite as being in it and as having truth only in
it; and because all philosophy seeks cognition's unity and connection and is
philosophy only so far as it pursues that unity. When thought sets about the
task of finding reason in existence, then everything to which thought attains
must form a unity, a connection in which no breach of immanence is found. For
every higher thought to which one ascends is precisely the higher only in
this, that it comprehends the lower ones, that it, as the richer and fuller,
has them within itself as its moments. And what holds of every higher thought
holds also of the highest: it cannot for thinking separate itself from the
thoughts which are enclosed in it as its moments and which find their final
grounding in its comprehensive wholeness. But this highest thought, the
thought of the idea, is precisely thought's expression for the Absolute, for
God. The Absolute, then, which comprehends and bears within itself the single
thoughts, the thoughts of rational existence, must be thought in essential
unity with that existence itself. In finding reason in existence, and in
having the Absolute in reason, thinking seeks nothing outside the universe,
outside the unity of existence, the corporeal and the spiritual. In this sense,
then, all philosophy is pantheistic: it must hold fast the Absolute's
immanence; it cannot think the Absolute outside that existence whose principle
it is; it cannot think an abstract transcendence's separation of the divine,
for thereby infinity would be lost. By its own nature thought is not led away
from immanence's unity; but it can, through consequences of one-sided
presuppositions (for example
% Emphasis on p. 82: three runs in which the connecting „og“ / „for“ ARE
% letterspaced with the rest --- unlike p. 78, where the „og“ between spaced
% nouns is set close. As printed in both places, and so in the English.
% „Væsen og Sandhed“ breaks across the line at „Sand-“/„hed.“, and the
% transcription flags that the continuation may in fact be tight, as at pp. 6
% and 15, while setting it as one emphasised phrase. The English follows the
% transcription and inherits the flag: if a later pass reads it as
% „\emph{Væsen og Sand}hed“, this run splits at the same joint.
% --- p. 83 ---
Aristotle, Descartes), and through foreign influences, be led to break the
unity; and within the pantheistic conception the unity of thoughts can take a
different form. But all real philosophy is as such pantheistic in the sense
indicated, as holding fast the Absolute's immanence and the untruth of finite
being when, as finite being, it is separated from the abstractly transcendent
divine.

In order to determine pantheism's \emph{true} form, we single out the main
forms of a \emph{one-sided} pantheistic conception whose possibility is given
in the relations of the concepts.

As the \emph{first} one-sided form we bring out that in which the fundamental
concept is the concept of \emph{the power of nature}, more closely apprehended
as the mighty one substance of nature in relation to the merely vanishing
accidents. In this form of pantheism the one-sidedness comes to light
\emph{immediately}, in that the existence of spirit, the human
self-consciousness, is to be set in relation to a substance which does not
explain it, and in this relation to come under the determination of the
vanishing accident; and \emph{mediately} the one-sidedness shows itself in
that the conception of the accidental as the merely vanishing, the powerless,
reacts upon the conception of the Absolute, the substance, which --- when it
posits only something vanishing and reveals itself only as power over the
merely powerless --- itself becomes the empty and the impotent.

As the \emph{second} one-sided form of pantheism we apprehend that in which the
Absolute acquires the determination of the substance's infinite unity, which
binds together within itself thought and extension. Here, indeed, through the
determination of thought, the determination of spirit has been brought into
the substance, the Absolute; but since the substance in its unity is to unite
the \emph{merely disparate} attributes, it becomes \emph{as} substance the
purely determinationless, that in whose differenceless unity everything merely
vanishes; and its accidents, the modifications, come to occupy the same
position towards it as the accidents of that substance of nature. The form of
self-consciousness is not so clarified in the single thinking \textit{modus}
that it could, reacting back, determine the concept of the Absolute;
self-consciousness becomes
% ANTIQUA LATIN, p. 83: „modus“ is set in upright antiqua against the Fraktur
% and is NOT letterspaced. \textit{} without \emph{}.
% „som“ in „bliver den som Substants“ is letterspaced; so is „som“ on p. 84 and
% on p. 88. The book emphasises this qua-„som“ repeatedly, and the English
% carries the run on "as" each time.
% --- p. 84 ---
in reality only a postulated accessory to natural causality, as it were the
expression of natural causality's sensing of itself.
% „Accessorium“ and „Imponderabile“ (p. 88) are set in FRAKTUR, not antiqua, so
% neither takes \textit{}. Only „modus“ (p. 83), „eo ipso“ (p. 89) and
% „implicite“ (p. 91) are antiqua in this stretch.

As the \emph{third} one-sided form we apprehend, finally, that in which the
Absolute is determined as \emph{subjectivity}, but \emph{in such a way} that
in this subjectivity's eternal process nature --- whose mode of being is the
untrue --- becomes an eternally abolished moment, and the finite spirits ---
whose individuality is apprehended as resting only upon a vanishing ground of
nature --- become only vanishing forms in the absolute unity's eternal
movement.

For the one-sidedness in these forms the corrective is to be sought in the
\emph{concrete} determination of the Absolute as \emph{absolute subjectivity},
and in such a way that the corrective is applied \emph{within} the pantheistic
conception.

The task becomes, with respect to the \emph{Absolute}: to maintain a moment of
transcendence \emph{within} immanence; and with respect to \emph{that which is
grounded by the Absolute}: to secure to real existence and to the finite
spirit a relative independence.
% „Transcendentsens“ here has ONE s, and „Transscendentsen“ two lines below has
% two; both readings confirmed at 400 dpi. The book varies; not normalised, and
% the variation has no English counterpart.

This twofold task is solved when the consequences are drawn of a distinction of
concepts which in the foregoing has only been touched upon, because what
mattered there was above all to maintain the validity of the ideal at all ---
of that determination which is common to the totality and to the principle by
which the totality is constituted. This distinction of concepts is the
distinction between \emph{the totality of existence and the principle of
existence}.

Between the principle determined as ideal-real and the whole of existence
grounded by it there appears that distinction by which a moment of
transcendence --- in the sense in which transcendence can have validity for
philosophy at all --- is brought in, the emphasis falling upon \emph{the moment
of otherness}: the real, and upon \emph{the real's form of being}. Since, as
was shown above, the Absolute \emph{as} ideal encloses in its concept the
determination of \emph{otherness} --- as energy, energy's opposite: the passive
--- it must be comprehended as that which in its eternal positing of itself
posits what is characterised by the determinations of otherness and of
passivity,
% --- p. 85 ---
and posits it as being in reality's \emph{essentially valid},
\emph{conceptually necessary for}ms, while it is at the same time thought as
that which, in this relation, is eternally in a relation of unity to itself.
There is therefore posited at one and the same time a twofold relation: a
relation of \emph{eternity} and a relation of \emph{time}, a \emph{gathering
of reality into the concrete ideal's unity} and an \emph{existence of reality
in the separation of externality}; and by this twofold relation, by this
twofold determination of reality, the possibility of solving the task set
above is conditioned.
% Three partial letterspacings on this page, all as printed: „begrebsnødvendige
% For“ is spaced and „mer,“ on the next line is not; „Evigheds“ and „Tids“ are
% spaced and „forhold“ close in both compounds. The English divides "forms" and
% hyphenates the two compounds so each run falls where the Danish one does.
% The „og“ joining the two long runs is set close here, unlike p. 82.

The separation of externality which, as essential to the real, is inseparable
from existence, makes itself felt as a pervasive determination in \emph{all}
forms of the real single existences. The primary single existences, the real
fundamental parts, become --- as determined by the opposition of ideality and
energy --- materially separated elements, self-less (totality-less) and
passive in the sense indicated above. And since within real existence the
ideal's determining reveals itself in this, that its material elements are
joined together into the relative totalities of the organic, in which a
principle of life with the determination of energy and of ideality is
realised, the material foundation for the realisation of this principle of
life preserves, under its determinateness by the ideal, its own essential mode
of being.

In the material elements, and in the relative totalities realised in temporal
determinateness and spatial separation upon a material foundation, the
separation of finitude in real existence comes into its right: by it the
members of real existence acquire their relative independence and can assert
their individual existence. The relative independence is conditioned by the
form of existence, by this, that each single member has stepped out into
separateness and works out from this separated being. Since all being is
working, all individual being --- and this has actuality only through the
separation of reality --- becomes a working out from an individual point of
action; and such a working is what lies abstractly at the foundation of
% --- p. 86 ---
the determination of freedom, and also of the concept of freedom in the
pregnant sense: of human freedom, of the will's freedom, which is the
expression of the universal being's ideality and is an energy that can with
consciousness be asserted against the Absolute itself.

But when the separation of finitude thus comes into its right, and the real
members' relative independence makes itself felt, then, in virtue of the
original relation of essence, there is through the separation a pointing back
to the unity, through the relative independence to the being worked through by
the universal principle.

The material elements of matter, whose being expresses the greatest opposition
to the principle's unity and ideality, lack in their self-less being energy's
actuality and originality and are determined by inertia, so that they express
the activity slumbering in them only in relation to an other and in virtue of
an impulse from outside which wakes the slumbering activity into life. But
this relation to the \emph{homogeneous} thing which wakes the activity, and
which itself in its singularity does not have the source of activity within
it, points back to \emph{the aggregate of the material elements}, which, as
already developed earlier, is to be determined as a \emph{totality} which as
such has the determination of energy and the conditions for the continuance of
activity within it.

The single members in real existence give their individually determined being
expression in an individually determined working; but this individual being is
drawn into the whole's connection and determined out of it, and the single
member's working, which is actualised by the working of the other members,
becomes enclosed by the whole's necessity. The single starting point for an
individual working therefore becomes a point of passage for a universal
working, and its working is converted, through the effect of totality
resulting from all the single individual activities, into the working of the
principle of totality. Since what results from the working of the single
acquires the determination of \emph{activity of totality}, this determination
expresses that what results is, according to
% --- p. 87 ---
the relation of essence, what grounds, the primitive; the end by which the
working of the single is determined.

When the infinity of time and of space is the form for the being of real
existence, then the relation of eternity to itself is the form for the being
of the ideal principle, and in the principle's being-for-itself the real, the
material side of the whole of existence, is idealised. But the principle's
ideality would become abstract if the determination of eternity did not
reflect within itself the determination of time, and thereby ideality reflect
the real's form of actuality. In that cycle of eternity in which the real
proceeds from the principle in order to give the principle's ideality fullness
and actuality, time and the determination of reality generally are so
reflected that reality preserves its significance as actuality. The real in
God is an actuality, not a fantasy-image of a`` Physis in God'', and in the
% PRINTER'S DEFECT, p. 87: the opening low quote is set TIGHT against the
% preceding word and LOOSE from the quoted matter --- „af en„ Physis i Gud““.
% Both OCR witnesses reproduce the fault. Reproduced here, the opening mark
% tight to "a" and spaced from "Physis". The quoted matter is Fraktur and is
% NOT letterspaced; the quote balance is unaffected.
eternal unity of real existence in God there is taken up as a moment that
unity in natural existence by which natural existence at \emph{every moment}
reveals its indwelling energy. The real, which in each element of time is
joined together into a whole and in that whole ideally determined, is
converted through this intermediate determination, in each element of time and
with reality preserved, into the absolute ideal. It is precisely upon this
conversion in each element of time (time-differential) that the emphasis here
falls. The totality through which it takes place becomes the intermediate
determination between the real's being in the form of externality and its
eternal unity in the Absolute: that intermediate determination through which
this eternal unity acquires concretion by the reflex of the determinateness by
externality which is preserved in it.

The relation in the Absolute between the real --- which even as idealised must
not be annihilated but must retain its essentiality --- and the ideal can
receive an elucidation through analogies from the finite, through
psychological and physiological facts. In the ensouled organism spatial
extension is idealised, in that in the unity of self-feeling the externality
of the points of space is abolished; an affection of a single part of the
organism is sensed in the ideally concentrated self; the physical,
spatially determined impression is trans\-%
% --- p. 88 ---
formed into an ideal moment in the self-sensing. But the organism's being in
space is not thereby annihilated, and the idealisable space subsists \emph{as}
space for sensing and sensation. --- Matter is idealised in the organism, in
that the sensing organism in its normal state is for itself an imponderabile;
but the organism does not thereby cease to be material, to be subject to the
laws of gravity. --- Consciousness binds the different elements of
representation together into a unity of representation, to which there
corresponds a unity of time out of the elements of time, in which the
externality of succession is abolished; although of the elements of
representation --- which, as bound to a nerve-act, each require its own
element of time --- only one can at any moment be at consciousness's focus, so
that consciousness's affection by the single elements must have the form of
succession. But while time is thus gathered into a unity of time, and in this
gathering --- in which succession is concentrated and thereby idealised ---
loses its externality, it continues to have its real validity for
consciousness. --- Thus these facts point to a preservation of the real's
fundamental determinations while these are idealised. And when we now, with
these facts as points of attachment, transfer by an inference from analogy
what holds of the individual to the absolute principle of the totality of
existence, then --- apprehending this principle as that which according to its
concept posits the real, which is its \emph{own} essential determination, as
something \emph{essential and valid}, as a system of existence being in
\emph{essentially valid forms of existence}, and at the same time seeing it as
that which, in asserting its ideality, joins this real system of existence
together into that content-rich unity of ideality by which the system is
grounded --- we shall obtain space, time and matter, the real's fundamental
determinations, as essentially valid for the Absolute, but at the same time as
so idealised that a living presentness of eternity concentrates within itself
the infinity of time, and that the Absolute's unity of life concentrates
within itself, in a point of ideality and subjectivity, the externality of
boundless space and the infinity of matter.
% --- p. 89 ---
The conception we have brought to bear of the Absolute's relation to real
existence --- existing with relative independence in its essentially valid
forms --- encloses, in the conceptual distinction between the absolute
principle and the whole of existence, and in the determination of the
ideal-real principle's relation of eternity to itself, a \emph{moment of
transcendence}; but in such a way that no breach of immanence is thereby made,
for the connection of essence is preserved, and existence is apprehended as a
totality only in virtue of the principle's indwelling. In virtue of the
original relation of essence to real existence, the Absolute, when it is
determined as \emph{spirit}, is \textit{eo ipso} to be determined as
\emph{principle of nature} (\emph{idea of nature}) --- by which expression it
is designated as the principle for natural existence's ideal unity, for its
universal totality.
% ANTIQUA LATIN, p. 89: „eo ipso“ is upright antiqua and is NOT letterspaced ---
% compare the letterspaced Fraktur „Aand“ immediately before it, on the same
% line. So this book's Latin is not uniformly letterspaced: cf. the
% letterspaced Latin at p. 68 against the close-set Latin at pp. 64, 71, 83
% and 91. The letterspacing of „Naturprincip, (Naturidee)“ runs through the
% comma and the parentheses; the emphasis is set on the two words only.

By the Absolute's relation of eternity to itself the Absolute is, under its
connection of essence with real existence in reality's fundamental forms,
withdrawn from becoming. Its relation to the single finite thing --- which,
seen from the finite thing's standpoint, is in temporal determinateness ---
becomes, in being taken up as a moment in the Absolute's relation of eternity,
an ideal-real determination in the concrete eternity, and as such freed from
externality. And with the Absolute's freedom from becoming its freedom from
finitising limitation stands in connection. By the idealisation of the real,
which precisely confirms the real in its essential validity as real, the
Absolute's infinity is maintained as concrete infinity; and in its relation of
transcendence, enclosed by the relation of immanence, that limitation is
removed which is not avoided in the ordinary consciousness's conception
resting upon the \emph{representation} of transcendence, nor in those
philosophical attempts which under the name of logical \emph{theism} conceal
theological \emph{theism}. --- In what follows, through the development of the
problem of evil, the Absolute's freedom from finitude's imperfection in spite
of immanence will be elucidated from a new standpoint.
% „ideal-real“: the hyphen is printed with a space on each side („ideal - real“),
% where „ideelt-realt“ (p. 84) and „idealt-reale“ six lines above are set
% tight. Neither word is letterspaced --- only the hyphen is loose. The Danish
% transcribes it as a plain compound and logs the spaced hyphen; the English
% does the same.

\begin{center}\rule{2cm}{0.4pt}\end{center}
% p. 89 is the last page of the sub-section „1. Det Absolutes Begreb“ and is a
% short page: the type block ends and a short centred rule closes the
% sub-section, as at pp. 29 and 63. Unlike p. 29 this one is a SINGLE rule.

% --- p. 90 ---
% The printed head is bold Fraktur on two lines, the „2.“ standing to the left
% of the first line. The printed head reads „Det menneskelige Væsens Begreb“;
% the book's own Indhold reads „De menneskelige Væsens Begreb“ --- one of the
% eight head/Indhold divergences in this book. The divergence is a Danish
% article form with no English counterpart; it is recorded here and at the
% Contents, and not manufactured in the English.
\section*{2. The Concept of the Human Being, Determined in Relation to the Absolute and to the Totality of Existence.}

After the concept of the Absolute has been developed in its connection with
the concept of the whole of existence, the task becomes to determine the
concept of the human being in connection with those developments of concepts.
In the determination of the concept of the human being the determination of
the human \emph{task of life} is enclosed, and from it and from those
determinations of concepts the attempt at a solution of the decisive
ethical-religious problems must proceed.

The \emph{universal} concept which has been won in the foregoing for the
determination of the human being is the concept of the \emph{relative
totality}: the concept which designates the living, the organic, in general,
and which distinguishes it on the one side from the concept of the
totality-less elementary, on the other side from the concept of the universal
totality constituted by the principle, the whole of real existence. Within the
universal concept of the relative totality the human being then stands as a
particular determination over against the lower forms of the living: over
against the plant and the animal. Plant life has for its fundamental
determination the determination of \emph{selfhood}, the \emph{abstract
universal determination} of the organic, which indicates ideality, the relation
of the principle of totality to itself through its other; but it lacks the
concentration in a \emph{psychical self}. Such a concentration is attained in
animal life; but animal life's highest determination is that
\emph{self-sensing} in which the individually-universal self is constantly
only for itself in a determinateness of singularity, in immediate unity with
that single determination, and by this relation betrays the limitedness of its
essence. Only in human life is \emph{self-consciousness} the fundamental
determination: the self which according to its ideal essence is universal is
as being for itself in the form of \emph{liberated} universality; and this
determination of self-consciousness is not an accessory determination, a
determination of property in a really presupposed subject --- or rather,
substratum --- but \emph{an essential determination in} the human. It
expresses the ideal fundamental essence itself, which
% A small ink speck stands between the letterspaced „i“ and „det“ in
% „Væsensbestemmelse i det Menneskelige“; it is not a sort and is not
% transcribed. The letterspacing stops after „i“; „det Menneskelige“ is
% close-set, and the English run ends at the same point.
% --- p. 91 ---
is for itself in its determinateness by reality. But in order that
self-consciousness may become possible, the particular essential determination
which characterises animal life must have given place to that of \emph{essential
universality}; only through it does the universal's being for itself \emph{as
universal}, which is self-consciousness's fundamental form, become possible,
and only by it do \emph{thought and will} become possible. And since
self-consciousness thus points back to essential universality, it points back
to \emph{the being's infinity}: only the being which according to its concept
is \emph{infinite} is a \emph{self-conscious} being.

This connection between the determination of \emph{self-consciousness} and the
\emph{infinite} was first indicated by \emph{Leibnitz}, in that he brought
self-consciousness into connection with the cognition of the ``eternal
truths'' which have their springing-forth from the divine. Only through the
possibility of cognising these eternal truths, with their necessity and
universality, does the self's perception have the possibility of
reflexibility; only by it is self-consciousness, in other words, possible. For
in that cognition, in which the object bears the mark of infinity,
consciousness does not relate itself to anything foreign, is not bound by the
particular, but relates itself in the object to itself, through the universal
to its own essence as determined by the universal.

Since self-consciousness is man's \emph{essential expression}, the form for
the being's self-manifestness, self-consciousness encloses a relation to the
human being's \emph{ground}; and since the being which makes
self-consciousness possible is the \emph{universal} being, that ground of the
human being to which self-consciousness has a relation is immediately
determined as \emph{existence's universal ground} (principle); and through the
relation to it a relation is posited to \emph{the totality of real existence}.
In the human spirit's relation of self-consciousness to itself there is
therefore enclosed \textit{implicite} the twofold relation: to the Absolute
and to that \emph{totality of existence} which is grounded in the Absolute.
This twofold relation must come forward for consciousness in the unfolding of
self-consciousness into concrete self-cog\-%
% ANTIQUA LATIN, p. 91: „implicite“ is upright antiqua against the Fraktur and
% is not letterspaced. A third close-set Latin word, after „modus“ (p. 83) and
% „eo ipso“ (p. 89).
% --- p. 92 ---
nition; it must show itself in concrete unity with the relation of
self-consciousness where abstract self-consciousness, through comprehending
cognition, has been unfolded into a reason-determined cognition of that which
lies as inner possibility in the ideally all-comprehending human being. ---

From the concept of self-consciousness --- as the essential expression of the
human, and as enclosing the self's relation to the Absolute and to the whole
of existence --- the \emph{human task of life} must be determined. If we look
to the kind of being which self-consciousness reveals, then the task in its
universality is to be determined by that \emph{unity of infinity and
individual being} which is expressed in self-consciousness, and it becomes
thus: \emph{a realisation of the being's infinity in the form of the given
individuality}. And since the human individual, as originally endowed for
self-consciousness and self-determination, is endowed for personality, the
task can in other words be determined as \emph{personality's realisation
through the unfolding of the being's inner infinity}. But personality must
then not be apprehended as the expression of a single side or a single form of
activity of the human being, but as \emph{the expression of unity} for the
\emph{total} human being in the aggregate of its fundamental functions.
% EMPHASIS, p. 92: the run breaks mid-phrase --- „Enhedsudtrykket“ is spaced
% across its line break (both halves), „for det“ is close-set, „totale“ spaced
% again, and the rest close-set. As printed, not regularised.

\emph{The universal} task must, as a task for a self which reveals itself in a
plurality of primitive forms of expression, articulate itself into a plurality
of essentially connected \emph{particular} tasks; but the universal expression
for the task must preserve its validity for these.

Thus the \emph{task of human cognition} becomes: infinity's realisation in
that cognition of existence and of its principle which is bound into unity
with self-cognition. By such a cognition --- in which the self ideally takes
up the reason-determined, infinite, real existence and joins itself in unity
with its own and existence's comprehended principle --- the self's inner
infinity is realised, and in the relation of infinity to itself freedom is
realised.

% --- p. 93 ---
The task of the human \emph{will} becomes the self's ethical willing of itself
according to its concrete infinity. In ethical willing the self asserts itself
according to its infinity and makes its freedom felt, in that it is
determining out of its own true infinity. But the condition for such a willing
is that no dualistic, finitising separation divides the will from spirit's
other forms of activity; for only so can the will will the being according to
its total determinateness. Just as cognition on its side must preserve the
character of universality, so that it does not, in relation to anything
belonging to spirit, become a merely mirroring hiding-place, but can really
make spirit's content transparent to spirit, so must the will too preserve this
character of universality, and be able to have all spirit's essential forms of
activity for its content. And this means, in other words: the ethical task
must become an all-comprehending task, must enclose human life as a totality,
so that the whole of human life, all the expressions of spirit, come under
ethical responsibility.

To these determinations of the task we come in taking our starting point from
the \emph{infinity of essence} revealed in self-consciousness. But when we now
accentuate the original \emph{relation to the principle and to existence}
which lies in self-consciousness, this leads us to a more concrete expression
for the task.

In determining cognition's task we reached the concept of the unfolded
self-consciousness's unity with the cognition of existence and of the
principle. In this cognition the self becomes conscious of itself \emph{in its
ordering as an organic member into the existence governed by the concretely
determined principle}; and since the self --- which through consciousness's
self-deepening finds the Absolute as that in which its self-consciousness is
first fulfilled --- has in \emph{self-surrender to the Absolute} the
consciousness of \emph{essential unity} with the Absolute, the Absolute's
essential determination acquires validity for the self and for the totality
posited by the principle. But the principle's essential determination was
\emph{the}
% --- p. 94 ---
\emph{concrete unity of the ideal and the real, which excluded all dualistic
separation}: this determination therefore acquires, since the relation of
immanence is not broken, its validity for that which is through the principle.
% EMPHASIS across the p. 93/94 joint: the spaced run begins with „det“ as the
% last word of p. 93 and runs to „Sondring“ on p. 94; the colon is close-set.
% Set as two \emph{} groups, as in the Danish, so the page marker can stand
% between them.

When, then, the self is that which is divided by no dualism and which reveals
itself in the essential unity of its fundamental functions, and when it is
ordered as an integrating member into the ideally-real organisation of
existence, then there lies in this that its task must become to develop the
undivided personality in its inner harmony, to develop it according to the
being's completeness as naturally-spiritual, in inseparable unity of the ideal
and the real, and to develop it as a peculiar member in the whole of existence
determined by the principle, by penetrating its peculiarity with the universal
and making it into a peculiar expression of the universally valid. But since
the human self becomes conscious of itself in its being ordered into the
totality of existence, it becomes at the same time conscious, in virtue of its
essence's universality, of its \emph{relative independence as an integrating
member}; and its task then becomes to will itself in this independence and to
assert that independence through that harmoniously complete development of its
essence's peculiarity. Man must will his essence as this concrete thing, in
which the real foundation of nature is not to be negated but taken up into
spirit, formed through with consciousness, so that the will acquires content
and concretion through the individual's determinateness by reality, while it
has what regulates it in the a priori, which in the domain of the ethical
acquires the same validity as in that of the intellectual. But since the self
wills itself \emph{ethically} in its individual peculiarity and relative
independence, this willing is in unity with its willing of itself as a member
in the comprehensive whole, as subordinated to the principle; and by its
willing it thus makes a movement of infinity, eternalises its acting while
that acting preserves time and actuality as its stage. ---

So apprehended, the task of life becomes a \emph{task of infinity}, in that
the self's willing is a willing of itself according to its essence's infinity,
according to its relation of infin\-%
% --- p. 95 ---
ity to existence and to existence's absolute principle, according to its
significance as an ideal moment in spirit's eternal unity of life. But since
the task of life becomes a task of infinity, a task of eternity, it acquires
the closer determination of a task of infinity which is to be solved in
finitude, of a task of eternity which is to be solved in time, under conditions
of existence, conditions of finitude. Hereby an emphasis falls upon the
\emph{ethical choice as the self's choice in a moment of time}, and thereby
upon the \emph{individual free will's choice} as distinct from the eternal
determining of metaphysical necessity; and an emphasis falls upon the
significance which time acquires for the solution of the task of eternity and
for the failure of that solution, upon the discord in the self which arises
through the self's twofold relation --- to the eternal and to the conditions
of time --- and upon the guilt which the self in this relation acknowledges
itself to have incurred.

From this \emph{standpoint} the determination of \emph{guilt} will \emph{not}
be volatilised into a \emph{metaphysical} moment; it will become an \emph{in
truth ethical} determination with decisive significance as \emph{actuality};
and, as \emph{the ideal's de}mand is maintained, it will be seen as a
\emph{universal} determination in the human as conditioned by finitude. The
ethical individual, which acknowledges the \emph{infinite} demand, must bow
under the feeling of \emph{guilt}, must acknowledge the necessity of
\emph{repentance} and the need of a \emph{supplementary assistance} in order
that guilt may be reconciled, restitution take place. But since the whole of
existence is seen as borne by the immanent principle, and its \emph{harmony as
a whole} can therefore not be thought as disturbed --- just as neither can the
single existence's \emph{essence} be thought as transformed --- the individual
must not seek the supplementary help outside existence, but must seek the
reconciling powers within this existence: in the other supplementary single
beings, with which the individual joins itself in love, in mutual
communication of essence; and, definitively, in that principle which is
indwelling in the individual itself and works through it, and which joins it
together with the other individuals in the being's unity. It will be shown in
what follows how
% EMPHASIS, p. 95: two more mid-word breaks in the letterspacing ---
% „Virkeligheds“ is spaced and „betydning“ close-set in the same word; and
% „Idealets For-“ is spaced while „dring“, after the line break, is close-set.
% The English hyphenates "de-mand" so the run can stop at the same point.
% --- p. 96 ---
redintegration is definitively to be sought through the absolute self-surrender
--- to this principle working through it --- of the self actualised according
to its real ground in its resistance to the ethical; and how through
redintegration the ethicising of the whole of life becomes possible.

Since the task of life is so determined that the personality in its free,
harmonious self-elaboration is constantly to live itself into the whole and
thus realise the relation to the whole's principle, the same thing is hereby
expressed which has been maintained earlier: \emph{that the religious relation
becomes the concluding one}, in which the particular relations first find
their completion. That the religious relation must become that through which
the solution of the task of life is accomplished lies in the human being's
fundamental relation; for the human being is grounded in the divine, and can
as a self-conscious being come to realisation only in so far as this relation,
by which the being is, comes to be for the being's consciousness, and comes to
be for it in such a way that the God-relation goes in unity with the self's
relation to itself in self-consciousness.
% PRINTER'S DEFECT, p. 96: the Danish prints „selvbevist“ for „selvbevidst“ ---
% the d simply absent. Orthographic, and so not reproducible in English.
% Note also „Livshelhedens“ (single e) four lines above, where the book
% elsewhere sets „Heelhed“. As printed.

But when the religious has thus once more shown itself as the concluding
thing, then it is clear from earlier developments that the religiousness which
rests upon the ground of knowledge and of the ethical will, and which has its
content in the ethical that comprehends the \emph{whole} of human life,
becomes such a religiousness as does not separate itself off as an exclusive
relation; rather its cultus becomes action, and its ideal is that the whole of
life \emph{as ethical} life is cultus.

For the ethical-religious self the Absolute becomes conscious in the
determinateness of the \emph{ethical}; for the self's being determines the
form of the Absolute's being for the self: as man is, so is his God. Religion
has its source in the human being itself: its content is conditioned by this
being's manifestness to itself; its God-consciousness keeps pace with its
self-consciousness. But religion has its source in the human being only
because the human being is grounded in the Absolute; and in the self, which is
posited by the Absolute and whose successive manifestation for consciousness
rests upon the communication of essence\-%
% --- p. 97 ---
\setcounter{footnote}{0}%
to the human spirit, the Absolute is also in the religious relation that which
works through and that which works primitively: in the \emph{last} instance,
therefore, we seek religion's source in the Absolute.

\begin{center}\rule{2cm}{0.4pt}\end{center}
% RULE, p. 97, the FIRST of the two on this page: a short centred single
% hairline standing mid-page with no head adjacent to it --- a thought-break in
% running text, like the rules of pp. 19, 28, 42 and 62. A small ink blot sits
% on its middle; it is not a second rule. Transcribed.

The conception developed in the foregoing will be more clearly elucidated in
its peculiarity through a comparison with theories with which it has striking
points of contact, while it nevertheless differs from them in something
decisive. With respect to the conception of the Absolute, of religion, and of
the relation between the ethical and the religious, the comparison with
\emph{Feuerbach}'s theory will be instructive; with respect to the ethical
life's relation to the determinations of actuality, the comparison with the
\emph{Hegelian} way of viewing things will be so.
% The SECOND of the page's two rules stands immediately below this paragraph
% and immediately above the run-in head: it is head ornament and is NOT
% transcribed, as at pp. 126, 128, 152, 164, 168, 177 and 182.

\textbf{Ludvig Feuerbach.} --- Already in the preface to ``Problemet om Tro og
Viden'' Feuerbach's standpoint has been characterised in general, and the
essential points in which I diverge from it indicated. It will however, for
the sake of the comparison and the critique, be necessary to enter more fully
upon a presentation of his conception in its connection and in its inner
development; and it will be so all the more since the utterances which have
appeared in this country about Feuerbach almost without
exception\footnote{As such an exception I must particularly bring out a
treatise which was printed many years ago in \emph{L. Helweg}'s
Fjerdingsaarsskrift for Literatur og Kritik (1845, 3rd number, pp.\ 1--17) and
whose author was my late friend, the schoolmaster \emph{Chr.\ Fenger
Christens}, known as a writer on pedagogy. But this treatise lies in time
before those writings of Feuerbach in which he has carried his theory through
to its final consequences.} betray not merely a highly incomplete acquaintance
with his writings, but a slight insight into what is characteristic of his
theory.
% HEAD, p. 97: the printed form is a run-in halbfette (bold) Fraktur head
% followed by an em-dash --- not letterspaced. The Indhold instead reads
% „1) Ludvig Feuerbachs“; the printed head is kept in the body, as in the
% Danish, and the divergence is recorded at the Contents.
% FOOTNOTE p. 97: printed mark *). Its two names are Fraktur and letterspaced,
% hence \emph{} and not \textit{}.

% --- p. 98 ---
What makes the apprehension and presentation of Feuerbach's view of the
religious difficult --- as of his philosophical view in general --- is partly
the aphoristic form of his philosophising, particularly at his last stage of
development, and partly the constant movement which takes place in his
thinking, and which lets every new writing mark a new stage in a development
whose beginning and end points lie at an enormous distance from one another
and belong to two wholly different directions of spirit. Only when the driving
force of this restless development is found, and the development's wholeness
surveyed, does the right understanding become possible.
% p. 98, l. 3: „hans Philosopherens aphoristiske Form“ as printed --- awkward
% but read the same way by both witnesses. The awkwardness is Danish-internal.

\emph{Feuerbach} began as a theologian, with a passionate striving to
penetrate faith's content with thought. The form of theology which drew him to
itself was therefore speculative theology; by it he was carried over to
\emph{Hegel's philosophy of religion}, and in apprehending this in the light
of the Logic he soon became a stranger to theology: the divine became for him
\emph{the essence of reason}, reason's unity the eternal in which the
transitory individuals are dissolved. In spite of beginning doubt he still
held fast the fundamental view of the Hegelian philosophy: the view of thought
as reality and truth existing in and for itself, so that \emph{the thought as
such} is actual, thought and being identical. But soon doubt of this view's
truth gained power, and as he moved away from it his conception of the divine
and of religion was transformed. \emph{Indirectly} the Hegelian philosophy had
acknowledged nature's essentiality and truth, by positing nature as a
necessary moment in the idea's eternal process; but this acknowledgement was
as it were revoked again, in that nature \emph{as} nature became something
merely untrue, vanishing. Feuerbach held fast that acknowledgement; he
attributed to ``the sensible'' first truth and essentiality \emph{within
nature's domain}, and then \emph{absolute reality}, since ``the being which is
opposed to sensibility as a heterogeneous being is nothing other than the
abstract being of sensibility''. What was primitive became for him no longer
spirit but na\-%
% --- p. 99 ---
ture, no longer thought but being: a proceeding of matter out of spirit he
declares unthinkable; therefore matter must precede spirit, nature be spirit's
ground. ``The true relation between thinking and being is only this: being is
subject, thinking predicate. Being is out of itself and by itself --- being
has its ground in itself, because only being has meaning, reason, necessity,
truth; in short, is all in all. Being is, because not-being is not-being, i.e.
nothing, meaninglessness. Being's essence as being is nature's essence''. And
being, nature, has now for Feuerbach not the significance of a universal
unity, but of \emph{the aggregate of the sensibly singular}, of the ``wholly
determinate, individual, sensible things and beings'' which become the one
actual thing; while the universal becomes only a subjective common
designation, just as ``spirit'' becomes only the unity of the senses.

This development of Feuerbach's fundamental philosophical view, from
\emph{idealism} to a peculiarly formed \emph{realism} --- which by the
acknowledgement of the human being's original activity, and by the
apprehension of the phenomena of spirit in which that activity manifests
itself, separates itself from ordinary sensualism and materialism --- had to
become decisive for his \emph{conception of the religious}. ``God was my first
thought'', he says, ``\emph{reason} my second, \emph{man} my third and last
thought. The subject of the Godhead is reason, but the subject of reason is
man''. The transcendent God's existence outside man transforms itself for him
first into an \emph{existence in reason}, i.e. an existence within reason, and
God's essence is identified in a \emph{pantheistic} way with \emph{reason's
essence}; then the religious representation is seen as an \emph{unconscious
projection} of the human self, which in God objectifies its own universal
essence as something different from itself; and finally --- as the reality of
the \emph{universal} is denied in a \emph{nominalist} manner, and the essence
of reason, which had been identified with the human race as nature's
conclusion, is absorbed by the single individuals --- the religious
representation becomes only a
% --- p. 100 ---
\emph{mirror-image of those individuals' craving}. In this process
\emph{pantheism gives place to atheism}, theology is dissolved, as metaphysics
was before it, into \emph{anthropology}, God vanishes, and only \emph{man}
remains as the true and the actual; and the \emph{religious relation} which
\emph{positive} religion posits becomes \emph{truth's utmost opposite}. Only
in a transferred sense is there still talk of religion as something valid.

What in the development of Feuerbach's philosophical standpoint --- which
conditions the development of his conception of the religious --- is the
moving motive, is \emph{the interest in actuality}. And while this interest,
as an interest in real \emph{natural actuality}, the foundation of concrete
human life, guides his polemic against abstract idealism, it guides, as an
interest in \emph{the ethical demand of actual human life}, his polemic
against the positively religious.

The writings which mark the stages of development in Feuerbach's
\emph{peculiar} conception of the religious are essentially the following:
1) Philosophie und Christenthum; das Wesen des Christenthums; (cf.\ P.\ Bayle
and the contemporary critical treatises) 2) Ueber das Wesen des Glaubens im
Sinne Luthers; das Wesen des Religion; Vorlesungen über das Wesen der
Religion; 3) Theogonie; Ueber Gottheit, Freiheit und Unsterblichheit. --- Of
these writings the fullest account must be taken of \emph{das Wesen des
Christenthums}, in which Feuerbach's conception of the Christian is developed
in connection. To this writing --- which precisely marks the transition in his
fundamental philosophical view, and is understood only when the stages between
which it forms the transition are clearly apprehended --- the writing Ueber
Philosophie und Christenthum is the polemical introduction, and the treatise
vom Wesen des Glaubens im Sinne Luthers the postscript. In the treatise on
religion's essence what in W.\ d.\ Christenthums was developed with particular
regard to Christianity is carried out with regard to the nature-religions, and
their relation to the religions of spirit is determined. The lectures on
religion's essence are a
% PRINTER'S DEFECTS in the list of Feuerbach's titles, p. 100, both confirmed
% at 300 dpi and read alike by both witnesses, both carried over as printed:
%   „das Wesen des Religion“ for „der Religion“;
%   „Unsterblichheit“ for „Unsterblichkeit“.
% The German titles are set in Fraktur, not antiqua, so none takes \textit{};
% only „das Wesen des Christenthums“ is letterspaced.
% --- p. 101 ---
further working out of the fundamental thoughts of the last-named treatise;
\emph{the Theogony}, and the parts belonging with it of F.'s last work,
Gottheit, Freiheit und Unsterblichheit, are a development of religion's
genesis and of its connection with the pathological in man --- a connection
which the treatise on religion's essence already strongly stresses. F.'s
polemic against \emph{theology} comes forward partly in various places in the
writings named, particularly in those cited under 1), and partly in his
critiques, as in his ``Beiträge zur Kritik des modernen Afterchristenthums''.

\emph{The fundamental features of Feuerbach's conception of religion}, as it
is given in ``\emph{das Wesen des Christenthums}'', are the following.

\emph{Since religion is found only in man, it must be grounded in that which
is the human being's peculiar} mark. This is, expressed generally,
consciousness; but consciousness in the strict sense is consciousness of one's
species, one's essentiality. The animal is an object for itself only as
individual; man is an object for himself also as species; therefore he has the
capacity for thinking, for knowing; he can make other things objects according
to their essential nature. Religion in general, \emph{as identical with man's
essence}, is identical with man's self-consciousness, with his consciousness
of essence; and when religion is at the same time consciousness of \emph{the
infinite}, it becomes this as man's consciousness of \emph{his own infinite
essence}. A really finite being has not the remotest inkling, still less
consciousness, of an infinite being; for \emph{the barrier of the being is
also the barrier of consciousness}. Consciousness of the infinite presupposes
the infinity of the conscious being; therefore religion points back to
\emph{the human being's infinity}.

Man is nothing without an object: through the object the subject becomes
manifest; but the object to which the human subject \emph{essentially and
necessarily} relates itself, the \emph{inner and chosen} object, is nothing
other than this subject's \emph{own}, but \emph{objectified}
% From p. 101 the long Feuerbach citations begin. They are run into the
% paragraph and marked by letterspacing alone, with no quotation marks; they
% are neither indented nor set as a quotation environment. The English does
% the same: the letterspacing is carried by \emph{} and no quotation marks are
% introduced where the printing has none.
% --- p. 102 ---
essence. Such an object is the \emph{religious} object; in it, therefore, man
becomes conscious only of himself.

In the case of the sensible object, consciousness of the object is separable
from self-consciousness; in the case of the religious object the two fall
immediately together: the object is man's manifest essence, his true objective
I, and its power over man is the power of his own essence.

\emph{Man cannot become conscious of a being higher than his own}. Whatever
object he becomes conscious of, he becomes constantly conscious of his own
essence; he cannot affirm another without affirming himself. Since the powers
of his essence, which constitute the essence, are realities, every working of
them --- and therefore also that by which they express the absolute object ---
is the immediate affirmation and confirmation of themselves; every expression
of them must be accompanied by consciousness of them as perfections. If man
could become conscious of himself, i.e. of his essence, as finite over against
another being known as infinite, then he would in consciousness limit his own
essence, i.e. negate it: \emph{his consciousness would therefore reach out
beyond his essence}. But consciousness is only the essence's
self-objectification, and therefore its self-affirmation, its realisation of
itself; it is precisely the highest form of self-affirmation; that negation is
therefore an impossibility. But where, then, does the distinction between the
finite and the infinite find its place? In the distinction between
\emph{individual and species}. \emph{The species' essence is the individual's
absolute essence}; beyond it the individual can only \emph{imagine} itself
able to reach. Every being is infinite for itself; its barriers exist only for
another being outside and above it: the being's measure is the understanding's
measure.

What holds of consciousness in general in relation to the object holds of
course of \emph{the single organs of spirit}. In that which is the essential
object for a single organ of spirit, the organ relates itself to itself,
objectifies itself. Thus \emph{reason's}
% --- p. 103 ---
\setcounter{footnote}{0}%
\emph{object is reason objective for itself, feeling's object} is
\emph{feeling} objective for itself. Therefore the \emph{difference of form}
between the religious and the philosophical consciousness becomes of such
\emph{decisive} significance: \emph{precisely the form}, which is the
expression of the organ's peculiarity, determines \emph{religion's essence}.
The Absolute is \emph{not} religion's object \emph{as such}, but \emph{as} it
is an object for the heart\footnote{``The heart'' F.\ sets on the one side in
opposition to \emph{reason}, as the individual to the universally valid; on
the other side in opposition to ``the Heart'' as the subjectively arbitrary,
limitlessly craving, over against the natural feeling and drive which
acknowledges its own limit.} and for the imagination. What transcendent
speculation about the religious posits as the merely secondary, as means and
organ, has precisely the significance of the primitive, of the essence, of the
object itself.
% FOOTNOTE p. 103: printed mark *). This page is NOT on the OCR-derived list of
% note-bearing pages --- found by eye. „Fornuften“ inside the note is Fraktur
% and letterspaced across its line break, hence \emph{} and not \textit{}.
% Feuerbach's „Gemyt“ / „Hjerte“ pair has no settled English equivalent; the
% note itself is what distinguishes them, so both are rendered "heart" and the
% distinction is left to Brøchner's own gloss, with the second capitalised as
% the printing capitalises „Hjertet“.

But when it is now said that man in the consciousness of the religious object
relates himself to his own objectified essence, so that \emph{consciousness of
God is man's self-consciousness}, then, in order that the \emph{religious
standpoint} may be designated in its peculiarity, a closer determination must
be added as essential: namely, that the religious man \emph{is not conscious of
this relation to himself}, and that the lack of such a consciousness precisely
constitutes religion's peculiar character. \emph{Religion is man's first and
indirect self-cognition}. Just as the child sees its essence, man, outside
itself, so that man as child is an object to himself as \emph{another} man, so
also, in the course of historical development, man first places his essence
outside himself before he finds it within himself; his own essence first
becomes an object to him as \emph{another being}. The historical development
in the religions therefore consists in this, that what counted for the earlier
religions as something objective is now intuited as something subjective ---
that is to say, that what was intuited and worshipped as God is now cognised
as something human. For the
% --- p. 104 ---
later religion the earlier one is idolatry: man has in it worshipped his own
essence. But every single religion of course excepts itself from this
judgement; because it has another object, another content, it imagines itself
raised above the laws which constitute religion's essence; it imagines that
its object, its content, is \emph{super}human. But philosophy sees through
religion's essence, hidden from religion itself; it shows that the opposition
between God and man is completely illusory, that it is nothing other than the
opposition between the human essence (the species) and the human individuals,
so that consequently the object and content of \emph{every} religion, even the
highest, are entirely human.

When it is conceded that God must \emph{unavoidably} be determined by man
through human predicates, and that consciousness therefore so far relates
itself in God to something humanly determinate, but when it is at the same
time said that God according to his essence \emph{in and for himself} is
infinitely higher than man, then this separation between what God \emph{is for
man} and what he \emph{is in and for himself} not merely disturbs religion's
peace but is in and for itself groundless and untenable. In this distinction
man sets himself beyond himself, beyond his essence, his absolute measure; but
this is an illusion. That distinction has significance only where an object
really \emph{can show itself} otherwise to me than it does, or where the
object is given immediately to the senses and therefore precisely also for
beings other than man; but it has none where an object which is exclusively
for man shows itself to me as it shows itself \emph{according to my absolute
measure}, \emph{as it must} show itself to me. If the representation of God
answers to \emph{the species' measure}, then that distinction falls away and
the representation is absolute; for the species' measure is man's absolute
measure, law and criterion. This \emph{every} religion also holds fast with
respect to \emph{its own} representation of God: that distinction is sceptical
and irreligious. Religion gives itself up when it gives up God's essence; it
is no longer a truth when it renounces
% EMPHASIS, p. 104: „er over“ is letterspaced and „menneskeligt“, in the same
% word, close-set --- the third mid-word break in this stretch. The English
% divides "super-human" so the run can stop at the same point.
% --- p. 105 ---
the possession of the \emph{true} God. What for man has the significance of
what is in and for itself, what for him is the highest being beyond which he
cannot represent anything higher: that is precisely for him the
\emph{divine} being. In the case of this divine object it can no longer be
asked what it is in and for itself. To ask whether God in and for himself is
such as he is for me is the same as to ask whether \emph{God is God}; it is to
raise oneself above one's God, to rebel against him. If one first doubts
\emph{the predicates' objective truth} and makes them anthropomorphisms, then
one must consistently doubt also \emph{the objective truth of these
predicates' subject}, and make the subject too --- that is, God's existence,
belief in a God generally --- into an anthropomorphism, a completely human
presupposition. \emph{The subject's necessity lies only in the predicates'
necessity}; for the subject is only the personified, the existing predicate.
The predicate is what is properly primitive. \emph{The certainty of God's
existence depends upon the certainty of God's quality}; only the predicate's
reality is the guarantee of the existence. For the Christian only the
Christian God's existence, for the heathen only the heathen God's, is a
certainty. \emph{God is man's essence, intuited as absolute truth}; but he is
man's essence \emph{apprehended in a determinate quality}, which, by being the
truth, becomes the highest existence; for what man intuits as true he intuits
immediately as existing.

\emph{If the divine predicates are human predicates, then the divine subject
too must be a human subject}. But here a phenomenon must be brought out which
characterises religion's inmost essence. The more human the divine subject is
in essence, the greater is apparently the difference between God and man, the
more does theology --- reflection upon religion --- deny the identity, and
lower the human as it \emph{as such} is an object for man's consciousness. The
ground
% PRINTER'S DEFECT, p. 105: „Prædicater ere menneskelige Prædikater“ --- the
% same word spelled with c and then with k within one sentence. Orthographic,
% and so not reproducible in English; recorded here.
% --- p. 106 ---
of this is simply that, since what is \emph{positive} in the intuition of the
\emph{divine} being is precisely the \emph{human} determinations of essence,
the intuition of the \emph{human}, as it \emph{directly} is an object for
consciousness, must be \emph{only} a \emph{negative} one.
% p. 106: the compositor spaces „kun“ and „negativ“ but leaves „være en“
% between them unspaced. As printed, and so in the English.

In order to enrich God, man must be made poor; he has \emph{his essence} in
God, how then should he have it in himself. Man negates with respect to
% TERMINAL MARK, p. 106: the printing closes „hvorledes skulde det saa have det
% i sig selv.“ with a FULL STOP, not a question mark, though the clause is
% interrogative in form. Reproduced as printed, per the rule that sentence-final
% marks follow the printing. Not in the Trykfeil leaf.
himself only what he posits in God. Religion abstracts from man, from the
world; but it can abstract only from the phenomenon, not from the essence, and
therefore religion must unconsciously posit in God everything it has
consciously negated, when that thing is something \emph{essential} and
\emph{true}, something therefore from which it cannot abstract, which it
cannot negate \emph{at all}: thus sensibility, reason, goodness, moral
activity.
% p. 106: „og“ between the two spaced words is itself unspaced; the same holds
% at pp. 108 and 114, though not at pp. 111, 115, 117, 119. As printed.
Through God man receives back, in infinitely richer measure, this which --- by
being posited in God, the essential object of spirit --- has already been
indirectly declared to be an essential determination in human nature.
Religion's secret is that man objectifies his essence, but then makes himself
an object for this objectified essence transformed into a subject. Man has
indeed in religion God for his end, but God has no other end than man's
welfare: the divine activity becomes a means for this, and therefore man has
in religion only himself for his end, in and through God. And what is posited,
as it were by an act of a religious force of repulsion, as an activity of God
--- man's separated-off subjectivity, from whom the religious consciousness
must let every impulse proceed --- becomes again, by an act of a religious
force of attraction, in that the divine activity is to work upon man, in man,
to become the principle of man's good disposition and of his good actions,
made into an activity human in its essence. Religion's course of development
therefore consists more closely in this, that man constantly appropriates more
of what he previously ascribed to God. What at the earlier stage was
% --- p. 107 ---
religion is so no longer at the later: what at the former counted as atheism
counts at the latter as religion.

But when religion is thus determined as man's indirect self-cognition and as
his unconscious objectifying of his essence, then the question must be
answered: in what relation does the \emph{religious} consciousness of the
essence stand to the consciousness present in other forms of the life of
spirit --- in other words, \emph{what part of man's essence} it is that
essentially gives the religious consciousness its content.
% p. 107: the letterspacing stops at „Væsen“; the rest of the clause is set
% tight. The English run ends at the same point.
The answer to this question will emerge from a consideration of religion's
character and essential standpoint.

\emph{Religion's essential standpoint} shows itself immediately as the
\emph{subjective, practical} one. Religion's end is the individual's welfare,
his salvation, his \emph{blessedness}. Man's relation to God is only his
relation to his own salvation. \emph{God is realised blessedness, or the
unlimited power to realise man's blessedness}. God as God, i.e. as he is an
object \emph{for religion}, is essentially \emph{only} an object for religion,
not for philosophy; for the heart, not for reason; for the heart's need, not
for thought's freedom --- in short, a being which does \emph{not} express the
peculiarity of the \emph{theoretical, objective} standpoint but that of the
\emph{practical, subjective} one. This also becomes plain from the fact that
religion attaches curse and blessing, damnation and blessedness, to its
doctrines, to belief in them. It appeals therefore not to reason but to the
heart, to the drive for happiness, to the affect of fear and hope. Doubt, the
principle of theoretical freedom, becomes here a crime, doubt of God the
highest crime. But what I cannot or dare not doubt without feeling disquieted
in my heart is not a matter for theory but for conscience, not a being of
reason but a being of the heart.

Now when religion's standpoint is the practical or subjective one, there lies
in this that \emph{only} the \emph{practical} man --- the man acting according
to his conscious physical or moral ends, who regards the world only in
relation to those ends
% --- p. 108 ---
and not in and for itself --- counts for religion as the \emph{whole}
essential man; and it is in relation to man in this determinateness that
religion's highest being is determined as a being of \emph{the heart}, as a
being which satisfies the need of the heart and of the imagination. All that
which lies behind the practical consciousness, but is the essential object of
theory --- of \emph{theory} in the most original and general sense, in the
sense of objective intuition and experience, reason, science in general ---
does not relate itself essentially to the religious consciousness. Religion,
being exclusively practical, subjective, lacks the true intuition of the
universe, of the really infinite, the true consciousness of the species. And
since now the true universal essence of nature and of man makes itself felt
for spirit as something that cannot be negated, that lack is indeed sought to
be made good in God, in a separate, personal being: the determinations of
essence of nature and of man, accessible only to the theoretical eye, are
involuntarily transferred into God --- but partly in such a way that they
constitute only the outermost point of attachment of the religious
consciousness, ``religion's mathematical point'', and partly in such a way
that they are transformed, and in this transformation appropriated by the
heart and the imagination. Since religion takes the apparent, superficial
essence of humanity and of nature for their true, inner essence, and
represents their true esoteric essence as one other being, it metamorphoses,
under the personification, the true essence of man and of nature into a
\emph{superhuman} and \emph{supernatural personality}, incomprehensible for
cognition, working through \emph{the miracle} and through it satisfying the
religious individual's practical need in contradiction with the law of reason.
% p. 108: „èt andet Væsen“ --- the accent is printed and is a GRAVE (è);
% 19th-c. Danish „èt“ = the numeral one, against the article „et“. Rendered
% "one other being" so the numeral force survives.

Religion's original separation of the divine and the human is involuntary and
naive; it originally identifies God with man just as immediately as it
separates him from man; it moves in the imagination's quantitative
differences, does not make the difference a qualitative and essential one. But
as reflection awakens, the sepa\-%
% --- p. 109 ---
ration becomes deliberate and studied, precisely in order to suppress that
consciousness of the identity which begins to dawn in the religious
consciousness. The first way, according to the concept, in which reflection
secures God as a being distinct from man is by \emph{the proofs of God's
existence}. Their end is precisely to separate God off from man, by
maintaining for him an existence as a thing in and for itself, which is not
merely to have being in our faith, our feeling, our thought, but is also to
have a real being distinct from these, an existence \emph{outside} us; and
this \emph{outside} must be taken \emph{literally}, for otherwise the
existence itself becomes an existence in an improper sense. Such an existence
outside us would however be nothing other than a \emph{sensible} being. But
God's being does not have the peculiarity of sensibly real being: that I am
involuntarily affected by it, that it is even when I am not, do not think it,
do not feel it. God is only for man, in so far as man raises himself above
sensible life, in so far as he feels, thinks and believes God. But since God
is \emph{only} an object for \emph{man}, this will simply mean: God is only in
so far as he is felt, thought and believed. The existence one wishes to
maintain for God would become a real being which yet was not real: a
\emph{spiritual} being, one calls it, as a way out. But spiritual being is
precisely only to be felt, to be thought, to be believed. God's being is
therefore a contradictory middle thing between sensible being and thought
being: a sensible being which, by being deprived of all the determinations of
sensibility, is reduced to a vague existence in general, and thereby abolishes
itself, since to existence there belongs full, determinate reality. \emph{A
necessary consequence of this contradiction is atheism}. God's existence,
which does not denote inner reality, truth, but outer existence, has the
essence of an empirical existence, but without having its marks; it is in and
for itself a matter of experience, and yet in reality no object of experience.
It itself calls upon man to seek it in actuality; and when man finds
experience in contradiction with the representations that existence has
awakened, he is completely justified in
% --- p. 110 ---
denying this existence. Only sensation, not reason, can give the proof of
empirical existence, and in the proofs of God's existence being has precisely
the significance of such an existence. For religion, belief in God's existence
without regard to the inner quality, the spiritual content, becomes at last
unavoidably the main thing. God's existence in and for itself establishes
itself as religious truth, and to it, as experience shows, the crudest and
most unhallowed representations can attach themselves.

At religion's standpoint it is \emph{the imagination} --- the place for the
absent existence, withdrawn from the senses but yet sensible in essence ---
that resolves the contradiction in an existence which is at once to be
sensible and non-sensible. Only the imagination preserves from atheism; for in
it the existence has sensible effects. Where God's existence is a living
truth, a matter for the imagination, there apparitions of God are believed in
as well. Where on the contrary the fire of the religious imagination is
quenched, and where the sensible effects or phenomena which are necessarily
bound up with an in-and-for-itself sensible existence cease, there the
existence becomes a dead and self-contradictory existence, which falls
irretrievably to atheism's negation.

The same contradiction which makes itself felt in the representation of God's
existence makes itself felt in the determination of \emph{God's essence
generally}. All the positive determinations in that essence are human
determinations; but in negating the determinateness in the determinations,
reflection makes them non-human. Yet by this negation the concept is
dissolved; there remains only a negative, contentless, self-contradictory
representation; for precisely that determinateness has been negated which made
the determination what it is, and only this is a real concept; whereas that
which is to make the determination something essentially other is objectively
nothing, subjectively a mere fancy.

This Feuerbach now carries out in ``das Wesen der Christenthums'' with respect
to the chief Christian dogmas.
% PRINTER'S ERROR, carried over as printed: „das Wesen der Christenthums“ for
% „des Christenthums“. The same error recurs at pp. 113 and 117, so it is the
% book's settled form. The title is set in FRAKTUR here, not antiqua --- as are
% all the German titles in this stretch --- so it takes no \textit{}.
The contradiction in Christianity, ``the religion of love'', he sees
% --- p. 111 ---
decisively concentrated in \emph{the opposition between faith and love}.

For in love he sees a revelation of religion's \emph{hidden} essence, in faith
that which constitutes its \emph{conscious} form. Love identifies man with God
and God with man, and therefore man with man; faith separates God from man,
but therefore also man from man; for God is nothing other than man's mystical
concept of the species: God's separation from man therefore becomes man's
separation from man, the dissolution of the common bond. Since faith makes God
into a \emph{separate} being and apprehends him as \emph{person}, the relation
to God is separated from the relation to men, and there must unavoidably arise
a conflict between the two relations, between religion and morality; and since
the relation to God, as the highest, lays claim to all man's powers and his
whole interest, morality must give way to religion, love be annihilated by
faith. Faith must raise itself above the laws of natural morality; there must
be actions in which faith comes to light in distinction from and in opposition
to morality; the duties towards God must collide with the simple human duties,
and faith privileges everything. Even when Christianity says God is love,
there remains --- since love becomes only predicate, and the proposition is
not at the same time reversed, so that it would read love is God --- a place
for the divine personality as such; and the subject, the divine personality,
becomes the main thing, the predicates, the ethical determinations, a mere
side-issue. But in thinking God as subject, in separation from the predicate,
a room is left over for unlovingness, and it becomes unavoidable that I give
now the subject, now the predicate the preponderance. This contradiction the
history of Christianity has sufficiently confirmed. --- Christian love is, as
Christian, a limited love. But a love which is limited by faith is an untrue
love. True love, free, universal love, knows no other law than
% --- p. 112 ---
its own and reason's. It is the universal law of intelligence and of nature;
it is nothing other than the realisation of the species' unity by the way of
disposition. It thus has a self-validity independent of the religious, and
does not need to be supported upon it. When universal love, grounded in
itself, becomes the principle, Christ vanishes before the consciousness of the
species; for Christ is himself only an image under which the popular
consciousness represented the species' unity, and true universal love can be
grounded only upon the unity of the species, of intelligence, upon humanity's
nature. Only then is it a thorough, principle-secured, free love; for it
supports itself upon love's origin, from which Christ's love too derived. That
love was itself a derived one: Christ did not love of himself, but in virtue
of humanity's nature. The law of the species, of intelligence, is that we
should love man \emph{for man's sake}; love tolerates no intermediary.

Although Christianity acknowledges love, it nevertheless comes to deny it
again, because it has not set love free, has not apprehended it absolutely;
and this it could not do, because it is religion, and love is therefore
subjected to the dominion of faith. Love is only Christianity's exoteric
doctrine, but faith its esoteric; love is only its morality, but faith its
religion.

In the contradiction between faith and love lies the necessity of going beyond
Christianity, and beyond religion's peculiar essence generally.
\emph{Religion's content and object is a completely human one: theology's
secret is anthropology.}
% p. 112 is a whole page of Feuerbach quotation carrying no quotation marks at
% all; the citation is marked by letterspacing alone and is run into the
% paragraph, so no quotation environment is used --- in the English either. The
% letterspacing of the thesis runs unbroken across the Danish line breaks.
Of this religion has no consciousness; but the task is to lift the illusion,
to come to the cognition that consciousness of God is nothing other than
consciousness of the species; that man can and should raise himself above the
barriers of his individuality, but not above his species' laws and positive
determinations of essence; that man cannot think, divine, represent, feel,
believe, will, love and worship
% --- p. 113 ---
any other being than the human species' being --- including nature, which
belongs to man's essence as man belongs to nature's --- as the absolute being.

\emph{Man, who in religion is in and for himself the first, must be posited
and declared as the first.} All man's moral relations to man must be posited
as in and for themselves valid, holy, and thereby as religious. As substantial
relations they can be grounded only by themselves: a grounding of morality by
theology is an illusion, a movement in a circle. The right, the true and the
good has everywhere its ground of holiness in itself, in its own quality.
Where right and morality are taken in earnest, there they are in and for
themselves divine powers. If they have no ground in themselves, then neither
is there any inner necessity for them: they are then given over to religion's
groundless arbitrariness.

In the self-conscious spirit's relation to religion, then, all that is at
stake is the annihilation of an illusion --- but of an illusion which works
fundamentally corruptingly upon humanity, and robs man both of powers for
actual life and of the sense for truth and virtue; for even love, in and for
itself the most inward, truest disposition, becomes through religion a merely
apparent, illusory one, since religious love loves man only for God's sake,
and therefore only apparently man, but in reality only God, and in God
egoistically itself. But that illusion is annihilated and the truth becomes
manifest as soon as we merely turn the religious relations round --- as soon
as we constantly apprehend what religion posits as means as the end, and what
it makes the subordinate thing, the side-issue and condition, as the main
thing and the cause.
% p. 113: a small raised speck stands immediately before „om“. The same speck
% recurs at p. 117 (twice) and p. 118. Four instances in six pages, all at the
% same height and all of the same round shape: read as set-off / ink specks in
% this copy, not as type, and therefore not transcribed in either file, though
% both OCR witnesses pick them up.
It is therefore in the interest of reason, of humanity and of morality that
the religious representations are given up; precisely that interest is
atheism's source, and this atheism therefore negates not the real, positive
content of the representation of God, but only the untrue representation of
God as a separate personal being, and what is inseparable from that
representation. ---

What in ``\emph{das Wesen der Christenthums}'' is
% p. 113: here the title IS letterspaced (it is not at pp. 110 and 117).
% --- p. 114 ---
developed about the religious with particular regard to the Christian is
completed in the treatise ``\emph{on religion's essence}'', in that closer
account is taken of the lower forms of religion, and in that the universal
relation between \emph{the nature-religions} and \emph{the religions of
spirit} is determined.

The starting point is formed by the proposition \emph{that man's feeling of
dependence is religion's ground} --- religion therefore being the conscious,
declared feeling of dependence. As the \emph{original object} for this feeling
of dependence \emph{nature} is then determined; and nature is posited not
merely as religion's \emph{original object} but also as its \emph{abiding
ground}, the lasting though hidden background for it, as the ground of belief
in God's \emph{objective} existence outside thought and reason, as the
foundation for those properties which express the divine being's difference
from the human. God, thought as the originator of nature, is indeed
represented as a being distinct from nature; but this being's real content is
only nature. All God's properties, which make him an objective, actual being,
are properties abstracted from nature, which presuppose it, express it, and
which therefore fall away when nature falls away. There does then indeed
remain a being, an aggregate of properties such as infinity, unity, necessity,
eternity; but this being is nothing other than nature's \emph{abstract}
essence, which as such presupposes nature's concrete, sensible, actual
essence.

But when it is now said that religion's first and essential object is nature,
this is, be it noted, to be understood thus: that nature, even in the
nature-religions, is \emph{not} an object \emph{as nature}, but as a
\emph{personal, living, sensing, man-like being}. Man does not originally
distinguish himself from nature, and consequently does not distinguish nature
from himself either: he therefore makes the sensations which a natural object
awakens in him immediately into properties of the object. \emph{The
nature-being} thus becomes a \emph{being of the heart}, i.e. a subjective,
human being.
% --- p. 115 ---
It is particularly nature's changeableness that brings man to apprehend it as
a human, arbitrary being. This conception reveals itself in \emph{sacrifice},
the most essential act of the nature-religions. And in sacrifice the whole
essence of religion concentrates itself; for sacrifice's ground is the feeling
of dependence, but its result and end self-feeling. \emph{And thus the feeling
of dependence upon nature is indeed religion's ground; but the abolition of
this dependence is religion's end.} The contradiction between willing and
being able, wishing and attaining, representation and actuality, is religion's
presupposition; the being in which this contradiction is abolished is the
\emph{divine} being, and this being is an object of prayer, determinable by
prayer and sacrifice. But since nature becomes religion's object as
metamorphosed into a being of the heart, of representation, of imagination,
the nature-religion becomes the striking contradiction between representation
and actuality, fancy and truth. This contradiction religion escapes only by
making its object an \emph{invisible} one, a wholly \emph{non-sensible} one, a
being which is an object only for \emph{faith}, for representation, for
imagination --- in short, for \emph{spirit} --- and therefore in and for
itself a \emph{spiritual} being. When man, from a merely physical being,
becomes a \emph{political} one, and generally a being which distinguishes
itself from nature and concentrates itself upon itself, then his God too
becomes, from a merely \emph{physical} being, a \emph{political} being
distinct from nature. Where man, through the life of society, becomes a being
which determines itself by principles, rules of wisdom, laws of reason ---
therefore a thinking, understanding being --- and where powers other than
powers of nature, namely political, moral, abstract powers, become objects of
his feeling of dependence, there the world, nature, also shows itself to him
as a being dependent upon and determined by \emph{understanding and will}.
Where man with will and understanding raises himself above nature, becomes a
supranaturalist, there God too becomes a \emph{supranaturalistic} being,
\emph{lord over nature}, and finally --- what is
% --- p. 116 ---
the condition of being lord in the strict sense --- \emph{creator of nature}.
The unity of the human consciousness and spirit, which comprehends the whole
manifoldness of the outer world, is what leads to laying down \emph{God's
unity}. Proper \emph{theism} arises only where man refers nature only \emph{to
himself} and makes this relation nature's essence, thus making himself
nature's final end. Where nature has \emph{its end} outside itself, it
necessarily has also its \emph{ground} and \emph{beginning} outside itself;
where it exists only \emph{for another being}, there it necessarily exists
also \emph{through another being} --- and through such a being as in producing
it had man in view as enjoying nature. The doctrine that God is the world's
creator has its ground and significance only in the doctrine that man is
creation's end. And nature's creator is only the condition of the creator of
man raised above nature; only he is God and creator \textit{sensu eminenti},
just as religion's proper providence is the \emph{special} providence
\emph{for man}.
% p. 116: „sensu eminenti“ is the only antiqua in this stretch, and it is set
% tight (not letterspaced), so \textit{} without \emph{}.

The \emph{spiritual} being which man, having advanced from the stage of the
nature-religions to that of the religions of spirit, posits above nature and
makes nature's presupposition, is nothing other than \emph{man's own spiritual
being} in the \emph{religious} consciousness's apprehension --- but which for
that reason shows itself to man as \emph{another} being, \emph{distinct} from
man, incomparable, because it is made into \emph{the cause of nature}, into
the cause of effects which the human spirit, the human understanding and will,
cannot produce; for the reason, therefore, that man combines with this
spiritual human being at the same time nature's being, which is distinct from
the human. While nature's essence is in this combination humanised, its being
in objective separation from man conditions the objectification of the human
essence as divine.

Belief in a God becomes therefore either \emph{belief in nature}, in the
objective being, \emph{as a human, subjective being}, or \emph{belief in the
human being as nature's being}. The former belief is
% --- p. 117 ---
\emph{nature-religion}, polytheism; the latter \emph{religion of spirit}
(``spirit-man-religion''), monotheism. Religion's first principle runs: I am
nothing over against nature, all is God over against me (the standpoint of
fetishism, the foundation of polytheism). Religion's concluding proposition,
on the contrary, runs: all is nothing over against me, all is only a means for
me. But this nature can be only when it is not through itself but through God.
In \emph{the miracle} religion's end is fulfilled in a sensibly popular way;
man's dominion over nature, his divinity, is in the miracle a striking truth.
\emph{The God-man is the revelation of the miracle's secret thought: that man
is God.}

God is a \emph{religious} word, a religious object and being, an object whose
existence is given only with religion's existence, whose essence is given only
with religion's essence, which therefore cannot exist outside religion, and in
which there is contained no more objectively than there is subjectively in
religion. God expresses only religion's essence: imagination and heart.
Religion's object, \emph{as} it is an object for religion, does not exist in
actuality, but is only an object for faith, which represents that which is not
but is craved to be, as being. --- As men's wishes are, so are their gods. The
Greeks had limited wishes and limited gods, the Christians an unlimited God,
exalted above all natural necessity, superhuman, transcendent, and absolutely
unlimited, fantastic wishes. Their fundamental wish is the enjoyment of an
infinite, unlimited, unutterable, indescribable blessedness. ---
\emph{Blessedness and divinity are one.} ---

In the treatise on religion's essence, as emerges from the foregoing, not only
is the horizon widened by the account taken of the nature-religions, but the
determination of the religious has undergone an essential modification. In
``das Wesen der Christenthums'' religion's standpoint was indeed apprehended
as the one-sidedly \emph{practical} one, and the determination of
``\emph{the heart}'' as religion's organ prepared the consequences later
drawn; but there lay
% --- p. 118 ---
nevertheless, in the apprehension of religion as man's indirect
\emph{self-cognition} --- when that was put in connection with the
apprehension of the species', the universally human's, relation to the
individual, and of its relation to reason --- a point of attachment for a
theoretical, rational element; and this also came to validity in so far as
God's essence acquired the determination ``being of understanding'' too, even
though this determination was to be only an outermost point of attachment for
the religious consciousness and not to express what was religion's proper
kernel. Now, on the contrary, in the apprehension of religion the accent falls
far more sharply and exclusively upon that side of human nature which, in its
separation from the theoretical, becomes the \emph{egoistically drive-like,
the pathological, the alogical.} As a connecting transition there appears, in
the treatise \emph{vom Wesen des Glaubens im Sinne Luthers}, the apprehension
of \emph{faith} as the expression of \emph{indirect self-love}. But still more
sharply than in the treatise on religion's essence, and with a still more
complete pushing back of the rational moment, the change in the apprehension
of the religious appears in ``\emph{Theogonie}'' and in Feuerbach's last
writing, ``\emph{Gottheit, Freiheit und Unsterblichheit}''.
% PRINTER'S ERROR, carried over as printed: „Unsterblichheit“ for
% „Unsterblichkeit“ --- the same error as at p. 100. The whole title, and
% „vom Wesen des Glaubens im Sinne Luthers“, are Fraktur and letterspaced; the
% Luther title carries no quotation marks, and none is supplied here.

In ``\emph{Theogonie}'' it is developed, with many turns of phrase, how
\emph{the wish} --- the irrational, egoistic craving --- is \emph{the gods'
first ground}, the source of belief in them; that the gods are \emph{not}
beings of reason but essentially expressions of craving and wish. God is the
supplement to man's defective essence, to his capacity, so finite and limited
in contradiction with his wishes; the certainty of God's being rests upon the
power of self-love, which over against actuality's painful limitation forces
out a representation of that limitation's abolition --- a representation so
rooted in man's craving that it unavoidably acquires actuality for him. Before
the divine power the tormenting feeling of limitation vanishes. God is and can
do what man is not and cannot do, but would gladly be and be able to do. The
wish has, as a \emph{theogonic} power, precedence over fear:
% --- p. 119 ---
it is not fear in and for itself, but fear's wish --- that there should be
nothing to fear --- that creates gods. But the wish which sums up all wishes
is \emph{the wish for happiness}; in this wish, therefore, the gods' origin,
religion's principle and fundamental essence are to be sought, and by the
nature and limit of the wish for happiness the nature and limit of the gods
are determined.

\emph{The wish itself as such} stems from \emph{man}; but the wish's
\emph{object} stems from \emph{outer nature}, from what is given by the
senses; for man has originally no empty, supranaturalistic, fantastic wishes.
No: the objects of his senses are also the objects of his wishes. Precisely
therefore the gods, as ``\emph{wish-beings}'', are at the same time
\emph{nature-beings}. But this is not to be understood as though the gods as
such were deified or personified powers of nature or bodies of nature; they
are expressions of \emph{nature's effects upon the heart and the imagination}
--- personified, made independent, objectified \emph{feelings, affects} --- but
affects which are bound to those bodies of nature by which they are awakened
or brought about. Thus the \emph{human} shines through in the determination of
the nature-being; and the \emph{human wish}, which draws its content
\emph{from nature}, is raised in power --- as man distinguishes himself as the
\emph{higher} from nature --- into the \emph{supranaturalistic wish}, into the
\emph{wish for blessedness}, in which the mystery of the Godhead is first
revealed. \emph{The wish for blessedness, self-love's most intensive
expression}, is \emph{Christianity's} fundamental wish; and the Christian God,
who fulfils this wish, is personified blessedness, absolutely raised above the
barrier of nature and of natural necessity.

That men have \emph{everywhere} represented gods, but as \emph{human} beings,
is the historical, factual proof that man in reality represents the
\emph{human} essence as divine, and that the human essence is the fundamental
presupposition for the representation of a Godhead --- that it is the
primal-being which lies at the foundation of
% --- p. 120 ---
\setcounter{footnote}{0}%
the gods. As \emph{highest} being I can worship only what expresses \emph{my
own} essence, but in a degree and perfection which I lack, and which precisely
therefore carries me away to worship and admiration. God is indeed the
\emph{superhuman}, the infinite being --- but, be it well understood, the
\emph{infinitely human}, the \emph{superhumanly human} being: a being which is
more, infinitely more man than man himself. The Christian God is just as fully
as the heathen one a human being, only of another kind than the latter,
because the Christian too is a man of another kind than the heathen. ---

With this result this writing joins itself to the earlier ones; and in taking
polemical account of the distorting reinterpretations of religion, it declares
itself against them in these words: ``There is only one true, one venerable
God; --- it is the immediate, self-active, self-speaking, self-shining,
self-lightning-flashing, self-thundering, self-raining God of the ancient
world. Either this one or no God! Where the old view of the world is
abandoned, where there is an independent doctrine of state and of law, an
independent doctrine of nature, which no longer sees in the wind God's breath,
in the thunder God's voice, in the lightning God's fire; there
\emph{atheism} is not merely a scientific but a moral truth and
necessity''. ---
% The Danish sets this as a marked quotation in „…“; the English keeps it as a
% quotation, as everywhere else in this file. Brøchner is quoting Feuerbach in
% Danish, so what stands here is a rendering of Brøchner's Danish, not of
% Feuerbach's German.

By means of the survey given of the development of Feuerbach's theory of
religion, what is unsatisfactory in that theory can easily be made plain.
There is in Feuerbach a great deal to which one cannot deny recognition. There
is in him a great clarity in the negative: in the polemic against theology and
against those philosophical mixed forms in which philosophy and Christianity
were sought to be fused into a unity. There is in him a sharp eye for what is
\emph{common to all the positive} forms of religion, and a sharp apprehension
of the phenomena which appear within the domain of the positive
religions\footnote{This \emph{phenomenal, factual} element F.\ then apprehends
as the \emph{necessary} expression of the positive religion's essence, without
letting that \emph{question} come forward which is \emph{indirectly} answered
by \emph{Kierkegaard}, in ``Kjærlighedens Gjerninger'' and in other places:
whether this factual, phenomenal element may not find its corrective at the
positive religion's own standpoint, and whether it is not to be explained by
an unconscious relapse from the sphere of the religious to lower stages of
life.}.
% FOOTNOTE p. 120: printed mark *). It runs over onto p. 121 in the print; set
% here in full at its call, as in the Danish. „Kierke-gaard“ is letterspaced
% across the page break, both halves.
% EMPHASIS p. 120: the run „fælleds for alle de positive“ stops at the line
% break; the following „Re-“/„ligionsformer“ is tight. As printed, and the
% English run ends at the same point.

% --- p. 121 ---
There are fine and deeply penetrating analyses of the positive religion's
determinations, and of the fundamental psychological relations which
illuminate the religious representations. There is a correct recognition of
religion's essential connection with human nature, and of the essential
connection between the religious consciousness's object and the apprehending
organ. There is a clear demonstration of the impossibility of reaching beyond
the human being's fundamental determinations, of the impossibility of negating
thought by thought. And to this comes, as something essential, that all F.'s
investigations bear the stamp of a high degree of honesty and love of truth,
of a living conviction, and at the same time of an energetic striving forward
which always sets his inquiry in motion anew and does not let it content
itself with results once won. But in spite of all this the Feuerbachian theory
has its great and essential defects. We bring out the most decisive.

\emph{The fundamental defect} is that with him a one-sided emphasis is laid
upon religion's proceeding \emph{from the human being}, and that it is
overlooked that the human being is itself something posited, something
derived, and that what posits and gathers together works through it. This
defect already becomes striking in F.'s earlier writings, as in \emph{Wesen
der Christenthums}, where nevertheless, by a reminiscence from idealism, the
concept of ``the species'' was still accentuated as the universal, the
essential, in which the individuals are absorbed and which realises itself in
them, and where the apprehension of \emph{the human species' relation to
reason and to nature} made it natural to seek back to a principle of existence
which within the unity of essence with the human had a transcendence over it.
But still more determinately must the defect come forward in the later
writings, where in a nominalist manner the individuals
% PRINTER'S DEFECT p. 121: the page prints „Relionens“ for „Religionens“ (the
% -gi- simply absent). Orthographic, and so not reproducible in English.
% EMPHASIS p. 121: „Slægtens“ inside the quotation marks is NOT letterspaced;
% the letterspaced runs are the four carried here.
% --- p. 122 ---
are posited as the properly real, the species only as a subjectively valid
universal expression for their aggregate, just as reason's identity is
apprehended only as ``an expression necessary for thinking and for language,
by which the universal is fixed and made independent, but in respect of which
it must not be forgotten that it is only a product of thinking''. Here religion can
be thought only as proceeding from the individual, and its source then
becomes, since the individual is not governed by the species' unconsciously
through-working rational power, merely \emph{the alogical in the in}dividual.

\emph{The rational} moment in religion is with F.\ pushed back in a one-sided
manner. Since religion's essential standpoint becomes the
subjective-practical, man in religion is to apprehend himself only from the
standpoint of the practical, and everything which is theory's object is to
fall outside the religious individual's consciousness of the human, and to
find its place in the divine only in a transformation necessitated by the
religious standpoint, or as an outermost point of attachment which is soon
lost wholly from sight. But in this apprehension the essential co-operation of
the different capacities, the different organs, in the human being --- even
under the sporadic development --- and the necessity of a striving after a
satisfaction of \emph{all} the fundamental capacities, are misapprehended. It
is indeed originally conceded that the theoretical too makes itself felt
involuntarily for consciousness; but this concession is in reality taken back
by the limitations added, and by the significance F.\ attributes to the
determination of \emph{personality} in the religious representation of God.
That, however, the theoretical \emph{as} theoretical must acquire a
determining influence upon the religious apprehension of the Absolute lies not
merely in this, that spirit in religion seeks reconciliation, satisfaction,
\emph{according to its wholeness}; it also emerges from the fact that it has
been possible for the human spirit to work its way out of the first, crudest
religious forms, before which as yet no independent --- i.e. independent of
the religious ---

development of thought had taken place. At \emph{higher} stages of religion's
development
% EMPHASIS p. 122: two runs stop inside a word at the line break and are kept
% as printed --- „Alogiske i In“ (with „dividet.“ tight on the next line) and
% „Personligheds“ (with „bestemmelsen“ tight on the same line). In „det
% Theoretiske som Theoretisk“ only „som“ is letterspaced; both occurrences of
% „Theoretisk“ are tight.
% --- p. 123 ---
\setcounter{footnote}{0}%
there is, in spite of the apprehension of the \emph{divine} will as
\emph{arbitrary}, a possibility that a \emph{content of thought} may lay
itself into the religious forms, and that in them the rational may assert an
\emph{essential} significance; and this content of thought, which is as it
were precipitated into those forms, has partly come about through a
development of thought which has not separated itself from the religious but
has been furthered by religion itself, and has partly, in virtue of an inner
need, been \emph{appropriated} by the religious consciousness from a
development of thought relatively independent of it.

The pushing back of the rational which shows itself in F.\ hangs together with
this, that the significance of the working of the \emph{a priori} in the
religious consciousness, as rationalising and ethicising, is overlooked. The
unconscious power is misapprehended which \emph{man's essence} exercises over
\emph{the individual}, in that the essence, working as a unity, governs what
is individually irrational and egoistic. But since this working of the a
priori is overlooked, there is a one-sided accentuation not merely of the
irrational but also of the \emph{egoistic} in the religions. Therefore neither
does the moment of religious self-surrender receive its proper significance.
F.\ is not blind to this moment's presence in the religious mind, which
demands a disinterested, self-surrendering love: \textit{diligere Deum propter
Deum et non propter se ipsum} (St.\ Bernhard). But just as the rational
``metaphysical'' determinations of the divine were only the ``outermost points
of attachment'', not the proper, specifically religious determinations, so
self-surrendering love is to designate only the moment of highest religious
enthusiasm, which is at once succeeded by that relation in which the subject,
as separating itself from the divine object, relates itself to itself as an
object for God, and therefore indirectly egoistically. And even where the self
negates itself in self-surrender, it is to do so only because it has in God
the enjoyment of a personality freed from all barriers; because God is the
human heart's highest self-feeling and feeling of pleasure\footnote{Cf.\ with
this F.'s apprehension of sacrifice: above, p.\ 115.}.
% ANTIQUA p. 123, the decisive page for the convention: the Latin sentence is
% set in upright antiqua and is TIGHT, not letterspaced --- hence \textit{}
% alone. Inside the same Latin sentence „(St. Bernhard)“ is Fraktur and stays
% so; not one proper name in the book is set in antiqua.
% EMPHASIS p. 123: „Menneskets Væsen“ and „Individet“ are two separate runs
% with „udøver over“ tight between them.

% --- p. 124 ---
This one-sided singling out of a single side in the religious relation is
strengthened in F.\ by this, that precisely through the accentuation of
personality's determination of form, in separation from its content, the
\emph{one-sided separation between the divine and the human} gains the
preponderance over that \emph{equilibrium} between unity and separation --- or
rather, over that \emph{preponderance} of unity over separation --- which was
originally laid down by the determination of the \emph{specifically religious}
predicates of the divine. And it is strengthened likewise by this, that the
\emph{difference of degree in the clarity and distinctness of the content of
consciousness} which takes place in the course of the religious consciousness's
development is as little sufficiently brought out as \emph{the relation between
the moments of the religious process of consciousness} is correctly
apprehended. The \emph{mystical} element in the life of faith does not come
into its right in F.'s apprehension. But precisely upon this mystical element
--- conditioned by the fact that in the unity of the religious relation of life
the moments of the relation of consciousness \emph{do not step sharply} apart
from one another --- rests the immediate religious consciousness's feeling of
unity with the divine: the possibility it possesses of a self-surrender of
enthusiasm, of a finding of itself again according to its essential
determination in the divine. \emph{It is unavoidable that reflection}, the
source of separation, should make itself felt in the religious domain. Where
human thinking, freeing itself from the religious, develops beyond the
determinate religious standpoint and reacts back upon it, or where the human
spirit, in the domain of the religious itself, begins by a spontaneous
development of thought to loosen itself from the earlier standpoint,
reflection cannot fail to make itself felt; but it is only where reflection
acquires an \emph{essential} significance in the domain of the religious, and
where it seeks to maintain religion's truth \emph{in opposition to thinking's,
that the moment of separation} gains preponderance in the religious
consciousness, and that the consequences F.\ draws acquire their validity.

With the one-sided bringing out of the separation there hangs
% p. 124 prints „F.s Opfattelse“ without an apostrophe, against „F.'s“ on
% pp. 121 and 123 --- kept as printed on both readings in the Danish; the
% English uses the apostrophe throughout, the distinction being purely one of
% Danish typography.
% EMPHASIS p. 124: the letterspaced runs carry the small words with them ---
% „Det er uundgaaeligt, at Reflexionen,“ is spaced throughout, and so are „at“
% and the opening „i“ of „i Modsætning“.
% --- p. 125 ---
together finally that defect in F., that \emph{the relation of opposition
between the religious and the ethical is one-sidedly brought out}, while the
unity which the mystical element in religion grounds is one-sidedly pushed
aside. ---

Through this indication of what I regard as the defects in Feuerbach's
apprehension of \emph{the positively religious}, its difference from mine will
have become sufficiently plain, and from the developments of concepts given in
the foregoing an essential difference in our whole fundamental view will
emerge.

The difference appears at once in the determination of the \emph{fundamental
concept}. For me it becomes the concept of the ideally-real Absolute as that
principle of unity which grounds existence in relative independence; for
Feuerbach it becomes, where his peculiar standpoint is \emph{fully} developed,
the individuals, the single sensible things and beings --- while the divine
becomes only a mirroring of subjective craving, and the universal only a
subjectively indispensable expression for the aggregate of the individual, and
as such without real validity.

From this difference in the fundamental concept there follows as a consequence
that, whereas at the standpoint I have indicated the conditions will be
present for maintaining the possibility of \emph{freedom in the} human, for
maintaining a \emph{cognition and an acting} which bears the stamp of the
\emph{universally valid} --- and therefore a true scientific cognition and a
true ethical acting --- and finally for maintaining a \emph{religious
relation} to the through-working Absolute, in such a way that this relation is
everywhere in harmonious unity with the relation to man: with Feuerbach, on
the contrary, the successive development of his fundamental view into
individualism and sensualism must lead to this, that not merely is a religious
relation to an absolute made impossible, but that the very foundation of the
universal is annihilated upon which a true scientific cognition must rest ---
and that even, through the annihilation of freedom's possibility and of the
concept of an ideal for ac\-%
% EMPHASIS p. 125: „Friheden i det“ is spaced through „det“, with
% „Menneskelige“ tight; in the next phrase „der har det“ is tight, so the page
% carries two runs, not one. The English divides at the same joints.
% --- p. 126 ---
tion, the ethical itself must be undermined --- the very thing for whose
benefit the religious was to be sacrificed.

When F.\ retains the expressions \emph{religion and religious}, it is only in
an \emph{improper} sense, as a designation for that to which man relates
himself as to something highest, and which also acquires practical
significance for him. They do not designate for him what they designate for
me: an absolute relation to the Absolute, a relation which is through this
Absolute and in which man relates himself to it with what is of his own
essence, but at the same time to that higher principle which grounds the human
being and with it the whole of existence, and gathers it together in its
eternal unity. For me, therefore, ``man's self-cognition'' in religion becomes
a self-cognition \emph{by and through his essence's ground}, his absolute
ideal.

\textbf{The Hegelian Philosophy.} --- This philosophy's apprehension of
\emph{the task of life} follows with necessity from its apprehension of
\emph{the idea}. Since the idea, as the highest expression of ``pure thought'',
comprehends within itself all

existence, is \emph{a unity of thought and being}, then \emph{the idea's
eternal process} --- in which nature and the finite spirits, whose
independence is only a semblance, are eternally taken up as abolished moments
--- comprehends \emph{all} actuality and \emph{all} truth; all true existence
is reduced to eternity's element; and what falls outside this in the forms of
time and space is the \emph{merely untrue}, which \emph{with eternal
necessity} is abolished in the Absolute. Therefore human life's forms of
nature and its natural content too become merely something unwarranted and
untrue; spirit's one absolutely true being becomes \emph{thought's
individuality-less, contemplative life in eternity's element}, in which spirit
knows the Absolute \emph{as the sole existent in its eternal process}, and
knows \emph{itself as a moment in this Absolute}, in whose knowing of itself
man's knowing of the Absolute, and of himself in the Absolute, is taken up.
Over against the view of the eternal metaphysical process --- which is not a
striving towards a goal but an eternal
% HEAD, p. 126: run-in halbfette (bold) Fraktur head followed by an em-dash,
% NOT letterspaced --- set exactly like „Ludvig Feuerbach.“ on p. 97. A short
% centred rule stands above it; not transcribed. The Indhold instead reads
% „2) den Hegelske Philosophies“; the printed head is kept in the body.
% EMPHASIS p. 126: of the three „al“ the first two are letterspaced and the
% third is tight. The English follows.
% --- p. 127 ---
existent's self-affirmation, in whose necessity therefore every opposition is
eternally \emph{abolished} --- the problems of actuality belonging to
personality lose their significance and their sharpness: the problem of the
finite spirit's freedom of will; of the free will's realisation of the being as
naturally-spiritual; of the human spirit's relation to time, within which it
relates itself to a task of eternity, to the ethical-religious; of its
self-determination towards that task's solution; of its failure of the task; of
its guilt and its reconciliation of guilt. Over against the necessity of the
\emph{eternal and eternally completed} process, \emph{finitude's} movement ---
even where it concludes in \emph{evil} --- becomes only a jest, eternally taken
back into the Absolute's process, in which all the seriousness falls away.
Through \emph{the thought of the Absolute} the individual can take itself out
of finitude's contradictions, and in the individuality-less contemplative unity
with the Absolute it has \emph{the absolution} for all imperfection and
\emph{the reconciliation} with its ground of essence. In nature-determined
existence one is to live in the abstraction of pure spirituality; in finite
existence one is to live only in eternity's element, so that time is in time to
be taken back into eternity, so that the temporal relation acquires no
significance as actuality, does not become a moment in the ethical-religious
resolution to action and self-surrender. \emph{Thought is action}; and to
\emph{comprehend} the Absolute is \emph{the absolute action}. Life becomes an
illusion of spirituality, \emph{ethics} is transformed into abstract
\emph{metaphysics, religion} into a \emph{comprehending of the idea}.

\emph{The difference} between this apprehension of the task of life and the one
indicated above is determined by the difference \emph{in the fundamental
concept}. The apprehension developed earlier of the Absolute's
\emph{essential} relation to the determinations of \emph{reality}, of the
\emph{essentially valid} real's \emph{essential forms of be}ing, leads with
necessity, on the one side, to maintaining the \emph{nature-determined} finite
spirit's \emph{relative independence}, and on the other side to laying a
decisive emphasis upon \emph{the conditions of reality's}
% EMPHASIS p. 127: two more runs stop inside the word --- „Realitets“ (with
% „bestemmelserne“ tight) and „væsentlige Være“ (with „former,“ tight on the
% next line). „Spøg“ is NOT letterspaced, though „Onde“ just before it is.
% --- p. 128 ---
significance for ethical acting, and generally for the determination of the
human task of life. This difference will appear in its concrete form through
the following developments of central ethical-religious problems. ---

From the critical comparisons we now pass over to these developments, in which
essential determinations of concepts that have been anticipated in the
foregoing are to find their justification, and what has only been indicated is
to acquire its more concrete determinateness. From the nature of the case the
problem of \emph{the will's freedom} is that whose solution must first be
attempted.

\section*{1. On the Freedom of the Will.}

The problem of the will's freedom leads us into the domain of metaphysics, of
physics, of psychology and of ethics, but must find its definitive solution in
the last.

The concepts which give the first, \emph{abstract points of attach}ment for the
determination of freedom are the concept of \emph{being as activity} and the
concept \emph{selfhood}. In the concept of being itself the determination of
activity lies enclosed, since all being asserts itself as being, preserves its
being, only by activity, and without this determination of activity would
immediately sink back into not-being. But this activity, which is inseparable
from being, expresses being's \emph{peculiarity}; through it, therefore, every
being makes itself felt \emph{according to its determinateness}, and posits in
the reciprocal action something proceeding from its own essence, expressing its
own essence.

In this determination of a peculiar \emph{expression} of essential activity we
have the first intimation of freedom; but a step nearer to the concept of
freedom we come through the concept \emph{selfhood}, subjectivity. The concept
of selfhood expresses a being-for-itself through its other; selfhood, as the
expression of the \emph{ideal}, is, according to the relation of the concepts,
what determines in relation to otherness and what converts
% HEAD, p. 128: centred halbfette (bold) Fraktur display head; a short centred
% rule stands above it, not transcribed. p. 128 carries NO footnote, contrary
% to the OCR-derived hint list.
% EMPHASIS p. 128: „abstracte Tilknyt“ is spaced and „ningspunkter“ tight ---
% another run stopping at the line break. The English divides "attach-ment" at
% the same joint.
% --- p. 129 ---
every determinateness by an other into self-determinateness. But the concept of
selfhood too still affords only \emph{abstract} points of attachment for the
concept of \emph{the will's freedom}; it is not sufficient as a foundation for
that concept, nor for maintaining the finite existent's freedom in relation to
the infinite and Absolute. Selfhood is the fundamental determination for
\emph{all} that is organic, for \emph{all} that is living, and therefore also
for plant life and animal life. But even animal life's highest, most developed
forms come under the determination of instinct: that abstract, metaphysical
determination of freedom which lies in selfhood therefore becomes, in real
actuality, in the existence of the single animal life, enclosed and hidden by
necessity. The finite selfhood's relation to the real which is not posited by
this self, and does not belong to its essence, is not immediately determined by
the universal metaphysical relation between selfhood and otherness. And what
metaphysics lays down with respect to this relation --- which is apprehended
rather as a relation between the fundamental moments in the absolute principle
and in the \emph{total} existence determined by the principle --- can still
less immediately secure the finite, individual selfhood in its relation to the
absolute subjectivity.

Through the points of attachment indicated there are thus reached indeed
moments in the concept of freedom, but not the concrete determination of
freedom as \emph{the will's freedom}, nor a maintaining of this concept's
reality. In order that a securer foundation may be found for the determination
of the human will's freedom, a new determination must come to the concept of
selfhood, namely the determination of \emph{essential universality}. In this
determination lies, as shown earlier, the determination of \emph{inner
infinity}, and thereby the possibility of \emph{a being-for-itself in the form
of self-consciousness and of comprehending cognition}. The working of the
selfhood-energy becomes, with essential universality for its presupposition, a
\emph{conscious} working --- which is the condition of its being able to be a
working of \emph{will}; and since it is the expression of the energy in a being
with \emph{inner infinity},
% --- p. 130 ---
there shows itself a possibility that it can be a \emph{free} will's
self-determining.

Freedom would then express, more closely, that the human spirit, in its working
through the will which expresses the being's infinity, takes the initiative for
the will with consciousness \emph{within itself}, and gives the will its
content out of \emph{itself}. What held generally of every existent --- that as
existing it must be active out of its being's peculiarity --- would then here
take the form that spirit, in virtue of its inner infinity, posits its content
of will in its self-determining, and, where the self-determination was the true
one, gave the will a content which expressed the self according to its truth,
i.e. according to its essence, and therefore an essentially valid,
reason-determined content. In this determinateness the will was then
essence-willing, i.e. self-determining out of its universal essence. But this
essence, which is seen as self-determining in its willing and as having itself
for the content of its self-determination, is according to its primitive
conceptual determinateness the self-conscious and cognising essence: in that
true form of self-determination there must therefore be a unity of the
fundamental determination of the theoretical and of the practical spirit, and
freedom thus became the highest expression of the human being according to its
totality, became this total being's true form of actuality, the τέλος of
cognition and of will.

In the determinations given, what conditions the possibility of freedom of will
has been indicated, and freedom's relation to the being intimated. But when the
human being and the human will are seen not merely according to the universal
concept of the being as it were in isolation, and according to the will's
fundamental relation to the being, but are considered under the conditions of
reality and in actuality's concrete connection, there arises a doubt whether
what showed itself as a possibility under abstraction from the concrete can be
maintained as valid under actuality's real conditions. Only when this doubt is
removed does the developed concept of freedom acquire more than a merely
hypothetical validity.

What awakens doubt about freedom's actuality is on the one side the
individual's relation to \emph{the Absolute},
% GREEK p. 130: τέλος, antiqua italic with the acute --- the only Greek in
% pp. 120--133. Neither OCR witness reads it; read off the image at 400 dpi.
% EMPHASIS p. 130: the run „i sig selv“ takes the preposition with it, while in
% „og ud fra sig selv“ the words „og ud fra“ are tight. The English runs fall
% on the same members.
% --- p. 131 ---
on the other its relation to \emph{the real conditions of existence}. The first
of these relations will receive its proper elucidation through the treatment of
the problem of evil's possibility; the second is here to be treated from the
various standpoints from which it shows itself.

The individual's relation to its conditions of existence is namely to be seen
partly under the standpoint of a relation to \emph{the outer}, partly under the
standpoint of a relation to \emph{the conditions of development}. Under the
determination of \emph{the relation to the outer} we gather together: the
individual's relation \emph{to the outer real}, into whose causal connection it
is ordered; the relation to \emph{the outer within the individual itself}, i.e.
to its \emph{determinateness by nature}; and the relation to \emph{other
individuals}.

\emph{The first} relation --- the relation to \emph{the real's causal
connection} --- has been accentuated above all as a ground of doubt against
freedom's actuality, and not merely from materialistic \emph{systems}, for
which the human will becomes only a member in the mechanically connected
infinite series of material causes and effects, but also from such systems as
separate a sphere of the noumenon from that of the phenomenon and concede an
intelligible freedom, but urge that man, as phenomenon in time and space, is
unconditionally subject to the necessity of the law of causality, and therefore
in his willing is with necessity determined by an \emph{outer motive}. Although
it is not here overlooked that the presupposition upon which the necessity of
all causes' effects rests is each thing's inner essence, and although there is
talk of \emph{inner} incitements (in the organism) and of an \emph{inner}
prompting (in animals, through instinct), it is nevertheless held fast ---
since the intelligible act of freedom is to have posited that fundamental will
which constitutes the essence, as an original, \emph{unchangeable} unity ---
that \emph{every} motive to a movement of will, to a determinate act of will,
that the whole ``matter of the will'', is to be given by \emph{the outer
world}. The abstract, merely thought motive too is to be just as fully an
outer cause determining the will as the intuitable one, which consists in a
% EMPHASIS p. 131: „Det første“ is spaced but the following „Forhold:“ is
% tight, so the sentence carries two runs. The English keeps them separate.
% --- p. 132 ---
really present object; in the last instance it is nevertheless always to be
something real, something material, in so far as it always rests upon an
\emph{outer} impression received at some time and in some manner --- and to
this is reckoned also transmission through word and writing from the remotest
times (Schopenhauer). For this conception it then becomes just as impossible as
for materialism to think the human will's freedom; in that freedom, as a
capacity for an \emph{absolute beginning}, it can see only what would
constitute a conflict with the universal world-order.

Now as regards this so strongly accentuated ground of doubt, it must at once be
pointed out that that mechanical conception of nature's connection --- which
acquires its validity also within the form of idealism characterised --- suffers
from the unavoidable defect that a \emph{primitive} mover, a \emph{definitive}
originator of the movement given \emph{as a fact}, \emph{cannot} be found
without an inconsistency. \emph{If all movement is only movement communicated
from outside, then that means that it is nowhere originally present}: an
infinite increase of the members becomes only a tautological repetition of the
same relation, and \emph{it becomes a contradiction to acknowledge movement,
activity, as factual}, since it can have no ground in the material and none
outside it; for that which was thought as standing outside the series of
material causes would, as essentially opposed to the material, be unable to
exercise an effect upon it. If movement, activity, in nature is to be
explained --- indeed even merely acknowledged as factual --- it becomes
unavoidably necessary to go back to the concept of a \emph{totality-activity},
and then to see the effect of the totality as resulting from the working of the
\emph{totality's} members, since the very concept of a \emph{totality's} working
points to \emph{relatively independent individual centres of activity}. When,
then, the concept of an \emph{active totality of nature} has stepped into the
place of the concept of a \emph{mechanical connection of nature}, then the
relation indicated above between the self-less and the self --- more closely,
the self-conscious self --- is to be ac\-%
% EMPHASIS p. 132: „Totalitets“ is spaced with „leddenes“ tight, whereas
% „mechanisk Natursammenhæng“ carries its letterspacing across the line break.
% --- p. 133 ---
centuated, as is also the concept of the human self's centralisation in itself
as a relative totality with inner infinity and therefore with a peculiar centre
for conscious activity. The concept which for mechanism showed itself
immediately as an absurdity --- the concept of a \emph{primitive beginning} ---
must \emph{necessarily} come forward when the individual being which has inner
infinity concentrates, centralises itself in itself as a relative totality, and
thereby, as taking possession of itself, in its constituting of itself posits a
new point of beginning, a peculiar relative centre in existence's total
organisation. How this working of the relative totality is to be apprehended in
relation to the \emph{absolute} principle will be elucidated in what follows.

In the self's essence, in its ideal doubling, in its character as energy, lies
the possibility of a self-determining whose motive falls \emph{within} the
self, in its concrete content of essence; in the relation between the self
which according to its foundation is \emph{eternal} and its unfolding in a
\emph{development in time} lies the possibility that a determining which from
the phenomenon's side appears as a temporally determined \emph{change}, and
which was therefore asserted to have to be effected by something \emph{outer},
can proceed from the self's eternal ground. And since the self exists as
ordered into a totality of existence, and at every moment receives influences
from \emph{the outer}, this relation too does not abolish, for the
self-conscious self, the possibility of a free self-determining.

For with respect to the influences from \emph{the outer real} --- which, as
influences upon the will, do not become mechanically working causes but
\emph{motives} --- there lies in the first place in the human being's
\emph{universality the possibility of abstracting from the determinate motive}.
In virtue of abstraction's inner infinity freedom's possibility is preserved,
even though for the present only in \emph{the form of negativity}, as a free
not-willing of what presses in from outside; and this freedom in the form of
negativity has as its outermost limit a free surrender of individual existence
in order to maintain freedom. But that uni\-%
% PRINTER'S DEFECT p. 133: the page prints „abstrabere“ for „abstrahere“ ---
% a b where an h is wanted. Orthographic, and so not reproducible in English;
% recorded here.

% --- p. 134 ---
\setcounter{footnote}{0}%
versality of the being --- which through the capacity for infinite abstraction
grounded freedom's possibility in the form of negativity --- encloses at the
same time a \emph{positive} determination of freedom. When the outer motive's
determining is warded off by the will, this happens in virtue of \emph{inner
motives, i.e. motives proceeding from the being}, which form the foundation for
abstraction's act; and what determines in the act of will therefore becomes a
determination in \emph{the self's essence}. The will is thus determined by what
belongs to that essence which the will \emph{itself} expresses; and since the
essence becomes transparent to consciousness in this determining, the
\emph{inner necessity of essence transforms itself into conscious freedom}.
Finally it lies in the being's \emph{universality} that the \emph{outer motive},
where an \emph{essential unity} takes place, \emph{can be transformed into an
inner one, i.e.} converted into a moment in the determination of essence.
Throughout, then, it is held fast \emph{that the act of will results with
necessity from motive and will}; but since the motive becomes an \emph{inner}
one, freedom is maintained --- not in the sense of a groundless and therefore
meaningless \textit{liberum arbitrium}, but in the sense of the true,
essence-determined freedom, that freedom through which the will, in being
determined by motives proceeding from the being, shows itself as an expression
of essence, in unity with the essence.

In the foregoing regard has been had to the will's free \emph{resolution, to
free acting in the inner, to freedom's inner realisation}. As regards the free
will's \emph{outer realisation}, this is --- when it is seen only \emph{in
relation to the individual} will --- to be determined as something essentially
independent of the resolution of freedom. The relation shows itself otherwise
when the relation between the will and its execution in the outer, when the
free ethical acting's relation to the course and order of nature and of
history, is seen as mediated by the \emph{universal} --- not as an ordering
third standing outside (Kant), but as the immanent harmonising
principle\footnote{The question how far the free will's non-realisation in the
outer, as hanging together with a culpable defect of cognition, can become an
object of ethical responsibility also from the standpoint that such
non-realisation may react upon freedom's \emph{inner} realisation and hinder
the development of freedom, falls outside the present investigation.}.
% printed mark: *) --- the only note on p. 134; it runs over onto p. 135 in the
% print and is set here in full at its call.
% ANTIQUA LATIN, p. 134: „liberum arbitrium“ is upright antiqua and CLOSE-SET,
% not letterspaced --- unlike the spaced antiqua of p. 68.
% NAME NOT LETTERSPACED, p. 134: „(Kant)“ is ordinary close-set Fraktur,
% neither antiqua nor spaced. So no \emph{} here, contrary to expectation.
% SPLIT SPACING, p. 134: the run „i Forhold til den individuelle“ stops at the
% line end; the „Villie,“ that completes the phrase is close-set. The English
% run ends at the same point.

% --- p. 135 ---
When the free acts of will which proceed from the primitive centres of activity
step out into outer action, the action comes \emph{as phenomenon under the
necessity of the law of causality}, and from the co-operation of the free acts
of will there results in the historical \emph{the necessity of universal
action}. Here a \emph{conversion of activities} takes place, like that which
takes place among the activities of nature among themselves; and under this
conversion the activity's goal is preserved unchanged.
% PRINTER'S DEFECT, p. 135: the page prints „Her findes en Virksomhedernes
% Omsætning Sted“, where the idiom requires „finder … Sted“ --- and „der finder
% Sted“ occurs correctly in the very next line, giving a control. Idiomatic and
% Danish-internal; recorded here rather than reproduced.

In the foregoing the will's \emph{relation to the outer} has been seen as the
relation to the real existing \emph{outside the individual}: the single real
and the totality of the real. It is now to be seen as a \emph{relation to the
outer within the individual itself, i.e. to the individual's side of nature}.
Setting aside what will be treated under the standpoint of the
indi\emph{vidual's relation to its conditions of development}, there are here
% ONE run in the Danish, beginning MID-WORD: „Indi-“ at the line end is
% close-set and „videts Forhold til sine Udviklingsbetingelser“ is spaced. The
% English begins the run at the same point in the word; an earlier draft split
% it into two and the parity check caught it.
to be considered the relation to \emph{the influences from the untransparent
natural element in the individual}, and the relation to those \emph{hindrances
to freedom} which may lie in original determinatenesses of nature: endowment,
temperament and the like. As regards those influences --- which hang together
with the abiding separation between \emph{the conscious} and \emph{the
unconscious in the individual's life}, and which may be designated as
\emph{promptings of drive}, if we take the concept \emph{drive in a wide}
sense --- the concept of \emph{essential unity} and of \emph{the totality of
essence} becomes with respect to them decisive for freedom's maintenance. Since
what proceeds from \emph{the unconscious} proceeds from the \emph{same}
undivided essence as what comes forward in \emph{the life of consciousness},
the necessity lying therein becomes --- so far as we do not reflect upon the
determinateness possibly working through it from outside --- an \emph{inner
necessity of es}\-%
% SPLIT SPACING, p. 135: „Indi-“ at the line end is close-set while the
% continuation on the next line is spaced --- the mirror of the usual case. The
% English divides "indi-vidual's" so the run can begin at the same point.
% --- p. 136 ---
\emph{sence}, and therefore a \emph{free necessity}. And when the determination
of \emph{the totality of essence} and of \emph{essential universality} is then
more closely brought to bear, what proceeds from the individual's side of
nature is seen as taken up, through the organic relative \emph{totality, into
the spiritual} being's unity of life, as gathered together in that being's
ideal concentration; and the inner necessity of essence is clarified into the
\emph{spiritual} being's inner necessity of essence, i.e. into
essence-determined freedom. Over against the single determinateness, \emph{the
totality of essence} and \emph{spontaneity} then make themselves felt, and the
consciousness of the totality preserves its significance even when it appears
in the form of abstraction.

As regards what lies as \emph{hindrances to freedom} in the individual
endowments, in the temperament and the like --- which as something determined
by nature, something given by nature, goes \emph{before} the individual's
conscious life of spirit and is presupposed by it --- there is, in relation to
these individual differences, however strongly they may appear, the human
being's \emph{abiding universality} to be held fast, and thereby the validity
of the universally human for \emph{every} human individual, in spite of the
intensity of the differences. But with the \emph{fundamentally human} there is
given that \emph{ideal relation to itself} which is freedom's fundamental
determination; and this relation, which is the human individuality's ideal
\textit{prius}, is realised for consciousness in the ideal concentration (see
below, pp.\ 139 f.), in which the individual constitutes itself as personality.

What holds of the natural differences in the individual generally holds also of
that peculiar form of them which comes forward in the \emph{abnormal physical}
and \emph{psychical} states. In the individual's relation to these states,
which hinder the will, there lies partly a relation to the outer within the
individual itself, partly a relation to the real outside the individual, in so
far as a disharmonious influence of this real upon the organism conditions
those states. With respect to these states the standpoint is to be brought to
bear that the continued existing, in spite of these states, of the human
individual endowed for a
% ANTIQUA LATIN, p. 136: „prius“ is upright antiqua, close-set.
% p. 136 carries no footnote and no Greek.
% --- p. 137 ---
life of spirit and of freedom involves the possibility of a liberation --- a
possibility therefore present \emph{so long as the individual exists as an
individual of this kind,} as a \emph{human} individual. ---

When finally the individual's relation to the outer is seen as a relation to
\emph{the other individuals} as acting, to the effects of the other
self-determining wills, then this relation comes on one side immediately under
the same standpoint as the individual's relation to the influences of the outer
real; on another side there shows itself a \emph{nuance in the rela}tion, in
that the other individual has, not merely through the unconscious power of
likeness of essence but through its cognition, the possibility of an influence
embracing the will far more firmly than the outer real has. But the difference
becomes nevertheless \emph{only a difference of degree}, which does not alter
the fundamental relation in what is essential.

On the other hand, in the individuals' mutual relation the \emph{supplementary
relation of unity} between them acquires an essential significance for
freedom's realisation in the individual --- which will be more closely
elucidated in what follows, under the development of the problem of
reconciliation. The supplementary relation of reciprocal action between the
individuals --- in which, as in the realisation of self-consciousness, the
relation to the other I \emph{as the species' representative} acquires its
significance --- will be shown as the condition for freedom's realisation in
the individual: a realisation which takes place in the \emph{ethical} relation
to those spiritual individualities by which the individual is integrated, and
which points back to a \emph{presupposing of freedom by itself}, and finally to
its last ground in that which gathers together the mutually integrating
individualities in spirit's unity: to its ground \emph{in the absolute
principle}.

When freedom's significance as actuality has thus shown itself to be preserved
under the relation to the outer, it is now, according to what was indicated
above, to be maintained \emph{in the relation to the conditions of
development}. With respect to this relation a twofold standpoint makes itself
felt, which leads to a twofold aporia. For since freedom --- in the pregnant
sen\-%
% SPLIT SPACING, p. 137: two places where the compositor drops the spacing at a
% line break --- „Nuance i For“|„holdet“, and the „som“ between „af denne Art,“
% and „menneskeligt“, close-set although both neighbours are spaced.
% --- p. 138 ---
se as \emph{spirit's freedom} --- cannot be immediately given to the
nature-determined individual, but must be \emph{acquired by the individual},
and therefore acquire a \emph{becoming} through the individual's development in
time, the difficulty shows itself first when this becoming is reflected upon.
According to spirit's concept freedom must be posited as that which is to
become, because it is to be acquired; but since it can be only as
\emph{self-acquired}, and there\emph{fore must be through itself, and therefore
must presuppose itself}, this seems to exclude a becoming, since such a
becoming, as freedom's becoming, would presuppose \emph{unfreedom} as that out
of which freedom was to become. And the difficulty shows itself \emph{secondly}
when the human spirit's development in time is seen \emph{as an approximation to
the being's ideal} --- when therefore freedom, as the expression and the
realisation of man's true essence, is seen under the standpoint of \emph{the
ideal}, as an object of \emph{an approximation}; for freedom then acquires the
stamp of that which is realised \emph{only approximately}, and which never in
the full sense acquires actuality --- which, if only the full and comprehensive
realisation of freedom rightly bears freedom's name, would be one with
\emph{freedom's acquiring no actuality at all}. This difficulty lying in the
approximation appears in a more determinate form when it is reflected that the
realisation of freedom of will has \emph{cognition} for its essential
condition, and therefore shows itself as conditioned in its perfection by
cognition's perfection --- which, when cognition is taken in its most concrete
sense, is only approximately attainable; so that freedom, on the presupposition
indicated above, would never acquire the complete condition for its
realisation.

As regards the difficulty which lies in \emph{freedom's presupposing of
itself}, its solution is to be sought above all in a concept which hovered
before Kant when he developed his doctrine of the timeless \emph{intelligible
act of freedom} --- but which with him, in virtue of his abstract separation of
the sphere of the noumenon from that of the phenomenon, received an almost
mythical clothing: namely in the concept of \emph{the spiritual}
% NAME NOT LETTERSPACED, p. 138: „Kant“, like the „(Kant)“ of p. 134, is
% close-set Fraktur. This book letterspaces most names, so the two Kants are a
% real exception and carry no \emph{}.
% SPLIT SPACING, p. 138: „alt“ at the line end is close-set while the
% continuation is spaced. The English divides "there-fore" at the same joint.
% --- p. 139 ---
\setcounter{footnote}{0}%
\emph{individual's concentration of the whole of its development in an ideal
point}, which abolishes the development's temporal externality and necessity.
This means, in other words: the solution is to be sought in the concept of
\emph{a breaking through}, in which \emph{the eternal} --- which in the
individual self has covertly determined that self's development --- joins
itself together in its unity as the presupposition of its own being. And this
moment of concentration coincides with \emph{the moment of the ethical choice},
in which freedom, through the breaking through of the eternal, first acquires
true actuality. But beside this decisive concept there acquires --- as was
already indicated above --- the single human individual's relation \emph{to the
supplementary personalities and to the divine}\footnote{Just as to that concept
of a concentration of the development in a point of eternity there is a
parallel in \emph{Aristotle}, in the apprehension of form (εἶδος) as that
which, at the goal of the development which it itself determines, can be said
only to \emph{be, not to have become}, so there is in him a parallel to the
thought of freedom's realisation through the relation to the supplementary
individualities and to the principle, in the determination that the moral's
realisation in the individual presupposes an \emph{already occurring actual
being of the moral}, just as all becoming generally presupposes something
actually existent. That determination makes itself felt in his well-known
apprehension of \emph{the state's} significance for \emph{individual}
morality.} its essential significance for the aporia's solution.

The solution of the difficulty which lies \emph{in what is approximative in
freedom's realisation} must be sought through this, that \emph{a twofold}
standpoint is brought to bear for the apprehension of freedom, in that freedom
is seen at one and the same time \emph{as task and as existing as an essential
determination}. That freedom is seen from this twofold standpoint is grounded
in its essence, and generally in the essence of the ideal. Since freedom, like
everything ideal, is \emph{existing through itself}, it is only as
\emph{constantly self-produced}. In the highest relation too it holds of the
divine actuality that it, as living, ideal actuality, is an eternally
self-producing actuality, at one and the same time existing and τέλος for a
working.
% printed mark: *) --- the only note on p. 139.
% GREEK, p. 139: τέλος in the body and (εἶδος) in the note, both antiqua italic
% with accents, read off the image at 4--5x. Neither OCR witness renders any of
% it. Copied verbatim.

% --- p. 140 ---
\setcounter{footnote}{0}%
But in the relation of human freedom there comes in that determination in which
the difficulty lies: that the self-production is a \emph{temporally determined}
self-development. When freedom therefore becomes a task for itself, it must as
task become something fuller than it was as presupposition; and since freedom
is a determination of \emph{eternity}, and since it is \emph{an ideal} not
merely for \emph{the individual} but for \emph{humanity}, there shows itself
\emph{an infinite possibility of development}. But when this has been laid
down, then it is on the one side to be held fast that the presupposition is
nevertheless \emph{in its essence} equal to the task --- that therefore the
will which is still realising its freedom \emph{is in principle free}, in spite
of its relative abstraction --- and on the other side the standpoint indicated
above is to be held fast, that the free will at every moment of realisation
joins its own past together in an \emph{ideal point of unity}, and thus takes
the development in time ideally \emph{back into a moment of eternity}, in order,
in the last instance, to concentrate itself in ideal eternal unity, as an
organic moment in the all-comprehending, to concentrate itself into the ``willing
of one thing'', understood in eternity's sense. Thus what characterises the
ideal is seen at every point of the development.

When it is therefore conceded that freedom's realisation, since it must take
place through a development towards an ideal which is not merely an individual
but a universally human one, must in virtue of this determinateness of the
ideal be subject to an infinite approximation, then at the same time --- since
the relation to the ideal is the relation to the will's τέλος --- what is at
work in the development, its entelechy, becomes \emph{the ideal itself as} an
energy of development, i.e. as the will's actual determination of
essence\footnote{Cf.\ for the rest what will be developed in what follows,
3, c, on the ideal's realisation.}.

Herein lies, then, that degrees \emph{in freedom} are indeed laid down; but in
the sense that freedom, in constituting itself, unfolds itself \emph{as}
freedom from the relatively more abstract form to the more concrete one; and
this is again to be determined more closely thus, that it unfolds itself from
the \emph{abstract willing}
% printed mark: *) --- the only note on p. 140.
% GREEK, p. 140: τέλος again, antiqua italic with acute.
% SPLIT SPACING, p. 140: three instances --- „Evigheds“|bestemmelse; „i sit
% Væsen er“ stopping at the line end while „liig med Opgaven“ is close-set; and
% „selve Idealet som“ stopping at the line end while the rest is close-set.
% The English runs stop at the same points.

% --- p. 141 ---
\emph{of the being} --- a willing which is only as \emph{ethical} willing ---
into a \emph{concrete essence-willing}, into a conscious, resolved willing of
the being according to its full content. This gradual development, through
which the \emph{prin}ciple of freedom revealed \emph{in the resolution} (in the
choice) wins concretion, is at the same time a development \emph{in intensity},
since freedom's energy increases in the same proportion as the will becomes the
expression of the concrete essence's energy, of the essence clear to itself, of
the essence conscious of itself in its fullness.

This leads us back to the difficulty specially brought out above, which lay in
this, that freedom presupposed as its condition a \emph{cognition} subject to
constant development, since only through cognition is the possibility given of
the essence's clarity to itself. Here reference must be made to the relation
between the \emph{a priori} cognition, which encloses the a priori element in
the ethical, and the \emph{empirical} one. The a priori cognition has in its
\emph{relative abstraction a completedness in itself}, and by its relation to
man's \emph{essence a real validity}; and it is therefore such a cognition as
\emph{a free} willing can rest upon. When I, with consciousness of the good's
validity, of duty's demand, will the good, will duty, then my willing --- as
determined by the a priori, and therefore by my own ideal essence's
determinations, as a conscious willing of what is essentially valid --- is
\emph{a free} willing; and this free, essentially valid willing acquires its
filling out, limited by no immovable barrier, through that cognition which,
resting upon the a priori and having in it the abstract essence of the self and
of the real, penetrates into the empirically given real --- here therefore above
all into the real within ourselves, into our own essence's determinations of
reality and relations of reality --- as an assimilating cognition. In the course
of this development of cognition, then, that cognition is present \emph{at every
point} which makes the will free, even where the will relates itself to the
concrete that is not cognised.

\begin{center}\rule{2cm}{0.4pt}\end{center}
% RULE, p. 141: the page ends short and a short centred SINGLE hairline closes
% the argument below the last line, the rest of the leaf blank. No head follows,
% so this is a section-closing rule of the kind at pp. 89, 143 and 201, not a
% head ornament. Transcribed.
% PRINTER'S DEFECT, p. 141: „paa det Aprioriske ag i dette“ --- „ag“ for „og“.
% Orthographic, and so not reproducible in English; recorded here.
% SPLIT SPACING, p. 141: „Friheds“ is close-set and only „princip“ is spaced,
% within the one word. The English divides "prin-ciple" at the same joint.

% --- p. 142 ---
Since in the foregoing freedom's reality was maintained against the doubt which
rested upon the individual's relation to the outer, the decisive concept was
the concept of the human individual's concentration in itself as a relative
totality, determined by essential universality, with inner infinity; and this
concept acquired its essential significance in the developments that followed
as well. Through that concept of self-concentration the unfreedom in the
relation to the outer was abolished, and converted into self-determination out
of the whole of essence, in which the ideal totality of spirit is in unity with
the organic totality. But since now the individual, in this joining together in
its relative totality, is seen as determining itself by freedom, the working of
\emph{the individual} which expresses this self-determining is at the same time
to be seen as ordered \emph{into the absolute totality}. This holds whether
freedom is apprehended according to its untrue or according to its \emph{true}
form. If freedom is only the \emph{formal} freedom, \emph{arbitrariness}, the
abstract expression of inner infinity, the \emph{semblance} of
essence-determined freedom, the apparently groundless determining, then there
are precisely here determinate, finite, though not clearly conscious motives
present, which have connection with the real's totality and draw the will's
working into it. If freedom on the contrary is the \emph{true, essential}
freedom, so that its ``absolute beginning'' has an essence-determined,
substantial content, then there is, in a timeless relation, a connection with
the totality of existence and its eternal principle --- a determining which goes
in unity with that principle's eternal determining.

Man's relation of freedom to the outer thus points back, through the concept of
totality, to \emph{the relation to the Absolute}. And it is only by this
relation that the individual being --- which as free is for itself in the form
of infinity, in which the principle's being reflects itself within it --- can
have the possibility of withdrawing itself from the dominion of outer causal
necessity, yet in such a way that at the same time, by its very positing of an
absolute beginning, it is ordered into the principle-determined tota\-%
% p. 142 carries NO footnote --- the OCR-derived hint list wrongly lists it ---
% and no Greek.
% --- p. 143 ---
lity and subjected to that totality's dialectic. For when the Absolute is the
human spirit's principle, then man's freedom is to be seen \emph{as grounded in
it}, the only thing which \emph{in an absolute sense is free}; and
\emph{freedom in its truth} can then be only one with the \emph{free absolute
surrender} to the Absolute as the power that works through all. But when
freedom as \emph{man's} freedom is then to be seen as \emph{derived}, and yet
is to be maintained as freedom; when it is to be seen as being through the
Absolute, while at the same moment, in expressing man's relative independence,
it shows itself to enclose the possibility of separating itself from the
Absolute --- then there appear, through the relation to the principle, new
grounds of doubt against freedom.

When these grounds of doubt are to be examined, they lead us into the
investigation of the other problems brought out above. The problem of
\emph{derived freedom}, and of its \emph{twofold} form of actuality --- as
asserting itself in \emph{self-will}, and as independently joining itself
together with the principle in \emph{self-surrender} --- coincides essentially
with the problems \emph{of evil's possibility and of evil's abolition in
reconciliation}.

The problem of human freedom therefore leads us, when the new grounds of doubt
are to be removed, once more into the domain of the ethical-religious, to which
we were referred above where the individual's self-concentration showed itself
as identical with the breaking through of the eternal in that choice by which
the ethical was constituted in the individual. We therefore point back at this
point to the earlier --- anticipatory and therefore for the present only
hypothetically valid --- development of the ethical's essence, and particularly
of the will's constituting of itself as ethical will in the ethical choice, in
which that breaking through takes place, in which freedom gives itself
actuality and, in giving itself actuality, presupposes itself as an essential
determination and posits itself as an infinite task.

\begin{center}\rule{2cm}{0.4pt}\end{center}
% RULE, p. 143: the sub-section closes with a short centred rule below the last
% line, the rest of the leaf blank. A SINGLE hairline, as at pp. 89 and 141 ---
% not the double rule of pp. 29, 63, 201 and 210. It closes a division and
% stands next to no head, and so IS transcribed. The head on p. 144 carries no
% rule of its own.

% --- p. 144 ---
\section*{2. On the Essence and Possibility of Evil.}
The concept of \emph{evil} must, as the concept of something \emph{negative},
be made plain through the concept of the \emph{positive} which it negates. The
concept of \emph{the good}, as the concept of this positive opposite, must
therefore form the starting point for the apprehension of evil.

The concept of \emph{the good} has already been touched upon in the foregoing
(p.\ 76) in the development of the Absolute's idea. When the Absolute was
determined as existence's principle, it was determined as that which posits
real being as determined by the ideal; and it was at the same time seen as the
\emph{teleologically} working, and in this working as its own end, as the goal
for its subjectively-ideal activity, and as comprehending and governing, in its
eternal process of self-realisation, the totality of the real moments. In these
determinations of the Absolute lay the point of attachment for the
apprehension of the Absolute as \emph{the good} (the idea of the good) --- a
concept which developed from the more abstractly metaphysical towards the
metaphysical-ethical.

In opposition to the \emph{good in its metaphysical} determinateness, evil
becomes the \emph{metaphysically evil}. If now the good in this determinateness
expressed \emph{the idea}, the \emph{Absolute, as its own end, as the goal} for
its \emph{eternal} process of realisation, then \emph{evil} --- since the
idea's process encloses all being, and evil can therefore express only
something falling within this process --- can, in that it must be something
standing in opposition to the goal, acquire only the general determination of
being \emph{the moment in its isolation}: which as such, in opposition to the
absolutely existent, becomes something merely \emph{relative}, something
determined by \emph{negation} (determination), which fixes

itself for thought in its determinateness by a μὴ ὄν.

But this most abstract universal determination of evil now takes on a nuanced
form according to the different apprehension of the Absolute.

If the Absolute is apprehended from an \emph{abstract standpoint of eternity},
so that its process is seen only \emph{\textit{sub specie}}
% HEAD, p. 144: bold (halbfette) Fraktur, centred; agrees with the Indhold. No
% rule stands under it; the rule belongs to the foot of p. 143.
% GREEK, p. 144: μὴ ὄν --- grave on the eta, smooth breathing with acute on the
% omicron. Read off the image at 4x; copied verbatim.
% ANTIQUA LATIN, pp. 144--145: „sub specie æterni“ is antiqua AND letterspaced,
% like the „causa sui“ of p. 68 and unlike the close-set antiqua of pp. 64, 71
% and 134. Encoded \emph{\textit{...}}; it breaks across the page, as here.
% --- p. 145 ---
\emph{\textit{æterni}}, then evil --- which from this standpoint is to be seen
as a \emph{necessary} moment in the process, because the process is a process
only through the development of the moments as \emph{such --- is eternally
abolished} just as it is \emph{eternally} posited. It becomes the idea's
instrument, which from the idea's standpoint cannot be seen \emph{as evil},
because it is eternally drawn into the process, but only as an eternally
governed moment. There then comes in a \emph{contradiction}, since it can be
seen \emph{as evil} only when it is as it were regarded from below, from a
standpoint \emph{outside and beneath} the idea, and from that standpoint held
fast in the abstraction of isolation. The metaphysically evil, so apprehended,
acquires actuality only \emph{for human thought} --- which itself, apprehended
from this abstract standpoint, becomes, like human existence generally,
something \emph{merely vanishing}, something incomprehensible in its
essencelessness (cf.\ Spinoza); and evil therefore acquires only \emph{an
illusory being for the thought of something vanishing and essenceless}. And
while evil becomes unreal, the \emph{metaphysically good} comes, in virtue of
the standpoint's abstraction, to suffer from the defect that in its abstract
eternity it becomes something abstractly existent, whereby the determination of
end becomes illusory.

If evil is to become more than a merely vanishing semblance in an eternal
process, then there is required an apprehension of the Absolute by which it
becomes possible that \emph{real actuality asserts a relative independence in
its relation} to the divine.

On such an apprehension an em\emph{phasis must fall upon the concept of the
really finite}, and the Absolute's process is apprehended more concretely, in
that it is seen as comprehending the real in its temporal determinateness, so
that the real in this determinateness acquires an \emph{essential} significance
for the Absolute. Here finitude is relatively fixed in virtue of the
essentiality of the real forms of existence. But even when the finite is thus
secured against immediate vanishing into the eternal, and in its relative
independence gives a place for evil, evil will nevertheless, so long as it is
seen \emph{outside}
% NAME, p. 145: „(jfr. Spinoza)“ is Fraktur and close-set --- neither antiqua
% nor letterspaced --- so it carries no \emph{}, though the book letterspaces
% „Aristoteles“ in the note to p. 139.
% SPLIT SPACING, p. 145: „Efter“ is close-set and only the rest of the word is
% spaced, the break falling inside it. The English divides "em-phasis" at the
% same joint.
% --- p. 146 ---
\emph{the will}, or in the will as brought about by a process of
\emph{nature}, still essentially retain the character of the
\emph{metaphysically} evil, and be burdened with that evil's unreality. Evil
then becomes real existence's \emph{finitude, defect, imperfection}, the
τὸ μὴ ὄν in the existent, and is seen as such in the light of mere
\emph{necessity}, as grounded in the law of the idea's movement but at the same
time as governed by that law and drawn into the harmony of the whole; and
although as the ``\emph{physically} evil'' it acquires a more concrete shape
than as the metaphysically evil, it acquires, even when it is apprehended as
\emph{privation}, only a being as evil for that consciousness which, forgetting
existence's necessity, judges the single thing by an abstract universal concept
or by its relation to the individual as sensing. And this being of evil
acquires in the last instance significance as actuality only in that \emph{the
will gives it significance for the personality}.

The more concrete apprehension of the metaphysical thus points, when evil is to
be comprehended in its actuality and in its full opposition to the good,
through the real \emph{towards the will and the ethical}. Only in the sphere of
the finite spirit, where we get \emph{an opposition to the Absolute posited by
the will with freedom}, do we get an evil in the \emph{proper} sense: only as
the \emph{ethically} evil has \emph{evil actuality}.

Just as the determination of the metaphysically evil was elucidated by that of
the metaphysically good, so must \emph{the determination of the ethically evil}
be elucidated by that of the \emph{ethically good}. And since now the good's
metaphysical determination is to be converted into an ethical one, the relation
must be that the former determination continues to be the foundation. Just as
in the development of the Absolute's concept (pp.\ 74 ff.) we saw the more
concrete metaphysical-ethical determination proceed from the more abstract
metaphysical one, so must we, when the good is to be determined as the
ethically good, the humanly good, see this determination resting upon the
metaphysical one. The good must everywhere, in order to be in truth the good,
bear \emph{the idea's} determination within it: all good is good only by
expressing the idea. But the ethically good in its \emph{concrete}
determinateness relates itself to the \emph{hu}\-%
% TWO runs in the Danish, split by the page break: „menne“ closes p. 146 spaced
% and „skelige Villie“ opens p. 147 spaced. The English splits "hu-man" at the
% same joint and keeps both halves emphasised; an earlier draft dropped the
% first half and the parity check caught it.
% GREEK, p. 146: τὸ μὴ ὄν, antiqua italic with accents and breathings; ABBYY
% drops it entirely. Copied verbatim.
% SPLIT SPACING, p. 146: „Natur“ is spaced and „proces,“ that completes the
% word is close-set. As printed.
% A round brown ink blot follows „Onde Virkelighed.“ --- a stain in this copy,
% not type; not transcribed.
% --- p. 147 ---
\emph{man will}; it must therefore become \emph{an expression of the idea
through the will, an ideal of the will}.

When we thus let the relation to the human --- more determinately, to the human
will --- be that by which the good determines itself from the metaphysically
good into the ethically good, then the Absolute is seen as ethical ideal in
virtue of this relation to the will as the will's ideal; but since the relation
to the human will, like the relation to the real generally, is grounded \emph{in
the Absolute's own essence}, the determination of ethical ideal does not become
a predicate added to the Absolute as it were from outside, but an
\emph{original determination of essence within it}.

The \emph{first} determination of the \emph{ethically} good becomes then, from
this starting point, \emph{the realisation of the idea (the idea of the good) by
the free human will}. But this determination then acquires --- as has been shown
in an earlier place (Section II, 3) --- \emph{the second expression}, that the
\emph{ethically} good, as the \emph{humanly} good, is the realisation of
\emph{man's idea}, i.e. of \emph{man's total, ideal essence}. This
transformation of the expression was grounded in this, that the \emph{absolute
idea} as the norm of human action \emph{is in the determinateness of the human
ideal}; wherefore to realise the idea is \emph{for man} the same as to realise
\emph{his own ideal essence} through the free resolution of will, through the
essence-willing clarified for itself by cognition. \emph{The Absolute} becomes
the end in that \emph{the human being} becomes the end; and it has appeared how
ethics could not reach beyond this determination of the human as the human's
task, when the task is determined \emph{positively}, and how the human will, in
virtue of the human individual's fundamental relation to the whole of existence
and its principle, can express the Absolute \emph{as such only in the form of
negativity}, in the form of unconditional resignation and self-surrender.

The specific determination of \emph{reality} by which the metaphysically good
becomes the \emph{ethically} good is therefore \emph{the human will as
essence-willing}, i.e. as de\-%
% PRINTER'S DEFECT, p. 147: the Danish prints „Bemmelse“ for „Bestemmelse“, a
% dropped „ste“. Orthographic, and so not reproducible in English; the Trykfeil
% leaf does not list it either.
% ANTIQUA NUMERAL, p. 147: the „II“ of „(Afsnit II, 3)“ is set in antiqua roman
% capitals, not in the Fraktur I/J sort. Left unmarked in both files, since the
% house \textit{} convention is for Latin words; recorded for a later pass.
% SPLIT SPACING, p. 147: „Realitets“ is spaced and „bestemmelse,“ close-set,
% within the one word. The English divides the phrase at the same joint.

% --- p. 148 ---
terminate, resolved realising of the being. This human being, which in virtue
of essential universality is unfolded into the being-for-itself of
\emph{self-consciousness} and of \emph{freedom}, we have seen as
\emph{naturally-spiritual}, as realising its determinateness by ideality
through and upon the foundation of a \emph{real organism}; and we have seen it
as a \emph{relative totality} within the universal totality. Herein lie the
closer determinations of the humanly good, of the ethically good. This humanly
good --- in which the self wills the Absolute, the idea of the good, in willing
its own ideal essence --- is by and in an act of will which in the human forms
the analogy to the teleological energy of subjectivity in the idea's process;
and in willing this good the will wills its own essence in the normal relation
between that essence's being as a relative totality and as a moment in the
universal totality, between the spirit-side and the nature-side within it.

Out of the concept of the ethically good thus determined there is now formed,
through the opposition, the concept of the \emph{ethically evil} --- which had
only an abstract prefiguration in the sphere of the metaphysical, in that in the
idea's process the necessary positing of the moments as moments in their
relativity was held fast in abstraction; or, what expresses the same thing, in
that the negation, the not-existent, fixed itself unresolved for abstracting
thought.

The determination of the ethically evil therefore leads back to what was
decisive for the ethical, \emph{to the free will}, and must be sought from that
will, \emph{not} from the determinateness by \emph{nature} in the human, still
less from anything absolutely \emph{outside} man, be it a \emph{higher good} or
a \emph{higher evil} being. The first determination of it becomes, then, that it
is a \emph{disproportion in the spiritual life posited by the will}; and how
this disproportion is to be apprehended lies implicitly in what has been
developed above. For if the good was determined by the normal relation between
man's being as a relative totality and as a member in the universal totality,
and between the nature-side and the spirit-side within him, and if the normal
form of this last relation is,
% SPLIT SPACING, p. 148: „Natur“ is close-set and only „bestemtheden“ is
% spaced, within the one word --- hence Natur\emph{bestemtheden} in the Danish
% and the same division here.
% EMPHASIS, p. 148: „Villiesact“ is TIGHT, though ABBYY renders it spaced.
% --- p. 149 ---
\setcounter{footnote}{0}%
as needs no closer development, conditioned by the normal form of the first,
then the disproportion which is constituted by the evil willing, and which
expresses it, is a disturbance posited by the will in these two relations, and
\emph{originally in the first} of them --- therefore, according to its
\emph{primitive form}, a \emph{disproportion between the individual's own will
and the demand upon the will declared by the Absolute through man's essential
determination.}

This disproportion can now also be expressed in the following way. If the good
is the \emph{fundamentally essential}, the \emph{in-and-for-itself valid, which
is posited as such by the will, in that it is what is unconditionally willed,}
then \emph{evil in the ethical sense is the opposite of this fundamentally
essential, of the in-and-for-itself existent --- but in such a way that it is
posited by the will as the in-and-for-itself existent, the absolute, the
unconditionally valid}\footnote{In these determinations the moment of
\emph{objectivity} in the good and in evil is accentuated: they are determined
as something \emph{realised, actualised} by the will. But from the determination
of objectivity there is precisely, through this ``by the will'', a pointing back
to that of \emph{subjectivity} as the \emph{primitive} one; and objectivity
becomes such as is borne by the will, has its reality in the will. Just as the
ethically good is only through the will, so the ethically evil is only in and
through the evil will. One may indeed speak of an untrue realised in the outer;
but it becomes evil only through the relation to the will, to which it stands as
effect or as incitement.}. The evil will then not merely determines the
not-in-and-for-itself-existent as an in-and-for-itself-existent, and in this way
makes it its end, but it determines \emph{itself} through this determining. For
in willing evil as \emph{its own end}, it has evil for \emph{its own}
determination. It is determined by the transitory, the essenceless, the
\emph{finite}, in that it itself determines itself for that; but in being thus
finitely determined it has the determination of \emph{infinity} in the form of
abstraction and of negativity, and precisely therein the inner contradiction
\emph{that by the formal infinity it determines itself finitely} --- that it
therefore gives itself a determinateness incommensurable with what determines,
and which must therefore be dissolved again by the will.
% FOOTNOTE, p. 149: printed mark *). The OCR-derived hint list does not name
% p. 149; the note is there.
% EMPHASIS, p. 149: within the long spaced run only „saa bliver“ is close-set.
% --- p. 150 ---
\setcounter{footnote}{0}%
The disproportion in the self which is constituted by evil then also acquires
the expression that the personality in the evil willing \emph{by an act of
infinity determines itself to finitude, that it wills itself in
finitude}\footnote{This determination holds, as will be shown below, also when
the self wills itself according to its abstract, empty infinity.}, but wills
this finitude \emph{as an infinite.} In virtue of the \emph{formal} freedom it
therefore wills itself, but \emph{not} according to its \emph{essential, ideal}
determination, but according to an \emph{untrue} one; it wills itself therefore
\emph{in virtue of freedom}, but \emph{not as free} --- rather, unconsciously,
as something unfree, something governed.

In the given conceptual determinations of evil --- which hold everywhere where
the will has got beyond the merely instinctive determinateness, beyond the form
of immediate drive, and where the consciousness of the ethical demand has been
awakened --- there lies the possibility of \emph{two main forms} of it, which
are nevertheless not simply co-ordinate, but of which the one must become the
proper fundamental form. They correspond to the two essential relations whose
normal form constituted the good: the relation between man as a member of the
totality and the principle of totality, the Absolute, and the relation between
the nature-side and the spirit-side in man. That form of \emph{evil} which
corresponds to the first, \emph{primitive} relation is \emph{self-will}
(defiance); that which corresponds to the second is \emph{degradation,
perdition. Defiance} is \emph{the fundamental form}, because the relation to the
Absolute is the primitive one, and because man \emph{as spirit in every} form of
evil asserts an \emph{active} willing --- so that \emph{degradation} too, even
in its \emph{maximum} as animal spiritlessness, can be only \emph{as willed},
and therefore as enclosing \emph{defiance}. Even when the will sinks back into
relative unconsciousness this standpoint holds good: not merely is \emph{the
state as a whole} incurred by the will, which has as it were by a single act
posited this state and thereby bears the responsibility for what happens within
it, but in \emph{every} actual expression of the state there is, because man is
\emph{spirit}, a \emph{renewed} willing; and even
% FOOTNOTE, p. 150: printed mark *); the hint list does not name p. 150 either.
% EMPHASIS, p. 150: „formelle“ is spaced but the following „Frihed“ is
% close-set, and „primitive“ is spaced but the following „Forhold“ is not.
% „sig“ in „vil den altsaa sig“ is close-set, though ABBYY renders it spaced.
% --- p. 151 ---
the apparent suffering is to be seen as an acting, as \emph{defiance's acting.}
That there is in the relation between evil's two main forms something
dialectical --- according to which, just as all degradation and weakness
encloses defiance, so conversely all defiance has a moment of weakness, of
degradation, within it --- does not contradict the fact that defiance is evil's
primitive form. As such it has also been apprehended by the religious
consciousness, which in the representation of Lucifer and of the Fall seeks
evil's source in pride and self-will. In the following section, on evil's
abolition, the development will therefore keep essentially to this fundamental
form of evil, which is accordingly apprehended as present in \emph{everything
unethical}, so that this must in \emph{all} its phenomena conceal
\emph{defiance, self-will, egoism, hatred.}

\emph{The possibility and the ground} of the ethically evil so determined we
seek in that very fundamental determination in the human being according to
which the human being, in its universality, as an ideally-real member in
existence's organic whole, is a relative totality subsisting in an essentially
valid separation of reality, which has the possibility of consciousness --- if
only in the form of relative abstraction --- of its inner infinity, and of a
free being-for-itself. But this \emph{possibility of the individually ideal for
an infinite being-for-itself is a possibility doubled within itself:} it is the
possibility of the \emph{true, filled-out being-for-itself in surrender to and
appropriation of the universal that works through all,} and the possibility of a
\emph{closing off within itself in its essential separation of reality from the
Absolute, in an inner infinity which, in the separation from the truly infinite,
becomes an abstract, empty infinity, finite through its being torn loose,} and
therefore a self-contradictory \emph{infinity}, whose illusory filling out must
be sought in what is for it contingent, which has only \emph{finitude's}
determinateness. In that \emph{first} being-for-itself the individual wills
itself as ordered into the whole, subordinated to the principle; it wills itself
therefore as that principle's peculiar

organ, and only so can it --- which lies in the fundamental relation --- will
itself according to its own essence. But since
% RULE, p. 151: the page carries NO rule at its foot, and none is expected ---
% the subsection does not end here but runs on into p. 152, and the page breaks
% in the middle of a sentence.
% --- p. 152 ---
the individual wills itself as an organ for the unconditioned, it wills
\emph{the good}; and this is through that will which is \emph{the divine's
willing in the individual in true unity with the individual's will.} But in that
\emph{other} being-for-itself, in which the individual wills itself through its
own self-will, it wills itself not as a member in the whole, but wills itself in
its finitude as an absolute; it therefore negates the principle, \emph{does not
will it}, but asserts its willing of itself against the principle --- but as the
\emph{evil} will; and by this evil self-will, which negates the \emph{good}, it
negates \emph{its own essence.}

In the \emph{selfsame} fundamental determination, then --- in the determination
of the individual relative totality conscious of itself in its infinity ---
there lies both the possibility of the \emph{good} and that of \emph{evil}: the
possibility \emph{of a being-for-itself in the Absolute, of the
self-affirmation of self-surrender,} and the possibility \emph{of the abstract
being-for-itself of isolation, of the self-affirmation of selfishness;} and the
self has according to its concept \emph{freedom of choice in relation} to this
twofold possibility. But through the determinations that have been given of
evil's essence and possibility it has already been indicated how evil \emph{in
its very essence} bears \emph{the germ} of a \emph{self-dissolution}; and
through the demonstration of that germ and of its unfolding
\emph{reconciliation's possibility} will reveal itself.

\section*{3. The Self-Dissolution of Evil as the Possibility of Reconciliation.}

In developing this point we start from the more abstract metaphysical-ethical
determinations of concepts, which will then in the following parts acquire a
more concrete form.

The \emph{self-affirmation of selfishness} by which the \emph{ethically evil}
was posited --- and which we have seen under the standpoint of a
\emph{displacement} of spirit's fundamental relation, as a willing of the finite
as the absolute and as a willing which negates the Absolute's willing --- showed
itself to have its possi\-%
% HEAD, p. 152: bold (halbfette) Fraktur, set full measure; agrees with the
% Indhold.
% RULE, p. 152: there is NO printed rule between the close of the previous
% subsection and this head. The mark standing in that gap is modern PENCIL ---
% a bracket with an oblique lead-in stroke, pale grey and quite unlike the
% printed rule of p. 159. Not transcribed in either file.
% --- p. 153 ---
\setcounter{footnote}{0}%
bility in the self's capacity for a \emph{willing of itself according to its
abstract infinity.} But this last determination encloses a peculiar dialectical
relation to the \emph{truly infinite}, the universal and absolute; and by this
relation the evil will bears its own negation within it.

The self-affirmation of selfishness is through the \emph{formal infinity} in
virtue of which the self closes itself off within itself as absolute. But this
infinity the self has, as a relative totality, only through the being within it
of the universal principle of totality, the \emph{Absolute}. In affirming
itself, therefore, in its isolation, within its really valid being as a member
in the totality, its self-affirmation is only through a power \emph{which is not
its own}, because it has that power only from something higher. In its
apparently greatest independence the self is working only through the
\emph{Absolute}. In positing itself as evil the will therefore posits itself
only apparently outside and in independence of the \emph{Absolute}: from a
twofold standpoint the continuing dependence shows itself, through which the
self falls into a self-contradiction. It shows itself \emph{first} in this, that
the willing self, which by an act of infinity wills itself in its isolation,
posits itself in that isolation in a finite relation of dependence upon and
opposition to the rest of the finite comprehended by the
universal\footnote{This shows itself from the various standpoints, whether
unfreedom is apprehended as the unfreedom of \emph{drive}, as the unfreedom of
\emph{arbitrariness}, or as the unfreedom of \emph{action} (of the will's
execution).}, comes under the law of efficient causes, and becomes --- while in
its isolated singularity it is powerless over against the sum of the finite ---
a member in an unfree causal connection which has its last ground in the
Absolute; which, from this standpoint, shows itself as a \emph{power of nature},
as blind, governing, overwhelming \emph{necessity}, as \emph{fate}. Thus the
independence aimed at is converted into dependence, freedom into unfreedom,
absoluteness into relativity. --- From a \emph{second} standpoint the selfish,
evil will's dependence upon the Absolute shows itself in the following way. In
order to will itself in a form which could
% FOOTNOTE, p. 153: printed mark *); the hint list does not name p. 153.
% EMPHASIS, p. 153: the spaced run opens at „Villen af sig selv“, not at „i
% Selvets Evne“ (the latter is close-set). The English run begins at the same
% point.
% --- p. 154 ---
answer to the infinity in the act of will which constitutes evil, the self would
have to withdraw into a hardened willing of itself in its empty infinity. The
self, which in the evil will wills itself selfishly, in isolation from the
principle, cannot will itself in the harmonious relation of its essence's
moments, and so give the abstract infinity a content by means of the concrete
--- precisely because the isolation, the tearing loose from the principle of
life, must make the selfish willing's infinity dead and unfit for inner, living
unfolding. But in that empty infinity, in which the self's determinateness by
actuality vanishes, the self cannot maintain itself. Abstraction is borne only
by the real concrete; abstraction's power rests upon the power of the real
foundation of nature which is at work in the will. The self must therefore ---
since its nature as active sets it in motion --- be led back from the empty
infinity to an \emph{unfree willing of the concrete finite}, to an arbitrary
positing of an \emph{isolated} determination of essence; and from there it must
again, in virtue of the \emph{incommensurability} between what is \emph{willed
and the determination of infin}ity, be cast back to the empty infinity, which
--- if the same oscillation is not to repeat itself restlessly and indefinitely
--- can definitively seek to assert itself only \emph{negatively}, through the
disturbance of the good, through a struggle against it, in which the egoistic
will of infinity is nevertheless constantly borne by the real foundation of
nature, whose promptings of drive meet in negative form in the drive to
annihilation. But this striving, in which self-will is raised in power, leads
not merely into the unequal struggle \emph{outwards}, but into an unequal
struggle of the self \emph{within its own interior.}

In the struggle \emph{outwards} the evil will fights powerlessly against the
overwhelming total sum of the other forces, and constantly stands, in a deeper
sense, \emph{alone}. For even under the co-operation of the evil, under their
mutual use of one another's forces, the single one who wills himself in
isolation --- and whose will's egoistic separation makes a relation of love and
concord impossible --- can never in the struggle against the good be borne by a
will of the whole; wherefore even the momentary victory of the outwardly
co-operating evil wills
% SPLIT SPACING, p. 154: „Uendelighedsbestemmelsen“ is spaced through
% „Uendeligheds“ and close-set from „bestemmelsen“ on, within the one word. The
% English divides "infin-ity" at the same joint.
% EMPHASIS, p. 154: „Sætten af en“ is close-set and only „isoleret“ is spaced.
% --- p. 155 ---
must bear the germ of dissolution within it. --- And the struggle against the
good then becomes at the same time a struggle of the self \emph{within its own
interior}, through which an \emph{inner contradiction} becomes manifest: a
contradiction in the relation to \emph{the self's essence}, and to the
\emph{Absolute} which is at work in that essence. The first contradiction comes
forward in this, that the self in this struggle must be \emph{divided}, and
\emph{must unconsciously will the opposite of what is consciously willed},
because it is to fight against that \emph{which is its own true nature and
essence and the ground of its own will.} The struggle must therefore become
fruitless, since the infinity which is used for the struggle has infinity as
\emph{content-filled} infinity over against it, and, having its own essence on
the side of what it opposes, is transformed into something \emph{finite} and
\emph{essenceless}, and sinks in its \emph{impotence}. In the self's struggle
against that which is its own essence, the raising in power of the evil will's
force becomes a raising in power of the force of \emph{resistance}: the evil
will's exertion of force therefore becomes a slackening of force, the highest
exertion the highest impotence; and the consciousness of this impotence breaks
the evil will, so that the essence-willing can come to a breaking through. ---
The contradiction in the relation to the \emph{Absolute} shows itself in this,
that the self --- since in its willing of evil it is willing only through that
power of infinity which it has from the Absolute --- must in its egoistic
willing of itself unconsciously will the infinite by which the possibility of
that willing is, in that its evil willing is broken by the Absolute which works
through it. The selfish will wills with the power which is not its own;
therefore it must, unconsciously, will what is not itself.

The self wills, in evil, itself in its separation of reality, in the form of
infinity, as an absolute. This will of infinity, which is only through the
Absolute, was that which --- when the self sought its will's content in
something finite --- consumed that content through the contradiction with
infinity; and it was that which, when the self then related itself to itself in
abstraction's emptiness, forced out the negative relation to the truly infinite,
to the good which expresses at once man's essence and the Absolute, in the
struggle against which
% PRINTER'S DEFECT, p. 155: the page prints „deu høieste Afmagt“ for „den“ ---
% a turned or wrong sort. And at l. 30 „forholdt sig til selv“, where „sig“
% appears to be dropped before „selv“. Both are Danish-internal and are
% recorded here rather than reproduced.
% EMPHASIS, p. 155: „Modsigelsen i Forholdet til det“ is close-set and only
% „Absolute“ is spaced, unlike the parallel clause earlier on the page.
% --- p. 156 ---
\setcounter{footnote}{0}%
the self hollows itself out, and in the infinite exertion of its force ends in
absolute impotence. \emph{That power of infinity by which evil is} is then in
\emph{evil as the power which abolishes evil;} it is in evil as \emph{evil's
inner negation}, as that which transforms the evil will --- an illusory willing
grounded upon an illusory cognising, as upon a false view of the essence ---
into an \emph{unconscious, indirect willing of the good}, of the essence's
truth. Through the evil will's self-dissolving willing of itself the good is
willed: the essence's truth forces its way through by the self-abolition of the
essenceless. The Absolute mirrors itself, but \emph{as transformed} by what
mirrors it, in the form of opposition, and thus inverted, in evil. It is present
in conscience, not as what reconciles but as what consumes, not as the loving
but as the wrathful\footnote{In this dialectical element in the evil willing
lies \emph{evil's punishment}, which is accordingly to be seen partly as an
\emph{inner} one, i.e. that which is given in the essence's inner discord and in
the relation to the Absolute in the evil conscience, and partly as the
\emph{outer} one, which lies in the reaction of the good, and of egoism of
another kind, against the egoistic willing. This outer punishment is essentially
the powerlessness of isolation, and being hated. As this outer punishment
becomes conscious to the individual according to its ground, the feeling of it
melts together with the feeling of the essence's inner discord. That this inner
discord also comprehends the disproportion in human nature given in degradation
is plain from the foregoing.}. And since the Absolute works as what consumes, as
what negates, as that which through the contradiction and the impotence leads
what is by its nature empty and essenceless to dissolution, and through the
consciousness of the dissolution and the impotence prepares the way for the
good's breaking through in \emph{repentance}, there becomes, in this breaking
through, the \emph{natural force actualised} in the self's infinite willing of
itself: the real ground of nature's \emph{power, taken up into the will, taken
into the good's service}, giving the good a \emph{heightened energy}, so that
\emph{evil} in conversion acquires a \emph{positive} significance for the
\emph{good.}
% FOOTNOTE, p. 156: printed mark *); the note carries an id-est mark.
% EMPHASIS, p. 156: „det Onde i Omvendelsen“ is NOT a spaced run --- only „Onde“
% is spaced. Likewise „i det Onde som“ before „det Ondes indre Negation“ is
% close-set. The English follows in both places.
% --- p. 157 ---
The concept of \emph{the actualising of the real ground} through the evil will
is to be determined more closely as follows. The evil will has, as an untrue
willing of infinity, for its conditions first a working-through of the
principle, the truly infinite, and secondly an \emph{energy in the foundation of
nature} (the real relative totality) \emph{which bears the individual will.} But
this second condition --- the foundation of nature's energy --- cannot, in
virtue of the fundamental relation between the real and the principle, be
thought without that working-through, which must therefore be covertly present
within it. The foundation of nature's energy makes itself felt in \emph{drive},
which, raised in power by infinity, becomes \emph{passion}, a willing of its
object \emph{as absolute}. Drive here wills itself infinitely; it is an
ἄπειρον. But under this striving of drive the will's form and its content are
incommensurable. From this follows the alternation in the content of will
demonstrated earlier. But the \emph{formal element: passion's infinity} remains;
and since the energy is now turned against the good in order to fight it, the
\emph{resultant of the workings of drive} becomes the \emph{formal, negative
infin}ity of the evil will \emph{fighting} for the good's annihilation. It is
this infinity of the will that is preserved when the will's direction is
reversed --- \emph{under the transformation of the force the sum of force
remains here too, as in the domain of the physical} --- and when
\emph{repentance}, which is the \emph{essence's} positive assertion through the
negation of the \emph{essenceless} willing, becomes victorious, the energy is
taken into the good's service, and is the more intensive as the real moments are
as it were singly penetrated by infinity. In this metamorphosis of energy too
the Absolute's working-through is to be accentuated; but here it is positive, in
unity with the essence-willing, whereas in the evil will it was negative, as
dissolving the essenceless willing.

According to what was developed earlier concerning \emph{evil's primitive form},
the dialectic demonstrated in the evil will --- which points directly to that
primitive form --- must acquire validity for \emph{every} phenomenal form of the
evil
% GREEK, p. 157: „ἄπειρον“ --- antiqua italic, smooth breathing with acute on
% the alpha. Neither OCR witness reads it; p. 157 is not in any Greek list.
% Copied verbatim.
% SPLIT SPACING, p. 157: „Uendelighed“ is spaced through „Uende“ (the end of the
% line) and close-set in „lighed.“ The English divides "infin-ity" likewise.
% EMPHASIS, p. 157: „kæmpende“ is spaced while the following „onde Villies“ is
% close-set.
% --- p. 158 ---
willing, because the peculiarity of that primitive form is a pervasive
determination in them all. Everywhere --- even where what is peculiar to the
fundamental form most of all conceals itself, as where evil is called forth by
an outer influence whose beginning goes back to earliest infancy, or where the
inclination is inherited --- there is, precisely because the evil will can be
only as \emph{an expression of spirit}, a point to be posited at which the will,
supporting itself upon a \emph{consciousness} whose degree of clarity makes
\emph{imputability} possible, has the form of \emph{defiance}, self-willedly
wills evil, and by that self-will constitutes itself as an \emph{evil} will.

But when what has been developed above is to be applied to evil's single forms,
a difficulty arises with such forms as those in which the individual's natural
force is broken and the organic conditions of an energetic acting are destroyed.
With these forms there seems just as little to be any question of an actualising
of the foundation of nature as of a preservation of force; and generally it
seems, where evil has the form of \emph{degradation} (weakness, wretchedness,
animality, lostness, ``in despair not to will to be oneself''), and where sin is
continued in \emph{torpor}, that through a life in such \emph{sin} nothing can
have been realised that could serve the good, and that a raising up must become
impossible. The relation is here to be apprehended as follows. In \emph{every}
unethical willing, resting upon an untrue working of drive, there is posited a
\emph{disharmony} in the essence, which must to a certain degree work
disturbingly upon the essence's side of reality: the organism (``death is sin's
wages''); and it is this \emph{everywhere} occurring disturbance whose degree is
in those forms so increased that the hidden sickness in the self becomes
manifest, and can even take on the character of animal torpor. Here, then, there
is not an opposition but a \emph{difference of degree}. Under the
\emph{disharmony} which occurs everywhere in evil, and which we designated as a
\emph{sickness in the self}, the will nevertheless has its \emph{abstract
infinity} upon the basis of the actualised real, because the essence as essence
\emph{cannot} be annihilated; and since the energy of will does not stand in a
\emph{direct} relation to the
% QUOTATIONS, p. 158: two pairs, one of them the Kierkegaard formula. Both
% balanced, and both kept as quotations.
% RULE, p. 158: the page carries NO rule at its foot; the subsection runs on
% into p. 159, where a rule does stand.
% --- p. 159 ---
physical force but to the individual relative totality's concentration (in the
contrary case the physically strongest individuals would have to be the
strongest in will), the \emph{intensity of will} can be preserved under the
disharmony, even where it takes on larger dimensions, in spite of the
disturbance; and therewith \emph{repentance's} starting point can be given. And
even where a disorganisation of the will's organic conditions takes place, the
possibility of ethical restitution remains, just as freedom's possibility
remains in the self so long as the self exists as a \emph{human} self. Since sin
is continued in torpor, but yet as willed (cf.\ pp.\ 150 f.), there is, under
the \emph{despairing} sinking down, in the very despair, an \emph{unconscious
struggle of anxiety} in which evil's power is undermined, and out of which a
saving flame of repentance may break forth. Even through such a life, then,
something may be thought to have been realised for the good.

In \emph{the rising of despair and anxiety}, which is only covered over by the
torpor and the lostness but is nevertheless present, there is work going on, as
it were in the depths of the essence, towards a breaking through; and when the
essence breaks through in repentance, it bears within it an infinite content
which has not been able to be annihilated, but which also can be set free only
through a far deeper suffering, through a painful burning out of evil. Here too,
then, the force can be preserved and the struggle bear its fruit in conversion.

\begin{center}\rule{2cm}{0.4pt}\end{center}
% RULE, p. 159: a short centred SINGLE hairline standing mid-page with no head
% adjacent to it --- a thought-break, as at pp. 19, 28, 42, 62 and 97.
% Transcribed. The display head lower on this page carries no rule of its own.

For a closer determination of the fundamental relations in evil demonstrated in
the foregoing there remains still to be developed: \emph{the relation between
the moment of necessity and the moment of freedom in evil's realisation and in
its abolition}, and \emph{reconciliation's essence, its principle and its goal.}

\section*{1. The Relation between the Moment of Necessity and the Moment of Freedom in Evil.}

\emph{a) in the realisation of evil.} --- The \emph{indifferentist} conception
shows itself at once to be untenable. It
% HEAD, p. 159: bold (halbfette) Fraktur, centred over two lines. The Indhold
% omits „i det Onde“; the PRINTED head stands in the body, here as in the
% Danish, and the divergence is recorded at the Contents.
% HEAD, p. 159: the run-in „a)“ head is letterspaced REGULAR Fraktur, not bold,
% followed by an em-dash --- hence \emph{} and not \textbf{}, unlike the run-in
% heads of pp. 97 and 126.
% --- p. 160 ---
is self-contradictory, in that it lets the action's, the choice's,
determinateness proceed from a will empty and undetermined in itself; it
excludes responsibility and imputability by annihilating the personality's
continuity; and it must for the same reason posit an ethical influence upon
character as impossible. But because indifferentism with its groundless
\textit{liberum arbitrium} is rejected, freedom is not thereby posited as
absorbed by necessity. In the section on freedom it was shown how the human
individual, in virtue of its essential determinateness as a relative totality
with inner infinity, must with consciousness concentrate itself in itself, take
possession of itself. In the fact that it \emph{must} do so, the
\emph{necessity} in the relation is posited; but since it constitutes itself in
its essential determinateness \emph{through itself}, posits itself through
itself, the necessity is in this conscious act of self-concentration converted
into \emph{freedom}. In the self's constituting of itself as a relative totality
there lies then \emph{a twofold relation}: a relation of the self to the
Absolute and to itself, and therefore \emph{a twofold possibility}: the
possibility of the ordering under the absolute totality and of isolation. Since
the twofold relation --- the relation to itself and to the Absolute --- is
conditioned by \emph{a twofold act of consciousness}, it must be posited in
\emph{different} moments of time; but the moments of time may be thought as
bordering upon one another so closely that there is nothing spacing them apart,
and therefore no separation in actuality of the two relations; and they may be
thought as drawn apart from one another to an indeterminate degree. In this
determination of time --- which is conditioned by the human consciousness's
peculiarity, and which acquires significance for the fundamental relations of
what is according to its essence eternal, in that these relations become an
object for consciousness --- lies the point of attachment for a moment of
\emph{contingency} in evil's realisation. For here the impulses coming
\emph{from outside} acquire their play; and evil's possibility, which is present
in every individual, can through the relation to actuality with its impulses
acquire an intensity heightened in indeterminable measure, so that it may be
thought to stand to the good's possibility as the infinitely
% ANTIQUA LATIN, p. 160: „liberum arbitrium“ is antiqua and CLOSE-SET, like the
% antiqua of pp. 64, 71 and 134 and unlike the spaced antiqua of pp. 68 and 144;
% \textit{} without \emph{}.
% PUNCTUATION, p. 160: no comma is printed after „Grundforhold“; both witnesses
% agree, and the English likewise runs the clause on.
% --- p. 161 ---
great to unity. But when that possibility is made actuality, the conversion from
possibility to actuality nevertheless remains always an \emph{act of freedom} ---
more determinately, an act, grounded in the human being, of freedom in its
\emph{infinity}, but in its \emph{abstract} infinity.

When, through periods of incalculable extent, a life has been lived by human
beings without what is in truth human having come to development in
consciousness, the question may be raised how far under such conditions there
can be talk of evil's realisation. The answer to this question --- which comes
under the more general question of evil's contingent or necessary relation to
\emph{the human being} --- will depend upon whether the factual lack of
spiritual development is seen as necessary, and as a real non-presence of what
is of principle in the content of human consciousness. The decision of this lies
in the relation between the determination of \emph{essence} and the influences
and conditions, incalculable in intensity. In this relation the determination of
essence must be held fast as what overreaches. If the life that is lived is a
life of \emph{human beings}, then the human's \emph{determination of essence}
and the \emph{demand} lying therein are present; and the non-fulfilment of this
demand constitutes evil's actuality, however strong the conflicting impulses may
have been, however weak the degree of consciousness.

If the moment of freedom and that of necessity in evil's realisation are seen
from the standpoint of the relation between the \emph{individual} will and the
working \emph{Absolute}, then the earlier developments contain what is necessary
for the apprehension. We point back here to what was developed about the
relation between the determination of ideality and that of reality, the
determination of eternity and that of time, in the Absolute; about this
relation's significance for the finite's relative independence; about the
Absolute's negative being in the evil will; and about this will's relation to
the essence's infinity. If the determinations freedom and necessity are applied
to the \emph{Absolute} in its relation to evil, then --- since in the
% SPLIT SPACING, p. 161: „Væsensbestemmelsen“ is spaced through „Væsens“ and
% close-set from „bestemmelsen“ on, within the one word --- the third such split
% in this stretch (cf. pp. 148, 154, 157).
% --- p. 162 ---
Absolute freedom and necessity must be thought in eternal unity --- evil is at
one and the same time to be apprehended as grounded in an act of freedom of the
Absolute (the absolute subjectivity), and as existing in virtue of a necessity
in the Absolute's essence. This relation of the moment of freedom and that of
necessity forms the parallel to the relation between the determination of
ideality and that of reality in the principle.

\emph{b) in the abolition of evil.} --- What has been developed in the foregoing
concerning the dialectic in evil, its inner self-dissolution, renders above all
a \emph{relation of concept}, a \emph{relation of essence}, a dialectic of
essence, which is therefore seen above all in abstraction from the temporal
relation and under the form of necessity. Here, on the contrary, evil's
abolition is to be seen as a \emph{relation of actuality}, in the
\emph{temporal relation} and in relation to the determination of
\emph{freedom}. The temporal relation's essential significance has been touched
upon in the determination of evil's first realisation in the will. From this
point of beginning time's significance shows itself in a twofold respect: in
respect of what is defective in the development, of the growing sum of what has
been neglected with regard to the demand's fulfilment (the growing guilt); and
in respect of the postponement or the failure of the will's turning away from
evil. The first will be elucidated in what follows, where reconciliation is to
be treated; the second is to be elucidated here at once.

The evil will's inner contradiction must prevent a standing still, a rest in
evil; it must force out a restless unrest, a ceaseless striving forward. Under
this striving --- which, whether it takes the form of a tumbling about in
pleasure or of a direct struggle against the good, constantly encloses defiance,
resistance to the ethical demand --- the evil will must more and more hollow
itself out; it must, in its negating of the demand of essence, while it wills to
make itself absolute, unconsciously strive towards a self-annihilation. But this
self-annihilation can only be approached, because the evil will --- which fights
against the essence, and wills to posit itself as the essence's expression ---
covertly preserves its connection with the essence,
% HEAD, p. 162: the run-in „b)“ head is letterspaced regular Fraktur, not bold,
% followed by an em-dash --- set exactly like its companion „a)“ on p. 159,
% hence \emph{} and not \textbf{}. The Indhold agrees word for word; no
% discrepancy to log.

% --- p. 163 ---
which as essence cannot be annihilated, so that it, while it
unconsciously works towards a self-annihilation, --- which, if the essence could
be absorbed by the evil will, would be the essence's annihilation, --- works
with a force borrowed from the merely hidden essence-willing, the will of the
good, and must always remain as a positive remainder beneath that doubling in
the evil will in which the will's force is consumed, but by a force which
belongs to the will. There therefore always remains a doubleness in the
relation; there is always as it were a flaring up of the flame which is to be
extinguished, and therefore there is always something merely approximative in
that striving which from an infinite distance strives towards a limit at which
the force reaches 0, where the will would sink away into annihilation, into an
eternal death. But in that \emph{dividedness} of the will which follows from the
evil will's fighting, as a consciously actualised will, against the
essence-willing, there is a discord posited in the essence which is like a
\emph{sickness of the essence}, a suspension of integrity, and the discord grows
as the evil will hollows itself out in fighting. In the evil will's
preponderance the essence becomes as it were invisible: its forces work
negatively in evil. As essence it cannot be lost, cannot be annihilated;
although this annihilation of it, which would be one with the will's
annihilation, is the goal towards which the evil will unconsciously works. But
in \emph{the discord in the essence}, in \emph{the sickness} which threatens
with an ever stronger \emph{dissolution}, there is an \emph{anxiety}, and in
this anxiety of sin is the reversal's possibility. In that double
movement in the will there shows itself, at the single points, the negation, the
self's lostness as the essence's disturbance, threatening for consciousness, and
in showing itself it lets the essence show itself, through the anxiety for
itself, as that which is to be maintained and restored, and herein lies the
possibility of repentance. Intermediate stages may now be thought, with an
oscillation in which the moments of repentance and of relapse alternate; but
there comes a point where consciousness, under the self's despairing working
towards a self-dissolution, is filled with the anxiety of spiritual death, of
the essence's
% RULE, p. 163: the page carries NO rule at its foot. It ends in mid-sentence,
% and the only mark below the last line is the gathering signature.
% ANTIQUA FIGURE, p. 163: „hvor Kraften naaer 0“ --- the zero is an antiqua
% figure (there is no Fraktur nought); set as the digit 0 in both files.
% EMPHASIS, p. 163: „i en evig Død“ is NOT letterspaced, though the parallel
% phrases around it are. „Spliden i Væsenet, i Sygdommen“ runs as one spaced
% phrase across its comma; the second „Angest“ is tight where the first is
% spaced. The English follows in each case.
% --- p. 164 ---
dissolution by the will, and therefore by itself, and at this point the
essence's struggle of anxiety becomes so strong over against the hollowed-out
evil will --- which is the divided will --- that the relation in the
possibility of evil turns about and comes to be designated by a formula opposite to that used
above: in that evil's possibility must now be posited as standing to the good's
as unity to the infinitely great. Here, then, is the moment of the reversal, of
the breaking through, and it is there in virtue of the essence's
\emph{necessity}, but through the resolution of \emph{freedom}, of the freedom
which asserts the essence and itself.

Two questions must however at once press themselves forward here. The first is
this: whether the case may not be thought in which evil comprehends the
\emph{whole} of the individual's existence, and in which the breaking through
therefore fails to come for this individual? The breaking through's relation to
freedom, and the contingency of the determinateness by time, seem to require an
affirmative answer. Nevertheless the answer will show itself to have to be
another, but this can first be grounded through the development of
reconciliation and of the relation between its factors. The \emph{second}
question is this: whether the breaking through --- the self's return to a
willing of its true essential determination --- is the self's \emph{own} deed, or whether
it is only through a relation in which the self stands to \emph{something
supplementary}, be that the other human individuals or the self's absolute
principle. This question too leads us to the problem of reconciliation, and in
the development of that problem what was touched upon above --- the possibility
of a making good of what has been neglected --- acquires a decisive significance.

\section*{2) Reconciliation, its Source and Goal.}

Reconciliation's concept we apprehend, in accordance with what has been
developed earlier, in the twofold sense of the individual's reconciliation
\emph{within itself}, i.e. of the inner harmony brought about in the individual
and its accord with the essence's demand, and of the individual's
% HEAD, p. 164: bold (halbfette) Fraktur, centred, on one line. TWO
% discrepancies with the book's own Indhold, both real and both left standing:
% (i) the printed enumerator is „2)“ with a closing PAREN, while the Indhold
% reads „2.“ with a full stop --- the reverse of the pattern at pp. 97 and 126;
% (ii) the Indhold places the entry at „167--201“, but the head is printed here
% on p. 164, and the Indhold's own sub-entry a) is given as 165--167. Neither is
% in the Trykfeil leaf, and the comment at the Contents records the same facts.
% RULE, p. 164: a short centred SINGLE hairline stands ABOVE the display head,
% as at pp. 97 and 126. Not transcribed, per the standing convention.
% EMPHASIS, p. 164: „Det første er det“ is TIGHT while „Det andet Spørgsmaal“
% has „andet“ spaced and „Spørgsmaal“ tight. Another split: „Selvets egen
% Gjerning“ spaces „egen“ only.
% PRINTER'S DEFECT, p. 164: „her strax paatænge sig“ --- „paatænge“ for
% „paatrænge“, a dropped „r“. Orthographic, and so not reproducible in English;
% not in the Trykfeil leaf either.
% --- p. 165 ---
reconciliation \emph{with its absolute principle}. Under both apprehensions,
whose inner connection has been shown earlier, the ethical disturbance is here
presupposed as going before, so that reconciliation becomes not the expression
of an undisturbed harmonious relation mirrored in consciousness, but of the
conscious \emph{restitution} of the disturbed normal relation.

We consider then first:

\begin{center}
\emph{a) that which is operative in reconciliation.}
\end{center}

Immediately this is of course \emph{the individual's own} \emph{essential will},
set free through \emph{repentance} and come to a breaking through in the
\emph{ethical resolution} for the good, the expression of the individual's
essential energy, of the personality's most intensive activity. But the
individual's ethical resolution is a resolution of that will in which the
personality wills itself as ordered into the totality and as subordinated to the
totality's principle. In this will, therefore, the relation to the living whole,
to its organic members and its principle, is posited from the beginning, and
this relation, which is here posited, shows itself,

in being posited, in that the individual becomes conscious of itself as
encompassed by the whole's life, to be the \emph{presupposition} of that will's
breaking through and the \emph{essential condition} of that development of the
will by which reconciliation becomes \emph{complete} as comprehending the whole
reality of life.

The relation to the members of the totality, to the \emph{other individuals}
who in union form a spiritual organisation, shows its significance through the
supplementary individuals' \emph{ethical} influence; through the conscious and
unconscious \emph{ethical communication of life}. The influence comes forward
in the education of the not yet ethical individual to moral willing by those
personalities in whom the spirit of morality is actualised; it comes forward in
the continued moral influencing, and, in negative form, in the determination of
the unethical individual by deterrence and punishment. It has its deepest, most
primitive and most universal expression in \emph{the relation of love}, in which
% --- p. 166 ---
the essence's unconscious energy goes in unity with the energy of the life of
consciousness. And just as in education and in deterrence it was the race, the
universal, reason-determined being of the human, that was at work through the
working of the ethical individuals, so in the relation of love this relation to
\emph{the race} --- the word taken in its \emph{full}, spiritual-bodily sense
--- comes forward generally as decisive. In the relation of love, for which the
love of the sexes is only a special and in and for itself not yet ethical
expression, the other individual is for the individual the race's
representative, the spiritual totality's representative, and becomes in the
ethical relation for the individual as it were its personified, objectified
conscience. The individual's relation to the other individual thus apprehended
becomes an \emph{ethical} relation \emph{to itself}: in the other individual's
consciousness it has, through that unconscious transposition of the individual
into a representative of the race, a determining \emph{norm} for its action.
Even the illusion which can attach to the relation of love, --- not merely to
the erotic relation ---: the illusion by which love's object is idealised, has
its ethical significance when it does not become an expression of the indirect
self-love; it gives the individual, which in this illusion relates itself to an
ideal, an \emph{ideal measure} for its action, a higher norm for it. And this
intuition of the ideal, of the universally valid, then fuses with the intuition
of \emph{the personalities' peculiarity}, of the ethical's peculiar forms of
realisation in the different individuals; and by this idealised ethical
peculiarity, which at one and the same time contains a corrective for what is
defective in the individual and an assurance of the ethical force in the human,
the individuals comprehended by the relation of love show themselves in a
twofold sense as supplementary. Thus love, whose germ is awakened to life in
the individual by the love already at work, becomes the true ethical source of
life and therefore a source of reconciliation.

Through the relation to the other individuals there is everywhere a pointing to
that which comprehends, \emph{the totality}, but
% RULE, p. 166: none, per the Danish; the page ends in mid-sentence.
% DASHES, p. 166: the pair round „Ordet taget i dets fulde, aandelig-legemlige
% Betydning“ and the „, --- ikke blot ved Elskovsforholdet ---:“ group are the
% PRINTING'S own, not editorial. The second is reproduced with its dash-comma
% opening and its dash-colon close exactly as set.
% --- p. 167 ---
through this there is a pointing back again to \emph{the totality's principle};
and since the totality's energy is grounded in this, as is the single
individuals' ethical working in the totality, reconciliation's source is in the
\emph{last instance} to be sought in \emph{the Absolute}. This has already in
essentials been developed in the foregoing. The principle's working has been
shown in every, positive and negative, individual relation of will to the
ethical. It has been shown in the the good will's self-surrender to the divine,
which is only by the latter's active indwelling in the will; it has been shown
in the principle's negating working in the evil will; it shows itself likewise
in that breaking through in which the good will breaks forth through the
dissolution of evil; for the essence's assertion in this breaking through can
only be by a working of that by which the essence is. It is finally also to be
posited as that which is definitively at work in the restitution of what has
been disturbed and neglected by the evil willing, in the attainment of the goal
in spite of the fact that in a span of time which cannot be repeated a deviation
from the goal has taken place. ---

But this last point of view leads us from the consideration of what is at work
in reconciliation to the consideration of:

\begin{center}
\emph{b) that which is to be attained in reconciliation.}
\end{center}

Apprehending reconciliation's concept in the twofold sense indicated above, we
find a comprehensive expression for reconciliation's goal in \emph{the
actualisation of the ideal}. The actualisation of the ideal becomes, when we
take the starting point for the determination in the individual's relation to
the Absolute, equivalent to the actualisation of \emph{the absolute will's
demand}; but this demand becomes again, in accordance with what was developed
earlier, identical with \emph{the essence's demand}. If we then take the
starting point from the human being itself, the ideal's actualisation acquires
a \emph{twofold} sense. On the one side it is to be apprehended as the
realisation of what is \emph{essentially} individual in the individual, of the
individual's \emph{peculiar} endowments, which are to be developed into an
ethical \emph{content of personality}; on the other side it is
% HEAD, p. 167: „b) det, der skal naaes i Forsoningen.“ --- centred, letterspaced
% REGULAR Fraktur, not bold, exactly like a) on p. 165; hence centred \emph{} and
% not \section*. The Indhold agrees word for word.
% DASH, p. 167: „har fundet en Afvigelse fra Maalet Sted. ---“ is the PRINTING'S
% own paragraph-closing dash and is reproduced.
% PRINTER'S DEFECT, p. 167: „paaviist i den den gode Villies Selvhengivelse“ ---
% the article „den“ is set TWICE. A dittography, and one of the few Danish
% defects that DOES travel: the English reproduces it as "in the the good will's
% self-surrender". Not in the Trykfeil leaf.
% --- p. 168 ---
to be apprehended as the realisation of what is \emph{universal} in the
individual, of \emph{the universal ideal self in man}, of \emph{man's idea}. But
the realisation of that individual and this universal must again show itself in
\emph{inner unity}, and this unity comes forward in \emph{positive} form in that
the individual in the unfolding of its peculiarity expresses the universal; it
comes forward in \emph{negative} form in self-surrender, in which
individuality's energy is at work in the subordination of the single self under
the totality, in resignation in the relation to the latter. In this unity, then,
the determination of the individual's individual τέλος and of the
universally-human τέλος determined by \emph{the human being} are so connected
that the human ideal, \emph{humanity's idea}, is individualised for the single
person, and \emph{the realisation of this idea made living in the single person}
becomes the ethical demand upon the single person: the \emph{ideal} which is
\emph{reconciliation's goal}.

The actualisation of this ideal, which becomes the ethical will's goal, now
acquires a \emph{threefold} condition. The \emph{first} condition is this,
\emph{that guilt can be made good}, that is, that in spite of the fact that the
ideal's actualisation is \emph{interrupted} by the unethical action, and that
therefore a time irreplaceable in an outward sense has been lost for the
actualisation, there is nevertheless a possibility of a making good of what has
been neglected, of a compensation for the lost time of working. The
\emph{second} condition, to which reference has already from one side been made
by the first, is that the ethical can become \emph{a determination of totality
for human life}, not merely in spite of that loss of time, but also in spite of
the fact that the individual's ethical willing, in accordance with the human
conditions of development, is necessarily limited to a definite stage of
maturity in the individual's life. The \emph{third} condition, finally, is that
the individual human existence acquires the \emph{determinateness by eternity}
which corresponds to \emph{the ideal's concept}.

\emph{α) The abolition of guilt.} --- We have above, from a rational starting
point, reached the determination of the individual will's incongruence with the
ethical de\-%
% HEAD, p. 168: run-in „α) Ophævelsen af Skylden.“ --- letterspaced regular
% Fraktur followed by an em-dash, the Greek enumerator α) in ANTIQUA ITALIC. Set
% with \emph{} like the run-in heads of pp. 159 and 162, and the em-dash after it
% is the printing's own. The Indhold reads „α) Skyldens Ophævelse“, the two nouns
% in the other order; the PRINTED head stands in the body, per the standing rule.
% RULE, p. 168: a short centred hairline stands above the α) head; head ornament,
% not transcribed, as at pp. 97, 126 and 164.
% GREEK, p. 168: τέλος TWICE in the one sentence, copied from the Danish.
% PRINTER'S DEFECT, p. 168: „ikke blot tiltrods for hint Tidstab, man ogsaa
% tiltrods for det“ --- „man“ for „men“. Orthographic and Danish-internal; it
% does not travel, and the English sets "but" without comment in the text.
% d. v. s. -> "that is", per the playbook.
% --- p. 169 ---
mand, and thereby the concept of \emph{guilt}. This concept is now to be
developed more closely.

The determination of the relation of incongruence must acquire its different
expression according as the incongruence is seen from the side of
\emph{essential necessity} or from that of \emph{determinateness of will}. Seen,
abstractly, from \emph{the essence's} side, it becomes, where it is admitted at
all as actual, apprehended in metaphysical determinations essentially
corresponding to those under which Greek ethics apprehended it. It is then
\emph{finitude's} μὴ ὄν which impresses itself in the \emph{imperfection} of the
essence's expression: the relation of incongruence with \emph{the ethical}
becomes determined \emph{negatively}, as a \emph{defect}. Seen, on the contrary,
from \emph{the will's} side, other categories assert themselves. There is then
talk of \emph{guilt}, still more determinately of \emph{sin}. The concept
\emph{guilt} can still be held hovering on the boundary of the metaphysical and
the ethical. Finitude's, ``the guilt of single existence'', which demands its
expiation, is the expression for the relation in the old Greek philosophy, which
essentially considers natural existences, and in this expression the oriental
fundamental view is mirrored. But in its stricter and proper sense guilt becomes
an \emph{ethical} category and designates \emph{the will's self-incurred
disproportion to the ethical}, whether now this disproportion be that existence
has not yet, though the possibility was there, been placed under the ethical's
determination, or that the break with the ethical is expressly posited through
the will's willed selfishness, which, asserting itself in its negative infinity,
posits a determination of finitude as the Absolute, thus in the Absolute's
place. Here \emph{the will} is accentuated as that which in \emph{self-wilfulness}
incurs guilt; and however much the significance of \emph{cognition} for the
resolution of will be emphasised, just as fully is the \emph{conversion} into
determination of will emphasised.

The concept of \emph{guilt} now acquires its more determinate expression in
\emph{sin}, whether now by this determination there be expressly brought out the
relation to the \emph{divine will}, or the affection of \emph{the essence} by
the unethical action,
% BROKEN WORD, p. 168/169: „For\-“ / „dring“ = Fordring, "demand"; the English
% divides "de\-" / "mand" at the same joint.
% GREEK, p. 169: μὴ ὄν, copied from the Danish.
% EMPHASIS, p. 169: „Enkelttilværelsens Skyld“ is set TIGHT inside its quotation
% marks --- no \emph{} inside the quotation, and none in the English either.
% PRINTER'S DEFECT, p. 169: „Naturexisteniserne“ for „Naturexistentserne“, the
% „t“ of the stem set as a dotted „i“. Orthographic; does not travel.
% --- p. 170 ---
the loss of its integrity. By both of these determinations there can be a
passing over onto the territory of the \emph{paradoxical}. For if the relation
to the divine will is apprehended as a relation to a will whose law is the one
heterogeneous with cognition's law of reason, incomprehensible in its
arbitrariness, then sin becomes the paradoxical, the consciousness of sin as sin
conditioned by a relation of faith of paradoxical nature, and sin's
reconciliation becomes conditioned by and an object for a relation of faith of
the same kind. And in the same way it stands with that apprehension of sin's
effect which in the doctrine of original sin lays down \emph{the corruption of
human nature} by the first sinful act of will. This corruption, as the Church's
doctrine apprehends it, becomes the incomprehensible. The human being is to be
thought as corrupted by an act proceeding from the essence itself, and, more
determinately, from the essence still in the state of integrity. The expression
of will which realises a possibility laid down in the essence is to transform
the essence, by one act to destroy it. As a parallel to this incomprehensible
there then stands the incomprehensible in the restitution by a divine miracle,
by God's paradoxical becoming-man, by the God-man's self-sacrifice, which is
appropriated through faith. In this apprehension sin and reconciliation withdraw
themselves from all cognition. But the determination of \emph{sin} as
comprehending both a strife \emph{against the divine will} and an effect
\emph{upon the essence} can be held fast without the territory of the rational
being abandoned. The determination sin can be held fast in the sense of
transgression of the divine will's demand, in that the concept of the
\emph{divine will} is apprehended in the way which was indicated in the
development of the Absolute's concept. And the concept of \emph{an affection of
the essence's integrity by sin} can be held fast, in that it is apprehended in
accordance with what was developed in the preceding sections. In accordance with
the twofold possibility present in the human spirit: to be and to will itself as
a consciously subordinated member in the totality, and to be and to will itself
in isolation according to its negative infinity,
% p. 170 carries no emphasis beyond the six runs marked, no footnote, no Greek,
% and no rule.
% --- p. 171 ---
there arises for the will a possibility, by positing itself, as evil, according
to its negative, abstract infinity, of \emph{finitising} itself; and as
finitised the will is not integral, and \emph{the expression of essence} the
will's lack of integrity touches \emph{the essence} by the discord which is
posited in it. But under this discord of the essence, under which integrity is
as it were suspended, the essential energy is according to its infinity
negatively present in the finitised will (cf.\ the earlier developments). The
discord does not acquire power to dissolve the essence; redintegration's
possibility and source is preserved in the essence itself, and the possibility
becomes actuality in the essence's self-acquisition, in its breaking through by
way of the evil will's self-dissolution. That this self-dissolution is brought
about only by the essence in that it is brought about by the indwelling
Absolute's power, just as the will's positing of itself as evil according to the
untrue abstraction of infinity was possible only by the true Infinite's
indwelling, has been shown above. But this relation does not annihilate the
essence's \emph{own} energy; for the through-working Absolute is such a one as,
in accordance with the relation between its side of ideality and its side of
reality, lets a place remain for the working of the finite.

Out of the given determination of the concepts of guilt and of sin, by which
these therefore become categories in a \emph{rational} ethics, the question is
then to be raised concerning \emph{the possibility of an expiation of guilt
through a factual making good of what was neglected in ethical action}. This
possibility was what was posited as the condition for the guilt-laden
individual's \emph{restitution in principle} being able to become an
\emph{actualisation of the ideal}, a \emph{full and true} reconciliation.

In order that a making good of what has been neglected, and thereby an ethical
redintegration of the individual in spite of guilt, should be thinkable, it is
required that in the \emph{transformation in principle} in the ethical's
breaking through a \emph{conversion} take place \emph{of a force actualised
through the evil willing}. Such an actualising of a force by the evil willing,
through this willing's infinity, was shown above,
% SYNTAX, p. 171: „og \emph{Væsensyttringen} Villiens Mangel paa Integritet
% berører \emph{Væsenet} ved Spliden“ --- the apposition is unmarked and the
% sentence will not parse as it stands; a preposition or a genitive appears to
% have dropped out. The transcription does not flag it and it is not in the
% Trykfeil leaf, so it is REPRODUCED rather than mended, and the English is
% broken at the same joint. Only the page image can settle whether the fault is
% the compositor's.
% EMPHASIS, p. 171: „berører Væsenet ved Spliden“ spaces „Væsenet“ ONLY ---
% „Spliden“ is tight here, unlike the spaced „Spliden i Væsenet“ of p. 163.
% --- p. 172 ---
and this actualised force is what in the good's breaking through in the will is
taken into the good's service, so that it becomes an energy whose earlier
development, in that through repentance the direction of the force is turned
about, acquires positive significance through a heightened activity of will. If
this actualising of force is seen in relation to \emph{the essence}, then the
emphasis falls on the essence's negative working under the disturbance, and this
hidden working of the essential energy is what makes it possible that, when the
breaking through occurs, the time filled out with the unethical can acquire
positive worth for the ethical development. It is then \emph{repentance}, as the
reaction of the essence come to a breaking through against the self-imputed evil
willing, that keeps the force actualised by this willing alive as a force at
work for the good. Therefore repentance acquires an \emph{abiding} significance,
and it must --- in opposition to what Fichte would have it --- be apprehended as
effecting an increase of force. --- If the actualisation of the force is seen in
relation to \emph{the objectively good and the Absolute's will}, then the
emphasis falls on that relation under the evil will's struggle against the good,
in accordance with which it is under this struggle \emph{unconsciously}
determined by the \emph{essentially valid}, against which it struggles, to a
straining of force, to an increase of force in a negative direction. In the
struggle itself there is thus posited a relation to the essentially valid as
that which governs the will; and by this relation, in that the direction of will
is transformed in the good's breaking through, i.e. in that the conversion
occurs, the struggle against the good can itself acquire positive significance
for the ethical development. In conversion the false, negative, expressions of
the will's energy enter, as transformed into positive determinations of force,
into the development, and the preceding life joins itself through repentance
together with the present into a continuous intensive unity of life. Striking
examples of such a relation come forward where an energetic struggle against the
higher truth is transformed into an energetic adherence to it. In the passion of
the struggle itself the truth's power is manifest: by the very passionate
resistance it is thus taken up into consciousness, penetrates
% DASHES, p. 172: both pairs are the PRINTING'S --- „den maa --- i Modsætning til
% hvad Fichte vilde --- opfattes“, and the paragraph-internal „Kraftforøgelse. ---“
% which opens the next movement without a paragraph break.
% NAME, p. 172: „Fichte“ is Fraktur and NOT letterspaced; no \emph{} on it.
% ID-EST MARK, p. 172: „i det Godes Gjennembrud ɔ: idet Omvendelsen skeer“ ->
% "i.e.", the second in this stretch after p. 164.

% --- p. 173 ---
it in such a way that in the reversal everything is prepared for an immediate,
collected working of the whole force of spirit in the service of that truth
which now becomes the principle of life and can at once work as such in all the
forces of life. One need only recall one example: Saulus's relation to
Christianity. ---

In the foregoing \emph{redintegration} has been seen essentially in relation to
\emph{the will itself}, but it is also to be seen in relation to \emph{that
which is effected by the evil will}, yet in such a way that this point of view
is not, as by the nature of the case it cannot and must not be, abstractly held
apart from the former.

What is effected by the evil will becomes partly a disturbance in the
individual itself, in its side of reality or in its spiritual development;
partly a disturbance in other individuals' physical or spiritual existence;
partly a disturbance in the objectively universal. The disturbance effected by
sin in \emph{the individual's side of reality} is in this connection to be seen
as a \emph{disturbance in the physical conditions of life for the will's
working}, and its significance to be apprehended from this point of view. With
respect to this the essential has already been indicated in the foregoing, in
the development of that possibility of freedom which was preserved under the
disturbances and defects in the physical organisation. The emphasis comes here,
as everywhere in these relations, to fall upon this \emph{abiding possibility
of freedom}, in which \emph{the energy for the transformation of will} is
present; and since this energy is there, the transformation of force
characterised in the foregoing becomes possible in spite of the physical
defect, and this defect becomes an ethical incitement for that will of
repentance which, in spite of it, is to assert itself in the ethical and is
relatively to be removed by the will's energy. --- If the disturbance in the
individual is apprehended as a disturbance in \emph{its spiritual development},
then the same fundamental point of view is to be asserted. The possibility of
freedom is to be seen as preserved; and since the \emph{whole} spiritual
development is here seen as concentrated upon the \emph{ethical} development,
there lies in the preserved possibility of ethical freedom at the same time the
possibility of an essential restitution of those functions of spirit which are
the conditions of the ethical.

% APOSTROPHE, p. 173: „Saulus's Forhold til Christendommen“ --- the printed form
% carries an apostrophe and an „s“; kept as printed, and the name is not
% Anglicised (proper names unchanged, per the playbook).
% DASHES, p. 173: both are the PRINTING'S --- „Christendommen. ---“ closing the
% first paragraph, and „ved Villiens Energie. ---“ opening the next movement
% without a paragraph break.
% SHORT PAGE, p. 173: the type block ends one line short of the measure, with
% white space below and NO rule; p. 174 opens a new paragraph.
% --- p. 174 ---
If the disturbance brought about by the evil will is seen as a disturbance in
\emph{other individuals' physical} existence, or in \emph{their spiritual}
conditions of life, then it falls, as a \emph{relatively} disturbing
intervention in these, under the same point of view as the disturbances in the
individual itself treated above. For the other individual the possibility of
freedom is, in spite of it, preserved, and the disturbance sets an \emph{ethical
task} for the individual, the possibility of whose solution is secured by that
possibility of freedom. If, on the contrary, the disturbance is the
\emph{total} annihilation of the whole physical existence, then we are
referred, as where the individual annihilates its \emph{own} existence, with
respect to the question of the ethical task's solution for the individual, to
the problem of the individual's eternity, which will be treated in what follows.

If the disturbance in the other individuals is a disturbance in their
\emph{moral} existence, then, under all the nuances of the relation, the other
individual's \emph{imputability} under the influencing is to be accentuated, so
that the reception of this influence is seen as \emph{willed}; and the point of
view for restitution which was the original one with respect to the acting
individual's will is then also to be asserted with respect to that of the
suffering one.

If finally the disturbance is a disturbance in the \emph{objectively universal},
in the \emph{moral order}, then the decisive thing here is what has already been
brought out above: the good forces' reaction against what disturbs, in that this
precisely by the struggle against the good calls forth an increased development
of force in those who fight for the good, so that in the whole's development a
compensation takes place, which holds also for the case that whole ages have a
direction towards evil; for then ages are to be seen in relation to ages as
before individuals in relation to individuals.

Since in all these relations the disturbance is seen as an object for the
consciousness of the individual from which it proceeded, it becomes, when the
principled breaking through of the
% EMPHASIS, p. 174: the opening line is NOT letterspaced; the spacing pass's flag
% on „tilveiebragt ved“ is an artefact of a loosely justified first line.
% --- p. 175 ---
good will has occurred, an object for repentance, an incitement for a working of
will in the individual which in relation to the other individuals and to the
moral order must come forward also as a striving after an \emph{outward}
restoration and compensation so far as this stands in the individual's power.
But the main accent falls everywhere on the restitution in \emph{the
individuals' inner being}.

\begin{center}\rule{2cm}{0.4pt}\end{center}

The foregoing developments have everywhere essentially taken account of that
form of the evil will which is posited as the \emph{primitive: that of defiance,
of self-wilfulness}. But in accordance with this primitive form's relation to
the particular forms of evil it finds application to these as well. Whether evil
comes forward as sensuality, animality, dullness, or as sloth and inactivity,
yet in all, even the most passive forms, this is to be held fast: that evil, as
a determination of will, must at every point be willed, that it has never come
to be purely passively nor is continued purely passively, but always comes to be
and is continued by a consent of the will. Therefore, where evil has the form of
\emph{weakness} in all its possible nuances, \emph{defiance} is constantly
present, and every willing of evil conceals within itself the evil will's
primitive determination: to strive to put itself in the Absolute's place. The
evil willing is thus everywhere an expression of an energy, but of an energy
which, broken in itself, works against itself, and therefore bears weakness,
powerlessness, concealed within itself, just as in its pride it conceals
self-abandonment, the direction downwards, towards that which is lower than
spirit.

Just as what has been developed holds of the different forms of evil, so it
finds application also to the different shapes under which the change of will
shows itself. Immediately it finds application to the form where the will to the
good energetically comes to a breaking through and energetically asserts itself
after the breaking through in life's concretion. But it finds application also
to those cases,
% RULE, p. 175: a short centred single hairline standing MID-PAGE with no head
% adjacent --- a thought-break, as at pp. 19, 28, 42, 62, 97 and 159, and so
% transcribed in both files. The batch pass's report that it sits off centre is
% wrong; it is centred within the scatter of the book's other rules.
% --- p. 176 ---
where what is sporadic in the self, even after the ethical's principled breaking
through, works against this, and where the will, which in repentance turns away
from evil, is weak and hampered and at every moment wavers over towards evil.
For when this wavering does not lead to a principled relapse, --- in which case
the general point of view for the evil will's working is to be asserted, ---
then the will's oscillation becomes through repentance an incitement for a
heightened straining of will, through the feeling of the threatening danger an
incitement for a more inward seeking back to those sources of reconciliation
which are found in the human and in the human's divine principle.

Where in the single individual conversion, the ethical's breaking through,
cannot be shown at all in any moment of the individual's life, the question of
reconciliation points back to the question of the relation of the individual's
time-determined life to the individual's determinateness by eternity, to its
eternal significance in existence's organisation. This question will be treated
in what follows.

\begin{center}\rule{2cm}{0.4pt}\end{center}

Since the individual as reconciled through its ethical willing relates itself to
\emph{the Absolute}, this shows itself decisive for it in the \emph{ethical}
determinations whose applicability to it was indicated in the foregoing (p. 78).
\emph{The Absolute} is seen as \emph{the good}, in that it, as through-working
the finite will, is this will's ethical τέλος, its ideal. It is seen as
\emph{the holy} in that its energy is determined by the law of essential
ideality. It is seen as \emph{the just} in that it lets the will's result be
determined according to its fundamental relation to the principle. It is seen
finally under the determination of \emph{love} as that which, communicating
itself and in the communication of essence revealing itself, joins the single
spirits together in its unity and comprehends them mutually in that harmony of
love which has spirit's unity for its ground.
% RULE, p. 176: a second mid-page thought-break hairline, transcribed. p. 176
% does NOT close with a rule; its type block runs the full depth.
% GREEK, p. 176: τέλος, tight (not letterspaced), copied from the Danish.
% DASHES, p. 176: the pair round „i hvilket Tilfælde …“ is the printing's own.
% PENCIL, p. 176: a modern reader has ruled a line through „lader“; marginalia,
% not transcribed in either file.
% --- p. 177 ---
\setcounter{footnote}{0}
In the developed ethical-religious relation of the individual to the Absolute
there lies a \emph{rational equivalent} for that which the \emph{positive}
religion seeks in the individual's turning to the \emph{personal} God through
\emph{prayer}: a relation in which, as shown elsewhere\footnote{Cf. Problemet om
Tro og Viden, pp. 95 ff.}, the fundamental peculiarity of positive religion is
concentrated. There is in that ethical-religious relation a \emph{turning of the
mind} to the through-working Absolute in order, through the unity with this, to
attain the lacking ethical force. There is in this turning a moment of
\emph{contemplation}, in that the Absolute is cognised according to its
fundamental relation to the human being and the human will. There is in it a
moment of \emph{action}, in that the will in self-surrender freely subordinates
itself under the divine will. There is a \emph{strengthening of the mind}, a
\emph{satisfaction of feeling} in this relation, in which the essence's sources
spring forth from the Infinite's depth. There is in it a \emph{sure attainment
of the goal striven for}; for in the full self-surrender itself is the full
self-affirmation, a winning of the Eternal's force, in that the individual
freely loses itself to the Eternal. There is finally in this relation a freedom
from the egoism which must remain when in prayer one prays to an arbitrary will,
even when one prays for the forgiveness of sins; for also in this prayer there
is, when it becomes the prayer for the remission of a punishment arbitrarily
determined by the divine personality's arbitrary will, a relation of externality
to the divine, in which finitude and precisely therefore egoism is preserved.

\section*{β) The Ethical as a Determination of Totality for Human Life.}

As the \emph{second} condition for the ideal's realisation it was demanded above
(p. 168) that the ethical must be able to become a \emph{determination of
totality} for human life, not merely in spite of what has been neglected and
dis\-%
% HEAD, p. 177: „β) Det Ethiske som Totalitetsbestemmelse for Menneskelivet.“ ---
% BOLD (halbfette) Fraktur, centred, set over two lines; hence \section*, unlike
% the letterspaced-regular run-in heads of pp. 159, 162, 165, 167 and 168. The
% enumerator β) is ANTIQUA ITALIC against the Fraktur of the rest; the Greek
% letter is typed directly, as in the Danish. The Indhold agrees.
% RULE, p. 177: a short centred hairline stands ABOVE the display head; head
% ornament, not transcribed, as at pp. 97, 126, 164 and 168.
% FOOTNOTE, p. 177: the only note on the page, printed *), called after „viist“
% and BEFORE the comma --- the anchor is kept at the same word in the English.
% EMPHASIS, p. 177: several split runs, all left as printed --- „Contemplationens“
% and „Handlingens“ are spaced but the following „Moment“ is TIGHT in both, so
% the English emphasises "contemplation" and "action" and not "moment"; likewise
% „Totalitetsbestemmelse“ is spaced only through „Totalitets“.
% --- p. 178 ---
turbed by sin, but also in spite of that limitation in time for the possibility
of an ethical willing which is given by the human conditions of development. How
far this demand can be fulfilled is now to be considered, chiefly from the last
point of view, since what concerns the first has already been treated in the
foregoing.

In the apprehension of the \emph{ethical will} as \emph{essence}-willing and of
the \emph{ethical action} as an expression of the \emph{essential} determination
there has been a pointing to the demand that the ethical shall be a
determination of \emph{totality}. The ethical demand is, as our \emph{essence's}
demand and as the absolute will's demand through our \emph{essence}, the
\emph{unconditioned} demand. It shows itself as such in a \emph{twofold} sense:
as that which as a categorical imperative has validity unconditionally for
\emph{all} human individuals, and as that which has validity unconditionally for
\emph{all} the individual's relations of life, for the individual human whole of
life. The demand for an ethical totality of life is thus justified by the
ethical's relation to the essence; but in order that this demand should be able
to be actualised under the individual's development in time, it is required on
the one side that the \emph{pre-ethical life}, i.e. the life which is led before
the development of consciousness necessary for the ethical choice has taken
place, can be \emph{taken over ethically}; on the other side, that the ethical
action can acquire a concretion corresponding to the personality's concrete
content of life.

That the \emph{pre-ethical life is taken over ethically}, in such a way that the
stage of life which, in accordance with the human spirit's psychological
conditions of development, \emph{must} lie prior to the ethical, is drawn in
under the ethical and becomes a moment in an ethical whole of life which has
unity and continuity, will more closely mean that the individual, in
constituting itself in the ethical choice as an ethical personality, in this
concentration comprehends its preceding life under a \emph{new principle of
life}, and that, in this concentration as it were beginning its life afresh, it
transforms the development which showed itself as a \emph{preceding condition}
for its \emph{true}, i.e. \emph{essence-corresponding}
% EMPHASIS, p. 178: three split runs of the kind recorded for pp. 6 and 15, all
% left as printed and all reproducible in English because the compound divides at
% the same joint: „Væsens=“/„villen“ spaced only in the first half -> "\emph{essence}-
% willing"; „Væsensbestemmelsen“ (whole, on the next line) again spaced only in
% „Væsens“ -> "the \emph{essential} determination"; „Totalitetsbestemmelse“
% spaced only through „Totalitets“ -> "a determination of \emph{totality}".
% Likewise „ubetingede“ and „dobbelt“ are spaced but their nouns are not.
% ID-EST MARK, p. 178: TWO ɔ: marks, both -> "i.e.".

% --- p. 179 ---
\emph{form of existence}, to \emph{determination within this existence}, and by
this transformation, \emph{ethically} seen, becomes \emph{\textit{causa sui}}.
That which in the pre-ethical life was \emph{natural determinateness,
endowment}, then becomes an \emph{ethical task of freedom}; and by the fact that
the preceding naturally determined development, the working of the not yet
ethicised forces, is ethically apprehended \emph{through repentance}, the result
of the pre-ethical life and of the pre-ethical action is metamorphosed, and in
this metamorphosis taken up into the sphere of the ethical principle of life, in
order to be worked through out of it, so that the principled transformation of
essence is impressed upon the concrete personality's ethical unfolding.

But since the personality is ethically to work through the result of its
pre-ethical life and as ethical to live a life determined all through by the new
principle of life, the demand for \emph{an adequate concretion of the ethical}
asserts itself. In order that the ethical should be able to form the
personality's natural foundation through and through, and out of the ethical
choice to form a totality of life in whose single actions it becomes the
determining principle, it is required that it be unfolded for consciousness in a
concretion corresponding to the concretion of life and of action, and that in
this concretion, in which it comprehends an \emph{empirically cognised}, it
preserve the \emph{ethical-apriori's certainty and necessity}. Since the
\emph{apriori} consciousness of the good, of duty, of virtue, is in the actuality
of action to become a consciousness of this determinate good, of this determinate
duty, of this determinate virtue, and since this consciousness is, through the
connection with the former, to preserve the former's \emph{unconditioned
certainty}, it is therefore required that the \emph{abstract consciousness of
essence} be unfolded into an \emph{essential consciousness of actuality}, out of
which action can proceed, so that in the concrete the decision does not, through
an \emph{immediate aesthetic tact's relative unconsciousness}, step into the
place of the ethical resolution borne by conscious reflection, which right up to
the last punctual sharpening towards actualisation rests upon a universal and
for the ethical subject's consciousness asserts its validity by letting itself
be led back to the laws of normative universal determinations.
% BROKEN RUN, p. 178/179: the emphasis begun at „til \emph{Væsenet svarende}“ on
% p. 178 continues here in „\emph{Existentsform}“, with a TIGHT „til“ dividing
% the two groups. Two runs in the Danish, two in the English --- rendered
% "\emph{essence-corresponding}" / "\emph{form of existence}" so that the joint
% falls at the page break exactly where the printing puts it.
% LATIN, p. 179: „causa sui“ is antiqua AND letterspaced --- \emph{\textit{…}},
% as at p. 68, not the tight \textit{…} of pp. 64, 71, 83, 89 and 91.
% --- p. 180 ---
But against such a demand upon the ethical consciousness there then sets itself
the \emph{twofold} difficulty, that the conditions of consciousness of human
life bring with them the periodic interruption of the clear life of
consciousness by the relatively and totally unconscious, and that human
cognition, though not hampered by a fixed, immovable limit, preserves a relative
abstraction. By this twofold difficulty the ethical continuity of life and that
punctuality in the ethical action which is conditioned by cognition's unfolded
concretion seem to be prevented.

As concerns the \emph{last} point, there lies indeed in the ethical-apriori's
relation to the metaphysical fundamental concepts, in the relation of human
nature's side of reality to the essence's ideal unity, and in the convertibility
of the naturally real coming forward in the drive's relative independence into
ideal moments of essence, a possibility of a constantly progressing development
towards concretion of the cognition of the ethical-apriori. But in the fact that
\emph{the relation of experience} becomes a moment in the cognition upon which
ethical action rests, there lies the unavoidability of a partial incompleteness
of cognition in every action's presuppositions, a persisting, though mobile,
limitation of cognition. In spite of this incompleteness and limitation, which
must work upon the clarity and compass of the consciousness out of which action
proceeds, the action can nevertheless not merely in general preserve its ethical
truth and its essential validity as ethical action, because it is in
\emph{principle} determined by the acknowledgement of the unconditioned, of the
absolute demand, by the ethical fundamental resolution, but it can preserve a
\emph{more concrete} correspondence to the punctual demand by that \emph{totality
of the conscious will} which corresponds to the totality of theoretical knowledge
of actuality grounded upon the apriori and as it were extensible from within, and
which, as determined by \emph{the apriori in the ethical}, has \emph{real}
validity and can not merely unfold itself from within to concretion through
cognition's progressing assimilation of the real, of the content of actuality;
% FOOTNOTE, p. 180: NONE, though the page stands in the OCR-derived hint list ---
% the hint is wrong, as at p. 166.
% --- p. 181 ---
but which in a joined-together unity works in the form of immediacy with the
sureness of genius.

By this last determination we are referred to that which removes the difficulty
first touched upon above, which lay in the periodicity of the clear life of
consciousness; namely to what we might call \emph{the working of the ethicised
nature}. By this, which has its first springing-forth from \emph{the self's
metamorphosis in the ethical choice}, there will be able, in the ethical life
constituted by this choice, in spite of that periodicity, for \emph{the totality
of what proceeds from the essence's activity} to have the character of \emph{the
ethical}; and the relatively and totally unconscious will, from this point of
view, as little abolish the ethical as on the theoretical's territory the
incessant relation of alternation between unconsciousness and consciousness,
even in the act of consciousness itself in the stricter sense, abolishes the
theoretically normal. \emph{In the ethicised nature's working of essence}, which
has \emph{immediacy's} form, but the form of \emph{resulted immediacy},
\emph{all} the acts of the advancing life will then be taken up as co-working
moments, after having gone back from \emph{consciousness's form} to \emph{the
unconscious}, and in this form enclosed by the essence's unity.

This working of the ethicised nature, by which the ethical continuity can be
preserved and cognition's incompleteness find a supplement, would nevertheless
still, like that unfolding of ethical cognition, leave us in the approximative
and in what is not in the last instance securing, if we did not go beyond
individuality's limit. But we are referred from it back again to what was
brought out in the section on what is at work in reconciliation: to the
influence, from the beginning fructifying and animating, thereafter upholding
and clarifying, which proceeds from the \emph{moral community} by which the
individual is spiritually enclosed, and from the \emph{ethically supplementary
individualities}. And definitively we are referred back to that which is the
ethical community life's and the individual ethical life's \emph{absolute
principle}; and what we have called the working of the ethicised nature we then
see at the same time as a working of the \emph{universal principle},
% EMPHASIS, p. 181: the long run „I den ethicerede Naturs Væsensvirken“ takes in
% its opening „I den“, and the „men den“ between „Umiddelbarhedens“ and
% „resulterede“ is tight; the English follows both.
% PRINTER'S DEFECT, p. 181: „ville da alle det fremskridende Livs Acter vare
% optagne“ --- „vare“ for „være“. Orthographic and Danish-internal; it does not
% travel, and the English sets "be taken up" without comment in the text.
% --- p. 182 ---
to which the self enters into relation in the ethical resolution and is
abidingly set in relation through \emph{love's self-surrender and appropriation
of essence}; and in this working of the Absolute we would then in the last
instance seek that which completes the ethical in the individual.

By this relation we are then referred back again to what was developed earlier
concerning \emph{the ethical's relation to the religious}. Since the ethical's
goal becomes the full actualisation of that human being which as self-conscious
must become conscious of itself according to its fundamental relation to the
divine, and since this goal can be attained only by the divine's working, the
ethical, in whose first coming-to-be the religious relation was present, cannot
be apprehended in its development otherwise than as everywhere penetrated by the
religious, in unity with which it first acquires its full actuality. \emph{The
single ethical relations} then become, since at \emph{every} point they are seen
as actualising a relation of the personality to the divine, and as having come
to be through this, to be apprehended as \emph{religious relations}, and
\emph{the whole of human life in its ethical unity} as \emph{religiously}
determined. In this unity of the ethical and the religious the \emph{ethical}
has its \emph{truth}, the \emph{religious} its \emph{actuality}.

\section*{γ) The Individual Human Life's Determinateness by Eternity Corresponding to the Ideal's Concept.}

As the \emph{last} condition for reconciliation's full actualisation it was
posited above that the individual's life can acquire that determinateness by
eternity which corresponds to the ideal's concept. The ideal, as an expression
of what is completed, has, also when it is seen under the point of view of
activity, energy --- and thus must spirit's ideal be apprehended --- the stamp
of \emph{the eternal}, while the individual life strives towards the ideal
through a \emph{development in time}, and more closely through such a one as,
when regard is had to the organic conditions of life, shows itself as concluded
by a final limit of time. In this indi\-%
% HEAD, p. 182: bold (halbfette) Fraktur, centred, over two lines; hence
% \section*, like β) on p. 177. The enumerator γ) is ANTIQUA ITALIC against the
% Fraktur of the head. The Indhold agrees.
% RULE, p. 182: a short centred hairline stands ABOVE the display head; head
% ornament, not transcribed, exactly as on p. 177.
% EMPHASIS, p. 182: the closing antithesis spaces only „Ethiske“, „Sandhed“,
% „Religiøse“ and „Virkelighed“, leaving „det“ and „sin“ tight both times; the
% English emphasises the four nouns and adjectives only.
% DASH, p. 182: the pair round „og saaledes maa Aandens Ideal opfattes“ is the
% printing's own.
% --- p. 183 ---
vidual's determinateness by time in its relation to the ideal's determinateness
by eternity lies the difficulty.

For the levelling out of the immediate incongruence there offers itself most
immediately the thought of the time-determined individual's \emph{immortality},
of a continued individual existence for it after the dissolution of its real
organism. But when this immortality, as an existence continued into infinity, is
apprehended as it is apprehended by the ordinary religious consciousness, and as
for example \emph{Kant} apprehended it, in that the dualism between nature and
reason, drive and freedom, led him to lay down an infinite progress, --- which
nevertheless, from another point of view, showed itself to him as concentrated
in a moment of eternity, --- then it does not help, even if it be thought as a
continued development, towards an actualisation of the ideal, but only towards
an infinitely continued approximation. One would attain eternity by a
\emph{quantitative} way through an infinite extension of time, and through a
development seen under time's point of view; but this is the impossible: such a
time infinitely extended from a definite point of beginning never becomes
eternity's correlate, its full unfolding. But when then the doctrine of an
immortality of the soul as such an infinitely continued individual existence
cannot remove the doubt concerning the ideal's realisation, then, from the
standpoint from which it has proceeded, the emphasis might be laid upon this,
that immortality should be represented as a being which is not affected by
finitude's conditions of existence, by its distinction of externality and its
form of time, generally as a being which, as belonging to a \emph{higher life},
has no conditions in common with the being of the life in time. Hereby, however,
partly the concept of a \emph{development} would fall away, partly
individuality's being as \emph{this} individuality, in continuity of
consciousness with its preceding development, would become incomprehensible; and
precisely upon this continuity the greatest weight is laid. But by that
modification a point of view might meanwhile be indicated which, conceptually
developed, might lead to seeing the individual's existence in time in another
light than that in which it showed itself only as an infinite approximation.
% EMPHASIS, p. 183: „Kant“ is letterspaced FRAKTUR, not antiqua --- proper names
% stay in the Fraktur fount and are emphasised by letterspacing, so \emph{} and
% not \textit{}.
% READING, p. 183: no dash before „Herved“, contrary to tesseract.
% The paragraph continues across the page break; no blank line before p. 184.
% --- p. 184 ---
There might be indicated a point of view from which the eternal shows itself as
that which can assert itself as energy ungoverned by determinateness by time, as
that which therefore will also be able to reveal itself as the eternal through
determinateness by time, to assert itself under the development in time as that
which it is according to its essence.

This then refers us back to what was developed earlier about \emph{the
significance for eternity of the ethical time}-determined life, and about
\emph{the concentration of the whole of life}, under the development, in a
\emph{moment of eternity}. \emph{The ethical life}, which rests upon a
\emph{relation of eternity}, becomes that which makes human existence a concrete
unfolding in determinateness by time of the eternal. In the life in love, in the
self-surrender to the principle posited by individuality's essential energy, a
life is lived with determinateness by eternity which in every moment of
development concentrates itself in an ideal point of unity, in which the
development in time is ideally taken back into a moment of eternity. And since
the individual's life in time ceases, since it as a whole is seen as comprehended
in a moment of eternity (cf. Kant), as concentrated in ideal eternal unity, this
individual existence, seen from the comprehending's point of view according to
its significance as a peculiar member in the eternal life's organisation, has
attained its τέλος, and has realised determinateness by eternity in spite of the
limitedness of existence in time.

What has here been said with respect to an individual life, seen in general as
an \emph{ethical} life, is now, when human existence generally is seen
\emph{from the universal totality's point of view}, to be extended not merely to
those forms of life which more weakly and more incompletely express the ethical,
but also to those forms of life which seem to be filled out only with an
unethical content and to be concluded or broken off without any breaking through
of the ethical. When the point of view is taken in that totality in which the
single human existences are comprised, then the goal for the individual
existence is seen as that which \emph{must} be attained, since \emph{the
totality's life cannot} be disturbed; and what the positive religion expresses
thus, that \emph{every} being attains
% BROKEN RUN, p. 184: the emphasis „Evighedsbetydningen af det ethiske tids-“
% stops INSIDE the broken word „tidsbestemte“, „bestemte Liv“ being tight. The
% English compound divides at the same joint --- "\emph{…the ethical
% time}-determined life" --- so the run count and the joint both hold.
% GREEK, p. 184: τέλος, tight, copied from the Danish. The page is NOT in the
% recon pass's Greek list; the survey is wrong in both directions here.
% EMPHASIS, p. 184: „ethvert“ is spaced but the following „Væsen“ is not, and
% „Totalitetens Liv ikke kan“ is spaced but „forstyrres“ is not; the English
% follows both splits.
% READING, p. 184: „(jfr. Kant)“ has a lower-case „jfr.“, unlike the „Jfr.“ of
% the note on p. 177; the English sets "cf." and "Cf." to match.

% --- p. 185 ---
that goal which is determined for it by the divine foreknowledge, acquires here
the form that \emph{every} individual, as a member in the universal existence's
organisation, must realise its τέλος, and in the realisation be ordered into
eternity's totality of life.

This conception is nevertheless not to be identified with that which from a
merely \emph{natural necessity's} point of view lets every individual attain its
goal, i.e. fulfil its \emph{fate}. The individual is seen from this point of
view as, in its determinateness by an \emph{outer} necessity, attaining what it
\emph{can} attain; and by this conception the ideal's concept is abolished. In
distinction from this we do indeed let that which the individual \emph{can}
attain, in accordance with its connection with the totality, become the measure
of what it actually attains; but we double the point of view in such a way that
what \emph{can} be attained is at the same time seen as that which \emph{shall}
be attained, thus as a \emph{task} which shall be realised by \emph{freedom},
and so first as an \emph{ethical ideal}. By this \emph{``shall''}, which is
\emph{directed to the will}, the moment of \emph{freedom} therefore comes to the
moment of necessity and enters into unity with it. The goal which shall be
attained becomes \emph{the ideal} which is realised with freedom, and this ideal
is \emph{the individual's essence according to its eternal validity}: the object
for the ethical choice, and that which realises itself through the concrete
ethical action. But this ideal, the \emph{individual} ideal, is at the same
time, as the source of that ``shall'', the \emph{universal} ideal. This is to
say: the essential demand upon the individual points out beyond the
\emph{individual} to the \emph{totality} into which the individual is ordered
according to its fundamental relation, and into which it shall with freedom
order itself through self-surrender. Thus there is through the moment of freedom
a pointing back again to necessity; but without freedom disappearing in it.

But when now the fact that the single individual, because its life is set within
the totality, must attain its determination, acquires the sense that it shall
with freedom realise its goal, by realising its ethical task, then it becomes
the question how what was said above about a
% GREEK, p. 185: τέλος, tight, copied from the Danish.
% EMPHASIS, p. 185: the two quoted „skal“ differ and both are left as printed ---
% the FIRST, „Ved dette „skal“, der er rettet til Villien“, IS letterspaced
% inside its quotation marks, so \emph{``shall''} with the quotes inside the run;
% the SECOND, „som Kilde til hiint „skal““, is TIGHT, so plain ``shall''. Also
% „Naturnødvendigheds“ is spaced but „Synspunkt“ is not, and „Frihedsmomentet“
% is spaced only in „Friheds“ --- a split inside an unbroken word, as on p. 178;
% the English emphasises "freedom" and not "moment".
% --- p. 186 ---
life which was posited as ethical without closer limitation can find application
to the different imperfectly ethical and unethical forms of life, and to human
lives whose development is broken off or disturbed.

If such an application is to become possible, there must be demanded a
\emph{moment of breaking through} and a \emph{moment of completion for all human
individuals}, whatever the apparent conclusion of life may be; and in this
moment of completion, whether the breaking through has gone before it or
coincides with it, the essence's force, divided in the preceding life but
preserved, will be joined together into an undivided energy in which the goal of
life is attained. The extension in time of this ``moment'' of completion becomes
the unessential. What psychological experiences show as possible in a moment
where life is threatened with annihilation: that the individual's whole
preceding life is in a few seconds drawn forth before consciousness in clear,
connected images, accompanied by an ethical judgement upon the content of life
which they present, can also find its application to the relation under
discussion. \emph{The moment of death} as the boundary between two forms of
existence for the individual, between that of existence in time (the phenomenal
form of life) and the form of eternal being, would then, however infinitely
small its determinateness by time may be, be able to become the time of
completion and give place for the thoroughgoing transformation. In this
completion the goal would then be attained, whatever the kind of the will may
have been in the preceding life; but in the good will's preceding working
individuality's essence-energy would have had direct, positive significance for
the goal's realisation, in the evil will's the energy would have worked
negatively, but in such a way that it, through the dialectical in the relation,
worked indirectly towards the goal's attainment. And in the good as in the evil
will's working, as the decisive and completing, as that which in the last
instance alone can actualise the ideal, the through-working \emph{divine will}
would be to accentuate.

But in order that this view of a moment of completion for all should not make
the ethical significance of the life in time
% READING, p. 186: „Individualitetens Væsens-“ / „Energie“ --- the second half
% opens with a CAPITAL E, so the hyphen is a compound hyphen and not a line-break
% artefact; the word is kept as „Væsens-Energie“, against the whole lower-case
% „Væsensenergie“ of p. 184. The English keeps "essence-energy" here.
% EMPHASIS, p. 186: „Gjennembruddets“ is spaced, the „og et“ between tight, and
% the run resumes at „Fuldendelsens Øieblik for alle menneskelige Individer“ ---
% two runs, and two in the English. „Øieblik“ inside its quotation marks lower on
% the page is NOT letterspaced.
% --- p. 187 ---
illusory, then it becomes \emph{partly} a matter of holding fast that the
difference in the individual's relation to the ethical essential demand through
the whole existence in time must mirror itself in an \emph{essential difference
of consciousness} which determines the individual's whole theoretical and
practical relation; \emph{partly} there must be thought a reflex of the
individual's ethical relation in its being as an eternal moment in the total
organisation of life. This last will nevertheless not be able to be thought in
such a way that in an existence of eternity there comes to be a disturbing
recollection of evil, or a ``metaphysical recollection'' of having been able to
become other than what this eternal existence is; --- for thereby the attainment
of the goal demanded by the totality, which can only be the good's goal, would
be made impossible, and the contradiction would be posited that the individual
moment's knowledge of itself in eternity's form acquired a finite content (the
discord, sin's punishment). But it can, when an eternal being is demanded for
all individuals in harmonious unity with the Absolute, only signify that the
\emph{intensity} of that life which is taken up as a moment in the organisation
of the eternal fullness of life expresses \emph{the ethical life in time's
content of eternity}.

Is there now hereby demanded a continued \emph{individual existence with
continuity of consciousness}? This question has been answered affirmatively by
supporting oneself upon \emph{the ethical obligation's character}. One infers:
individuality must have eternal validity, i.e. immortality as this conscious
individuality, as this individual will, because the ethical demand as
unconditioned has eternal validity, and the individual, if it were perishable,
could not be eternally obligated. But in this reasoning it is overlooked that,
even if the individuals' conscious being as these individuals is posited as
perishable, there remains the individuals' \emph{eternal essence} in which the
demand is grounded, and that the demand would therefore just as fully preserve
its \emph{unconditioned} validity for the individual as an expression of this
essence, as a moment in its process of life, and that it would preserve its
validity for it in every moment of its existence. It also becomes essential to
hold fast that the ethical de\-%
% DASH, p. 187: „…er; --- thi derved vilde“ is the printing's own.
% READING, p. 187: no comma between „er“ and „det“ in „Men i dette Raisonnement
% er det overseet“; the mark the ABBYY layer reads is an ink speck.
% --- p. 188 ---
mand, as a demand of eternity upon the \emph{time-determined individual},
becomes a demand which shall be actualised in \emph{time}, i.e. in a
\emph{determinate time}, so far as it relates itself to \emph{the individual as
individual}, and that it relates itself to time generally, and to the eternity
comprehending time's totality, only in that the individual, while it is an
ethical end in itself, becomes a moment for humanity's and existence's whole.

But when thus the kind of the ethical demand and of the ethical obligation does
not enclose the necessity of an individual immortality in \emph{the sense
indicated}, there yet lies in this, that the individual's life in time acquires
through it a determinateness by eternity in itself, since the individual relates
itself to a task of eternity, that there must be posited an \emph{eternal
essentiality in the individual}. When the eternal essence is seen as the source
of a demand whose fulfilment acquires time's form and the stamp of the
time-determined individuality, then the individual to which the eternal
essence's demand is directed and which actualises this demand cannot be a
\emph{merely perishable}, nor its mode of being a merely vanishing and
insignificant one. This then leads us back to the question of \emph{the form}
under which the individual human spirit's eternity can be thought.

The starting point for the answer must be \emph{the psychic-organic conditions
for individuality's existence}. The real human organism in its peculiarity
shows itself, in accordance with a relation of essence, as the
\emph{indispensable condition} for the existence of the only spiritual
individuality we know, of the human individuality in its peculiar mode of being
and peculiar development. There is therefore nothing which entitles us to posit
the possibility that this individual existence of spirit can be without the
organism. But we do not for that reason seek individuality's \emph{last ground
in this}. The last ground of the organisation itself could indeed not be sought
in the real, but had to be sought in that principle of ideality which, in
accordance with its \emph{essential} relation to the real, comprehends the real
elements into unity. And thus individuality's last ground becomes
% BROKEN WORD, p. 187/188: „For\-“ / „dring“ = Fordring; the English divides
% "de\-" / "mand" at the same joint, as at pp. 168/169.
% READING, p. 188: „psychisk=“ / „organiske“ breaks across a line but the hyphen
% is a real one --- the same compound stands WITHIN a line on p. 189. Kept as
% „psychisk-organiske“ in both places, "psychic-organic" in the English.
% --- p. 189 ---
to be sought in \emph{the principle's determinateness by ideality}, in the
self-articulating ideal. This relation, by which there is therefore a pointing
back to a higher principle which also grounds the organism, encloses meanwhile
in its generality no guarantee for the human \emph{individual's immortality}.
The eternal principle reveals itself also through the individualised selfless,
and this is nevertheless in its individuality something changing and vanishing.
Immortality's guarantee would then have to be sought in the determination of
\emph{selfhood} in the human individual, and more closely in \emph{that form of
selfhood} which makes \emph{freedom} possible. In freedom there lies in reality
a guarantee that the self cannot be annihilated; in every expression of freedom,
be it positive or negative, the self reveals its determinateness by eternity,
shows itself as that which cannot vanish as something essenceless. But on the
other side freedom points precisely back to the through-working divine, and the
highest act of freedom is that of self-surrender, of union of essence with the
principle. The question then repeats itself concerning \emph{the form of the
free self's} preservation when the organic conditions of individuality have
vanished, concerning its eternal being after the conclusion of this phenomenal
life in time. When on the one side that form of life which is lived in this
organism cannot be lived when the organism is dissolved, but on the other side
the self which has unfolded its life of spirit on the basis of the
psychic-organic, and which has freely revealed its determinateness by eternity,
cannot be annihilated, then it can only be thought as taken up \emph{as a moment
in the Eternal's life}, in that unity out of which the articulation of life
eternally proceeds. But can there in this being be thought a \emph{consciousness
of a preceding development}? To this it must be answered that when
consciousness's organic conditions are dissolved, consciousness \emph{as such}
cannot be. But what \emph{metamorphosis} can then be thought to take place with
it? Such a one that the individual spirit in its concentration of eternity is as
a moment in God's knowledge of himself in his eternal kingdom of spirit, in his
eternal fullness of life's organisation,
% EMPHASIS, p. 189: „Selvhedsbestemmelsen“ is spaced only in „Selvheds“, within
% an unbroken word --- the third such split in this stretch (cf. pp. 178 and
% 185); the English emphasises "selfhood" and not "determination". „Formen for
% det frie Selvs“ is spaced but „Bevaren“ is not.
% The paragraph continues across the page break; no blank line before p. 190.
% --- p. 190 ---
and --- as was already said above --- a moment whose intensity expresses the
ethical life in time's content of eternity, its eternal significance in the
organisation of the whole of life.

\begin{center}\rule{2cm}{0.4pt}\end{center}

For a clearer elucidation of the conception of reconciliation developed in the
foregoing we set it beside, on the one side, the \emph{abstractly metaphysical}
conception, on the other side \emph{Christianity's} doctrine of reconciliation,
in which we see the most perfect expression of the \emph{positive religion's}
conception.

\begin{center}
1. The Abstractly Metaphysical Conception of Reconciliation.
\end{center}

As the most prominent representatives of this we bring forward \emph{Spinoza}
and \emph{Hegel}. We leave out of account in this connection what is peculiar to
the first of them, that existence, when it is seen from the point of view of the
single, limited, modifications, shows itself as an infinite chain of causes and
effects, governed by the law of blind natural necessity, so that the ethical
must disappear in the physical, ethics become a doctrine of nature about action.
We lay the weight rather upon the \emph{essential likeness} between the two
thinkers' conception, which lies in the \emph{abstract conception of eternity}
common to both, which has its source in a one-sided idealism unconsciously
affected by dualism. This likeness, which becomes clear when in \emph{Spinoza}
existence is seen from \emph{substance's} point of view, that of the eternal and
the infinite, is not disturbed by the apparent equal entitlement which the
natural, the real, acquires in substance with the spiritual and the ideal, and
which might point to a corrective for the abstract conception of eternity, and
to an essential significance of the real's fundamental forms for the ethical
life. For in reality the \emph{ideal mo}ment, which originally formed the
starting point for the conception of substance and was the source of the
concept of substance, a prepon\-%
% RULE, p. 190: a short centred SINGLE thin rule, standing between two paragraphs
% and NOT under a display head (the head follows a further paragraph below);
% transcribed, as at pp. 19, 28 and 42. Contrast the DOUBLE rule at the foot of
% p. 201.
% HEAD, p. 190: „1. Den abstract metaphysiske Opfattelse af Forsoningen.“ ---
% centred on one line in ORDINARY Fraktur: neither halbfette nor letterspaced.
% Hence a bare centred group, with NEITHER \emph nor \section* --- the only head
% of this shape in the book. The Indhold reads „1) den abstract metaphysiske
% Opfattelse af Forsoningen“: a different enumerator („1)“ for the printed „1.“)
% AND a lower-case „den“ for the printed „Den“. Both divergences are real and are
% left standing; the printed head governs.
% NAMES, p. 190: „Spinoza“ and „Hegel“ are letterspaced Fraktur with a TIGHT „og“
% between them --- hence two \emph groups and not one, in both files.
% EMPHASIS, p. 190: „Christendommens“ is spaced but „Forsoningslære“ is tight;
% „positive Religions“ is spaced across its line break but „Opfattelse“ is not;
% and „ideelle Mo“ is spaced with the „ment“ completing the word tight --- the
% English divides "\emph{ideal mo}ment" at the same joint.

% --- p. 191 ---
\setcounter{footnote}{0}%
derance, and just as it becomes the overreaching in relation to the other
attributes, it \emph{definitively} determines the point of view for existence as
absorbed in substance's eternity.

\emph{The consequences} which the \emph{abstract conception of eternity} and the
\emph{one-sided idealism} out of which it proceeds bring with them are
essentially the following.

Out of that conception nothing that takes place in time can acquire
significance, i.e. \emph{essential significance}, because time becomes only an
\emph{untrue} form of existence, and does not stand in an essential relation to
the Absolute, or to what exists, seen according to its truth; for thus it is
seen only \textit{sub specie æterni}. Only when \emph{time}, as the necessary
form of existence for the \emph{essentially valid} real, is acknowledged as
having \emph{essential} significance for \emph{all} that is, can also the
\emph{ethical} existence, as an existence in time which relates itself to a task
of eternity, be apprehended according to its \emph{significance for actuality}.
Only then can the relation of time to the ethical, the ethical choice's
determinateness by time, the ideal's realisation through a totality of time, the
loss of time as grounding a guilt, and the making good in time of this guilt
acquire decisive significance. Only then can it become the task to see the
determination of eternity in the ethical individual as realising itself in
concretion in determinateness by time.

Out of that conception it becomes impossible to acknowledge \emph{freedom} in
the individual in its relative distinction from the Absolute --- a distinction
which is precisely conditioned \emph{by the essentiality of the determinations
of reality for the principle}. Everything must be drawn in under the universal
and eternal's metaphysical process of necessity, in which the individual must
become something merely powerless and vanishing\footnote{When thus \emph{Strauß}
--- in his Dogmatik --- seeks to make a place for freedom, for evil and for a
reconciliation, the movement becomes purely illusory, since the fundamental
point of view is that one-sided idealism's abstract conception of eternity and
of necessity.}. Only when the form of reality is ascribed essential
% LATIN, p. 191: „sub specie æterni“ is antiqua and TIGHT here --- \textit{…}.
% Contrast p. 192, where the same phrase stands twice on one page, spaced the
% first time and tight the second.
% FOOTNOTE, p. 191: the only note on the page. „Strauß“ is letterspaced Fraktur,
% with the ß as printed; „Dogmatik“ is tight and stays in Danish, per the rule
% that Danish work titles in notes are not translated. The OCR-derived hint list
% gave p. 192 as note-bearing; it is p. 191 that carries the note.
% --- p. 192 ---
validity, and the real's significance for freedom's determination by ideality is
discerned, can individual freedom be asserted.

Out of that conception, then, \emph{reconciliation} can be seen only as
determined by the \emph{eternal} process's \emph{necessity}, and it must be
posited as realised in the individual \emph{thought's} elevation above the
finite to a \emph{thinking} of \emph{the Absolute} (the Eternal) and of itself in
the Absolute. In this thinking the individual's \emph{thought}, which is its
\emph{essence's} expression, as determined by the Absolute, goes up into this as
an abolished moment, and the conditions of time and the distinction of time
become insignificant, unessential, whether the relation be apprehended in
Spinozistic or in Hegelian fashion. Since the essential relates itself
exclusively to the determination of eternity, which expresses the \emph{one true
being}, and since the \emph{true cognition}, which is reconciliation's
\emph{form}, is \emph{\textit{sub specie æterni}}, there is essentially
\emph{nothing} neglected and therefore in the word's strictest sense nothing to
restitute (redintegrate). In that true cognition, being \textit{sub specie
æterni}, there vanishes, as the relation of time loses its significance, that
content of life which is unfolded in a \emph{time prior to the true} cognition,
as a \emph{semblance} with the vanishing form of time; it is abolished in
eternity's eternally present actuality, in which not even a reflex can be
preserved of the finite individual life which, as \emph{only} borne by a
difference of nature, is \emph{only} borne by something vanishing, something
essenceless.

\section*{2. The Christian Doctrine of Reconciliation.}

That the \emph{Christian} doctrine of reconciliation is chosen as
\emph{representative} of the \emph{positive} religions' conception of
reconciliation has its justification in this, that in Christianity the
consciousness of sin, of the discord in the human spirit and of its distinction
from the divine, has attained an intensity and seriousness as in no preceding
religion, and that therefore the doctrine of reconciliation in it could acquire
the richest development and give the sharp expression for what in the other
religions was only unclearly and uncertainly indicated.
% LATIN, p. 192: „sub specie æterni“ stands TWICE on this page --- SPACED antiqua
% at its first occurrence (\emph{\textit{…}}) and TIGHT nine lines later
% (\textit{…}). The clearest instance in the book of the antiqua's spacing
% varying within a single page; both are carried over mechanically.
% HEAD, p. 192: BOLD (halbfette) Fraktur, centred, on one line --- hence
% \section*, against the ordinary-weight centred head of p. 190. The Indhold
% reads „2) den christelige Forsoningslære“: a different enumerator and a
% lower-case „den“. Left standing; the printed head governs.
% EMPHASIS, p. 192: „Tid forud for den sande“ is spaced and the „Erkjendelse“
% completing the phrase on the next line is tight; the English follows.
% --- p. 193 ---
There is in the present connection first to be taken into account
Christianity's conception of sin and its effect, and thereafter its conception
of reconciliation and of the abolition of sin's consequences.

\begin{center}
\emph{a) The conception of sin and of its effect.}
\end{center}

Decisive for the conception here becomes at once the determination of sin as
\textit{ἀνομία}, (1. Joh. III, 4) as transgression of God's law, as disobedience
to God's commandment. For since the divine commandment which is transgressed in
sin is apprehended as the commandment of an arbitrary divine will, the
commandment's content must show itself as the contingent and in a deeper sense
indifferent, and the emphasis must fall upon the demand for a blind obedience to
the arbitrary will (cf. Augustin). The transgression of this demand and of the
arbitrary commandment is then not set in relation to man's and more closely to
the human will's determination of \emph{essence}; only upon the \emph{formal}:
upon the disagreement with the commanding divine will, does the emphasis fall.
And since this will stands as that which by transcendence is distinguished in
essence from the human, and since its arbitrary determination, which is declared
in the commandment, cannot express a determination of essence in the human
\emph{will}, the transgression, as not touching a relation of essence, cannot be
comprehended as having an essence-determining effect upon the human will and
human nature. When therefore the Christian doctrine lays down such an effect, it
comes under the determination of the paradoxical.

The \emph{paradoxical} comes forward already where sin, as something contingent,
is for the first time brought into relation with the human. Prior to the
\emph{first sin}, which was to acquire decisive significance for the human
race's whole, there is represented a primal state in which the still sin-free
first man, who was created in God's image, had an intellectual and moral
perfection which makes the fall an \emph{incomprehensible}. And the
\emph{paradoxical} then continues, in that in the Christian Church, under the
development of the Church's doctrine, the
% HEAD, p. 193: centred, letterspaced REGULAR Fraktur, not bold --- same
% treatment as the centred heads of pp. 165 and 167. Absent from the book's own
% Indhold, which jumps from „2) den christelige Forsoningslære“ (192--201)
% straight to „Slutning“: this and its fellow at p. 195 are body-only heads,
% with no Indhold reading to compare against.
% GREEK, p. 193: „ἀνομία“, antiqua, the only Greek on pp. 190--201.
% CITATION, p. 193: „(1. Joh. III, 4)“ is kept VERBATIM, with the Danish stop
% after the numeral and the Danish abbreviation of the name --- scripture and
% work references are not Anglicised, per the playbook's rule for citations.
% „Augustin“ likewise stands as the author sets it, not as "Augustine".
% EMPHASIS, p. 193: „Væsens“ is spaced and the „bestemmelse“ completing the word
% is tight --- the same split as at pp. 178, 185 and 189; the English emphasises
% "essence" and not "determination".
% --- p. 194 ---
first sin's effect is determined as the corruption of human nature and of the
human will by sin --- an effect which had to be seen as a punishment, imposed by
the divine wrath. And the fact that human nature and the human will are
corrupted by the first sin --- the first man's sin --- becomes, in relation to
the \emph{other indi}viduals in the race, apprehended in such a way that the
corruption in nature propagates itself to the whole of the human race, spreads
from generation to generation as original sin, and that the transgression is
imputed to the generations as inherited guilt after the forefather. Original sin
becomes the first sin's \emph{punishment}, and the guilt which rests upon the
race must constantly be increased by the race's new offence. The effect of sin
which is laid down comes under the \emph{paradoxical} both when it is considered
in relation to the individual who is represented as sin's beginner, and when it
is seen in relation to the other individuals. It is a \emph{paradox} that the
actualisation of a possibility laid down in the human being should be able to
recreate the essence itself, that the expression of will, which is an expression
of essence, should be able to bring with it the human being's corruption,
irreparable by man. It is a \emph{paradox} that sin, the individual free will's
deed, and the guilt which it brings with it, should be transferred to the still
will-less individuals existing only as possibilities, and that it should
predetermine their action in that that forefather of the race's sin corrupts the
race's nature in such a way that sin becomes for the succeeding individuals the
unavoidable. And the \emph{paradoxical} finally also asserts itself in this way,
that, under the presupposition of original sin, of the human being's
disturbance, \emph{the consciousness of sin} must consistently be posited as
lacking in man, so that it can be awakened only by a divine \emph{revelation};
but when revelation generally, as \emph{miracle}, shows itself to belong under
the paradox, then it is particularly striking that \emph{such} a revelation about
sin to men existing in sin, transformed by sin, must do so.

When the \emph{positive religion's} determinations of \emph{sin's con\-%
% READING, p. 194: „deres Handling idet hiin Slægtens Stamfaders Synd“ is printed
% with NO punctuation between „Handling“ and „idet“; reproduced.
% EMPHASIS, p. 194: „andre In“ is spaced and the „divider“ completing the word is
% tight --- the English divides "\emph{other indi}viduals" at the same joint.
% --- p. 195 ---
cept} and \emph{sin's effects} are held together with those of the
\emph{rational ethics}, the opposition shows itself everywhere. It shows itself
in the conception of sin's essence, in that sin becomes for the \emph{positive
religion} the transgression (or the inclination to transgression) of an
arbitrary divine commandment, formal disobedience, a breach of the feeling of
dependence upon omnipotence; for the \emph{rational ethics}, on the contrary, a
violation of the essence's demand, and precisely therefore of the Absolute's
demand expressed through the human being. It shows itself in the conception of
\emph{sin's punishment}, which for the \emph{positive religion} becomes
determined by the divine will's \emph{arbitrariness}, for the \emph{rational
ethics} by the essence's law, by the consequence of egoism's isolation. It shows
itself in the conception of \emph{the essence's disturbance} by sin, which for
the \emph{positive religion} becomes an incomprehensible \emph{transformation of
essence}, for the \emph{rational ethics} a discord in the essence, a suspension
of its integrity, under which the essence is present as negatively working, as
preserving the possibility of redintegration. It shows itself finally in
\emph{the consciousness of sin}, which according to the \emph{religious}
conception must consistently be awakened by a supernatural influence, according
to the \emph{rational ethics} proceeds naturally out of the essence's negative
working in the evil will. Under the \emph{opposition} to the \emph{rational
ethics'} determinations, \emph{all the positive religion's} must come under the
\emph{antirational}, the \emph{paradoxical}.

\begin{center}
\emph{b) The conception of reconciliation and of the abolition of sin's effects.}
\end{center}

Since \emph{sin} and \emph{sin's punishment} are for the \emph{positive
religion} knit together by the representation of the \emph{divine will},
\emph{not} by a \emph{relation of essence}, the abolition of what is posited by
\emph{sin} is to be seen from a \emph{twofold} point of view: as abolition of
\emph{sin's effects} (\emph{the punishment}), and as abolition of \emph{sin's
actualised principle}, of the sinful determinateness of will; and these two
points of view meet only in so far as \emph{sin it\-%
% BROKEN RUN, p. 194/195: the emphasis „\emph{Syn\-“ / „dens Begreb}“ is ONE run
% carried across the page break inside a broken word. The English breaks
% "\emph{sin's con\-“ / „cept}" at the same place, so the run count, the page
% joint and the broken word all hold together.
% PRINTER'S DEFECT, p. 195: „det Absolutes gjennm Menneskevæsenet udtrykte
% Fordring“ --- „gjennm“ for „gjennem“, a dropped „e“. Orthographic and
% Danish-internal; it does not travel.
% HEAD, p. 195: centred, letterspaced regular Fraktur, NOT bold, set on TWO
% printed lines with the turn after „Ophævelsen“; like p. 193's head it is absent
% from the Indhold.
% --- p. 196 ---
self}, apprehended as \emph{original sin}, becomes \emph{sin's punishment}.

Since \emph{the punishment} is posited by the \emph{divine will}, the
\emph{first condition} for \emph{the punishment's} abolition becomes, in
accordance with this will's nature, the \emph{objective} one: \emph{God's
forgiveness}. In the conception of this a doubleness shows itself in the
Christian doctrine. On the one side the divine forgiveness is represented, as
corresponding to the arbitrary determination of the punishment, as
\emph{arbitrary}, as \emph{independent} of determinate conditions of
reconciliation; on the other side the divine \emph{justice} is accentuated, and
this justice must demand an \emph{adequate compensation}. But since those two
opposed ways of seeing stand against one another, \emph{that} representation of
the divine \emph{arbitrary will} shows itself as the \emph{overreaching} one,
and it determines the opposed way of seeing in such a way that the divine
justice is thought \emph{to let itself} be satisfied by a compensation which
only by the divine will becomes \emph{taken for} a compensation. The
arbitrariness in the choice of the means of reconciliation is not merely
strongly brought out by the Scholastics, but in the different nuances which the
theory of satisfaction receives under the development of the Church's doctrine,
the moment of arbitrariness, which is expressly acknowledged in the
\emph{acceptilation theory}, is throughout preserved. The means of
reconciliation by which the divine will \emph{lets itself} be satisfied is
\emph{the God become man's self-surrender, Christ's death}; by this redemption
and reconciliation are brought about, and reconciliation is the incarnation's
end.

To the \emph{objective} condition for the abolition of the punishment of sin
there then comes the \emph{subjective}: \emph{faith}, as appropriation of the
\emph{objective satisfaction}. Faith as the \emph{justifying faith} is the faith
that Christ's \emph{death} is satisfying for \emph{this} believing individual.
By this faith, not by works, which, in order to have truth, are to have faith as
their presupposition, it is that reconciliation is won.

By faith is abolished that \emph{effect of sin} which
% BROKEN RUN, p. 195/196: „\emph{Syn\-“ / „den selv}“ --- again one run across
% the page break inside a broken word; the English breaks "\emph{sin it\-" /
% "self}" at the same joint.

% --- p. 197 ---
consists in the divine wrath, and the blotting out of guilt is won. But since
the blotting out of guilt and the remission of the punishment enclose the
abolition of that punishment which is \emph{sin itself}, then in faith sin also
as \emph{determinateness of will}, as a power of will, is seen as abolished in
principle, and the beginning of an action in agreement with the divine will, of
\emph{sanctification}, as given; and regeneration's principled transformation
acquires actuality of life through a \emph{mystical union} with God. Here, then,
the two points of view distinguished above meet; but they meet only to be
separated again.

That they are separated again has, apart from the fact that sin's punishment is
originally posited as \emph{arbitrarily} determined, its ground in \emph{faith's
relation to the human} and in its \emph{content}. Since faith, as will be shown,
has \emph{not} its \emph{origin} in the \emph{human}, in the being out of which
the sinful act has proceeded, and whose true transformation can only be thought
as proceeding out of its own centre, it must retain a relation of externality to
the human inner being which makes the transformation inexplicable. And to the
same result we are led from the point of view that faith, as justifying faith,
with respect to its \emph{content}, and thereby with respect to its
\emph{essence}, does not become so determined that it can become an actual
\emph{principle of life}, wherefore it must again come into a relation of
externality to the true action.

This peculiarity in faith in religion's conception stands in the closest
connection with the fact that religion in a \emph{dualistic} manner lays
\emph{all good} over into \emph{God} and deprives \emph{man} of it, in that it
lets man have possessed the good as a gift only in an imaginary past. It will
become clear when the relation between the divine and the human activity in
reconciliation is more closely considered.

When human nature is posited as corrupted, God's image as lost and the will as
bound to sin's necessity, then, consistently, no activ\-%
% --- p. 198 ---
ity can be thought as proceeding from human nature and the human will which can
contribute to reconciliation. All such activity must then proceed from
\emph{God}; and just as reconciliation's \emph{objective} condition was given by
him alone, in that he became man, died for men's sins and received this death as
satisfaction, so must also the \emph{subjective} condition, \emph{faith}, be
given by \emph{God} by an \emph{act of grace} of the divine will, and be given
\emph{as a reflection of the arbitrary forgiveness}, as a \emph{faith} that
\emph{this} believing individual is reconciled with God. When then some men
become believing, it is not in virtue of their own merit, but in virtue of the
divine grace's choice; when others remain as non-believers, thus as
unreconciled, and therefore as forfeited to damnation, it is because the divine
grace is withdrawn from them, because they are left to their own forces. In this
representation of the election of grace, which is the \emph{unavoidable}
consequence of the Christian conception of the fundamental relation, the
conception of God's will as the \emph{absolute arbitrariness} comes forward in
its greatest sharpness, and the opposition between the positive religion and the
rational ethics becomes the greatest possible. With right does \emph{Calvin}
call that divine resolution to salvation and damnation, incomprehensible for
man's thought, \emph{\textit{et horribile decretum}}; and it sounds like a
terrible mockery when \emph{Luther} with respect to those not elect interprets
the Gospel's words: ``do this, and ye shall be saved!'' thus: ``only try it, and
ye shall see that it is not possible!'', and when he from God's \emph{revealed
resolution} to \emph{all} men's salvation, as it were with an unconscious irony
upon \emph{revelation}, distinguishes a \emph{hidden resolution} to damnation.
And yet Calvin and Luther are, when they draw the most frightful consequences,
in agreement with the presuppositions which religion gives and, according to its
fundamental way of seeing, must give.

In the Christian doctrine of the abolition of the punishment of sin we meet, as
in the doctrine of sin's essence and consequences, everywhere the
\emph{paradoxical}. We meet it in the representation of \emph{the incar\-%
% ANTIQUA, p. 198: „et horribile decretum“ is antiqua AND letterspaced, and the
% antiqua run begins at the DANISH article „et“ --- the article is inside the
% \textit group, as printed. Contrast p. 199, where „hiint horribile decretum“
% has „hiint“ in Fraktur and the Latin in TIGHT antiqua.
% QUOTES, p. 198: two pairs, both balanced; the closing mark of the second
% stands BEFORE its comma, as printed, and the English follows.
% --- p. 199 ---
nation}, of the eternal, infinite God's becoming and living in the individual
human nature's limitation, under a human life's development in time. We meet it
in the doctrine of \emph{satisfaction}: in the presupposition of the necessity
of a satisfaction of the divine will by a divine deed, and of a possibility of
this deed's acceptance as satisfying. We meet it in the representation of a
\emph{transference} of righteousness, of that which, as a determination by the
free will, can only be thought as acquired by the free will's own working, while
now man is represented as relating himself purely passively in reconciliation.
We meet it finally in \emph{the arbitrariness in the divine election}, in the
representation, annihilating of all ethical force, of an absolutely groundless
determination to salvation and to damnation.

If we compare reconciliation as it becomes an object for consciousness on the
\emph{positive religion's} territory and within the \emph{rationally ethical's}
sphere, then there shows itself, as concerns the positively religious, a defect
both with respect to reconciliation's \emph{certainty} and to its
\emph{completeness}. As concerns the \emph{certainty}, it is evident that so
long as that \textit{horribile decretum} stands as the background for the
religious consciousness, no rest can be found in a faith in the reconciliation
brought about. What comes forward for consciousness as effects of grace cannot
with certainty be cognised as such: the possibility of a delusion constantly
remains, precisely because these effects are thought as proceeding from
\emph{another being} which by \emph{transcendence} is distinguished from the
human. And the will which groundlessly elects and groundlessly damns can at
every moment take its deluding effect of grace back, let man sink down into
damnation. Thus is anxiety constantly kept awake; there is no rest, no peace.
--- And as concerns reconciliation's \emph{completeness}, the defect in the
positive religion's conception shows itself \emph{first} in this, that the
abolition of \emph{sin's punishment} is again separated from the abolition of
\emph{sin's ground} in the individual. The
% BROKEN RUN, p. 198/199: „\emph{Men\-“ / „neskevordelsen}“ --- one run across
% the page break inside a broken word; the English breaks "\emph{the incar\-" /
% "nation}" at the same joint.
% LATIN, p. 199: „horribile decretum“ is TIGHT antiqua here and „hiint“ before
% it is Fraktur --- \textit{…} with the demonstrative outside, against the
% spaced \emph{\textit{et horribile decretum}} of p. 198.
% DASH, p. 199: „ingen Fred. --- Og hvad“ is the printing's own.
% --- p. 200 ---
\emph{restitution} which, according to the positive religion's doctrine as it is
developed in the Church's greatest representatives, in reality takes place, is
the \emph{remission of the arbitrarily imposed punishment}, thus the
\emph{represented} restitution through the reconciliation with the
\emph{arbitrary} divine will, in accordance with the \emph{arbitrary}
forgiveness of the transgression of the \emph{arbitrary} commandment; but the
restitution is \emph{not the redintegration of the essence and the will},
\emph{not} the \emph{actual} restitution, the restitution of the force for the
good, of that state which was represented as the first men's. Let one compare
\emph{Augustin's}, \emph{Luther's} and \emph{Calvin's} pronouncements, and one
will find that they, in spite of man's regeneration as believing and in spite of
the believer's mystical union with God, posit the human will as constantly drawn
to sin, as constantly entangled in evil, and that they hope for the true
sanctification only in a \emph{future} state, in another life in which human
life's real conditions are abolished. And they do this consistently in
accordance with the dualistic distinction between God and man with respect to
the good touched upon above. On account of the distinction lying in the
\emph{abstract transcendence}, reconciliation with God on the positive
religion's territory cannot become the \emph{full} reconciliation in the
consciousness of \emph{the unity of essence} with the principle. \emph{To this}
form of reconciliation there is a pointing only through the \emph{mystical}
element in religion; but \emph{reflection}, which awakens in religion,
unavoidably presses it back again by giving weight to transcendence's
distinction. --- If reconciliation is seen as the individual's reconciliation
\emph{within itself}, then its \emph{incompleteness} shows itself not merely
through the struggle of the will entangled in evil with the better cognition,
but also in what was brought out earlier, that the religious consciousness,
which has its point of support in what is inaccessible to reason, can live its
life only by negating the rational, thus by negating an expression of spirit's
essence, and must therefore preserve a discord in its inner being.

On the \emph{rationally ethical's} standpoint reconcil\-%
% DASH, p. 200: „Vægt. --- Sees Forsoningen“ is the printing's own.
% NAMES, p. 200: „Augustins“, „Luthers“ and „Calvins“ are letterspaced Fraktur;
% the genitive is Danish and is rendered with the English apostrophe.
% --- p. 201 ---
\setcounter{footnote}{0}%
iation's \emph{certainty} is given in spirit's consciousness of its \emph{unity
of essence} with the \emph{indwelling} principle, in the cognition of the
reconciling force as \emph{spirit's expression of essence in its unity with the
through-working Absolute}. And reconciliation's \emph{completeness} is made
possible by the fact that reconciliation becomes not merely the individual's
\emph{principled reconciliation with its origin} and in \emph{its own inner
being}, but that there is, as has been shown in the earlier developments, in
spite of that breach in the ethical life's development which occurs by the evil
action, and in spite of the will's wavering and relapse, preserved a possibility
of bringing about a \emph{redintegration}, of abolishing discontinuity in
continuity, of making the ethical a determination of totality for human life,
and thus of realising the ideal ethical demand\footnote{Both in the
ecclesiastical conception and in the usual attacks upon it there is lacking ---
even in \emph{Anselm}, who has come nearest to it --- the point of view that the
actualisation of the task, of the ideal demand, is conditioned by a continuity's
being brought about in the ethical striving, or by what has been neglected in
time being able to be made good again in time, which seems to be made impossible
by the fact that every moment in time has its task, and therefore seems unable
to yield place for the positive filling out and ordering into the whole of an
earlier moment not filled out or filled out in a negative direction. Thus there
is for example in \emph{Strauß}, --- in whom there is talk only of an abstract
necessity in the single individual's development, in that each single one
becomes what the whole's necessity of development brings with it, which is
apprehended as freedom only because the whole actualises itself only in the
individuals, --- also only talk of again \emph{in general} being able \emph{to
act well} after the individual's cognition of the right has led the will in the
right direction, not of acting \emph{so} well as could have been acted without
that interruption by the unethical. In reality by such a conception the single
action becomes loosened from the essence and seen as purely \emph{accidental}.
--- How the difficulty with respect to the relation of time is to be removed has
been developed in the foregoing.}.

\begin{center}\rule{2cm}{0.4pt}\end{center}
% FOOTNOTE, p. 201: the only note on the page. „Anselm“ and „Strauß“ are
% letterspaced Fraktur. Inside the note „Almindelighed“ is spaced, the „at“ after
% it tight, and „kunne handle godt“ spaced again --- two runs, and two in the
% English, with "being able" tight between them.
% RULE, p. 201: the division closes below the footnote with a short centred
% DOUBLE rule --- the only double rule so far besides those at pp. 4, 29, 63 and
% 210. Set with the single \rule the skeleton uses throughout, exactly as the
% Danish does; the printed weight is recorded here and in transcription.tex.
% p. 201 is a SHORT page: text, then the note, then the closing rule.

\chapter*{Conclusion.}
\addcontentsline{toc}{chapter}{Conclusion}
\markboth{Conclusion}{Conclusion}

% --- p. 202 ---
\setcounter{footnote}{0}%
% ENUMERATION 1.--3., pp. 202--208. In the print the numeral stands at the
% PARAGRAPH INDENT and the following lines run flush to the left margin, not the
% hanging setting of pp. 1--4. Set here with the same enumitem call as pp. 1--4,
% for consistency across the edition, exactly as transcription.tex does; the
% divergence in the printing is recorded rather than reproduced.
\begin{enumerate}[label=\arabic*., leftmargin=*, labelsep=1em,
                  topsep=0.35\baselineskip, itemsep=0.35\baselineskip]
\item From all the earlier developments it follows as a result that there can
  come to be an essential unity between the religious consciousness and the
  human consciousness, ``the consciousness of culture''; and this unity can and
  will, because it is grounded in human nature, assert itself in spite of the
  fact that the religious, where it still has the form of the \emph{positively
  religious}, comes into \emph{opposition} to the human as the essence of both
  begins to be clarified for consciousness. In this opposition the religious
  consciousness holds fast the dogma of the, supranatural, the supranatural
  principle of cognition and the supranatural principle of will, while the human
  consciousness holds fast nature's essentiality as the foundation of human
  life, and the untransformed, complete human nature's capacity to attain out of
  its own inner being to the truth of cognition and of action, whose source the
  human being conceals within itself. This opposition, which, within
  Christianity, came forward for consciousness at the close of the Middle Ages,
  was clarified, under the unfolding of philosophical thought, more and more in
  the succeeding centuries\footnote{Cf.\ Problemet om Tro og Viden, pp.\ 21
  ff.}, and there developed, from realistic as well as from idealistic
  presuppositions, an opposition to the Christian view of the world and of life,
  and philosophy's strife against the religious either directly determined the
  ordinary consciousness or at least indirectly affected the development by
  which in modern times this
% PRINTER'S DEFECT, p. 202: the page prints „Dogmet om det, Supranaturale“, with
% a COMMA between „det“ and „Supranaturale“. It is a comma and not a speck, and
% it is one of the few Danish defects that DOES travel: the English reproduces it
% as "the dogma of the, supranatural". Not in the Trykfeil leaf.
% EMPHASIS, p. 202: „positivt Religiøses“ spaced, „Form,“ tight; the FIRST
% „Modsætning“ spaced, the second, eight words later, TIGHT --- the English
% emphasises only the first.
% --- p. 203 ---
\setcounter{footnote}{0}%
  consciousness removed itself from the Christian and from the positively
  religious generally.

  There could indeed, under the development of the opposition, and while this
  development in part led to open struggle, still for a time \emph{partially} be
  preserved an uncertainty. A particular form of philosophy could\footnote{Cf.\
  Problemet om Tro og Viden, p.\ 23.}, in that by its one-sidedly idealistic
  principle it acquired an abstract likeness to the Christian view of the world,
  deceive itself and seek a unity where the opposition was nevertheless what was
  of principle. And since the Christian had constantly had beside it a
  development which was affected by it, even if it had its root in another
  principle of life and was borne by the view of nature and of human life
  inherited from antiquity, which the new age independently reacquired, then
  that which had developed contemporaneously with Christianity, in part under
  its influence, but out of another principle, could be apprehended as something
  which had proceeded from Christianity itself, so that this was accordingly
  seen as the principle which out of its inner fullness of life could unfold a
  content of actuality, in science, in art, in the forms of moral and political
  life. And thus there could still in our time stand, over against the view
  which, in laying down a relative unity of the Christian and the human, sees in
  Christianity the principle for the shaping of a world-actuality in religious
  spirit, that view which apprehends Christianity as relating itself only
  negatively to the natural and to human life's unfolding in real actuality, and
  which has its representatives just as fully within the positively religious as
  within the human. But there will not be able to subsist an uncertainty about
  which of the two opposed conceptions is the true and the historically secured
  one, when one does not let oneself be confused by the unavoidable
  inconsistencies into which a religion which posits the denial of the world,
  the denial of nature, as its principle must come, when it is to
  \emph{continue to subsist in the alien} world whose destruction within a short
  term was the original expectation.
% EMPHASIS, p. 203: only two runs. In the second, „vedblive at bestaae i den
% fremmede“, the following „Verden,“ falls back to TIGHT type in the middle of
% the phrase --- another half-executed run, reproduced as printed.
% --- p. 204 ---
  Let one but analyse the Christian representations of the purely spiritual God,
  of the creation out of nothing, of the God-man, of the state of blessedness
  and of the destruction of the world, and Christianity's \emph{abstractly
  spiritualistic} character, its \emph{denial of the natural} and thereby also
  of \emph{human life's content of reality}, will become indubitable. And since
  there is now united with this what Christianity must share with \emph{every}
  positive religion: \emph{the placing of truth's, of theoretical and practical
  truth's, principle in a transcendent being}, over against which man stands as
  the merely finite, as that which by its nature is powerless to find the truth,
  and since this conception is seen as sharpened by the Christian representation
  of the state of sinfulness and of the first sin's consequences, then the
  opposition between the Christian and the human view of the world and of life
  must show itself as complete.

  But under this opposition between the positively religious in its highest form
  and the human consciousness's principle and content there shows itself
  nevertheless not merely the possibility, but the necessity, of a
  reconciliation in spirit, by which the human and the religious are brought to
  unity; and to such a reconciliation the development of the modern age itself
  has led.

  When the modern consciousness became clear about its opposition to the
  Christian-mediaeval one, the opposition to the transcendent supranatural could
  by its abstraction lead to a \emph{conception of nature and of human life in
  terms of finitude}; but this conception ends in consequences which must lead
  spirit back to the acknowledgement of the religious according to its
  \emph{essential} significance, according to its \emph{content of infinity}.
  The human consciousness is forced by the ``real actuality's'' finitude and by
  its own finitude to seek an infinite, and finitude conceals within itself the
  trace of that infinity by which it has being. Consciousness finds the infinite
  again, in that the \emph{nature} which for materialism became something full
  of contradiction and inexplicable shows itself
% p. 204 opens FLUSH, not indented: it continues the paragraph begun on p. 203,
% although p. 203 ends on a full stop.
% EMPHASIS, p. 204: the run „Fornægten af det Naturlige“ stops at „og
% derigjennem ogsaa af“, which is tight --- INCLUDING the „af“ that governs the
% next spaced run. The English keeps that "of" outside both runs.
% --- p. 205 ---
  for spirit in \emph{the light of life and of ideality}, and \emph{in that
  consciousness sinks into itself}, and cognises its \emph{relation} to the
  \emph{living} actuality's \emph{ideal-real principle}. In this sinking-in the
  self finds again the \emph{religious relation}, finds it again by
  \emph{knowledge's way}, and finds it again \emph{without giving itself up} and
  without denying a \emph{fundamental activity} of its essence. The
  \emph{human} consciousness becomes, in seeking down to its essence's ground,
  in truth \emph{ethical} consciousness, and as such \emph{religious}
  consciousness. Only by this relation of unity does the human consciousness
  acquire its \emph{truth}, as surely as the human spirit first in the deepest
  sense becomes \emph{conscious of itself} when it becomes \emph{conscious of
  itself as grounded in the divine}, and as freely subordinated to this. And the
  religious acquires in the unity with the human \emph{concretion and
  actuality}, and it has in this its concrete actuality its \emph{truth} as
  \emph{principle for the actual human life's totality}. In this unity the
  religious reconciliation becomes \emph{complete}, because religion's life is
  not lived by a flight from human life's actuality, by a denial of essential
  forms of its expression: a denial which must always become self-contradicting,
  since it becomes possible only by the forces of human nature itself; --- but
  is lived in \emph{every} essential form of spirit's life, and has its
  \emph{positive} expression in these forms of life, so that the religious
  relation becomes an \emph{absolute} relation, comprehending the \emph{whole}
  of life, and connecting the human spirit with its \emph{essential} Absolute,
  with \emph{the in truth infinite}. And in this unity is, with the
  \emph{absolute relation} and with \emph{reconciliation}, --- spirit's
  reconciliation with \emph{its principle} and in \emph{its own inner being},
  --- \emph{what is essential in the positive religions} preserved, while the
  historical form is dissolved, the transcendence of representation abolished in
  the relation of immanence in which transcendence becomes a moment.

\item The \emph{essential task} of the present work can accordingly be
  determined as this, \emph{to point to that form of religiosity which is
  present}, though often
% DOUBTFUL READING, p. 205: the print gives „dens væ=/senslige Absolute“, with s
% and not t --- „væsenslige“, not „væsentlige“. The transcription keeps the
% printed form. The distinction is invisible in English, both readings giving
% "essential"; it is recorded here so the pair stays comparable.
% EMPHASIS, p. 205: „frit underordnet dette“ is TIGHT, and the „som“ between
% „sin Sandhed“ and „Princip for …“ is likewise tight between two spaced runs.
% DASHES, p. 205: the „; ---“ before „men leves i enhver“ and the pair round
% „Aandens Forsoning med sit Princip og i sit eget Indre“ are the printing's own.
% --- p. 206 ---
  \emph{unclearly and half-consciously, in our age's human consciousness},
  which, more clearly or more darkly, has proclaimed itself in this
  consciousness \emph{from the beginning of the modern age}, which has asserted
  itself in the social and moral relations, in the endeavours after humanity, in
  the form of the life of spirit generally, and which needs only to be \emph{set
  free and clarified} through consciousness's self-deepening in order to show
  itself in its essentiality, as one with what is deepest and most moving in the
  positive religions and as one with the human spirit's inmost essence. It is
  \emph{the religion of actuality}, which is \emph{not} to be founded or
  created, but only to be shown as being and as grounded in the human spirit's
  law of development; --- which is \emph{not} to be a rationalising
  transformation of a positive religion or itself a new positive religion, but
  \emph{the expression of religiosity's eternal spirit as the whole of human
  life's principle of life}. Therefore it has been the task to show this true
  and eternal religion of actuality as that in which the \emph{total} human life
  is in unity and concord, and in which therefore the \emph{human consciousness}
  can find \emph{religious reconciliation}. The solution of the problem of
  faith's and knowledge's relation is therefore sought out of that which in the
  indivisible human being has the principate of concept, but in such a way that
  what is peculiar in the human spirit's essential forms of expression is
  acknowledged, that the human spirit's working according to its concrete
  determinateness has come to its validity. What in the introduction to the
  work ``Problemet om Tro og Viden'' was indicated has received its closer
  execution. In the human spirit's \emph{essence} is \emph{the religious
  relation's necessity} sought, and it has shown itself as the necessity of an
  in truth \emph{absolute relation}, of a \emph{true relation of infinity}, of
  such a relation in which \emph{the full reconciliation: reconciliation within
  actuality}, can be found, in that the human spirit, in the \emph{absolute
  self-surrender to the principle}, comes into \emph{inner unity with itself}.
  In this relation are \emph{religion, science and morality} seen in inner
  unity. Out of the acknowledg\-%
% p. 206 is set almost throughout in letterspaced type; the tight passages are
% the exception. „Væsen“ in „I Menneskeaandens Væsen“ is spaced while the
% following „er det“ is tight.
% DASH, p. 206: „Udviklingslov; ---“ is the printing's own.
% --- p. 207 ---
  ment of the human spirit's natural ground as something essential by the
  relation to the absolute principle, there has been a pointing to an
  \emph{ethical action} which acquires a \emph{content of actuality} by this
  natural ground, and which, in acquiring \emph{the religious's stamp, fills out
  the religious with actuality}, and more closely, with a
  \emph{thought-determined, rational actuality}.

\item The standpoint which is asserted in this work does \emph{not}, then,
  \emph{set an abstract metaphysical} in the place of the \emph{religious
  actuality of life}. It finds in the principle itself the guarantee for the
  real's essentiality, and by accentuating this it makes place for a
  \emph{relation of freedom} in opposition to the abstractly metaphysical
  relation of necessity, and secures the essential significance of the real's
  necessary mode of being for the ethical-religious relation; and hereby, and by
  letting this relation comprehend a concrete totality of life, it asserts it
  precisely as a \emph{relation of actuality}. Against this it does not tell
  that recourse is everywhere had to metaphysical determinations; for the
  metaphysical preserves its \emph{universal} validity; but it must only be so
  determined that it points \emph{to the real}, to the finite
  \emph{actuality's conditions of life}.

  The standpoint which has been asserted cannot either with right, as some have
  wished, be designated as a \emph{negative} one, and as such be brought
  together with the standpoints from which those attacks upon the Christian have
  been made which belong to our age and to the immediately preceding one. Its
  relation to the most considerable representative of the negative tendency in
  our age, to \emph{Feuerbach}, has above been set forth in detail. Through this
  presentation, and through all the earlier developments generally, it will have
  shown itself that it relates itself negatively to the positive religions only
  in that it negates the determinateness by \emph{finitude} in them, but that it
  precisely through this negation strives to assert \emph{the essentially
  religious's infinity}. From this point of view the standpoint is precisely to
  be designated as \emph{positive}: the negative
% READING, p. 207: „til den endelige Virkelighedens Livsbetingelser“ --- printed
% so, with no comma; as printed in both files.
% EMPHASIS, p. 207: in „ne=/gerer Endeligheds bestemtheden“ the compositor spaces
% „Endeligheds“ and falls back to tight type for „bestemtheden“ inside the same
% word; the English emphasises "finitude" and not "determinateness".
% READING, p. 207: „fra hvilket de Angreb … ere gjorte“ --- the antecedent is
% plural but „hvilket“ is printed; not emended, and the English follows the
% printed sense.
% NAME, p. 207: „Feuerbach“ is letterspaced FRAKTUR, not antiqua --- \emph{}.
% --- p. 208 ---
  becomes a determination for such a one as, in relating itself polemically out
  of a \emph{one-sided} starting point, \emph{can only tear down}, and for such
  a one as \emph{consciously holds fast the determinateness by finitude}. ---

  Should now nevertheless anyone, from the negative relation I take to the
  moment of finitude in the positive religions, while I strive to assert what in
  the religions is essential, eternal and infinite, wish to take occasion, in
  the manner of sycophants, to throw out the accusations well calculated upon
  ignorance, and always welcome to fanaticism, of atheism, smouldering hatred of
  Christianity, and whatever else the catchwords may be, then it shall on my
  side be answered only with the silence which is that of contempt. I have once
  in another work expressed myself about my personal relation to the positive
  religion, and an utterance of that kind it is meaningless to repeat.
\end{enumerate}

\begin{center}\rule{2cm}{0.4pt}\end{center}

Herewith I close this work and hand it over to the judgement of those who are in
a position to follow a scientific investigation with insight, and to be able to
subject it to a \emph{scientific} criticism. Such a criticism will always be
welcome to me, and however remote the standpoint from which it proceeds may lie
from mine, I shall always be ready to acknowledge whatever the criticism may
bring of what is really valid and true. But also only with such a criticism,
which rests upon a real understanding of what I have given, which has the matter
itself for its end, and which fights with grounds that science must acknowledge,
shall I be able to engage. When a recently published, so-called
``\emph{Religionsphilosophie}'' instead of refuting my proof against the
possibility of a self-limitation of knowledge, of knowledge's acknowledgement of
the validity of a principle \emph{absolutely} heterogeneous with it, confines
itself to showing the fact:
% RULE, p. 208: a short centred SINGLE hairline between the close of the numbered
% argument and the valedictory paragraphs; transcribed, as at pp. 89 and 143.
% The enumeration is closed above it, nothing after it being numbered.
% TITLE, p. 208: „Religionsphilosophie“ is a Danish work title (R. Nielsen's) and
% is kept in Danish, per the playbook's rule for work titles; here it is
% LETTERSPACED inside its quotation marks, on p. 209 it is tight.
% DASH, p. 208: „Endelighedsbestemtheden. ---“ is the printing's own.
% --- p. 209 ---
that science declares the positive religion's principle \emph{incompatible} with
its own --- (precisely what I have urged!) --- but makes no serious attempt
whatever to ground that knowledge must acknowledge this principle \emph{as
valid}, or that it can at all \emph{cognise} something as falling
\emph{outside} its limits; when this ``Religionsphilosophie'' by a simple
postulate makes the difference between abstract and concrete knowledge into an
\emph{absolute} opposition between two kinds of knowledge, and by similar
postulates would reach the conclusion that only the personal satisfaction which
the positive religion gives is the true one; and when it finally first
trivialises and then conjures away the ``absolute heterogeneity'' which was to
give it its character; --- then a discussion with such an inverted science
hovering in the air is superfluous; there is simply to be a referring back to
the criticism I have given of it in my earlier work; nothing needs to be added
to that.

From \emph{theology's} side there has hitherto been made no attempt to meet what
were the main points in my criticism of it: my \emph{denial of theology's right
as a science and of the validity of its concept of God}. I direct an appeal to
those among our theologians who have a sense for and a capacity for a scientific
investigation, to make such an attempt, and then to make it, not by the
\emph{indirect} way, by showing possible defects and contradictions in my
theory, but by the \emph{direct} way, by refuting the grounds of my criticism.
By that indirect way no final result could be attained; for even if it were made
good that I had erred ever so much in what I have \emph{positively} set up, it
could nevertheless not become of any advantage for theology, so long as it were
not shown that those errors in the positive stood in inner connection with my
criticism; but this very showing would indeed transform the indirect attack into
a direct one. Only by a direct criticism of my criticism can a result be
attained, and it seems to me that the matter has importance enough that there
ought
% DASHES, p. 209: all three are the printing's own --- the pair round „(netop
% hvad jeg har urgeret!)“, and the „; ---“ before „saa er en Forhandling“.
% --- p. 210 ---
to be striving after it. What has recently by one of our theology's most
excellent representatives been reproached to a certain theological tendency,
that it, instead of refuting serious attacks, talks them away or keeps silent
about them, ought not to be able to be made a complaint against theology
generally, also against that theology which is represented by men who really are
men of science, even if they are not so \emph{as} theologians.

\begin{center}\rule{2cm}{0.4pt}\end{center}
% p. 210 is the LAST page of the book's text and is a short page: seven lines of
% type, then the closing rule. The rule is a DOUBLE rule, as at pp. 4, 29, 63 and
% 201, and unlike every other rule in the book; set with the same single \rule
% the skeleton uses throughout, exactly as the Danish does.
% EMPHASIS, p. 210: „som“ alone is letterspaced, „Theologer.“ tight --- the only
% emphasis on the page, and the English emphasises "as" alone.

\chapter*{Errata:}
\addcontentsline{toc}{chapter}{Errata}
\markboth{Errata}{Errata}

% --- p. 211 ---
\begin{center}
\emph{Errata:}
\end{center}

\noindent Page 62, line 24: fra det [from that], read: for det [for that].

\begin{center}\rule{2cm}{0.4pt}\end{center}

\begin{center}
in ``Problemet om Tro og Viden''.
\end{center}

\noindent
\begin{tabbing}
Page \=\ 114,\ \ \=\ \ 26: \=\kill
Page \> 14, \> Line\ \ 7: \> sat, read: satte.\\
--- \> 29, \> --- \ \ 2 f.\ b.: \> findes. r.: finder.\\
--- \> 42, \> --- \ 20: \> afgivende, r.: afvigende.\\
--- \> 51, \> --- \ \ 1: \> om sine, r.: og sine.\\
--- \> 71, \> --- \ 22: \> Indhold, r.: Forhold.\\
--- \> 114, \> --- \ 16: \> dobbeltlydige, r.: dobbelttydige.\\
--- \> 119, \> --- \ 22: \> man, r.: men.\\
--- \> 176, \> --- \ 26: \> viser, r.: vises.\\
--- \> 223, \> --- \ 15: \> som at, r.: som i at.
\end{tabbing}

\begin{center}\rule{2cm}{0.4pt}\end{center}

\begin{center}
in ``Et Svar til Prof.\ R.\ Nielsen'' (2det Oplag).
\end{center}

\noindent
\begin{tabbing}
Page \=\ 22,\ \ \=\ \ 14: \=\kill
Page \> 6, \> Line\ 14: \> fører, r.: føres.\\
--- \> 22, \> --- \ \ 6: \> lønner sig, r.: lønner.
\end{tabbing}
% THE ERRATA LEAF. Handled as in the sibling volume's translation
% (../svar-nielsen/translation.tex): the RUBRIC is translated --- „Trykfeil:“ ->
% "Errata:", „Side“ -> "Page", „Lin.“ -> "Line", „læs:“/„l.:“ -> "read:"/"r.:",
% „f. n“ (fra neden) -> "f. b." --- while the WORDS BEING CORRECTED stay in
% Danish, since each erratum is about a Danish sort in a Danish line and has no
% referent in the English. A bracketed gloss is given on the one entry that
% belongs to THIS book; the other two blocks are errata for „Problemet om Tro og
% Viden“ and for „Et Svar til Prof. R. Nielsen“, not for this volume, and are
% reproduced unglossed exactly as the leaf prints them --- including the „findes.“
% that takes a FULL STOP where every other entry takes a comma.
% NONE of these corrections is applied to the body, here or in the Danish: the
% edition is diplomatic, the printed readings stand, and the leaf is set at the
% end as the author printed it. Note in particular „Side 119 … man, l.: men“,
% which is why the „man“ for „men“ at p. 168 of THIS book is NOT covered by the
% leaf: that entry belongs to the „Problemet om Tro og Viden“ block.
% The printed spelling of the heading is „Trykfeil:“, with -ei-.
% NOTE the printed spelling of the heading: „Trykfeil:“, with -ei-. The
% leaf carries three blocks, only the first of which belongs to this book;
% the other two are errata for „Problemet om Tro og Viden“ and for „Et Svar
% til Prof. R. Nielsen“, and are reproduced as printed, as in the Danish.

\end{document}
